\chapter{人体化身与动画}
\label{ch:16}

\paragraph*{本章阅读锚点}
本章的锚点是可动画的人体表示。原始 3DGS 不知道骨骼、表情和衣物；数字人方法把高斯绑定到身体模型、骨骼场或变形场上。


数字人体化身（digital human avatar）是影视特效、虚拟社交和沉浸式通信等领域的核心技术。其核心挑战在于：从有限输入（单目视频、稀疏图像或文本描述）重建高保真、可动画驱动的人体三维表示。3DGS 的实时渲染能力和显式点云式表示为数字人提供了理想的底层表示；它既能实现照片级面部和身体细节，又支持高效的姿态驱动和形变控制。

3DGS 的出现为数字化身领域带来了根本性变革。高斯基元的显式特性使其天然支持变形和动画驱动，而基于 tile 的可微分光栅化器（differentiable rasterizer）则实现了实时渲染。这一技术组合催生了大量创新工作：从单张图像生成可动画的全身化身，从单目视频重建可重光照的人脸模型，从语音信号驱动说话人头像的唇形同步，到利用物理模拟实现衣物动态——3DGS 已经成为数字化身领域最活跃的研究方向之一。

本章系统梳理了 258 篇相关论文，涵盖三个核心方向。第~\ref{sec:16_1}~节聚焦面部与头部重建，包括参数化模型、单图像重建、生成模型、重光照、头发建模和人脸分析；第~\ref{sec:16_2}~节讨论身体与运动，涵盖全身化身、铰接物体重建、人体-物体交互、场景中的人体动画和人体新视角合成；第~\ref{sec:16_3}~节介绍语音驱动、手部与局部建模及服装试穿与编辑。

\begin{figure}[htbp]
\centering
\resizebox{\textwidth}{!}{\begin{tikzpicture}[
  route/.style={rectangle, rounded corners=5pt, draw=#1, fill=#1!10,
    text width=3.2cm, minimum height=0.9cm, align=center, font=\small},
  method/.style={rectangle, rounded corners=3pt, draw=#1, fill=#1!6,
    font=\scriptsize, text width=3.0cm, align=center, minimum height=0.6cm},
  arr/.style={-{Stealth[length=4pt]}, techgray!60},
  >=Stealth
]

  \node[route=techgray, text width=4cm, font=\small\bfseries] (root) at (24.0, 5.6) {人体化身与动画方法分类};
  \node[route=techblue] (b0) at (10.4, 3.1) {人脸重建与动画};
  \node[method=techblue] (b0_l0) at (1.7, 0.6) {参数化化身\\{\tiny GaussianBlendshapes, NPGA, GraphAvatar}};
  \node[method=techblue] (b0_l1) at (6.2, 0.6) {单图/生成重建\\{\tiny FaceLift, FastAvatar, SEGA, GGHead}};
  \node[method=techblue] (b0_l2) at (10.4, 0.6) {头发建模\\{\tiny MeGA, 3DGH, HairGaussian}};
  \node[method=techblue] (b0_l3) at (14.0, 0.6) {眼动控制\\{\tiny ScaffoldAvatar, HyperGaussians, STAvatar}};
  \node[route=techgreen] (b1) at (24.0, 3.1) {身体与运动};
  \node[method=techgreen] (b1_l0) at (18.1, 0.6) {化身表示驱动\\{\tiny GaussianAvatars, GauHuman, HUGS}};
  \node[method=techgreen] (b1_l1) at (22.2, 0.6) {单图创建\\{\tiny HumanSplat, EliGS, ExAvatar}};
  \node[method=techgreen] (b1_l2) at (25.8, 0.6) {物理衣物\\{\tiny PhysAvatar, ClothGS, DressCode-GS}};
  \node[method=techgreen] (b1_l3) at (29.4, 0.6) {压缩流式\\{\tiny CompactAvatar, StreamAvatar}};
  \node[route=techorange] (b2) at (34.8, 3.1) {语音驱动与局部};
  \node[method=techorange] (b2_l0) at (33.0, 0.6) {音频驱动\\{\tiny GaussianTalker, Talk3D, GSTalker}};
  \node[method=techorange] (b2_l1) at (36.6, 0.6) {手部建模\\{\tiny HandGS, Dextrous-GS}};
  \node[method=techorange] (b2_l2) at (40.2, 0.6) {服装试穿\\{\tiny GarmentGS, TryOn-GS}};
  \draw[arr] (root) -- (b0);
  \draw[arr] (root) -- (b1);
  \draw[arr] (root) -- (b2);
  \draw[arr] (b0) -- (b0_l0);
  \draw[arr] (b0) -- (b0_l1);
  \draw[arr] (b0) -- (b0_l2);
  \draw[arr] (b0) -- (b0_l3);
  \draw[arr] (b1) -- (b1_l0);
  \draw[arr] (b1) -- (b1_l1);
  \draw[arr] (b1) -- (b1_l2);
  \draw[arr] (b1) -- (b1_l3);
  \draw[arr] (b2) -- (b2_l0);
  \draw[arr] (b2) -- (b2_l1);
  \draw[arr] (b2) -- (b2_l2);
\end{tikzpicture}}
\caption{人体化身与动画方法分类}
\label{fig:ch16_avatar_taxonomy}
\end{figure}

%=====================================================================
\section{人脸重建与动画}
\label{sec:16_1}

人脸是数字化身中最重要的组成部分，也是 3DGS 应用于人体领域最早、研究最密集的子方向。本节梳理 81 篇相关工作，涵盖人脸化身的表示架构、单图像/少样本重建、生成模型、重光照与可编辑性、头发建模、眼动控制和人脸分析等主题。本节首先讨论基于参数化模型和高保真策略的人脸化身，然后逐步深入到单图像重建、生成模型、细节建模等方向。

\begin{figure}[htbp]
\centering
\resizebox{\textwidth}{!}{\begin{tikzpicture}[
  box/.style={rectangle, rounded corners=4pt, draw=#1, fill=#1!10,
    text width=3.2cm, align=center, font=\small, minimum height=1cm},
  sbox/.style={rectangle, rounded corners=3pt, draw=#1!60, fill=#1!5,
    font=\scriptsize, text width=2.4cm, align=center, minimum height=0.5cm},
  arr/.style={-{Stealth[length=6pt]}, thick, techgray},
  darr/.style={-{Stealth[length=6pt]}, thick, techgray, dashed},
  >=Stealth
]

  \node[box=techblue] (p1) at (3.2, 0.6) {输入(视频/单图/多视角)};
  \node[box=techgreen] (p2) at (11.6, 0.6) {3DMM参数估计(FLAME/BS)};
  \node[box=techorange] (p3) at (20.9, 0.6) {高斯化身绑定(面部/头发/牙齿)};
  \node[box=techpurple] (p4) at (29.4, 0.6) {表情驱动(blendshapes)};
  \node[box=techblue] (p5) at (37.4, 0.6) {实时渲染(可控表情/姿态)};
  \node[font=\scriptsize, text=techgray] at (11.6, -0.6) {中性脸建模};
  \node[font=\scriptsize, text=techgray] at (20.9, -0.6) {逐顶点高斯};
  \draw[arr] (p1) -- (p2);
  \draw[arr] (p2) -- (p3);
  \draw[arr] (p3) -- (p4);
  \draw[arr] (p4) -- (p5);
\end{tikzpicture}}
\caption{人脸高斯化身重建与动画 pipeline}
\label{fig:ch16_face_pipeline}
\end{figure}

\subsection{参数化与高保真人脸化身}
\label{ss:16_1a}

本节讨论~2~条技术路线：（1）基于参数化模型的人脸化身；（2）高保真人脸化身。

\subsubsection{基于参数化模型的人脸化身}

本路线介绍将~3DGS~与参数化人脸模型结合实现表情驱动的混合形状方法。

如何让 3D 高斯"听懂"表情指令并做出正确的变形响应，是人脸化身的基础问题。早期工作将 3DGS 与成熟的参数化人脸模型（3D Morphable Model, 3DMM，一类通过 PCA 基向量的线性组合来控制面部形状和表情的统计模型）结合，利用已有的表情参数化能力来驱动高斯基元。

\paragraph{Gaussian Blendshapes}

Gaussian Blendshapes（2024.04）~\cite{2404_19398} 是该方向的先驱工作，提出了基于混合形状（blendshape）的 3D 高斯表示。

\textbf{Baseline 与问题}：传统参数化人脸模型（如 FLAME）虽然能通过线性混合形状控制表情，但其基于网格的表示在渲染速度和细节保真度之间存在根本矛盾：网格分辨率限制了高频表面细节的表达，而 NeRF 虽然能表达细节却无法实时渲染。

\textbf{改了什么}：Gaussian Blendshapes 将 FLAME 的混合形状机制从网格域迁移到 3D 高斯域，用一组 3D 高斯基元代替网格顶点来表示头部几何和外观。

\textbf{实现原理}：具体而言，该方法定义一个中性基础模型 $B_0$ 和 $K$ 个表情混合形状 $\Delta B_1, \Delta B_2, \ldots, \Delta B_K$，全部表示为 3D 高斯集合。每个高斯具有位置 $x$、不透明度 $o$、旋转 $q$、缩放 $s$ 和球谐系数（Spherical Harmonics, SH，用于表示视角依赖颜色的球面基函数，下文简称 SH 系数）$SH$ 五种属性。混合形状定义为各高斯属性相对于中性模型的偏移量 $\Delta B_k = B_k - B_0$。

任意表情的化身通过线性混合生成：$\hat{B} = B_0 + \sum_{k=1}^{K} \alpha_k \Delta B_k$，其中 $\{\alpha_k\}$ 为 PCA 表情系数。该模型采用 FLAME 的 50 个基表情，同时利用 FLAME 的关节和姿态参数 $\theta$ 通过线性混合蒙皮（LBS）变换高斯位置：$B^* = \mathrm{LBS}(\hat{B}, \theta)$。此外，该方法为嘴部内部定义独立的固定高斯 $B_m$，仅随下颌关节运动而不受表情影响，解决了 FLAME 网格无法覆盖嘴内区域的问题。

\textbf{设计动机}：将混合形状应用于高斯域而非网格域，其核心优势在于：高斯基元的半透明核函数天然适合建模头发、透明边界等离散表面难以处理的区域；同时基于 tile 的可微分光栅化器实现了实时渲染。

\textbf{公式与参数}：训练时使用泊松盘采样在 FLAME 中性网格上初始化约 50k 个高斯，表情偏移通过光度损失和感知损失联合优化。该方法在保持 FLAME 参数化控制能力的同时，实现了比当时最优方法更高的渲染质量和帧率。

\paragraph{NPGA}

NPGA（2024.05）~\cite{2405_19331}（Neural Parametric Gaussian Avatars）在架构上提出了根本性改进。

\textbf{Baseline 与问题}：Gaussian Blendshapes 使用基于 PCA（Principal Component Analysis，主成分分析）的线性表情空间，其表达能力受限于基表情的数量和线性假设，难以捕捉个体特异性的细微表情变化（如特定皱纹模式、血流量变化引起的肤色变化）。

\textbf{改了什么}：NPGA 将高斯基元的动态行为条件化于神经参数化头部模型（NPHM）的丰富表情空间，而非传统的基于网格的 3DMM；同时引入逐高斯潜在特征来增强化身的表达能力。

\textbf{实现原理}：NPGA 由两个核心组件构成。(1)~规范高斯点云 $A_c$，除了标准的高斯属性（位置、缩放、旋转、不透明度、球谐系数）外，每个高斯还附加一个可学习的潜在特征 $f \in \mathbb{R}^8$。该特征与空间坐标无关，提供语义信息描述各基元的动态行为。(2)~动力学模块 $D$，分解为粗粒度先验网络 $D^p$ 和细节网络 $D^a$。

$D^p$ 以空间坐标 $x$、逐高斯特征 $f$ 和从 MonoNPHM 蒸馏得到的表情编码 $\tilde{e}$ 为输入，预测高斯中心的偏移量 $\delta x^p$，作为 NPHM 后向变形的逆操作。$D^a$ 则预测所有高斯属性的残差偏移 $\delta a^p$，捕获表情依赖的皱纹和外观变化（如环境光遮蔽、血流量变化）。两个网络均使用 MLP（Multi-Layer Perceptron，多层感知机）架构。

\textbf{设计动机}：将 NPHM 的表情先验蒸馏到高斯动态模块中，既继承了 NPHM 丰富的表情空间（远超 PCA 的线性基），又利用了 3DGS 的高效渲染能力。逐高斯特征充当类似位置编码的角色，具有无限空间分辨率且不需要额外数据结构。

\textbf{公式与参数}：关键蒸馏策略是将 NPGA 的前向变形场 $D^p$ 优化为 MonoNPHM 后向变形 $D^{-1}_{\text{NPHM}}$ 的逆操作，通过拉普拉斯正则项约束动态行为的平滑性。在 NeRSemble 数据集上的自重演任务中，NPGA 比之前最优方法高出 2.6 dB PSNR。

NPGA 通过神经参数化表情空间和逐高斯潜在特征实现了高保真的化身渲染。但其动力学模块 $D$ 仍以 3D 高斯作为唯一图元，而 3D 高斯天然具有椭球体积，在从渲染结果提取几何表面时会产生多视角不一致和表面噪声问题，几何精度受限。

\paragraph{MixedGaussianAvatar}

MixedGaussianAvatar（2024.12）~\cite{2412_04955} 针对这一盲点提出了混合 2D-3D 高斯表示，与 NPGA 纯 3D 高斯的策略形成互补。

\textbf{改了什么}：该方法放弃 NPGA 的逐高斯潜在特征和 MLP 动力学模块，转而采用两阶段混合表示。先用 2D 高斯（surfels）贴附于 FLAME 网格三角形面片重建几何表面，再通过基于瓦片损失的自适应筛选机制识别渲染质量不足的区域并附加 3D 高斯补偿外观。具体而言，渲染图像与真实图像被分割为 $w \times h$ 大小的瓦片，计算逐瓦片 MSE 损失 $L_{ij} = \|o_{ij} - x_{ij}\|_2^2$，选取损失最大的前 $k$ 个瓦片并利用高斯树结构回溯其关联的 2D 高斯父节点（同一三角形面片上的兄弟节点一并纳入），经 $n$ 次不同 FLAME 参数和相机参数下的迭代取并集 $U = \bigcup_{i=1}^{n} S'_i$ 后，为 $U$ 中的每个 2D 高斯生成一个继承其属性的 3D 高斯子节点。

\textbf{混合渲染}：两种高斯通过局部-全局坐标变换统一由 FLAME 参数驱动，其中位置变换公式为 $\mu_g^{2D} = R \mu_l^{2D} + T + p^{2D}_*(t)$（$T$ 为三角形三顶点平均位置，$p^{2D}_*(t)$ 为可学习扰动项以自适应修正位置误差）。

混合渲染时，遍历排序后的 2D 高斯序列进行 $\alpha$-混合，当遇到具有 3D 子节点的 2D 高斯 $g_k^{2D}$ 时，将其子节点 $g_j^{3D}$ 的颜色贡献 $c^{3D}(x) = c_j \alpha_j G^{3D}_j(x) \prod_{t=1}^{k}(1 - \alpha_t G^{2D}_t(x))$ 紧跟插入，最终像素颜色为 $c^{mixed}(x) = c_p^{2D}(x) + c^{3D}(x) + c_b^{2D}(x)$。以 2D-3D 距离正则化 $L_{dis} = \|(\mu_l^{3D} + p_l^{3D}, \delta_{dis})\|_2$（阈值 $\delta_{dis}=1$）约束子节点在全局空间中与父节点的邻近性，确保两者投影落入同一瓦片以最大化补偿效果。

\textbf{设计动机}：2D 高斯具有 surfel 的表面贴附特性，能保证几何精度，但其颜色表达能力弱于 3D 高斯。混合表示在保持几何准确性的同时用 3D 高斯弥补渲染质量。训练分两阶段：2D 高斯阶段使用 $L_{rec} = L_c + \lambda_1 L_{pos} + \lambda_2 L_{sca} + \lambda_3 L_d + \lambda_4 L_n$（$\lambda_3=1000$, $\lambda_4=0.05$）约束局部位置和缩放；混合阶段冻结 2D 高斯的几何属性，仅优化 3D 高斯全参数和 2D 高斯的球谐系数。

与 MixedGaussianAvatar 将高斯绑定到 FLAME 三角形面片的策略不同，以下方法从顶点绑定的角度探索了另一条路线。

\paragraph{Topo4D}

Topo4D（2024.06）~\cite{2406_00440} 将高斯中心直接绑定到 FLAME 网格顶点而非三角形面片中心，利用预定义拓扑信息（8280 个顶点、16494 个面片）实现跨帧稳定的顶点级追踪。

\textbf{改了什么}：不同于 MixedGaussianAvatar 的两阶段 2D+3D 混合策略，Topo4D 采用交替几何-纹理优化方案：在每帧先执行 1000 次几何优化迭代，再做 300 次纹理优化迭代。几何优化阶段引入物理损失和拓扑损失的双重约束。物理损失 $L_{phy} = \lambda_{rot} L_{rot} + \lambda_{iso} L_{iso}$（$\lambda_{rot}=\lambda_{iso}=20$）强制相邻帧间旋转一致性（$L_{rot}$，基于四元数的逐边旋转差异）和长程等距约束（$L_{iso}$，惩罚边长变化 $|d_{0,j} - d_{0,i}|_2 - |d_{t,j} - d_{t,i}|_2$）；拓扑损失 $L_{topo} = \lambda_{pos} L_{pos} + \lambda_{flat} L_{flat}$（$\lambda_{pos}=10^{-3}$, $\lambda_{flat}=2 \times 10^{-4}$）通过顶点-邻居平均位置偏差（$L_{pos}$）和相邻面片夹角余弦变化（$L_{flat}$）维持网格规则性和表面光滑度。纹理优化阶段采用 UV 空间致密化技术，在每帧的稀疏高斯网格上通过双线性插值插入 $N \times N$（$N=30$）的密集子网格，学习高频纹理细节，最终通过高斯法向扩展（沿顶点法线方向偏移投影缩放值）提取更贴近真实表面的网格。

\textbf{设计动机}：Topo4D 的目标不是化身驱动渲染，而是从多视角时序图像直接优化密集对齐的 4D 头部网格和 8K 纹理图，因此需要保持严格的拓扑一致性。在自采集的 Light Stage 数据集（16 视角、$4096 \times 3000$ 分辨率、60 fps）上，该方法在面部重建精度上比次优方法多出 52.8\% 的顶点落在 0.5 mm 精度范围内，生成的 8K 纹理可忠实还原毛孔级细节。

\paragraph{GraphAvatar}

GraphAvatar（2024.12）~\cite{2412_13983} 从存储效率的角度重新审视了高斯化身的设计。NPGA 通过逐高斯潜在特征和 MLP 动力学模块实现了高保真渲染，但其存储开销仍与高斯数量直接相关（数万高斯的所有属性参数）。GraphAvatar 将"存储高斯参数"替换为"用 GNN（Graph Neural Network，图神经网络）生成高斯参数"，将存储开销从高斯参数本身降低到 GNN 模型权重（仅约 10 MB），同时在渲染质量上保持了竞争力。

\textbf{锚点回顾}：与 NPGA 用 MLP 从逐高斯潜在特征 $f \in \mathbb{R}^8$ 预测属性偏移的策略不同，GraphAvatar 转而采用双路 Graph U-Net 架构，直接从 FLAME 跟踪网格 $M = (V, A)$ 的顶点位置和法线中生成完整的高斯属性，利用图卷积捕获面部网格固有的拓扑结构信息。

\textbf{改了什么}：GraphAvatar 由三个核心改进组成。

(1)~双路 Graph U-Net 生成器：几何 U-Net $U_{\text{geo}}$ 和外观 U-Net $U_{\text{app}}$ 均基于 Chebyshev 图卷积构建，共享相同的编码器-解码器对称结构（两层图卷积编码→网格下采样至瓶颈特征 $z \in \mathbb{R}^8$→与表情系数 $e$ 拼接→两层图卷积解码上采样）。$U_{\text{geo}}$ 输出高斯中心偏移 $\mu'$、旋转 $R$ 和缩放 $S$，$U_{\text{app}}$ 输出颜色 $c$ 和不透明度 $\alpha$。由于 FLAME 模板固定为 $n = 5023$ 个顶点，输出高斯数量恒定，消除了 NPGA 中高斯数量的动态波动问题。(2)~神经高斯扩展：受 Scaffold-GS 启发，GraphAvatar 将 Graph U-Net 生成的 5023 个高斯视为锚点 $G_{\text{anc}}$，为每个锚点生成 $k$ 个神经高斯 $\{G_0, \ldots, G_{k-1}\}$，其属性通过独立 MLP 从锚点特征 $f_{\text{anc}}$、视角方向 $v_c$ 和表情系数 $e$ 联合解码。这一机制覆盖了 FLAME 无法建模的头发和配饰区域，弥补了跟踪网格的拓扑缺失。

(3)~3D 感知增强器与图引导优化：为缓解 GNN 的过度平滑效应，3D-aware Enhancer 将粗渲染的深度图 $D$ 通过 $1 \times 1$ 卷积预测逐块缩放 $\gamma_{i,j}^k$ 和偏移 $\beta_{i,j}^k$，以 $\hat{F}_{i,j}^k = \gamma_{i,j}^k(D) \cdot F_{i,j}^k + \beta_{i,j}(D)$ 的形式调制 UNet 每个中间块的激活，引入几何先验引导细节恢复。同时，图引导优化模块（GGO）通过交叉注意力机制融合 Graph U-Net 的瓶颈特征 $f_g$ 与时序特征 $f_t$，端到端地预测表情系数偏移 $\delta_e$ 和相机姿态偏移 $\delta_{\text{pose}} \in SO(3)$，降低了跟踪参数误差的累积。

\textbf{设计动机}：GraphAvatar 的核心洞察是"以生成替代存储"。与其为每个化身存储数万高斯的完整参数，不如训练一个轻量 GNN 按需生成。图卷积天然适配网格拓扑，Chebyshev 多项式近似使计算复杂度与顶点数呈线性关系，而固定 5023 个锚点高斯 + 可扩展神经高斯的分层策略在覆盖完整性和存储效率之间取得了平衡。

\textbf{公式与参数}：训练分两阶段进行：先用 vanilla 3DGS 对每个目标进行 10,000 次静态场景优化生成伪高斯，再以 MSE 损失预热 Graph U-Net；随后以 $L = \lambda_f L_f + \lambda_c L_c + \lambda_w L_w$（$\lambda_f = 1.0$, $\lambda_c = 0.1$, $\lambda_w = 0.1$）联合监督最终渲染、粗渲染和前景分割。在 INSTA 数据集上，GraphAvatar 以约 10 MB 模型尺寸达到了与 GaussianAvatars 等大模型相当的 PSNR，在 NBS 数据集上 PSNR 超越了 FlashAvatar 和 Gaussian Head Avatar。

\paragraph{MeGA}

MeGA（2024.04）~\cite{2404_19026}（Hybrid Mesh-Gaussian Head Avatar）则从根本上质疑了上述所有方法共享的一个假设：能否用同一种表示同时高质量地建模面部的皮肤和头部的头发？\textbf{锚点回顾}：Gaussian Blendshapes 及其后续工作（NPGA、GraphAvatar 等）均用统一的 3DGS 表示建模整个头部，但人头的不同组件具有截然不同的特性：面部皮肤近似光滑曲面且可用低维参数空间驱动，而头发则是包含大量细丝状结构的体积型对象。

\textbf{改了什么}：MeGA 放弃了"单一表示"策略，转而采用混合网格-高斯解耦架构：用增强的 FLAME 神经网格建模面部（皮肤、眼睛、嘴内），用静态规范空间中的 3DGS 独立建模头发，两者通过遮挡感知混合渲染（occlusion-aware blending）合成最终图像。

\textbf{设计动机}：这一解耦基于一个核心观察：纯 3DGS 方法（如 GaussianAvatars）虽然在头发等体积结构上表现出色，但在面部区域存在抗锯齿缺陷和穿透伪影（如张大嘴时高斯交错），而基于网格的延迟神经渲染恰好擅长产生高质量的面部纹理细节。与其用一种表示的弱点去妥协另一种表示的优势，不如让每种表示各司其职。

\textbf{具体差异}：在面部建模方面，MeGA 对 FLAME 网格进行四路细分并添加牙齿面片，生成约 16k 顶点的增强网格 $\mathcal{T}$，并预测以 FLAME 表情参数 $\psi$ 和姿态参数 $\phi$ 为条件的 UV 位移图 $\hat{G}_d$ 来补充个性化几何细节（如皱纹深度）。纹理方面，将神经纹理解耦为漫反射分量 $\hat{T}_{di}$、视角依赖分量 $\hat{T}_v$ 和表情依赖分量 $\hat{T}_{dy}$ 三部分，分别建模基础肤色、镜面反射和动态纹理（如皱纹、酒窝），通过延迟神经渲染生成面部颜色。在头发建模方面，从选定帧中构建规范空间中的静态 3DGS 头发，动画时施加全局刚体变换和基于 MLP 的非刚性变形场来捕捉动态发丝运动。

在渲染合成方面，MeGA 提出遮挡感知混合策略：先用增强 FLAME 网格渲染面部深度图，再在头发高斯渲染时使用"近端 z"深度（near-z depth，而非常用的积分深度）进行遮挡测试，配合提前停止策略跳过被遮挡的高斯和软混合（soft-blending）技术实现发际线处的平滑过渡。这种解耦表示使 MeGA 在 NeRSemble 数据集上超越了此前最优方法，同时自然支持发型更换、纹理编辑等下游应用——正是得益于头发在规范空间中的独立表示。

\subsubsection{高保真人脸化身}

本路线介绍通过高维扩展、局部条件化和自适应正则化等策略提升人脸化身保真度的方法。在参数化模型解决了"表情可控"的基础问题之后，如何进一步提升化身的视觉质量成为新的研究焦点。

当基础的"表情可控"问题被解决后，下一个瓶颈在于保真度：如何让化身在 3K 分辨率下仍无破绽？如何让一副眼镜的金属镜腿在不同角度呈现正确的镜面反射？这些问题的核心是 3D 高斯表示本身的表达能力上限。

\paragraph{ScaffoldAvatar}

ScaffoldAvatar（2025.07）~\cite{2507_10542} 引入局部表情条件化策略。

\textbf{Baseline 与问题}：此前的人脸化身方法（如 NPGA、Gaussian Blendshapes）对所有面部区域使用统一的表情驱动机制，忽略了不同面部区域（如嘴角、额头、鼻翼）的表情响应具有显著不同的空间模式。

\textbf{改了什么}：ScaffoldAvatar 将面部划分为多个块（patch），每个块独立建模表情驱动的动态外观和运动，并耦合到 Scaffold-GS 的锚点层级结构上。

\textbf{实现原理}：借鉴 Scaffold-GS 的锚点-层级组织方式，该方法在面部表面分布一组可学习锚点，每个锚点负责其邻域内的高斯基元。关键创新在于表情条件化的局部化：每个锚点关联的 MLP 分支独立接收表情编码作为输入，预测其管辖区域内高斯的属性偏移。这意味着嘴角附近的锚点可以学习到与嘴角运动高度相关的变形模式，而额头锚点则专注于额头皱纹的表达。

\textbf{设计动机}：面部表情本质上是局部化的——不同面部区域对同一表情编码的响应截然不同。全局表情条件化迫使单个 MLP 同时学习所有区域的不同动态模式，增加了学习难度并降低了精度。局部条件化则让每个模块专注于局部模式，显著提高了表达能力。

\textbf{公式与参数}：这种局部化的动态建模策略能更精确地捕捉皮肤褶皱等微观表情细节，同时实现了 3K 分辨率训练图像的高质量重建。

\paragraph{HyperGaussians}

HyperGaussians（2025.07）~\cite{2507_02803} 从根本上重新思考了 3D 高斯的表示能力。

\textbf{Baseline 与问题}：标准 3D 高斯的位置 $\mu \in \mathbb{R}^3$ 和协方差 $\Sigma \in \mathbb{R}^{3\times 3}$ 定义在三维空间中，其表示能力受限于三自由度的均值和六自由度的协方差，难以捕获眼镜框架、镜面反射等高频外观变化。

\textbf{改了什么}：HyperGaussians 将标准 3D 高斯扩展为高维 $(m+n)$ 多变量高斯，其中 $m=3$ 为原始空间维度，$n$ 为额外引入的潜在维度，通过可学习的局部嵌入增加表达能力。

\textbf{实现原理}：形式化地，随机向量 $\xi = (\xi_a, \xi_b)^T$ 被分为属性分量 $\xi_a \in \mathbb{R}^m$ 和潜在分量 $\xi_b \in \mathbb{R}^n$，其联合分布遵循 $(m+n)$ 维高斯。在渲染时，通过边缘化潜在变量 $\xi_b$（积分消去），得到仅关于 $\xi_a$ 的边缘分布，这恰好恢复了标准 3D 高斯的渲染管线。关键创新是"逆协方差技巧"（inverse covariance trick）：利用分块矩阵求逆公式 $\Sigma^{-1}_{aa|b} = \Sigma^{-1}_{aa}$，避免了直接计算高维 $(m+n)\times(m+n)$ 协方差矩阵的逆，将计算复杂度从 $O((m+n)^3)$ 降低到与标准 3DGS 相当的水平。

\textbf{设计动机}：从贝叶斯视角理解，高维联合分布的边缘化相当于在所有可能的局部嵌入上积分，等价于无限多个标准 3D 高斯的混合。这种概率解释赋予了单个 HyperGaussian 远超标准高斯的表达能力，而逆协方差技巧确保了计算效率不受影响。

\textbf{公式与参数}：在 19 个被试和 4 个人脸数据集上的实验表明，HyperGaussians 在眼镜框架、牙齿、复杂面部运动和镜面反射等高频细节上显著优于标准 3DGS。

\paragraph{STAvatar}

STAvatar（2025.11）~\cite{2511_19854} 与上述 ScaffoldAvatar 的局部锚点策略形成互补：ScaffoldAvatar 在三维空间中用锚点层级结构组织高斯并在每个锚点邻域内局部驱动表情，而 STAvatar 将高斯绑定到二维 UV 特征空间，在 UV 域中预测偏移。

\textbf{锚点回顾}：ScaffoldAvatar 的局部条件化解决了"不同面部区域对同一表情编码响应不同"的问题，但其锚点和高斯的绑定关系在训练中保持固定，且致密化策略对所有帧采用统一的时间无关准则。

\textbf{改了什么}：STAvatar 在空间和时间两个维度分别提出了改进——空间维度上，用 UV 自适应软绑定框架替代锚点绑定，高斯通过 UV 坐标从特征偏移图 $\delta_{\text{map}} \in \mathbb{R}^{256 \times 256 \times 13}$ 中采样 13 维偏移（位置 3 维 + 缩放 3 维 + 四元数旋转 4 维 + 颜色 3 维），且 UV 坐标在每次致密化后动态更新以兼容高斯数量的变化；时间维度上，用 FLAME 条件化时间聚类（FTC）替代均匀帧采样策略，将视频帧按 FLAME 参数（表情、姿态、平移权重分别为 0.3、0.6、0.1）聚为 $K \in [5, 12]$ 组（通过最大化轮廓分数自适应选择），在组内计算融合感知误差（FPE-AP）指导致密化。软绑定框架采用双分支 UV 空间网络：全局分支 $\phi_g$ 编码参考图像的整体纹理特征并输出全局偏移，局部分支 $\phi_l$ 按 UV 空间掩码 $M_i \in \{0,1\}^{256 \times 256}$ 将特征分为眼、嘴、鼻、额四个区域，由共享解码器和区域特有头 $H_i$ 分别处理（$\delta_l = \sum_{i=1}^4 H_i(M_i \odot \phi_l(T, UV_{\text{disp}}', c))$），最终通过融合函数 $F(\delta_g, \delta_l)$ 生成偏移图。这种 UV 域操作的优势在于：特征偏移图天然具有空间局部性和连续性，无需显式维护锚点邻域关系。

时间密度控制方面，FPE-AP 准则融合 $L_1$ 像素误差和 D-SSIM 损失（权重 $\lambda_1=0.2$），对每个高斯计算平均融合感知误差 $\epsilon_i = \frac{1}{C_i} \sum_{p \in P_i} E(p)$，并追踪峰值误差 $E_i^{\text{peak}} = \max_t \epsilon_i^{(t)}$，选取前 3\% 高斯构成 $S_{\text{peak}}$ 集合，当 $\epsilon_i > \delta_{\text{avg}} \cdot \mathbf{1}_{i \in S_{\text{peak}}}$ 时触发克隆（$\delta_{\text{avg}}=10^{-3}$），分割仍使用位置梯度。

\textbf{设计动机}：ScaffoldAvatar 的锚点绑定在致密化过程中需要重新计算锚点邻域归属，而 UV 坐标采样天然支持高斯数量变化——只需为新增高斯分配 UV 坐标即可。时间聚类则解决了瞬态可见区域（如张大嘴时突然暴露的口腔内部）在均匀采样下被稀释、得不到充分致密化的问题。

\textbf{具体差异}：STAvatar 在 NeRSemble 数据集上的实验表明，FTC 聚类策略使瞬态可见区域的 PSNR 提升约 0.5 dB，而 UV 软绑定相比硬绑定（LBS 双线性映射）在细纹和发丝细节上有明显改善。

\paragraph{GeoAvatar}

GeoAvatar（2025.07）~\cite{2507_18155} 从正则化粒度的角度补充了 ScaffoldAvatar 的设计。

\textbf{锚点回顾}：ScaffoldAvatar 的所有锚点采用统一的 MLP 架构和优化策略，对不同几何复杂度的面部区域施加同等强度的约束——但面部并非各区域同等刚性：鼻梁应紧贴网格，而头发和嘴角需要更大的偏移自由度。

\textbf{改了什么}：GeoAvatar 引入自适应预分配阶段（Adaptive Pre-allocation Stage, APS），在初始训练 $N$ 步后根据各面部区域高斯到网格的平均距离 $\mu^{(k)} = \frac{1}{|G_k|} \sum_{i,j \in G_k} \| \mu_{i,j} \|$ 将 FLAME 三角面片无监督地分割为刚性集 $F_r$（$\mu^{(k)} < \bar\mu$）和柔性集 $F_f$（$\mu^{(k)} > \bar\mu$），对两组施加差异化的位置正则 $\delta_r=0.1$（刚性，严格约束高斯贴近网格）和 $\delta_f=2.0$（柔性，允许较大偏移以建模头发等非刚性区域）。此外，GeoAvatar 还引入角度正则 $L_{\text{angle}}(\theta) = \mathbf{1}_{r > \sigma_r}(\theta - \theta_0)$（$\theta_0 = 45°$），防止柔性集中距离网格较远的高斯偏移到错误的相邻面片区域。

在嘴部建模方面，GeoAvatar 扩展了 FLAME 的嘴部几何结构——在 GaussianAvatars 唇环网格的基础上增加了磨牙、上颚和口腔底部的面片（构成嘴部集 $F_m$），并设计了上下分离的嘴部变形网络 $\phi_{\text{upper}}$ 和 $\phi_{\text{lower}}$，分别处理上齿+上颚和下齿+口腔底部的联动运动，输入为 FLAME 表情参数 $\psi$、下颌姿态 $\phi_j$ 和经位置编码的时间步 $T$（$\Delta v = \phi(\psi, \phi_j, \gamma(T))$），推理时置 $T=0$ 即可。

\textbf{设计动机}：统一的正则化强度是次优的——对刚性区域施加过弱的正则会破坏动画精度，对柔性区域施加过强的正则则限制了几何表达能力。自适应分割让模型自动发现哪些区域需要松绑、哪些需要紧箍。

\textbf{具体差异}：刚性集正则 $\delta_r=0.1$ 与柔性集正则 $\delta_f=2.0$ 相差 20 倍，这一数量级的差异使得面部刚性区域（鼻、额）的高斯始终贴近 FLAME 网格，而头发和耳部区域的高斯可以自由偏离网格以捕捉细丝状结构，在 NeRSemble 和 iPhone 数据集上均优于均匀正则化策略。

\paragraph{ArchitectHead}

ArchitectHead（2025.10）~\cite{2510_05488} 首次实现了头部化身中连续的细节层次（Level of Detail, LOD）控制。

\textbf{Baseline 与问题}：HyperGaussians 通过高维扩展大幅提升了单个高斯基元的表示能力，但所有现有头部化身方法（包括 HyperGaussians 本身）在训练后高斯数量即固定——无法根据渲染场景的算力约束动态调节精度。已有的 LOD 方案（如 LoDAvatar）仅支持 3-4 个离散级别，且依赖合成多视角数据训练，无法用于单目视频个性化化身。

\textbf{改了什么}：与 HyperGaussians 在表示维度上做加法不同，ArchitectHead 转向效率维度——将高斯参数化在 2D UV 特征空间中，使得调节高斯数量如同缩放一张 2D 图像般简单。核心创新是可学习的多级 UV 特征场（multi-level UV feature field）：由 $N=3$ 个不同分辨率（$256\times256$、$128\times128$、$64\times64$）的特征图构成金字塔，每个像素存储 $D_f=64$ 维的潜在特征。

渲染时，连续 LOD 参数 $l \in [0, 1]$ 先通过 $S = S_{\max} - l(S_{\max} - S_{\min})$ 计算目标分辨率，再以 softmax 加权混合所有级别的特征图：$w_i = \frac{\exp(-|r_i - r|/\tau)}{\sum_j \exp(-|r_j - r|/\tau)}$，其中 $r_i = \ln S_i$，温度 $\tau=0.35$ 控制插值软度。该公式赋予分辨率最接近目标 $S$ 的级别最高权重，同时保证所有级别权重不为零以维持梯度传播。

\textbf{设计动机}：直接缩放单张 UV 特征图会面临两难——高分辨率图下采样时相邻高斯的旋转和尺度差异被平滑掉，损失关键局部细节；低分辨率图则先天不足以表征精细外观。多级金字塔让每个分辨率级别独立学习其尺度下的最优特征表示，避免了跨尺度表征冲突。消融实验证实，单张 $256\times256$ 特征图必须同时兼顾所有 LOD 级别，导致低 LOD 表示反向干扰高 LOD 质量，出现更多视觉伪影和模糊边缘。

\textbf{实现原理}：解码器接收三路信息的拼接——位置编码后的 UV 坐标（12 个频率，$D_{PE}=75$ 维）、采样所得的 UV 潜在特征、以及 FLAME 表情参数经 MLP 映射网络压缩为 $K=20$ 维的驱动码与 LOD 值 $l$。五个独立 MLP 分别解码高斯的五类属性（位置偏移 $\Delta\mu$、尺度 $s$、旋转 $q$、不透明度 $\alpha$、颜色 $c$）。

训练分为两阶段：第一阶段固定 $l=0$（最高 LOD）联合优化全部参数 15K 步以稳定学习；第二阶段随机采样 $l \in [0, 1]$ 训练 30K 步，使解码器适应任意细节层次。

\textbf{公式与参数}：在最低 LOD（$l=1.0$）时，模型仅使用 6.2\% 的高斯基元（约 3.2K 个），渲染速度提升 96\%，而 L1 损失仅增加 7.9\%，PSNR 下降 0.97\%，SSIM 下降 0.6\%。训练仅需单张 RTX 4090 约 3 小时（阶段一 30 分钟 + 阶段二 2.5 小时），训练后可即时切换任意 LOD 而无需重训练。

\subsection{单图像与生成模型人脸重建}
\label{ss:16_1b}

本节讨论~2~条技术路线：（1）单图像和少样本人脸重建；（2）基于生成模型的人脸化身。

\subsubsection{单图像和少样本人脸重建}

本路线介绍从前馈生成、扩散优化到先验学习等策略实现单图像或少样本人脸重建的方法。这些方法与前文基于多视角视频的化身方法形成互补，面向更具挑战性的少输入场景。

上述方法通常需要多视角视频或数十张输入图像，但实际应用中最常见的场景是用户仅有一张正面照片。从极少输入（单张图像或少量视频帧）创建高质量、可动画的人脸化身是该领域的重要目标，核心难点在于模型需要"脑补"所有不可见区域。

\paragraph{FaceLift}

FaceLift（2024.12）~\cite{2412_17812} 提出了前馈式的 360 度头部重建方法。

\textbf{Baseline 与问题}：单图像 3D 重建的核心困难在于输入图像仅提供正面信息，侧面和背面完全不可见，需要模型具备强先验知识来"脑补"不可见区域。

\textbf{改了什么}：FaceLift 采用两阶段策略：首先用多视角潜在扩散模型从单张人脸图像生成一致的侧面和背面视角，然后通过基于 Transformer 的重建器从多视角图像生成 3DGS 表示。

\textbf{实现原理}：第一阶段训练一个条件扩散模型 $f_D$，输入单张正面图像 $y$ 和对应 CLIP 文本嵌入 $e^{1:N}$，联合生成 $N=6$ 个覆盖 360 度的视角 $x_0^1, x_0^2, \ldots, x_0^N$（正面、左/右 45 度、左/右 90 度、背面）。扩散模型采用 DDPM 框架，通过多视角注意力机制确保生成视角之间的多视角一致性。第二阶段使用基于 Transformer 的重建器 $f_R$，将生成的多视角图像及其 Plücker 射线坐标作为输入，预测一组像素对齐的 3D 高斯。选择 3DGS 而非 NeRF 或网格的原因在于：高斯的半透明核函数天然适合建模头发等复杂结构，显式体素基元比密度场或离散表面更具表现力。

\textbf{设计动机}：两阶段解耦的设计将"视觉想象"（多视角生成）和"3D 重建"（高斯预测）分工到各自擅长的模型中，避免了单阶段方法中 3D 先验和 2D 纹理先验的冲突。

\textbf{公式与参数}：训练数据通过 Blender 渲染管线生成，包含 200 个身份、每个 50 种外观（不同发型、肤色、表情、服装），在 $512^2$ 分辨率下渲染。该方法在合成数据上训练，但展现出对真实世界图像的强泛化能力。

\paragraph{FastAvatar}

FastAvatar（2025.08）~\cite{2508_18389} 实现了约 3 秒内的单图像 3D 人脸重建，与 FaceLift 的两阶段扩散+重建策略形成对比。FaceLift 依赖扩散模型先"想象"不可见视角再重建，虽能处理任意视角但生成过程较慢且可能引入扩散模型的伪影。FastAvatar 的核心思路转变是：与其让模型从零"脑补"缺失视角，不如直接利用数据驱动的平均模板作为强先验，只需预测个体与模板之间的偏差即可。

\textbf{与锚点的核心差异}：FaceLift 通过扩散模型在 2D 图像空间生成缺失视角、再由 Transformer 重建 3DGS；FastAvatar 完全绕过视角生成步骤，转而在 3D 参数空间建立显式模板先验，通过残差预测直接从单张图像回归完整 3DGS 模型。

\textbf{模板构建}：对 $N=410$ 个训练身份各自优化一个 3DGS 模型（每个高斯放置在 FLAME 网格三角面片中心，加上 168 个牙齿锚点，共 $K=10{,}144$ 个语义一致的高斯），然后对对应高斯的参数取平均得到模板 $T$。这种简单平均策略相比联合优化模板或随机初始化能产生更稳定的先验。

\textbf{残差预测}：前馈编码器-解码器不直接回归绝对参数值，而是预测相对于模板的残差。解码器 $h_\theta$ 接收姿态不变的恒等嵌入 $w \in \mathbb{R}^{|w|}$ 和高斯位置编码 $e_k \in \mathbb{R}^{|e|}$，输出每个高斯的残差 $\Delta \mu_k, \Delta s_k, \Delta q_k, \Delta \alpha_k, \Delta c_k$。最终参数为：位置 $\mu_k = \mu_k^{\text{FLAME}} + \Delta \mu_k$（相对于 FLAME 标准位置偏移，而非模板位置），缩放 $s_k = s_k^T \cdot \exp(\Delta s_k)$，旋转 $q_k = \text{norm}(q_k^T + \Delta q_k)$，不透明度 $\alpha_k = \sigma(\alpha_k^T + \Delta \alpha_k)$，球谐系数 $c_k = c_k^T + \Delta c_k$。这种残差化设计使网络仅需预测小幅度偏差，极大地降低了回归难度。

\textbf{两阶段训练策略}：直接端到端训练编码器-解码器会导致网络坍缩到预测平均人脸的退化解。为此，FastAvatar 采用分阶段训练：第一阶段冻结编码器，将所有训练身份的恒等编码 $w_i$ 作为可学习变量，与解码器联合优化（使用 $L_{\text{dec}} = \lambda_1 L_{\text{LPIPS}} + \lambda_2 L_1 + \lambda_3 L_{\text{SSIM}} + \lambda_4 L_2(\Delta \mu) + \lambda_5 L_2(\Delta s)$ 损失），迫使解码器学习平滑且可泛化的潜在空间；第二阶段冻结已训练好的解码器，训练编码器 $g_\phi = g_{\text{MLP}} \circ g_{\text{FR}}$（基于预训练人脸识别骨干网络）将输入图像映射到该潜在空间，用 $L_{\text{enc}} = L_2(w, w^*) + \lambda_{\text{cos}} L_{\text{cos}}(w, w^*)$ 监督。

\textbf{轻量外观精修}：推理时，单次前馈编码器-解码器传递即可产生稳定的 3D 一致几何（base 模型达到 21.17 dB PSNR），随后仅优化潜在编码 $w$ 和解码器的外观输出约 300 次迭代（约 3 秒，A100），所有调整均通过编码器-解码器先验进行而非直接修改高斯参数，确保不引入视角相关的伪影。精修后全模型在 Nersemble 数据集上达到 24.01 dB PSNR、0.91 SSIM、0.81 身份相似度，在极端输入姿态（超过 90 度侧脸）下仍保持强身份保持性，且质量-速度权衡大幅优于扩散方法（Arc2Avatar、DiffusionRig）和优化方法（GaussianAvatars、FlashAvatar）。

与 FaceLift 和 FastAvatar 的前馈生成策略不同，

\paragraph{SEGA}

SEGA（2025.04）~\cite{2504_14373} 将重建问题拆解为``静态身份''和``动态表情''两个独立子问题。FaceLift 用扩散模型生成多视角再重建，SEGA 则完全绕过多视角生成：静态分支用 DINOv2（自监督视觉Transformer，通过自蒸馏学习强大视觉特征）编码器提取 2D 身份特征，经 UV-Alignment Transformer（10 层）映射到 FLAME UV 空间，再由静态解码器直接预测全头部高斯属性（$1024 \times 1024$ UV 分辨率），其中位置偏移量 $M_{\text{offset}}$ 捕捉身份特有的几何细节（发型、面部轮廓等）；动态分支用 VQ-VAE（Vector Quantized Variational AutoEncoder，向量量化变分自编码器）编码器提取表情特征，通过位移 VAE 预测表情驱动的形变图（$400 \times 400$ UV 分辨率），仅覆盖嘴唇、眉毛等表情敏感区域。推理时静态分支仅计算一次并缓存，动态分支逐帧更新，通过预定义二值掩码 $M_d$ 将两路输出在 UV 空间中融合，并用线性插值权重掩码确保过渡区域无缝拼接。这种分层解耦的设计动机在于：有限 3D 训练数据的身份多样性不足，将静态身份（可预计算复用）与动态表情（轻量实时）分开建模可最大化数据利用效率——静态分支在 FFHQ 等 2D 数据集上联合训练以弥补 3D 数据的身份覆盖不足，动态分支则仅需学习表情相关的局部变化。

\paragraph{Arc2Avatar}

Arc2Avatar（2025.01）~\cite{2501_05379} 则走了一条完全不同的路线：不训练任何前馈重建网络，而是利用预训练的人脸基础模型 Arc2Face（在 WebFace260M 上训练）通过基于分数蒸馏的迭代优化直接塑造 3D 高斯。与 SEGA 的显式编码器-解码器架构不同，Arc2Avatar 的核心是 ISM（Interval Score Matching，区间得分匹配）优化——相较于标准 SDS（Score Distillation Sampling，分数蒸馏采样）的单步去噪伪真值不一致问题，ISM 通过 DDIM（Denoising Diffusion Implicit Models，去噪扩散隐式模型）反演构建可逆噪声轨迹，使伪真值始终与原始渲染视角对齐，再用多步去噪提升质量，收敛速度更快。

为实现 360 度生成，作者用 PanoHead 合成多视角数据（10K 身份，每人 26 个视角），通过 LoRA（Low-Rank Adaptation，低秩适配，rank $r=16$）微调 Arc2Face 以扩展其视角覆盖范围，推理时 LoRA scale $s_L = 0.45$ 平衡新颖视角生成与身份保持。条件化方面，ArcFace 身份嵌入 $v \in \mathbb{R}^{512}$ 通过占位符注入文本提示，并与方向性文本嵌入（``front/side/back view''）按权重 $b = 0.85$ 混合，兼顾身份一致性与视角控制。在 3DGS 优化层面，Arc2Avatar 采用掩码策略：面部区域的高斯与 FLAME 顶点建立一一对应并禁用自适应密度控制，非面部区域则正常执行分裂和剪枝以建模头发等复杂结构；同时引入拉普拉斯连接性正则化约束相邻高斯的拓扑关系，防止优化过程中面部结构崩溃。表情驱动时，仅需 500 次迭代即可从变形后高斯出发精修出包含牙齿和舌头的极端表情。

在少样本设定方面，上述方法（FaceLift、FastAvatar、SEGA、Arc2Avatar）均依赖真实 3D 数据训练前馈网络或迭代优化，而以下三个方法探索了不同的先验学习与个性化策略。

\paragraph{HeadGAP}

HeadGAP（2024.08）~\cite{2408_06019} 与 SEGA 的静态-动态双分支解耦思路相似，但采用了不同的先验学习范式。SEGA 用 DINOv2 编码器从 2D 图像直接提取身份特征再经 Transformer 映射到 UV 空间；HeadGAP 则构建基于自动解码器（auto-decoder）的部件级特征点云表示（part-based feature point cloud），从 FLAME 网格 UV 空间初始化 $n$ 个高斯基元，每个语义部件（如面部、头皮、耳朵、口腔）拥有独立的身份码 $z_j \in \mathbb{R}^{c_2}$ 和部件专属 MLP $f^M_p$，而非 SEGA 的单一全局编码器-解码器架构。这一架构选择的设计动机在于：auto-decoder 无需显式编码器即可在多身份联合优化中自然形成身份潜空间，使先验学习不依赖大规模配对数据，且通过共享点特征编码 $f$ 跨身份传递通用几何先验，同时以部件级身份码 $z_j$ 保留各语义区域的独立表达能力——面部和头皮的身份特征解耦度更高，更容易捕获发型等大变异部件的身份信息。

在动态建模上，SEGA 通过 VQ-VAE 编码表情特征后由位移 VAE 预测形变图；HeadGAP 则直接利用 FLAME 形变信号 $e = T(\hat{\theta}) - T(\hat{\theta}_{\text{neutral}})$（当前姿态全局位置减去中性姿态全局位置的差值）驱动每个部件 MLP 预测动态局部属性，包括位置偏移、旋转、缩放和不透明度。与 GaussianAvatars 使用静态局部属性不同，HeadGAP 的动态局部属性建模使每个高斯基元的属性能随表情显式变化，从而更好地捕捉动态细节。此外，HeadGAP 引入屏幕空间 CNN 精修网络 $f^C$，不执行超分辨率而是同分辨率精修：先渲染 RGB 图像 $I_{\text{rgb}}$ 和潜变量特征图 $I_h$，再由 $f^C$ 输出最终图像 $I = f^C([I_{\text{rgb}}, I_h])$。这一设计旨在利用大规模数据学习 3DGS 表示难以直接建模的结构化外观先验（如细密纹理和微妙的光照变化），在少样本设定下显著提升渲染真实感。

先验学习阶段，HeadGAP 在大规模多视角动态数据集上联合优化 $k$ 个身份的 GAPNet 参数（身份码 $z_{1 \ldots k}$、点特征编码 $f$、部件 MLP $f^M$、CNN 精修网络 $f^C$），使用 $L = L_{\text{rec}} + \lambda_\alpha L_\alpha + \lambda_s L_s + \lambda_\delta L_\delta + \lambda_{\text{arap}} L_{\text{arap}}$ 等多目标正则化（其中不透明度正则化 $L_\alpha$ 约束高斯不透明度累积为 1 且限制在头部区域内，ARAP 正则化约束静态偏移的刚性）。个性化阶段仅通过两步完成：先用部分线性组合权重 $w \in \mathbb{R}^{k \times p \times 1}$ 从身份码本中反演最接近目标身份的编码 $z^* = (w)z \in \mathbb{R}^{k \times p \times c_2}$，再以小学习率微调所有网络参数（排除嘴部区域并施加视角正则化防止过拟合），仅需少量帧即可完成个性化。反演策略利用了 auto-decoder 先验中已学到的 $k$ 个身份码构成的身份码本——通过线性组合已有编码来逼近目标身份，避免了从零搜索潜空间的不稳定性，这也是 auto-decoder 相比 encoder-decoder 在少样本场景的核心优势。

\paragraph{SynShot}

SynShot（2025.01）~\cite{2501_06903} 与 SEGA 和 HeadGAP 均采用身份-表情解耦策略，但核心差异在于先验数据来源：SEGA 依赖真实 3D 扫描数据补充 2D 数据训练，HeadGAP 从真实多视角动态数据学习先验，而 SynShot 完全在合成数据上训练先验模型（约 2000 个独特身份，使用 Blender Cycles 渲染器生成 1400 万张 $768 \times 512$ 图像）。选择纯合成数据的设计动机在于：真实数据受限于采集设备可达到的身份多样性和表情覆盖度，而合成数据可以程序化地生成任意身份-表情组合，确保先验模型在身份空间上的均匀覆盖，从根本上避免了真实数据的分布偏差问题。

在表示上，SEGA 用分离的静态/动态分支在 UV 空间预测高斯属性，HeadGAP 用部件级 MLP 从特征点云回归属性；SynShot 采用双分支 VQ-VAE 编码器，分别将身份纹理图 $T$（RGB 纹理 + XYZ 顶点位置图）和表情偏移图 $E$（FLAME 表情基的 UV 光栅化结果）编码为独立的量化潜在空间 $q(z_{\text{id}})$ 和 $q(z_{\text{expr}})$，通过 Transformer 解码器在量化码本空间建模长程依赖，预测 UV 空间的高斯参数纹理图。与连续潜空间相比，VQ-VAE 的离散码本将身份和表情信息约束在有限词汇表中，迫使编码器学习更具代表性的特征模式，同时为下游 Transformer 解码器提供更结构化的输入表示。解码器输出三个分支：特征图解码器 $D_f$ 输出 $C$ 维特征向量、身份图解码器 $D_t$ 输出纹理图、表情解码器 $D_e$ 输出表情偏移图，高斯位置则通过 $\hat{P} = \hat{V} + \hat{E}$ 将预测的顶点位置图与表情偏移图相加获得。

在基元数量灵活性方面，与 SEGA 固定 UV 分辨率不同，SynShot 采用逐部件双线性采样机制：从解码器的特征图和参数图中按任意 UV 坐标 $(u, v)$ 采样高斯属性，面部分辨率不同于头皮区域，实现自适应基元致密化。这一设计的关键在于将高斯数量与 CNN 输出分辨率解耦——训练和推理时可使用不同数量的采样点，无需重新训练网络即可调整渲染质量。每个部件区域的初始高斯参数也有差异：最近邻距离初始化缩放 $s_r = \min_j \|p_{r_i} - p_{r_j}\|_2$，初始不透明度设为 0.7，旋转通过 UV 图像空间梯度 $\mathbf{T} = \partial p_r / \partial u$、$\mathbf{B} = \partial p_r / \partial v$、$\mathbf{N} = \mathbf{T} \times \mathbf{B}$ 构建切线框架计算。精修阶段使用轻量级四层卷积残差块（带跳跃连接）作为回归器 $R_r$ 和 $R_c$，分别预测球谐系数偏移和位置/缩放/旋转/不透明度偏移。个性化阶段，SynShot 采用两阶段反演：先优化身份编码器 $E_{\text{id}}$ 冻结其余网络以恢复身份潜在码 $z_{\text{id}}^*$，再冻结 $z_{\text{id}}^*$ 精调解码器和回归器，辅以基于 ArcFace 的身份损失 $L_{\text{id}}$ 和 $L_{\text{arc}}$ 约束身份保持，最多 20 视角的反演优化在单张 H100 上不到 10 分钟。

其中 $L_{\text{id}}$ 直接约束渲染图像的 ArcFace 嵌入与输入图像嵌入的余弦相似度以保持身份一致性，$L_{\text{arc}}$ 则在反演阶段进一步强化身份编码器的特征空间与 ArcFace 预训练空间的对齐，防止优化过程偏离身份敏感区域。反演训练还引入了早停策略配合指数移动平均损失判断收敛，以及根据输入帧数动态缩放优化步数的启发式规则（常数因子 10 确保每个样本至少被采样一次），使少样本场景下的反演更加稳定。

\paragraph{Gaussian Deja-vu}

Gaussian Deja-vu（2024.09）~\cite{2409_16147} 与 SEGA 在 UV 空间参数化高斯的思路一致，但采用了``先通用化后个性化''的两阶段流程而非 SEGA 的单阶段前馈预测。这一流程的设计动机在于：通用化阶段通过大规模数据学习到的单图重建能力提供了稳健的初始化，使后续个性化阶段仅需在此基础上做小幅修正，而非从零开始优化整个 UV 高斯图——这与 HeadGAP 利用身份码本反演作为初始化的思路异曲同工，都是``先获取合理起点再精修''的策略。在 UV 参数化上，SEGA 使用标准 FLAME UV 对应关系；Gaussian Deja-vu 修改了 FLAME UV 映射以引入口腔内部形状并移除颈部区域。这一修改的设计考量在于：标准 FLAME UV 缺少口腔内部的几何表示，导致牙齿和舌头无法在 UV 空间中被编码；而颈部区域在头部动画中运动模式与面部差异过大，将其排除可减少无关区域对 UV 图空间的浪费，使有限的 UV 分辨率集中用于面部和头皮细节。

定义 UV 高斯图 $U \in \mathbb{R}^{H \times W \times 13}$（包含位置、缩放、旋转、不透明度和颜色五个通道组），考虑到头部表面材质以简单漫反射为主，颜色通道仅使用零阶球谐函数（3 通道 RGB）而不建模复杂各向异性反射，初始缩放值设为 $-8.3533$（经指数激活后为正）、初始不透明度为 0.1。通用化阶段，卷积编码器-解码器网络从单张头部图像预测 UV 偏移量 $\Delta U$，用 $L = L_{\text{photo}} + \lambda_\delta L_\delta + \lambda_s L_s + \lambda_{\text{view}} L_{\text{view}}$ 训练，其中位置正则化 $L_\delta$ 使用预定义权重图 $M_\delta$ 对面部区域施加更强约束而允许头发区域更大位置自由度，视角一致性损失 $L_{\text{view}}$ 惩罚同一身份多视角估计的 UV 偏移差异。训练中还采用同批次同身份的多视角图像随机配对策略进一步增强视角一致性。个性化阶段采用两步优化：先用单图重建对视频中每帧估计 UV 偏移并取均值 $\bar{\Delta}$ 作为基础，再学习全局 UV 修正 $\Delta U_{\text{global}}$（可学习张量，初始化为零）；第二步学习表情驱动的混合权重图 $\Delta U_{\text{blend}} \in \mathbb{R}^{D \times H \times W \times 13}$，以第一阶段的 $\Delta U_{\text{global}}$ 作为初始化（而非从零开始），通过从 FLAME 表情系数导出的 softmax 混合权重 $b = \text{softmax}(W \beta_{\text{exp}})$ 加权混合，使最终 UV 高斯图为 $U_{\text{blend}} = U_{\text{init}} + \bar{\Delta} + \sum_{i=1}^D b_i \Delta U_{\text{blend}}^i$。

softmax 函数确保混合权重为正值且总和恰好为 1，避免权重爆炸或退化。这种表情条件混合图设计使个性化过程既能捕捉全局身份特征，又能建模表情相关的细节变化——全局修正处理与表情无关的身份特征（如发型、面部轮廓），混合图处理表情驱动的细节变化（如皱纹、嘴唇形变），两阶段分工将单目视频的训练时间缩短至现有方法的四分之一。

\subsubsection{基于生成模型的人脸化身}

本路线介绍基于~GAN~和扩散模型的~3DGS~人脸化身生成方法。

3DGS 与生成对抗网络（GAN，一类通过生成器和判别器的对抗训练来学习数据分布的模型）和扩散模型（Diffusion Model，一类通过逐步去噪来生成样本的生成模型）的结合，为人脸化身的多样性和可控性开辟了新路径。从本节开始，讨论重点从"重建"转向"生成"——不再要求与特定个体一致，而是追求在身份、表情和视角等维度上的多样性和可控性。

\paragraph{GGHead}

GGHead（2024.06）~\cite{2406_09377} 在 3D GAN（Generative Adversarial Network，生成对抗网络）框架中采用 3DGS 表示。

\textbf{Baseline 与问题}：3D GAN 训练的核心挑战在于如何用 2D CNN（Convolutional Neural Network，卷积神经网络）高效生成无结构的 3D 点云表示。3D 高斯作为点基表示本身是无序的，直接生成面临顺序模糊性和大量空白区域的困难。

\textbf{改了什么}：GGHead 将 3D 高斯绑定到模板头部网格的 UV 空间中，利用 2D CNN 生成器在结构化的 UV 空间中预测高斯属性，巧妙地将 3D 生成问题转化为 2D 纹理生成问题。

\textbf{实现原理}：使用 StyleGAN2 架构作为骨干网络 $B$，将潜在编码 $z \in \mathbb{R}^{512}$ 映射到多通道 UV 图 $M \in \mathbb{R}^{256 \times 256 \times 14}$，分别对应位置、缩放、旋转、颜色和不透明度五种高斯属性。对于位置属性，额外添加一个权重初始化为零的可学习 CNN 层 $Z$，确保训练初始时高斯精确落在模板网格上，这是稳定对抗训练的关键设计。3D 高斯通过在 UV 图上均匀采样获得——每个高斯有固定的 UV 坐标 $x_{uv} \in [0,1]^2$，通过双线性插值从属性图中查询其属性。生成器可以与采样方案无关：训练和推理时可使用不同数量的高斯。

\textbf{设计动机}：UV 参数化利用了头部网格的固定拓扑结构，为无序的点云引入了空间组织，使 2D CNN 的局部感受野能有效建模面部的局部几何和纹理变化，同时避免了 3D 卷积或 Transformer 的计算开销。

\textbf{公式与参数}：该方法首次实现了 $1024^2$ 分辨率的 3D 一致性头部实时生成和渲染。

\paragraph{CGS-GAN}

CGS-GAN（2025.05）~\cite{2505_17590} 解决了 GGHead 和 GSGAN 共有的视角条件化缺陷——二者在训练中以 50\% 概率将相机信息输入生成器映射网络，导致推理时正面条件渲染正常但侧视角质量严重退化。

\textbf{锚点回顾与核心差异}：与 GGHead 依赖 FLAME UV 参数化将高斯绑定到固定网格拓扑不同，CGS-GAN 完全脱离网格约束，采用无网格的层级点云生成策略：从 512 个可学习特征点出发，通过自适应 Transformer 块（包含 AdaIN 层、自注意力层和 MLP）逐级处理，每级通过点上采样将高斯数量扩展为上一级的四倍，同时利用前一层的锚点高斯提供位置和朝向参考，形成粗到细的层级生成策略。这一设计的核心优势在于骨干网络中点数固定，可灵活添加额外 Transformer 层而仅线性增长计算量（GSGAN 因在骨干内上采样导致高内存开销，只能改用低效的图卷积）。

\textbf{训练稳定性保障}：为实现完全移除视角条件化后的稳定训练，CGS-GAN 引入三项关键改进：（1）多视角正则化——每步从多视角渲染同一场景并平均梯度，利用 3DGS 超过 200 FPS 的高效渲染将额外开销降至可忽略；（2）移除 GSGAN 的对比损失，改用 StyleGAN 条件训练机制将相机参数嵌入判别器图像特征，避免了对比损失在高度相似视角下产生过大损失值导致的训练不稳定；（3）以随机颜色背景替换 GGHead 的白色背景，阻止生成器利用背景颜色制造白色空洞伪影。

\textbf{公式与参数}：通过添加 Transformer 层和加倍点上采样因子实现分辨率扩展，对应高斯数量为 $256^2$: 109k、$512^2$: 240k、$1024^2$: 502k、$2048^2$: 1M，支持高达 $2048^2$ 的输出分辨率。

\paragraph{DGH}

DGH（2025.06）~\cite{2506_20875} 从组合性角度出发，提出了具有可组合头发和面部组件的无条件生成模型。

\textbf{锚点回顾与核心差异}：与 GGHead 和 CGS-GAN 将头部作为单一整体建模不同，3DGH 将头发和面部解耦为两个独立分支：各自拥有独立的模板网格和 UV 纹理图（共 131K 高斯），分别通过独立的 StyleGAN2 生成器 $G_{\mathrm{hair}}$ 和 $G_{\mathrm{face}}$ 生成高斯属性。

\textbf{可变形头发几何先验}：由于头发的几何变化远大于面部（从短发到长发、从蓬松到扁平），单一模板网格无法覆盖所有发型。3DGH 从多视角捕捉数据中拟合 283 个不同头发网格，通过 PCA 降维得到 32 个混合形状系数的线性变形模型 $M(\boldsymbol{\beta})$，在推理时由生成器的几何映射网络预测系数以适配不同发型，确保头发高斯分布在几何合理的表面上。

\textbf{头发-面部关联建模}：现实中发型与面部特征（如性别、种族）存在内在关联。3DGH 通过交叉注意力层将面部分支的隐编码 $w_{\mathrm{face}}$ 注入头发生成器的每个合成块，训练时以 10\% 概率随机丢弃条件（替换为零向量），推理时通过无分类器引导（CFG，Classifier-Free Guidance）因子 $\eta$ 混合条件和无条件特征，实现关联强度的连续可控调节。

\textbf{训练策略}：以 PanoHead 生成的前后视角图像作为监督信号，除标准对抗损失外，引入分割损失强化头发-面部的清晰分离——为头发和面部高斯分别分配不同 one-hot 标签渲染分割图，并用交叉熵损失（而非 $L_1$）缓解 $\alpha$ 混合在边界处的误标签问题。

\paragraph{AGORA}

AGORA（2025.12）~\cite{2512_06438} 将 GGHead 的无条件 3D GAN 框架扩展为可动画化（animatable）的化身生成器，核心问题是 GGHead 只能生成静态身份，无法根据表情参数驱动变形。

\textbf{锚点回顾与改了什么}：与 GGHead 仅用潜在编码 $z$ 生成单一规范（canonical）化身不同，AGORA 在其 UV 生成器基础上新增三个关键组件。（1）\textbf{空间形状条件化}（spatial shape conditioning）：直接将 FLAME 形状参数 $\beta$ 注入映射网络会导致 $w$ 被 $\beta$ 主导、多样性坍缩，因此 AGORA 改为计算形状孤立位移场 $V_{\text{shape}}(\beta) = M(\beta, \epsilon_0, \theta_0) - M(\beta_0, \epsilon_0, \theta_0)$，将其烘焙为 UV 对齐位移图并做逐样本方差归一化后拼接到生成器特征图中——这一空间化策略在几何敏感区域（如儿童较小的颅面比例）注入形状偏置，同时保留 $z$ 驱动的随机性。

（2）\textbf{轻量变形分支} $G_D$（仅 3M 参数，对比身份分支的 30M）：从主分支取 $64 \times 64$ 低分辨率特征图，通过两个连续 StyleGAN2 块上采样至 $256 \times 256$ 预测表情依赖的高斯属性残差 $\Delta\mu, \Delta s, \Delta q$，与规范属性通过后激活求和（位置）、四元数乘法（旋转）和求和（对数缩放）进行组合——由于 FLAME 表情参数 $\epsilon$ 和颌部姿态 $\theta_{\text{jaw}}$ 仅通过 style modulation 注入变形分支，表情控制与身份生成分离，动画时仅需运行轻量分支。（3）\textbf{位移着色判别器}：朴素图像判别器无法保证表情精度，Next3D 的 UV 着色条件化在强表情下欠表达；AGORA 改用表情孤立顶点位移 $V = V_{\text{posed}} - V_{\text{neutral}}$ 作为着色信号合成判别器输入，在 AED 和 AED-jaw 指标上显著提升表情一致性。

\textbf{设计动机}：GGHead 的高斯绑定在模板网格 UV 空间使其天然具备动画驱动的基础架构，但原始框架缺乏表情控制通路；AGORA 的双分支设计将身份（一次性计算并缓存）和表情（逐帧运行 3M 参数网络）解耦，实现了高质量动画与实时推理的兼得。

\textbf{AGORA-M 变体}：离线阶段从训练好的 AGORA 模型采样 $N{=}10000$ 组参数元组，计算 posed-minus-neutral 高斯属性残差矩阵 $B \in \mathbb{R}^{N \times (M \cdot D)}$，对 $B$ 做 SVD 保留前 $K{=}64$ 个右奇异向量作为共享高斯混合形状（Gaussian blendshapes），再训练轻量系数回归器 $f_{\text{mlp}}(w, \epsilon, \theta_{\text{jaw}})$（仅两层 MLP）；推理时高斯属性更新退化为一次线性组合 $A^M = A_{\text{neutral}} + \sum_{k=1}^{K} c_k U_k$，在单 RTX A6000 上以 560 FPS 渲染，在 VIVO X200 Ultra 手机上通过 WebGL 实现以 60 FPS 渲染，全部动画计算在端侧完成。

\textbf{公式与参数}：训练分两阶段——$256^2$ 分辨率 + 65K 高斯训练 6.5M 图像，再切换至 $512^2$ + 262K 高斯训练 14M 图像，总训练耗时 4 天（4$\times$ A6000）。正则化系数对主分支设 $\lambda_{\text{pos}} = 0.25$、$\lambda_{\text{scale}} = 0.5$，对变形分支设 $\lambda^D_{\text{pos}} = 1.5$、$\lambda^D_{\text{scale}} = 1.5$（更严格约束变形幅度以保持局部化）。最终模型在 FFHQ 上取得 FID 3.17，同时支持实时动画与移动端部署。

\paragraph{AvatarPointillist}

AvatarPointillist（2026.04）~\cite{2604_04787} 保留 GGHead 基于 FLAME 网格驱动的动画框架，但将点云获取方式从 UV 空间的 2D CNN 单次前向生成替换为仅解码器 Transformer 的自回归逐点生成。

\textbf{核心差异}：GGHead 的高斯数量由 UV 图分辨率固定（如 $256^2$ 即 65K），AvatarPointillist 则将规范空间点云表示为长度 $4N$ 的 token 序列——$N$ 个点各含 4 个量化 token（$x, y, z$ 坐标经 1024 级量化 + FLAME 三角面片绑定索引），由 Transformer 依条件概率 $p(T_n | T_{<n})$ 逐 token 预测，点数量 $N$ 不预设而是由自回归过程的结束标记决定，可根据主体复杂度自适应调整。训练数据通过 GaussianAvatars 为每个身份拟合高斯-绑定表示，坐标按 $y$-$z$-$x$ 字典序排列后量化为离散序列。

\textbf{两阶段训练}：第一阶段训练 AR 模型用标准交叉熵损失做 next-token 预测（16 块 H20 GPU、50K 步、窗口大小 12000 的截断训练策略处理长序列）；第二阶段冻结 AR 模型，单独训练基于 Transformer 的高斯解码器，该解码器将去量化后的坐标经位置编码生成几何特征 $P_n$，同时提取 AR 模型最终隐状态并将每点 4 个 token 的隐特征经 MLP 组装为语义特征 $F_n^p$，两者拼接后预测完整高斯属性（颜色、不透明度、缩放、旋转及精细位置偏移 $\Delta p_k$），渲染损失为 $L_{\text{total}} = \lambda_{L1}L_{L1} + \lambda_{\text{SSIM}}L_{\text{SSIM}} + \lambda_{\text{LPIPS}}L_{\text{LPIPS}} + \lambda_{\text{Reg}}L_{\text{Reg}}$，其中偏移正则化约束位置修正幅度。

\textbf{动画机制}：AR 模型为每个点预测的绑定索引将点锚定到 FLAME 网格三角面片，通过重心坐标插值获取 LBS 权重和表情混合形状，再经标准 FLAME 变形流程驱动——与 GGHead 的 UV 图隐式绑定不同，AvatarPointillist 显式输出每个点的绑定关系，消融实验表明静态 FLAME 顶点方案（跳过 AR 模型直接用 FLAME 顶点作为输入）无法捕捉主体特定几何（特别是头发等复杂区域），证实了自回归逐点生成的必要性。实验在 NeRSemble 数据集 419 个身份上进行，自重演和交叉重演任务上全面超越 AvatarArtist、Portrait4D-v2、LAM 和 GAGAvatar 等基线。

\paragraph{SplitAvatar}

SplitAvatar（2026.05）~\cite{2605_25751} 同样基于 FLAME 网格驱动动画，但采用图神经网络（GNN）的层级分裂策略取代 GGHead 的 UV 参数化和 AvatarPointillist 的 Transformer 自回归。

\textbf{核心差异}：GGHead 在 UV 空间用 CNN 一次性生成所有高斯，AvatarPointillist 用 Transformer 逐点自回归生成，SplitAvatar 则从 FLAME 网格顶点初始化 5023 个基础表情高斯（由 MLP 在每个跟踪顶点位置预测一个高斯），通过 $L{=}3$ 层自回归图分裂网络（GSN）逐级扩展——每层以图注意力网络（GAT，4 头）为骨干，将每个父高斯分裂为 $k{=}2$ 个子高斯（总高斯数按 $5023 \times k^L$ 指数增长），同时预测精细高斯属性。

\textbf{拓扑扩展}：GGHead 依赖 FLAME UV 的固定拓扑隐式维持高斯空间关系，SplitAvatar 需要显式处理分裂后的图结构。

其设计将边集 $E^i$ 分为两类——拓扑继承边 $E^i_{\text{topo}}$（在每个 $k$ 倍副本内保留原始 FLAME 网格拓扑）和内部连接边 $E^i_{\text{conn}}$（同一父顶点分裂出的 $k^i$ 个子节点之间全连接），使图注意力机制在拓扑邻近的三角面片之间传播信息，而非对无序点云做全局注意力。

\textbf{门控剪枝}：多层分裂导致高斯冗余和聚集，SplitAvatar 设计了软掩码网络 $N_{\text{mask}}$（256 维 MLP）为每个高斯预测二值掩码，采用延迟过滤策略——训练期间不实际删除高斯（避免重新计算图索引和批内序列长度不一致），仅用逻辑与运算 $M_{i+1}^{\text{logical}} = q(m_{i+1}) \wedge q(M_i)$ 标记无效高斯，在最终渲染时忽略；梯度反传通过直通估计器（STE）绕过量化函数和逻辑与的不可微问题。最终渲染高斯为身份高斯与各层掩码后表情高斯的叠加：$G = G_{\text{identity}} + \sum_{i=0}^{L} M_i G_i$。

\textbf{设计动机}：与 GGHead 的固定密度 UV 图和 AvatarPointillist 的全局 Transformer 相比，SplitAvatar 的 GNN 利用 FLAME 网格拓扑将面部分区域（如眼部、嘴部）的信息传播局部化，使关键面部区域的高斯获得更精确的优化信号；门控机制则允许不同区域拥有不同密度——消融实验表明 $k{=}2, L{=}3$ 的中等配置在重建质量和表情精度上最优，过度分裂（$k{=}3, L{=}3$）反而导致高斯冗余和泛化退化。

\textbf{公式与参数}：在 VFHQ 数据集（636,359 帧、14,513 视频片段）上训练 200K 次迭代（77 GPU 小时，单块 A100），推理速度 41 FPS、GPU 显存仅 3.1 GB，在 VFHQ 和 HDTF 上的 PSNR、LPIPS、AED、AKD 指标均优于 CVTHead、LAM、GPAvatar、Portrait4D-v2 和 GAGAvatar。

\paragraph{COSY}

COSY（2026.05）~\cite{2605_24114} 从组件可控性角度出发，解决了 GGHead 和 CGS-GAN 等单一统一生成器的根本局限——无条件潜在空间中所有语义属性（发型、肤色、配饰）高度纠缠，编辑某一属性时不可避免地改变其他区域。

\textbf{锚点回顾与核心差异}：GGHead 和 CGS-GAN 使用单一生成器将所有语义区域混合在同一个潜在空间中；COSY 则将头部合成显式分解为四个独立的 CGS-GAN 风格子生成器——头发、皮肤、眼镜和躯干，各自拥有独立的潜在编码，编辑时仅修改目标组件的潜在向量而保持其余不变。

\textbf{无监督组件分离}：组件间的分离完全通过隐式信号实现，无需分割掩码或几何先验。具体地，永久组件（皮肤、头发、躯干）通过稀疏的 10-bin RGB 颜色直方图（每个通道独立统计 10 个 bin，三通道拼接为 30 维向量）调制子生成器中高斯基元的颜色属性来分离——该信号仅输入对应子生成器及其判别器，使生成器隐式学会仅负责指定区域。

可添加组件（如眼镜）通过布尔标志（0 或 1）控制存在性：存在时正常生成并加入合成场景，不存在时将高斯不透明度减去一个大值以隐藏；判别器同步接收该标志，惩罚错误缺失或错误出现的生成结果。

\textbf{全局一致性机制}：多个独立子生成器的核心挑战是确保组合后几何与光照的匹配。COSY 提出三项关键设计。（1）\textbf{潜在混合}（latent mixing）：训练时以 20\% 概率随机替换某一子生成器的潜在向量（其余保持不变），使判别器产生训练信号鼓励不同潜在组合的兼容性——实验表明 20\% 混合率在 FID（4.42）与编辑 FID（8.32）之间取得最优平衡。（2）\textbf{形状上下文}：所有子生成器共享一个低维形状潜在编码，该编码从不被混合，仅在最终合成层（将特征解码为高斯位置坐标时）注入，确保各组件对头部整体大小和形状达成一致——改变该编码仅改变头部形状而不改变身份。

（3）\textbf{光照上下文}：所有子生成器共享一个低维光照潜在编码（同样不参与混合），通过条件化输入潜在向量来控制全局光照方向，并用 LPIPS 正则化约束同一潜在在不同光照编码下的渲染差异——正则化权重设置为使平均 LPIPS 距离约为 0.3，恰好允许光影变化而阻止身份偏移。此外，为解决面部生成器在发际线附近几何覆盖不足的问题（隐式依赖头发生成器补全），COSY 引入 EMA 几何一致性损失，惩罚生成面部几何相对于指数移动平均面部几何的偏离，推动生成更完整的面部表面。

\textbf{判别器设计}：与 EGG3D 和 3DGH 等基于分割掩码监督的组合式 GAN 不同，COSY 的判别器接收完整渲染图像（而非遮罩区域），将所有颜色直方图和眼镜标志拼接为单一条件向量乘以判别器特征——语义分离完全由生成器侧的稀疏条件化实现，避免了分割掩码精度不足和早期训练阶段生成器输出结构化程度低时掩码监督失效的问题。

\textbf{公式与参数}：在 FFHQC 数据集上，COSY 在 $512^2$ 分辨率取得 FID$_\text{3D}$ 4.42，在 $1024^2$ 取得 5.43，与 CGS-GAN（4.53 / 5.25）和条件化 CGS-GAN 基线（4.52）持平。在编辑质量上，COSY 在添加眼镜任务中达到近 100\% 的召回率（对比 GANSpace 约 50\%、InterFaceGAN 约 60\%、StyleFlow 约 95\%），且在无关区域的 LPIPS 距离显著低于所有基于潜在空间编辑的方法。

\subsection{人脸细节建模}
\label{ss:16_1c}

本节讨论~4~条技术路线：（1）重光照与材质分解；（2）头发建模；（3）眼球控制与注视重定向；（4）人脸分析与人脸编辑。

\subsubsection{重光照与材质分解}

本路线介绍将~3DGS~人脸化身的光照与材质属性解耦以实现可重光照的方法。前文讨论的化身方法将光照效果"烘焙"到球谐系数中，无法在新光照条件下重新渲染。重光照的目标是：给定任意新光照环境，化身应产生正确的阴影、高光和间接照明效果。

使化身在不同光照条件下呈现真实外观是实现实用化的关键一步。这一问题被称为重光照（relighting）：给定任意新光照环境，化身应产生正确的阴影、高光和间接照明效果，而非仅在训练光照下看起来合理。

\paragraph{HeadLighter}

HeadLighter（2026.01）~\cite{2601_02103} 提出了基于光阶段（light stage）捕获数据的监督式光照分解框架。

\textbf{Baseline 与问题}：此前的 3DGS 人脸化身将光照效果"烘焙"到球谐系数中，无法在新光照条件下重新渲染。换一个光源方向，面部阴影和高光就完全错误。

\textbf{改了什么}：HeadLighter 将 3D 高斯的属性分解为与光照无关的固有属性（几何、反照率、粗糙度）和物理渲染组件（法线、环境光遮蔽），实现任意光照条件下的物理基着色。

\textbf{实现原理}：核心是双分支 StyleGAN 生成器架构。基础分支 $G_{base}: z_{base} \to (A_{geo}, a)$ 从潜在编码 $z_{base} \in \mathbb{R}^{512}$ 预测几何属性 $A_{geo} = \{\mu, s, q, o\}$ 和反照率 $a$；重光照分支 $G_{relit}: z_{lit} \to (n, r, T)$ 从独立潜在编码 $z_{lit} \in \mathbb{R}^{128}$ 预测表面法线 $n \in \mathbb{R}^3$、标量粗糙度 $r$ 和间接光照传输 $T$。两个分支共享 UV 参数化但使用不同潜在编码，实现几何与光照属性的解耦。训练数据通过光阶段系统采集，包含 46 个标定点光源和 28 台同步相机，记录均匀光照、方向光照和随机光照三种配置下的多视角图像。

\textbf{设计动机}：光阶段数据提供了逐光源的控制能力，使得监督信号可以精确区分"哪个光源照亮了哪个点"。这是从普通视频中无法获得的珍贵监督信息。

\textbf{公式与参数}：通过渐进式解耦训练逐步注入头部外观先验，所有高斯属性嵌入模板头部网格的 UV 图中，通过高斯在 UV 空间的固定坐标采样获得。

\paragraph{RelightAnyone}

RelightAnyone（2026.01）~\cite{2601_03357} 与 HeadLighter 依赖逐身份光阶段采集的思路截然不同，改为融合多个已有数据集实现跨身份泛化的可重光照头部化身。

\textbf{改了什么}：将 HeadLighter 的单身份、双分支 StyleGAN 架构替换为两阶段流水线——Stage 1 在四个不同光照配置的多视角数据集（共 929 个身份）上联合训练均匀光照下的多身份高斯化身（基于 RGCA 架构），Stage 2 训练一个 UNet 重光照网络将均匀光照高斯颜色映射为可重光照 RGCA 参数（漫反射反照率 $\rho_k$、固有辐射传输函数 $d_k$、镜面粗糙度 $\sigma_k$ 和法线 $n_k$）。两阶段顺序训练使得 Stage 1 可以利用仅有均匀光照的普通数据集（D2--D4），仅 Stage 2 需要带 OLAT 标注的数据集（D1，127 个身份），大幅降低了对昂贵光阶段设备的需求。

\textbf{实现原理}：Stage 1 引入可学习身份码 $z_* \in \mathbb{R}^{256}$ 和数据集特定光照码 $z_l \in \mathbb{R}^{4}$，通过将 $z_l$ 与 $z_*$ 拼接后送入颜色解码器 $D_c$，将不同数据集的光照差异从身份外观中解耦——这使得重光照网络可以通过简单替换 $z_l$ 将任意数据集的高斯颜色对齐到训练 Stage 2 所用的特定均匀光照条件。Stage 2 的 UNet 包含共享编码器 $E$ 和两个解码分支：视角无关分支 $D_{ci}$ 预测 $\{\rho_k, d_k, \sigma_k\}$，视角相关分支 $D_{cv}$ 预测 $\{v_k, \delta_k\}$（法线残差），其中相机方向 $\mathbf{o}$ 在瓶颈层拼接到特征图上。为使网络学习形状相关的阴影变化（如自遮挡），3D 高斯位置 $t_k$ 作为额外三个通道拼接到编码器输入中。为避免反照率与光照参数的非物理解耦，引入两个正则化项：$L_\rho$ 鼓励预测反照率接近均匀光照下的平均纹理，$L_{mono}$ 鼓励漫反射球谐系数 $d_k$ 接近单色（$d_k^r \approx d_k^g \approx d_k^b$）。对新身份的拟合采用两步优化：先反演（固定网络，优化 $z_*$ 和逐点光源 RGB 强度），再微调 Stage 1 权重（Stage 2 冻结以保持重光照先验），支持单张 in-the-wild 图像输入。

\textbf{设计动机}：两阶段解耦的核心动机是数据可用性——高质量 OLAT 数据集稀缺且采集成本高昂，但普通均匀光照的多视角人脸数据集相对丰富。通过将身份先验学习和重光照学习分离，Stage 1 可以充分利用所有可用数据（929 个身份 vs. 仅 127 个 OLAT 身份），Stage 2 则专注在有限 OLAT 数据上学习从颜色到物理参数的映射。

\textbf{公式与参数}：高斯数量 $M = 1024 \times 1024$，所有高斯嵌入模板网格的 UV 纹理图；身份码 $z_* \in \mathbb{R}^{256}$，光照码 $z_l \in \mathbb{R}^4$；Stage 2 损失 $L_2 = \lambda_1 L_1 + \lambda_s L_s + \lambda_c L_{c-} + \lambda_n L_n + \lambda_\rho L_\rho + \lambda_{mono} L_{mono}$。

\paragraph{HRAvatar}

HRAvatar（2025.03）~\cite{2503_08224} 从单目视频重建可动画、可重光照的高保真头部化身。该方法解决三个核心问题：(1)~表情跟踪误差——通过端到端训练的表情编码器替代传统的拟合式跟踪，直接以光度损失为监督学习更精确的表情参数；(2)~个体变形差异——为每个高斯点独立学习形状基 $S \in \mathbb{R}^{N \times 3 \times |\beta|}$、表情基 $E \in \mathbb{R}^{N \times 3 \times |\psi|}$ 和姿态基 $P \in \mathbb{R}^{N \times 39K}$，以及逐点的混合蒙皮权重，替代 FLAME 的全局共享基；(3)~光照可分离性——将头部外观分解为反照率、法线、粗糙度等物理属性，整合基于物理的着色模型。

\paragraph{GTAvatar}

GTAvatar（2025.12）~\cite{2512_09162} 将 2DGS 与 UV 纹理映射结合，与上述 HRAvatar 的逐高斯独立属性不同，GTAvatar 将规范高斯基元的局部帧嵌入 FLAME 模板网格 UV 空间，使每个高斯的物理材质属性（反照率、粗糙度、法线）从 UV 纹理图中动态采样，实现连续可编辑的材质纹理和实时基于图像的光照（image-based lighting），支持单目视频重建可重光照化身。

\textbf{锚点回顾}：HRAvatar 为每个 3D~高斯点独立存储一组材质属性（反照率、法线、粗糙度），属性值离散分散在各个基元上，无法构成连续的纹理图，因此不具备 UV~空间的可编辑性。

GTAvatar 从三个维度突破这一限制：(1)~将 HRAvatar 的 3D~高斯替换为 2D~高斯 surfel，利用 2DGS 固有的切平面局部坐标系作为 UV~映射的桥梁；(2)~将逐基元独立属性替换为全局共享的 5~通道可学习 UV~纹理图（反照率 $\rho \in [0,1]^3$、粗糙度 $r \in [0,1]$、镜面反射率 $f_0 \in [0,1]$）外加 2~通道法线扰动图，材质属性通过双线性插值在渲染时动态采样；(3)~引入 UV~正则化损失确保纹理连续性，而非 HRAvatar 式的逐基元优化。

\textbf{改了什么}：核心创新是高效的切空间到 UV~空间映射算子 $\mathcal{M}$。给定光线-splat 交点 $(s,t)$，将其正交投影到绑定三角形平面，再经重心坐标插值得到 UV 坐标。关键是利用映射的线性性质，将整个过程压缩为单次矩阵乘法：$\mathcal{M}(s,t) = \text{uv}_0 + J_{\text{st} \to \text{uv}} [s, t]^T$，其中 $J_{\text{st} \to \text{uv}} = J_{\text{uv}} J_{\text{v}}^{\dagger} J_{\text{st}}$ 将 splat 切空间直接映射到 UV 空间（$J_{\text{v}}^{\dagger}$ 为三角形世界坐标 Jacobian 的伪逆），$\text{uv}_0$ 由高斯基元的重心坐标直接计算。

光照模型采用 Cook-Torrance BRDF 的 split-sum 近似，漫反射项 $L_{0,\text{diffuse}} = \rho \cdot \mathcal{L}(\mathbf{n}, 0)$ 与镜面项 $L_{0,\text{spec}} = R(f_0, r, \boldsymbol{\omega}_0 \cdot \mathbf{n}) \cdot \mathcal{L}(\boldsymbol{\omega}_r, r)$ 分离，其中 $\mathcal{L}$ 为可学习的 RGB cubemap 环境光（$6 \times 32^2 \times 3$），$R$ 为预计算查找表。法线由 splat 几何法线与 UV 法线扰动图合成：切空间扰动 $\mathbf{n}_t = [n_x, n_y, \sqrt{1 - n_x^2 - n_y^2}]^T$ 经旋转矩阵 $R = [\mathbf{s}, \mathbf{t}, \mathbf{n}_g]$ 变换到世界空间。训练引入三组损失：UV 扭曲损失 $\mathcal{L}_{\text{uv\_dist}} = \sum_{\text{rays}} \sum_{i,j} \omega_i \omega_j \|uv_i - uv_j\|_2$ 聚合同一光线上所有 splat 交点的 UV 坐标，FLAME PCA 反照率先验 $\mathcal{L}_{\text{stat\_albedo}} = \|\mathcal{T}_\rho - \tilde{\mathcal{T}}_\rho\|_1$ 引导物理分解，UV 边界掩码防止采样纹理图集间隙。

\textbf{设计动机}：GTAvatar 的核心动机是弥合高斯泼溅的表达力与 UV~纹理映射的可编辑性之间的鸿沟。HRAvatar 虽能重光照，但其离散的逐基元属性无法支持直观的局部编辑（如添加纹身、修改妆容、替换材质），也无法在 UV~空间施加物理先验约束。GTAvatar 利用 2DGS 的切平面结构实现高效 UV~映射，使渲染出的纹理图天然连续、可直接用标准图像编辑工具修改，且编辑后仍与物理光照模型兼容——这是此前可重光照高斯化身均不具备的能力。

相比之下，同期工作 FATE 需要额外的神经网络烘焙阶段才能获得连续纹理，且编辑后丢失高频细节。

\textbf{公式与参数}：UV 纹理分辨率 $1024 \times 1024$；材质通道 5 维（$\rho \in \mathbb{R}^3, r \in \mathbb{R}, f_0 \in \mathbb{R}$）+ 法线扰动 2 维；环境光 cubemap $6 \times 32^2 \times 3$；总损失 $\mathcal{L}_{\text{total}} = \mathcal{L}_{\text{photo}} + \mathcal{L}_{\text{PBR}} + \mathcal{L}_{\text{uv}} + \mathcal{L}_{\text{geom}}$；训练 15 epochs，NVIDIA RTX A5000 上耗时 1--2 小时；INSTA 数据集上 PSNR 30.52 dB、SSIM 0.9537、LPIPS 0.0552（HRAvatar 分别为 30.36 dB、0.9482、0.0569）；静态几何渲染 175 FPS、动态几何 83 FPS（10k 高斯），同等 LPIPS 下模型大小显著小于 HRAvatar。

\subsubsection{头发建模}

本路线介绍基于~3DGS~的头发几何重建、发丝表示和动态头发模拟方法。

头发的高几何复杂性和各向异性反射特性使其成为人脸化身中最难处理的组件之一。一根头发的直径仅约 50--100 微米，而人类头部平均拥有约 10 万根发丝，这意味着发丝级别的重建需要处理极其密集的细长几何结构。

\paragraph{Gaussian Haircut}

Gaussian Haircut（2024.09）~\cite{2409_14778} 提出了双重表示方案。

\textbf{Baseline 与问题}：头发重建面临独特的挑战——发丝具有极端的各向异性（细长形状），传统的等向 3D 高斯无法有效表示发丝的定向几何，而直接优化发丝参数则缺乏可微分渲染路径。

\textbf{改了什么}：Gaussian Haircut 设计了"非结构化高斯→结构化发丝"的两阶段管线，先用 3D 高斯建立场景级别的重建和方向场，再将非结构化高斯转化为发丝级别的表示。

\textbf{实现原理}：第一阶段（3D 线提升）使用修改版的 3DGS 进行场景重建，核心创新是将方向图（orientation map）信息嵌入到高斯的协方差矩阵中——每个高斯的协方差矩阵最大方差方向对齐局部发丝方向，正交方向的方差提供不确定性度量。渲染时同时输出颜色图和方向图，两者都具有高保真度。第二阶段（发丝拟合）采用粗到细策略：先用潜在纹理图参数化发丝，用第一阶段的高斯作为光度损失和几何损失的伪真值；粗优化后解码为显式发丝图并进行精细拟合。发丝通过在其每个线段上附着 3D 高斯来实现可微分渲染——这些高斯的缩放仅有一个自由度（与线段长度成正比），另外两个维度固定为小的预定义值。

\textbf{设计动机}：两阶段设计利用了 3DGS 在非结构化场景重建中的优势（不需要拓扑先验），同时通过发丝扩散模型正则化确保内部结构的物理合理性。

\paragraph{CGHair}

CGHair（2026.04）~\cite{2604_03716} 引入发卡聚类（hair card clustering）将发丝聚类为代表性发卡并分组为共享纹理码本，与 Gaussian Haircut 保留完整发丝表示不同，CGHair 先用生成式 PCA 先验加速发丝重建（4 倍加速），然后将重建的发丝参数化为头皮 UV 空间的潜在纹理图，通过聚类和码本量化将纹理从独立发丝压缩为共享表示，实现了超过 200 倍的内存压缩。

\textbf{锚点回顾}：Gaussian Haircut 采用"非结构化高斯→结构化发丝"的两阶段管线，最终输出的是逐根独立参数化的发丝——每根发丝由圆柱高斯链构成，拥有各自独立的球谐（SH）外观系数。这种逐发丝独立表示虽然精度高，但存在严重冗余：一个典型发型含约 10 万根发丝，导致数百万个高斯基元和庞大的 SH 参数存储开销。CGHair 正是观察到发型内部的结构相似性——相邻发丝在几何走向和视觉外观上高度一致——从而提出了发卡引导的层次化压缩表示。

\textbf{改了什么}：CGHair 提出了三阶段紧凑管线：(1)~基于生成式先验的快速发丝生成；(2)~发卡几何构建与纹理提取；(3)~基于发卡聚类的共享高斯纹理码本压缩。在阶段(1)中，CGHair 采用 PERM 参数化几何发丝模型作为先验，将头发表示为头皮 UV 空间的二维潜在纹理图（由引导图 $\mathbf{G}$ 捕获全局结构和风格图 $\mathbf{S}$ 捕获局部细节），通过预训练的 StyleGAN2 引导解码器 $\mathcal{D}_{guide}(\alpha)$ 和 VAE 风格解码器 $\mathcal{D}_{style}(\beta)$ 从图像优化潜在码 $\alpha$、$\beta$，再经 PCA 解码器 $\mathcal{D}_{\mathrm{PCA}}$ 输出完整 3D 发丝。圆柱高斯沿发丝线段附着，联合优化潜在码和微调解码器实现图像驱动重建——相比 Gaussian Haircut 依赖扩散模型正则化的慢速潜在空间优化，CGHair 利用 PERM 的强几何先验实现了 4 倍加速。在阶段(2)中，重建的发丝通过 k-means 聚类为 $N_c$ 个簇，每个簇用 B 样条拟合的引导发丝 $\bar{\mathcal{S}}$ 作为发卡的中央轴；法向量 $\mathbf{n}_k$ 通过受约束优化求解（$\mathbf{n}_k \cdot \mathbf{t}_k = 0$，最小化发丝点到发卡平面的距离之和），副切向量 $\mathbf{b}_k = \mathbf{n}_k \times \mathbf{t}_k$ 确定发卡宽度，构建条带网格；发丝通过优化反映射（可学习的 UV 坐标和法向位移 $\delta_i$）投影到发卡表面生成二维发丝纹理。在阶段(3)中，基于发丝纹理特征将发卡进一步聚类为 $N_T$ 个外观组，每组分配一个共享高斯外观码本（$K$ 个条目），每个码条目编码一根完整发丝的外观特征（维度 $D=64$，远低于直接存储 SH 系数的数千维），通过全局共享 MLP 解码器 $\phi_{dec}$ 将码本特征解码为各高斯的独立 SH 参数；每根发丝通过 Gumbel Softmax 产生的软分配权重 $\pi_k$ 对码本条目进行加权组合 $\mathbf{F}_{\text{strand}} = \sum_{k=1}^{K} \pi_k \cdot \mathbf{F}_k$，实现外观建模。

\textbf{设计动机}：发卡是游戏和 VFX 工业中广泛使用的紧凑头发表示——它将一束几何相似的发丝压缩为共享几何和纹理的多边形条带。CGHair 将这一工业实践引入 3DGS 管线：层次化聚类（发丝→发卡→发卡外观组）逐层消除冗余，从"每根发丝独立参数化"压缩到"每组发丝共享码本"，实现了超过 200 倍的内存压缩。此外，通过将 SH 参数的冗余压缩到低维特征空间（$D=64$），而非直接量化高维 SH 系数，CGHair 避免了码本存储爆炸问题。全局共享 MLP 解码器进一步降低了参数开销，使整个外观建模在发丝级别而非单个高斯级别进行——这意味着优化过程中梯度从渲染损失直接回传到发丝级特征，而非数十万个独立高斯参数，提高了优化的稳定性和效率。

\textbf{与 Gaussian Haircut 的具体差异}：(1)~发丝生成路径不同——Gaussian Haircut 用潜在纹理图配合扩散模型正则化从零优化发丝参数，CGHair 利用 PERM 生成式先验（StyleGAN2 + VAE + PCA 解码器链）快速初始化并从图像微调，实现 4 倍加速；(2)~外观建模粒度不同——Gaussian Haircut 为每个高斯独立存储 SH 系数，CGHair 通过发卡聚类和共享码本将外观压缩到发丝级特征（200 倍以上内存压缩）；(3)~表示范式不同——Gaussian Haircut 保留逐发丝独立表示，CGHair 引入发卡层次结构（发丝→发卡→外观组），显式建模发型内部的结构冗余；(4)~共同局限——两者均需要多视角输入，且重建结果为静态发丝，不具备动画驱动能力。

\paragraph{Im2Haircut}

Im2Haircut（2025.09）~\cite{2509_01469} 结合全局发型先验和局部优化，在合成数据（PERM 数据集，21054 种发型）和真实数据（Hairstep 数据集）上联合训练 Vision Transformer 先验模型，通过粗到细的 PCA 系数预测（先 10 个主成分捕获基本结构，再 54 个高频成分精修细节）实现从单张照片的基于高斯泼溅的发丝级重建。

\textbf{锚点回顾}：CGHair 和 Gaussian Haircut 均依赖多视角输入进行发丝重建——它们需要多张不同视角的图像来建立场景级别的 3D 几何。Im2Haircut 面向一个更受限也更具实用性的设定：仅从单张照片（通常是正面肖像）恢复完整的三维发丝几何。这一设定面临独特的挑战——照片只捕捉到头发可见的外层结构，而内部几何和背面几何完全被遮挡，必须有强大的发型先验来"补全"不可见区域。

\textbf{改了什么}：Im2Haircut 提出了"全局先验预测 + 局部高斯泼溅精修"的两步方案。全局先验模型采用粗到细的 Vision Transformer 架构：粗分支从输入图像提取方向图、分割图和深度图（4 通道特征），经卷积编码器和 ViT 处理（注意力层中丢弃头发区域外的 token 以提高效率），预测前 10 个 PCA 主成分 $Z_c$（捕获发型的基本形状和走向）及秃发掩码 $M_c$；细分支以 2 倍分辨率的特征图工作，将粗分支输出上采样后与自身编码特征拼接，预测剩余 54 个 PCA 系数 $Z_f$（高频细节和曲率变化）及精修秃发掩码。训练采用三阶段混合策略：(1)~粗分支在合成数据（PERM，168032 张渲染图像）上以 3D 几何损失 $\mathcal{L}_{\text{coarse}} = \mathcal{L}_{\text{MAE}} + \lambda_{\text{PCA}} \|Z_c - \hat{Z}_c\|_2 + \lambda_{\text{mask}} \|M_c - \hat{M}\|_1$ 训练，其中 $\mathcal{L}_{\text{MAE}}$ 包含点位置、方向和曲率监督；(2)~细分支在合成数据上以类似损失训练，但引入可见性加权 $\mathcal{L}'_{\text{MAE}}$（可见点权重为 3，不可见为 1），使高频细节聚焦于输入视角可见区域；(3)~混合训练将合成 3D 损失与真实数据（Hairstep，1250 张"野外"照片）的渲染自监督损失 $\mathcal{L}_{\text{real}} = \lambda_{seg}\mathcal{L}_{seg} + \lambda_{dirmap}\mathcal{L}_{dirmap} + \lambda_{pen}\mathcal{L}_{pen}$ 联合优化，其中分割损失和方向损失通过 3DGS 可微分光栅化将预测发丝渲染到 2D 并与标注对比，穿透损失 $\mathcal{L}_{pen}$ 惩罚发丝穿透头部网格。推理时对单张输入图像进行 400 步微调（前 20 步仅优化粗成分，后续优化全部 64 个成分），整个过程约 10 分钟。

\textbf{设计动机}：Im2Haircut 的核心创新在于合成-真实数据混合训练策略。纯合成数据训练的先验模型存在严重的域迁移问题——现有合成数据集的发丝数量少（远低于真实头发的 10 万根）且渲染纹理不真实，导致模型在真实图像上的泛化能力差。Im2Haircut 通过在合成数据上学习内部 3D 结构（有真值监督），在真实数据上通过 3DGS 可微分渲染学习外层可见结构的真实外观（无需 3D 真值），两者互补：合成数据确保不可见区域的几何合理性，真实数据缩小域迁移差距并提高发型多样性。梯度传播策略也体现了这一权衡——真实数据的梯度仅回传到粗分支（10 个成分），因为渲染损失信号过于嘈杂，直接训练细分支会导致不真实的高频伪影。

\textbf{与 CGHair 的具体差异}：(1)~输入需求不同——CGHair 需要多视角图像作为输入，Im2Haircut 仅需单张照片即可完成重建；(2)~发丝生成路径不同——CGHair 先用 PERM 先验从多视角图像优化潜在码初始化发丝，再经发卡聚类压缩；Im2Haircut 用 ViT 直接从单图回归 PCA 系数（预测而非优化），再通过 400 步微调精修；(3)~压缩策略不同——CGHair 通过发卡层次化聚类和共享码本实现存储压缩，Im2Haircut 未引入显式压缩机制，其紧凑性来自 PCA 参数化本身（64 个系数即可表示完整发丝）；(4)~数据利用不同——CGHair 的 PERM 先验仅在合成数据上预训练，而 Im2Haircut 显式引入真实数据参与训练，有效缩小了域迁移差距；(5)~共同基础——两者都使用 PERM 数据集的 PCA 发丝表示作为底层参数化，且都用圆柱高斯附着发丝实现可微分渲染，但用途不同：CGHair 用于多视角重建中的联合优化，Im2Haircut 则用于单图先验预测后的微调精修。

\textbf{从静态到动态——头发的动画驱动}。以上讨论的 Gaussian Haircut、CGHair 和 Im2Haircut 所解决的都是静态头发的重建问题：它们从图像中恢复发丝的几何形状和空间位置，但重建结果本身并不具备任何动画能力。当头部运动时，头发需要随头皮牵引、受重力下垂、因惯性摆动——这要求在静态重建的基础上增加物理驱动的变形机制。

\paragraph{One-shot Compositional 3D Head Avatars}

One-shot Compositional 3D Head Avatars（2026.04）~\cite{2604_14782} 正是面向这一需求，它不仅重建头发几何，更进一步赋予头发可动画的物理变形能力。

\textbf{锚点回顾}：Gaussian Haircut 将非结构化 3D 高斯转化为发丝级别的结构化表示，但发丝本身是静态的——它们被固定在重建时的空间位置上，不随头部运动而变形。One-shot Compositional 从根本上改变了这一范式：它不再追求发丝级别的精确重建，而是将头发视为一组非结构化高斯基元的集合，通过笼结构（cage）和基于位置动力学（PBD）的物理模拟驱动这些基元产生整体性的动态变形。

\textbf{改了什么}：One-shot Compositional 提出了"单张图像 → 面部/头发解耦重建 → 独立变形驱动 → 统一渲染"的组合式头部化身方案。核心改动有三方面。(1)~头发/面部解耦与重建：给定一张正面肖像，先用 2D 去发滤镜生成无发图像，再用 FaceLift 将原始图像和无发图像分别提升为高保真 3DGS 表示 $\mathcal{G}_{head}$ 和 $\mathcal{G}_{bald}$。通过 SAM2 多视角语义分割监督，为每个高斯添加可学习的 2 维特征 $\mathbf{f} \in \mathbb{R}^2$，经 softmax 分离出头发高斯 $\mathcal{G}_{hair}$；随后引入边界感知重分配策略（boundary-aware reassignment），利用 2D 边界与深度信息将 $\mathcal{G}_{hair}$ 划分为边界和非边界子集，以非边界头发高斯和 $\mathcal{G}_{bald}$ 为两个类中心，按颜色和尺度特征的相似度重新分配边界高斯，有效去除残留的皮肤基元。

(2)~面部绑定与驱动：$\mathcal{G}_{bald}$ 通过联合重建与注册损失 $\mathcal{L}_{r\&r}$（包含 Chamfer 距离项 $\lambda_1 \mathcal{L}_{chamfer}(\{\boldsymbol{\mu}_{prior}\}, \{\boldsymbol{\mu}_{bald}\})$ 和多视角光度重建项）与 FLAME 参数化网格建立绑定关系，每个高斯基元定义在其宿主三角形的局部坐标系中，随网格变形自然运动——缺失的牙齿和眼球则由预训练先验模型补充。(3)~笼-PBD 头发变形：这是该方法的标志性创新。将头发高斯点云体素化为水密源笼 $\mathcal{C}_s$（少于 500 个顶点），每个高斯椭球用中心加六个轴端点共 7 个控制点表示，并预计算这些控制点对笼顶点的均值坐标（MVC）权重 $\mathbf{W} \in \mathbb{R}^{N \times 7 \times M}$。动画时，笼顶点作为 PBD 粒子进行物理模拟：靠近头皮的顶点设为运动学粒子（零逆质量），通过线性混合蒙皮（LBS）跟随头部运动；自由粒子则受重力驱动进行半隐式欧拉积分预测位置 $\mathbf{p}_j^{(t)} \leftarrow \mathbf{c}_j^{(t)} + \Delta t \cdot \mathbf{v}_j^{(t)}$，再经约 15 次迭代约束投影（拉伸、弯曲、体积保持、碰撞）获得最终位置。变形后高斯的中心由主轴两端点的变形位置取中值 $\boldsymbol{\mu}_d = \frac{1}{2}(\mathbf{x}_d^{i^*+} + \mathbf{x}_d^{i^*-})$ 确定，旋转由主轴变形引起的相对旋转 $\Delta\mathbf{R}$ 更新 $\mathbf{r}_d = \Delta\mathbf{R} \mathbf{r}_s$，缩放按轴长比调整 $\mathbf{s}_d = \mathbf{s}_s \cdot \|\mathbf{x}_d^{i^*+} - \mathbf{x}_d^{i^*-}\| / \|\mathbf{x}_s^{i^*+} - \mathbf{x}_s^{i^*-}\|$。

\textbf{设计动机}：笼-PBD 方案的设计基于三个关键观察。(1)~One-shot 场景下没有多视角时序视频可供学习变形模式，因此不能依赖数据驱动的变形场，而 PBD 作为纯物理模拟方法只需几何约束即可产生可信的动态效果，无需任何训练数据。(2)~直接对数万个高斯基元做物理模拟计算代价过高，笼结构将高维变形降维到数百个控制顶点的运动上，再通过 MVC 插值传播给所有高斯，实现了实时性能（整个 PBD 模拟在 Taichi 框架中运行）。(3)~笼与高斯之间不可避免存在间隙，直接对笼顶点做碰撞检测会导致过度排斥——因此提出代理碰撞约束策略：为每个笼顶点预分配其最近邻高斯的 MVC 权重，在碰撞检测时先用权重插值出代理粒子 $\mathbf{p}_{proxy,j}^{(t)} = \sum_j w_j \mathbf{p}_j^{(t)}$，仅当代理粒子（而非笼顶点）穿透 FLAME 网格时才施加碰撞修正，确保修正力度与实际高斯位置一致，避免头发与头皮之间出现不自然的间隙。

\textbf{与 Gaussian Haircut 的具体差异}：Gaussian Haircut 重建的是静态发丝级几何（每根发丝由多个线段高斯组成），追求的是几何精度和拓扑正确性，但不考虑动画；One-shot Compositional 放弃了发丝级精度，转而用非结构化高斯表示头发区域，换取了物理驱动的动态变形能力。前者是"重建得准"，后者是"动得真实"。两者代表了头发高斯建模的两条互补路线：静态高保真重建 vs.~动态物理驱动。此外，Gaussian Haircut 需要多视角图像作为输入，而 One-shot Compositional 仅需单张正面肖像——输入需求的差异决定了它们面向不同的应用场景。

\paragraph{HairGS}

HairGS（2025.09）~\cite{2509_07774} 回到静态重建路线，但与 Gaussian Haircut 的策略形成鲜明对比。

\textbf{锚点回顾}：Gaussian Haircut 的核心思路是"先建非结构化高斯 → 再拟合结构化发丝"，发丝是从高斯几何中间接恢复的。HairGS 同样追求发丝级重建，但走了不同的路径：它直接将 3D 高斯作为发丝几何的基元，通过合并（merging）将短段拼接为长发丝，而非从高斯拟合到独立的发丝参数化。

\textbf{改了什么}：HairGS 提出三阶段优化管线，每个阶段都围绕 3DGS 的可微分光栅化器构建，全程无需预训练数据。(1)~几何重建阶段：以 FLAME 模型顶点（而非 COLMAP 稀疏点云）初始化高斯，配合 Gabor 滤波器提取的 2D 方向图及置信度图 $C_\theta$ 监督，引入双向方向损失 $\mathcal{L}_\theta = \frac{1}{HW}\sum \Delta_\theta(\hat{\theta}, \theta) C_\theta$（其中 $\Delta_\theta(\theta_1, \theta_2) = \min(|\theta_1 - \theta_2|, \pi - |\theta_1 - \theta_2|)$）加速高斯朝正确发丝方向收敛，同时用二值交叉熵掩码损失 $\mathcal{L}_m$ 约束高斯仅出现在头发区域，总体目标为 $\mathcal{L}_{first} = (1 - \lambda_{DSSIM})\mathcal{L}_1 + \lambda_{DSSIM}\mathcal{L}_{DSSIM} + \lambda_\theta \mathcal{L}_\theta + \lambda_m \mathcal{L}_m$。标准 3DGS 自适应致密化机制在此阶段确保几何充分覆盖。(2)~发丝生成阶段：第一阶段的高斯被转化为由两个关节点构成的短发丝段，每段用圆柱体高斯表示——中心取端点中值，旋转用 Rodrigues 公式将 $x$-轴对齐到段方向 $R = I + K + K^2/(1 + \mathbf{v} \cdot d)$，缩放为 $(\|p_{j+1} - p_j\|_2, \tau_j, \tau_j)$，其中 $\tau_j$ 是可学习的厚度参数。关键创新是合并方案（merging scheme）：将短发丝端点构建为二部图节点，边代价由空间距离和角度差异定义，用基于 K-D 树的贪心匹配（初始阈值 $d_m = 2$ mm、$\theta_m = 20^\circ$）而非全局最优的匈牙利匹配（因规模达数十万节点，匈牙利算法的计算和内存开销不可接受）逐步合并短发丝——每次合并在新端点的中点创建新关节并连接原始端点。优化过程中阈值逐步放宽至 $d_m = 4$ mm、$\theta_m = 40^\circ$，使更多远距离端点得以合并，平均发丝长度近乎翻倍。(3)~生长与精修阶段：恢复图像监督并引入角度平滑损失 $\mathcal{L}_{smooth} = \frac{1}{|\mathcal{C}|}\sum_{(a,b) \in \mathcal{C}} \theta_{a,b}^2$（当相连段夹角超过阈值 $\theta_s$ 时惩罚），防止弯曲发型中出现不自然的锐角；对超长线段进行分裂插入新关节，增加自由度以更好地拟合曲线。最终通过学习到的掩码值 $m$ 过滤非头发几何。

\textbf{设计动机}：HairGS 的设计围绕一个核心洞察展开——3D 高斯天然适合作为发丝的基元表示，因为高斯的各向异性椭球形状与发丝的细长几何高度吻合。与 Gaussian Haircut 先建无结构高斯再拟合发丝的间接路径不同，HairGS 直接在高斯空间中构建发丝拓扑，避免了中间表示转换带来的信息损失。合并方案避免了预定义每根发丝的点数（这在 Gaussian Haircut 和学习型方法中是固定超参数），而是由优化过程的梯度反馈动态决定发丝长度——简单发型形成较少关节的长发丝，复杂卷发则分裂为更多关节以捕捉曲率。此外，HairGS 提出了发丝一致性（strand consistency, SC）评估指标：对于每根真值发丝，计算其被正确匹配到单根预测发丝的点比例最大值并取均值，填补了现有评估指标只度量点云级几何精度、忽略发丝拓扑连通性的空白。

\textbf{与 Gaussian Haircut 的具体差异}：(1)~发丝生成路径不同——Gaussian Haircut 用潜在纹理图参数化发丝，HairGS 直接将高斯段拼接为发丝折线；(2)~方向监督不同——Gaussian Haircut 将方向嵌入协方差矩阵，HairGS 用独立的 2D 方向图损失 $\mathcal{L}_\theta$ 引导高斯对齐；(3)~初始化不同——Gaussian Haircut 用标准 COLMAP 点云，HairGS 用 FLAME 模型顶点，后者对头发的覆盖更充分；(4)~无需数据驱动先验——Gaussian Haircut 依赖发丝扩散模型正则化，HairGS 全程使用纯优化策略（约 1 小时完成），相比 Neural Haircut 的约 120 小时显著更高效。HairGS 的局限在于不处理动画驱动，重建的头发是静态的；而 One-shot Compositional 的笼-PBD 方案恰好弥补了这一缺口，两者在"静态精确重建"与"动态物理驱动"上各司其职。

\subsubsection{眼球控制与注视重定向}

本路线介绍基于~3DGS~的眼球注视重定向和眼周肌肉联动建模方法。

\paragraph{GazeGaussian}

GazeGaussian（2024.11）~\cite{2411_12981} 修改了 3DGS 头部化身的变形机制——用双流分离表示（面部变形分支 + 眼球旋转分支）替代单流统一变形——光栅化管线和混合渲染公式保持不变。

\textbf{Baseline 与问题}。此前基于 3DGS 的头部化身方法（如 Gaussian Head Avatar（2023）~\cite{xu2023gaussianheadavatar}）使用单一 3DGS 分支统一建模整个头部：所有高斯（包括眼球区域）通过同一组 MLP 接收表情潜码 $\psi$ 和头部姿态 $\theta$，输出统一的变形偏移。这种做法忽略了眼球运动的独特物理特性——眼球是一个近似刚体，其运动本质上是绕固定轴的旋转，而非自由变形。将眼球与面部一起用变形场处理，导致注视重定向时眼球形状畸变、注视方向不准确。此外，2D 图像级的注视重定向方法（如 DeepWarp）缺乏 3D 空间一致性，而 NeRF 类方法（如 GazeNeRF）虽然建模了 3D 结构，但隐式地修改特征图来实现注视变化，计算开销大且精度有限。

\textbf{改了什么}。GazeGaussian 将单流统一变形拆分为两个独立分支：面部分支（face-only branch）继续使用 MLP 变形策略，而眼球分支（eye branch）引入显式刚性旋转机制。拆分从初始化阶段就开始了——先从 SDF 场提取中性网格（neutral mesh），再用 3D 面部标志点将顶点分为面部区域和眼部区域，分别初始化两组规范高斯。关键区别在于：面部分支通过 MLP 预测位置偏移量 $\Delta\mu^f$，以表情潜码 $\psi$ 和头部姿态 $\theta$ 为条件；眼球分支则直接用目标注视方向 $g$ 作为旋转信号（注意：眼球分支省略了 $\theta$，因为眼球旋转独立于头部姿态），将规范空间中的眼球高斯旋转到目标位置。渲染时，两组高斯分别光栅化为面部特征图 $M_f$ 和眼球特征图 $M_e$，同时拼接后统一光栅化为完整头部特征图 $M_h$，三条特征图共同送入神经渲染器。

\textbf{实现原理——面部变形分支}。面部分支的规范高斯具有标准属性 $\{\mu_0^f, z_0^f, R_0^f, S_0^f, o_0^f\}$，其中 $z_0^f \in \mathbb{R}^{128}$ 是逐点特征向量，用于解码颜色和其他属性的变形。核心设计是基于距离的加权策略：计算每个面部高斯中心到 3D 面部标志点（排除眼部）的最小距离 $d$，据此分配表情潜码权重 $\alpha_\psi$ 和头部姿态权重 $\alpha_\theta = 1 - \alpha_\psi$。直觉上，靠近面部标志点的高斯应主要响应表情变化，远离标志点的高斯应主要响应头部运动。具体而言：

\[
\alpha_\psi = \begin{cases} 1 & d < d_1 \\ (d_2 - d)/(d_2 - d_1) & d_1 \leq d \leq d_2 \\ 0 & d > d_2 \end{cases}, \quad d_1 = 0.15,\; d_2 = 0.25
\]

即距离阈值 $d_1 = 0.15$ 以内完全由表情驱动，$d_2 = 0.25$ 以外完全由姿态驱动，中间区域平滑过渡。变形公式为：

\[
\mu^f = \mu_0^f + \alpha_\psi E_\mu^f(z_0^f, \psi) + \alpha_\theta P_\mu^f(z_0^f, \theta), \quad c^f = \alpha_\psi E_c^f(z_0^f, \psi) + \alpha_\theta P_c^f(z_0^f, \theta)
\]

其中 $E^f$ 系列 MLP 处理表情条件，$P^f$ 系列 MLP 处理姿态条件。

\textbf{实现原理——眼球旋转分支}。眼球分支的关键创新是显式刚性旋转。眼球高斯的缩放被约束为球形 $S_0^e \in \mathbb{R}^{N \times 1}$（单一标量），与眼球的旋转物理特性一致。由于注视标签存在噪声，直接用归一化注视方向 $g$ 旋转高斯会导致数值优化错误，因此 GazeGaussian 训练两个独立的 MLP 组来预测旋转偏置：

\[
\mu^e = E_\mu^e(z_0^e, \psi) + G_\mu^e(z_0^e, g) \cdot \mu_0^e
\]

这里 $G_\mu^e(z_0^e, g) \cdot \mu_0^e$ 表示以注视方向为条件的旋转操作——将规范空间的眼球高斯位置 $\mu_0^e$ 按照目标注视方向进行旋转，$G_\mu^e$ 学习修正噪声标签带来的旋转偏差。$E_\mu^e(z_0^e, \psi)$ 则捕获表情引起的微小眼球运动（如眯眼时眼球的轻微位移）。其他属性（颜色、旋转、缩放、不透明度）通过类似的 MLP 结构变形。

\textbf{表情引导神经渲染器}。两个分支的高斯分别被独立光栅化为面部特征图 $M_f = \mathcal{R}(\{\mu^f, c^f, R^f, S^f, o^f\})$ 和眼球特征图 $M_e = \mathcal{R}(\{\mu^e, c^e, R^e, S^e, o^e\})$。与此同时，将两组高斯拼接在一起形成 $\{\mu^h, c^h, R^h, S^h, o^h\} = \{\mu^f, c^f, R^f, S^f, o^f\} \cup \{\mu^e, c^e, R^e, S^e, o^e\}$，统一光栅化为完整的头部特征图 $M_h$。三条特征图 $M_f, M_e, M_h$ 共同输入 UNet 风格的神经渲染器 $U$，解码为最终的面部图像 $I_f$、眼球图像 $I_e$ 和头部图像 $I_h$。为实现跨主体泛化并注入主体特定信息，GazeGaussian 在渲染器的瓶颈层引入了切片交叉注意力（slice cross-attention）模块：

\[
z'_b = z_b + \mathrm{CrossAttn}(q=\psi, k=z_b, v=z_b)
\]

其中 $z_b$ 是 UNet 编码器的瓶颈特征，而表情潜码 $\psi$ 作为查询向量从瓶颈特征中检索主体特定的面部细节，引导渲染器合成具有正确身份特征的面部图像。这解决了标准渲染器只做图像后处理、丢失主体特定细节的问题。

\textbf{设计动机}。眼球与面部的运动物理本质不同：面部是非刚性变形（肌肉驱动），眼球是刚性旋转（绕固定轴），因此必须用不同的参数化方式建模。将两者分离不仅提高了注视精度，还使得眼球高斯可以用球形缩放约束来保持形状一致性。距离加权策略避免了"一刀切"的变形方式对面部不同区域的干扰。表情引导渲染器则弥补了纯 3DGS 渲染在跨主体泛化上的不足——同一组高斯变形 MLP 对不同人可能产生不同的颜色映射，交叉注意力通过注入主体信息解决了这一问题。

\textbf{公式与参数}。训练损失 $L = \lambda_I L_I + \lambda_G L_G$，其中图像合成损失 $L_I$ 包含面部、眼部和头部的渲染图像及特征图共六项 L1 + SSIM + LPIPS 损失之和，注视重定向损失 $L_G$ 使用预训练注视估计器计算渲染图像与真值图像之间的角度误差 $E_{\text{ang}}$。损失权重 $\lambda_I = 1.0$，$\lambda_G = 0.1$，SSIM 和 LPIPS 权重均为 $\lambda_{\text{SSIM}} = \lambda_{\text{VGG}} = 0.1$。

\paragraph{Roll Your Eyes}

Roll Your Eyes（2025.08）~\cite{2508_06136} 继承了 GazeGaussian 的双流分离架构，但沿三个维度做出了关键改进。

\textbf{锚点回顾}。GazeGaussian 用自由高斯直接表示眼球，通过注视方向 $g$ 旋转这些高斯实现注视重定向——但眼球高斯的初始化没有显式的 3D 几何约束，球形缩放 $S_0^e \in \mathbb{R}^{N \times 1}$ 仅是一个正则化手段，并非真正的眼球形状；同时，面部分支的表情-姿态双权重机制 $\alpha_\psi / \alpha_\theta$ 不包含注视条件，意味着注视变化不会驱动眼周面部高斯变形。与 GazeGaussian 的双流分离策略不同，Roll Your Eyes 将改进聚焦在三个具体问题上：眼球结构的几何精度、注视方向的真实性、以及眼周肌肉的联动建模。

\textbf{改了什么}。(1)~\textbf{显式 3D 眼球结构}：取代 GazeGaussian 的自由高斯眼球表示，从 3DGazeNet 引入专用的 3D 眼球网格——该网格在虹膜区域分配更高顶点密度（与角膜凸起一致），而非简单球体。高斯直接在该网格顶点上初始化，使眼球形状具有解剖学约束。Roll Your Eyes 发现，GHA 的动态生成器无法学习正确的眼球位置，因此引入 eyeoffset MLP $D_e$ 预测眼球中心的解剖位移 $\Delta P_{\text{eyeball}} = D_e(\text{exp}, \text{pose})$，根据表情和头部姿态调整眼球在眼眶中的位置。位移预测从几何引导初始化阶段开始训练，并在化身训练阶段进一步精修。(2)~\textbf{注视引导变形场}：GazeGaussian 的面部变形仅由表情 $\alpha_\psi$ 和姿态 $\alpha_\theta$ 驱动，不包含注视条件。Roll Your Eyes 观察到注视方向改变时眼周肌肉会微妙运动（如视线偏移时眼睑的轻微张合），因此引入第三个影响因子 $\alpha_{\text{gaze}}$——与 GazeGaussian 的 $\alpha_\psi$ 距离衰减机制类似，$\alpha_{\text{gaze}}$ 基于每个高斯到 3D 眼部标志点的最小距离 $d_e$ 计算：

\[
\alpha_{\text{gaze}} = \begin{cases} 1 & d_e < t_3 \\ (t_4 - d_e)/(t_4 - t_3) & t_3 \leq d_e \leq t_4 \\ 0 & d_e > t_4 \end{cases}, \quad t_3 = t_1 \times 0.5 = 0.075,\; t_4 = t_2 \times 0.5 = 0.125
\]

其中 $t_3, t_4$ 取 GazeGaussian 阈值 $t_1, t_2$ 的一半，使注视影响的范围比表情更窄——仅影响紧邻眼球的区域。这使得靠近眼部标志点的高斯被注视方向 $g$ 调制，远离眼球的区域则不受注视影响。(3)~\textbf{伪真值注视方向}：多数注视数据集仅提供焦点坐标而非真实眼球方向，受 kappa 角（视觉轴与光轴的夹角）影响，直接用焦点坐标旋转眼球会产生方向偏差。Roll Your Eyes 改用 3DGazeNet 输出的方向性注视作为伪真值，使模拟眼球方向与解剖学真实方向对齐——这也使方法能在无注视标注的数据集上训练。

\textbf{设计动机}。三项改进对应三个观察：(a)~GHA 的动态生成器无法学习正确眼球位置，必须用专门的 MLP 显式建模眼球偏移；(b)~注视变化伴随眼周肌肉运动，忽略这些运动会导致渲染的眼部缺乏真实感；(c)~数据集注视标注是焦点坐标而非方向向量，直接使用会因 kappa 角引入系统性偏差。三者的共同目标是让双流架构中的眼球流不再只是"旋转高斯"，而是具有解剖学精度的可驱动 3D 结构。

\textbf{具体差异}。与 GazeGaussian 的三路渲染+UNet 神经渲染器（$M_f, M_e, M_h \to U$）不同，Roll Your Eyes 的渲染管线是双流各自光栅化后拼接，再通过超分辨率模块生成最终图像。训练损失引入了三个 GazeGaussian 没有的专用项：低分辨率眼球损失 $L_{\text{eyeLR}}$ 和高分辨率眼球损失 $L_{\text{eyeHR}}$（用眼部掩码 $M_e$ 约束眼球流只在眼球像素区域渲染），以及面部空白损失 $L_{\text{blank}}$（惩罚面部分支在眼部区域的渲染输出，将面部高斯推离眼眶，避免双流互相干扰）。损失权重：$\lambda_{\text{lr}} = 0.5$，$\lambda_{\text{hr}} = 0.5$，$\lambda_{\text{blank}} = 1.0$。

\subsubsection{人脸分析与人脸编辑}

本路线介绍利用~3DGS~可微分渲染进行人脸姿态估计和微表情光流数据生成等分析任务。

3DGS 也被应用于人脸与人体分析任务。与前面侧重"重建与渲染"的方向不同，本节关注如何利用 3DGS 的可微分渲染能力来辅助分析任务。

\textbf{Baseline 与问题}。多视角 3D 人体姿态估计的主流方法依赖学习式融合策略，即在标注数据集上训练跨视角特征聚合网络。这种范式有两个根本局限：其一，融合策略与特定相机配置、姿态分布和遮挡模式耦合紧密，跨场景部署性能急剧下降（例如在 Human3.6M 上训练的方法直接迁移到 CMU Panoptic 时误差增幅可达数倍）；其二，对遮挡的鲁棒性完全取决于训练数据的覆盖范围，遇到未见过的遮挡模式便无能为力。因此，理想的多视角融合应当是场景无关的：无需重训练即可适应任意相机配置和遮挡条件。

\paragraph{SkelSplat}

SkelSplat（2025.11）~\cite{2511_08294} 从一个根本性的视角转换出发：不做外观重建，而是用 3D 高斯直接表示关节位置——将多视角姿态估计重新定义为关节高斯的可微分渲染优化问题。

\textbf{改了什么}。标准 3DGS 用 RGB 通道渲染密集场景外观，高斯的颜色属性 $c_j$ 通过球谐函数编码视角相关辐射。SkelSplat 对渲染管线做了三处关键改造：(1)~将颜色属性替换为 one-hot 关节身份编码 $c_j[k] = \mathbf{1}_{k=j}$（第 $j$ 个关节仅在自身通道激活），渲染函数从 3 通道 RGB 扩展为 $N$ 通道张量（$N$ 为关节数），使每个关节拥有独立的渲染-监督通道，即使多个关节在图像平面上重叠也能独立优化；(2)~用 2D 关键点检测生成的伪真值热图替代 RGB 图像作为监督信号，无需 3D 真值标注；(3)~跨视角梯度累积替代逐视角更新——将所有视角的渲染损失梯度累加后统一更新高斯参数，而非标准 3DGS 的逐视角独立更新，提高了多视角一致性。

\textbf{实现原理}。SkelSplat 的优化管线包含三个阶段。

(1)~\textbf{骨骼高斯构建}。给定初始 3D 姿态估计 $\hat{\mathrm{SK}} = \{\hat{sk}_0, \dots, \hat{sk}_N\}$（可通过三角化或单目估计获得），为每个关节 $sk_j$ 构建一个各向异性 3D 高斯 $gs_j = (\mu_j, \Sigma_j)$，其中 $\mu_j \in \mathbb{R}^3$ 是关节位置，协方差初始化为 $\Sigma_j = 3 \cdot \mathbf{I}_3$。整个骨架表示为高斯集合 $\mathrm{GS} = \{gs_0, \dots, gs_N\}$。

(2)~\textbf{伪真值热图生成}。对每个相机视角 $i$，用预训练 2D 姿态估计器检测关节的 2D 位置 $sk_{ij}^{2D}$，并以该位置为中心、将 3D 协方差 $\Sigma_j$ 通过透视投影映射为 2D 协方差来渲染热图：

\[
\Sigma_{ij}^{2D} = J_i W_i \Sigma_j W_i^T J_i^T,
\]

其中 $W_i \in \mathbb{R}^{4 \times 4}$ 是相机外参矩阵，$J_i \in \mathbb{R}^{2 \times 3}$ 是透视投影在关节相机坐标处的雅可比矩阵。这种基于透视投影的协方差映射确保了近处关节的热图更大、远处更小，几何一致地反映 3D 空间中的不确定性。

(3)~\textbf{可微分渲染优化}。将骨骼高斯通过改造后的 $N$ 通道渲染函数投影到所有相机视角，用带掩码的 L2 损失驱动优化：

\[
\mathcal{L}_{\text{render}} = \sum_{j=1}^{N} \sum_{i=1}^{M} \| \mathcal{M}_{ij} \odot (I_{ij}^{\text{render}} - I_{ij}^{\text{pseudo}}) \|_2^2,
\]

其中 $\mathcal{M}_{ij}$ 是掩码（仅选取热图非零像素，避免大量空白背景影响收敛）。对称肢体长度正则项 $\mathcal{L}_{\text{sym}}$ 约束左右对称肢体（如下臂、下腿）长度一致，以恢复解剖学上合理的姿态结构。优化采用 Adam，最多 125 次迭代，并在连续两个窗口内最小损失差低于 $10^{-6}$ 时提前停止。与标准 3DGS 不同，这里禁用了高斯增删（致密化与剪枝），因为关节数量固定，不允许改变 $\mathrm{GS}$ 的基数。

\textbf{设计动机}。核心洞察在于：高斯泼溅本质上是一种可微分的 3D 到 2D 投影机制，投影的对象不必是颜色，也可以是关节身份。one-hot 编码使每个关节拥有独立的渲染通道，解决了标准 3DGS 在 RGB 通道中无法区分不同关节的根本问题——消融实验表明这一改动将 MPJPE 降低了 10\%。跨视角梯度累积则利用了多视角互补性：某个视角中因遮挡而缺失的关节信息，可由其他可见视角的梯度补偿。伪真值热图的设计使得整个系统完全不需要 3D 标注，仅依赖现成的 2D 关键点检测器即可运行，实现了场景无关的即插即用部署。

\textbf{公式与参数}。总优化损失为：

\[
\mathcal{L}_{\text{opt}} = \mathcal{L}_{\text{render}} + \lambda_{\text{sym}} \mathcal{L}_{\text{sym}},
\]

其中对称损失权重 $\lambda_{\text{sym}} = 10^{-5}$。SkelSplat 不需要训练，仅需优化：在 Human3.6M 上无需 3D 真值监督即超越所有非学习方法，跨数据集迁移时将误差降低高达 47.8\%；在遮挡基准 Human3.6M-Occ 上达到最优性能，且无需场景特定的微调。对易遮挡关节（肘、腕、膝、踝），将其初始协方差放大 1.25 倍以提供更大的优化搜索范围。推理时单帧 4 视角估计约需 3 秒，主要瓶颈在于每个关节独立渲染一个通道。

\paragraph{SRFlow}

SRFlow（2026.01）~\cite{2601_06479} 从另一个方向利用 3DGS 的分析能力——不做化身重建，而是用 3DGS 的可微分渲染生成高分辨率面部光流训练数据。

\textbf{Baseline 与问题}。微表情识别受限于光流训练数据的匮乏——传统光流数据集（FlyingChairs、Sintel 等）由合成场景或刚体运动构成，缺乏面部特有的微妙、低振幅、空间局部化的非刚性变形。手工标注面部光流几乎不可能（微表情的像素位移常小于 1 个像素），而现有多视角人脸数据集（如 NeRSemble）虽有多视角视频，却未提供帧间光流标注。

\textbf{改了什么}。SRFlow 利用 GaussianAvatar 的动态 3D 高斯表示作为"光流生成器"——从 16 视角视频中重建参数化头部化身后，引入 Flow Rasterizer 直接将高斯中心的帧间位移光栅化为像素级光流，无需 2D 插值或光流估计。旧机制（从 2D 图像对估计光流，受纹理缺失/遮挡/光照变化干扰）被替换为新机制（从 3D 高斯的已知 3D 运动投影光流，几何精确且稠密）。在此基础上进一步提出 SRFlowNet，基于 SKFlow 骨干网络并设计了四种面向高分辨率面部光流的正则化损失。

\textbf{实现原理}。SRFlow 的数据生成管线包含三个阶段。

(1)~\textbf{动态高斯化身重建}。基于 NeRSemble 数据集（222 个主体、16 个同步相机、4700+ 序列），选取 27 个主体的多样化运动序列。先用 VHAP（基于 FLAME 的参数化先验）对多视角视频进行网格对齐，再用 GaussianAvatar 重建每个主体的动态高斯化身——每个化身由绑定到 FLAME 网格三角形的 3D 高斯集合表示，通过可微分渲染实现精细的面部几何与外观重建。

(2)~\textbf{Flow Rasterizer 光流渲染}。对相邻帧 $t$ 和 $t+1$，每个高斯中心 $P(X_w, Y_w, Z_w)$ 经外参矩阵和透视投影映射到像素坐标 $(u_t, v_t)$ 和 $(u_{t+1}, v_{t+1})$。Flow Rasterizer 的核心是利用 3DGS 的 $\alpha$-混合 机制将帧间位移累积为逐像素光流：

\[
O_{\text{optical}} = \begin{bmatrix} \Delta U \\ \Delta V \end{bmatrix}^T = \sum_{i=1}^{n} \begin{bmatrix} \Delta u_i \\ \Delta v_i \end{bmatrix}^T \alpha_i' \prod_{j=1}^{i-1}(1 - \alpha_j'),
\]

其中 $(\Delta u_i, \Delta v_i)$ 是第 $i$ 个高斯在两帧间的投影位移。这种设计直接在 3D 高斯域中计算运动并投影，避免传统 2D 插值带来的误差，确保几何一致性和空间稠密性。

(3)~\textbf{运动增强与数据集构建}。为增加运动多样性，通过修改外参矩阵随机旋转相机和主体轴，使每段序列包含正面视角和变化视角。最终数据集包含 11{,}161 对高分辨率光流标注（分辨率 $2200 \times 3208$，来自 27 个身份的多样表情序列），划分为 6791 训练对、1212 验证对和 3158 测试对。相比此前唯一的面部光流数据集 FacialFlowNet（$512 \times 512$），SRFlow 的分辨率提升约 27 倍，并包含基于体积表示的头发和精细面部细节。SRFlowNet 在此数据集上训练，引入四种正则化损失抑制高频伪影、保留微妙面部运动：全变差正则化（TVR，归一化 Sobel 算子约束光流空间平滑性）、均值图像梯度激活正则化（MIGAR，自适应指数加权的逐像素约束）、图像梯度方差激活正则化（IGVAR，用有效区域的梯度方差替代全图均值作为指数基底）等。

\textbf{设计动机}。3DGS 的关键优势在于每个高斯都有明确的 3D 位置——帧间位置差即为精确的 3D 运动向量，投影到图像平面即为稠密光流真值。相比传统立体匹配或多视角光流方法，这种基于泼溅光栅化的方案无需匹配或搜索，且天然保持了多视角几何一致性。

\textbf{公式与参数}。SRFlowNet 以 SKFlow 为骨干，正则化权重 $\lambda_N = 0.05$，衰减因子 $\beta \in (0,1)$ 逐步降低早期预测阶段的权重。在 SAMM、CASME II、SMIC 三个微表情数据集上，使用 SRFlow 数据训练的模型将 F1 分数从 0.4733 提升至 0.6947（提升 48\%），端点误差降低 42\%。

在人脸编辑方面，3DGS 的显式表示为高斯级别的操作提供了天然的接口，使面部部件的分割、交换和编辑成为可能。

\paragraph{FaceParts}

FaceParts（2026.05）~\cite{2605_13853} 提出了无监督的 3D 高斯化身分割与编辑框架，直接在高斯域中操作（而非图像域），将化身分解为语义一致的面部部件（胡须、眉毛、眼睛等），支持跨化身的部件交换——例如将一个人的胡须高斯替换到另一个人的化身上。

\paragraph{GaussianSwap}

GaussianSwap（2026.01）~\cite{2601_05511} 将 3DGS 与人脸交换结合，通过 FLAME 模型绑定高斯泼溅体来实现动态面部控制，利用三个人脸识别模型（ArcFace、CosFace、AdaFace）的复合身份嵌入确保在交换源面部身份的同时保持目标的面部运动。

\paragraph{GRMM}

GRMM（2025.09）~\cite{2509_02141}（Gaussian Residual Morphable Model）是首个支持全头部生成与实时渲染的高斯 3D 变形模型。

\textbf{Baseline 与问题}。传统 3D 变形模型（3DMM）以网格为基础，参数化地表示身份和表情变化，但其几何分辨率受限于网格密度，无法恢复皱纹、胡茬、发际线等高频细节。GaussianHeads 等基于 3DGS 的方法虽然能以高斯基元表示精细外观，但它们是针对单个主体设计的重建方法，无法跨身份泛化——为每个人建模都需要从头训练，失去了变形模型"一次训练、无限生成"的核心优势。因此，核心挑战在于：如何将 3DGS 的表示能力与 3DMM 的跨身份泛化能力结合起来。

\textbf{改了什么}。GRMM 不像传统 3DMM 那样在网格顶点上直接参数化外观，也不像 GaussianHeads 那样让高斯属性完全由网络预测，而是采用"粗到细"的残余策略：先用 FLAME 网格提供粗几何（粗），再用残余潜码驱动的高效 CNN 解码器预测高斯属性的精细偏移量（细）。旧机制（GaussianHeads 的 $D_{GH}(z_{exp}, v) \to (\Delta_r, \Delta_p, \Delta_s, o, c)$，仅以表情码为输入）$\to$ 新机制（GRMM 的多解码器 $D_{\Theta} := \{D_v, D_T, D_o, D_c\}$，以残余身份码 $z_{id}$、残余表情码 $z_{exp}$、姿态和视角为联合输入，分别预测网格变形、高斯变换、不透明度和颜色）。

\textbf{实现原理}。GRMM 的架构由四条解码管线和一层屏幕空间精修网络组成，形成一条从身份/表情编码到高分辨率渲染图像的完整生成管线。

(1)~\textbf{残余潜码与表情对齐}。每个主体用可学习的残余身份码 $z_{id} \in \mathbb{R}^{512}$ 编码，每个表情用全局残余表情码 $z_{exp} \in \mathbb{R}^{256}$ 编码，且同一表情在所有身份间共享同一个 $z_{exp}$——这种设计借助 GRMM 新发布的 EXPRESS-50 数据集（50 个身份、每人 60 个手工对齐的高强度表情）实现身份与表情残余的清晰解耦。

(2)~\textbf{网格解码器} $D_v$。一个共享 MLP 先融合身份与姿态信息：

\[
f_{base} = \phi_{shared}([z_{id}, \theta_{neck}, \theta_{jaw}]),
\]

身份位移由 $v_{\Delta,id} = \phi_{mesh,id}(f_{base})$ 预测。表情调制采用 FiLM（Feature-wise Linear Modulation）机制，将 FLAME 表情系数 $\beta_{exp}$ 映射为逐通道缩放 $\gamma$ 和偏移 $\beta$：

\[
(\gamma, \beta) = \phi_{FiLM}(\beta_{exp}), \quad \tilde{f}_{exp} = \gamma \odot [f_{base}, z_{exp}] + \beta,
\]

表情位移由 $v_{\Delta,exp} = \phi_{mesh,exp}(\tilde{f}_{exp})$ 预测。最终变形网格为 $v_d = v_{rec} + v_{\Delta,id} + M_{face} \cdot v_{\Delta,exp}$，其中 $M_{face}$ 掩蔽牙齿顶点以防止表情偏移影响齿列。

(3)~\textbf{高斯属性解码}。在变形网格的 UV 空间中以 $N_g \times N_g$（$N_g = 512$）分辨率均匀采样 $N_{prim} = 262{,}144$ 个高斯基元，初始位置通过重心插值确定。三个高效 CNN 解码器分别预测：
\begin{itemize}
\item 变换解码器 $D_T$：输出 $N_g \times N_g \times 10$ 的偏移图（位置 $\Delta_p$、旋转 $\Delta_r$、缩放 $\Delta_s$），最终位置 $p = p_{in} + \Delta_p$；
\item 不透明度解码器 $D_o$：输出 $N_g \times N_g \times 1$，仅以 $z_{id}$ 为条件；
\item 外观解码器 $D_c$：输出 $N_g \times N_g \times 35$，包含 RGB 颜色和 32 维学习特征向量 $f_i$，以 $z_{id}$ 和视角 $\omega$ 为条件。
\end{itemize}

(4)~\textbf{屏幕空间精修网络} $N_{ref}$。光栅化得到 RGB 图像 $I_{rgb}$、特征图 $I_{feature}$ 和深度图 $I_{depth}$ 后，CNN 精修网络将其融合为最终图像 $I = N_{ref}([I_{rgb}, I_{feature}, I_{depth}^{norm}])$，分辨率保持 1K 不变。该网络不改变底层 3D 表示，仅增强难以被高斯模型捕获的外观先验（如高频纹理）。

\textbf{设计动机}。残余参数化的核心思想是"站在巨人的肩膀上"：FLAME 已提供可靠的粗几何（面部形状、下巴旋转、颈部姿态），学习残余量比从零学习绝对量更容易收敛且泛化更好。将身份和表情残余分离为独立潜码（$z_{id}$ 与 $z_{exp}$），配合表情对齐数据集，使模型在推理时可以自由组合未见过的身份和表情——这正是变形模型区别于单主体重建方法的核心能力。UV 空间的 CNN 解码器设计兼顾了效率（卷积运算的并行性）和局部性（相邻高斯共享相似的偏移模式），使得单次前向传播即可解码全部 26 万个高斯属性，实现 75 FPS 的实时渲染。

\textbf{公式与参数}。端到端训练损失为：

\[
L = \lambda_{rec}\,L_{rec}(I, I^*) + \lambda_{rec}\,L_{rec}(I_{rgb}, I^*) + \lambda_{depth}\,L_{depth} + L_{reg},
\]

其中 $L_{rec} = \lambda_{l1} L_{l1} + \lambda_{ssim} L_{ssim} + \lambda_{perc} L_{perc}$ 包含 L1、SSIM 和 VGG 感知损失，$L_{depth}$ 是 L2 深度损失，正则项 $L_{reg} = \lambda_s L_s + \lambda_z L_z + \lambda_{lap} L_{lap}$ 分别约束高斯尺度范围、潜码范数和网格拉普拉斯平滑。具体权重为 $\lambda_{l1} = 0.8$、$\lambda_{ssim} = 0.2$、$\lambda_{perc} = 0.04$、$\lambda_z = 0.001$、$\lambda_{lap} = 0.01$、$\lambda_s = 0.1$、$\lambda_{depth} = 0.5$。推理时采用两阶段适配：第一阶段固定解码器优化潜码 $(z_{id}, z_{exp})$，第二阶段固定潜码精调解码器（PTI 策略，局部性正则 $\lambda_{loc} = 0.1$），在保持先验的同时适配新身份。GRMM 以 75 FPS 实时渲染 1K 分辨率的全身头部图像。

\paragraph{Mind-to-Face}

Mind-to-Face（2025.12）~\cite{2512_04313} 将 3DGS 化身的驱动信号源从图像/音频扩展到脑电信号（EEG）——化身本身的 3D 高斯表示和可微分光栅化渲染管线不变，改变的是从信号到面部几何的映射方式。

\textbf{Baseline 与问题}：前述所有人脸化身方法（ScaffoldAvatar、GGHead、GRMM 等）均从视觉数据（图像或视频）学习化身，驱动信号也来自视觉或音频模态。然而，在脑机接口（BCI，Brain-Computer Interface，一种在大脑与外部设备之间建立直接通信路径的技术）场景中，用户可能无法产生可见的面部运动（如面瘫患者），但仍能通过大脑活动"意图"表达情绪。EEG 信号是否包含足够的面部运动信息来驱动高保真化身，是一个未被探索的问题。

\textbf{改了什么}：Mind-to-Face 设计了从 EEG 信号到 3D 面部几何的端到端映射管线，最终通过修改的 GaussianAvatars 渲染照片级化身。

旧机制（GaussianAvatars 从多视角视频重建化身，驱动信号为 FLAME 参数）被替换为新机制（EEG 信号 $\to$ CNN-Transformer 编码器 $\to$ Stable Diffusion 解码器 $\to$ 3D 位置图 $\to$ GaussianAvatars 渲染）。核心创新在于将 EEG 到面部的映射问题转化为图像生成问题，用 Stable Diffusion 解码器生成 UV 空间的 3D 位置图。

\textbf{实现原理}：管线分为三个阶段。\textit{信号采集与预处理}：16 通道 EEG（125Hz 采样，OpenBCI Cyton-Daisy）与 16 个同步相机（最高 8K，120fps）同步采集。EEG 信号经 6 阶 Butterworth 带通滤波（4--40Hz）和逐通道 z-score 标准化后，切分为重叠时间窗口 $E^{(t)} \in \mathbb{R}^{W \times C}$（$W=375$，$C=16$）。

\textit{EEG 到位置图解码}：每个窗口先经过两层 CNN（第一层沿时序轴建模短期模式，第二层沿通道轴捕获电极间空间关联），经 batch normalization + ELU + 平均池化降冗余后，通过 Transformer 编码器捕获长程时序依赖，输出紧凑潜嵌入 $z^{(t)}$。然后，初始化自 Stable Diffusion 2.1 预训练权重的解码器将 $z^{(t)}$ 映射为密集 3D 位置图 $\hat{P}^{(t)} \in \mathbb{R}^{256 \times 256 \times 3}$（每个像素编码 UV 空间中的 $(x, y, z)$ 面部坐标），二元掩码 $m$ 将监督限制在内面部区域。位置图损失为 $L_{pm} = \lambda_{rec} L_{rec} + \lambda_{smooth} L_{smooth}$，其中 $L_{rec}$ 为逐像素 MSE，$L_{smooth}$ 为自监督拉普拉斯平滑正则。\textit{化身渲染}：预测的位置图通过 UV 布局采样为顶点坐标形成网格，输入修改的 GaussianAvatars 管线——每个三角面片绑定一个高斯，经刚性变换后可微分渲染。与原版 GaussianAvatars 不同的是，Mind-to-Face 不从 FLAME 追踪获取面部几何，而是直接使用多视角光度测量重建的几何（提供更大灵活性以适配极端表情）。

\textbf{设计动机}：将 EEG 解码目标设为 3D 位置图而非 2D 图像或 FLAME 参数，是因为位置图在保留几何连续性的同时具有图像般的规则结构。这让 Stable Diffusion 的 2D 卷积先验得以复用，无需设计全新的 3D 解码架构。使用预训练的 Stable Diffusion 解码器而非从零训练，弥补了 EEG-面部配对数据极为稀缺（每人约 25 分钟录制）带来的训练信号不足。

\textbf{公式与参数}：EEG 窗口大小 375（3 秒@125Hz），16 通道。位置图分辨率 $256 \times 256$。GaussianAvatars 训练超参与原论文一致，渲染损失 $L = (1-\lambda)L_1 + \lambda L_{\text{SSIM}} + \lambda_\epsilon L_\epsilon + \lambda_\sigma L_\sigma$，其中 $L_\epsilon$ 和 $L_\sigma$ 分别约束高斯位置和尺度。该方法证明了 EEG 信号中确实编码了可用于驱动高保真 3DGS 面部化身的面部运动信息，开辟了 3DGS 在脑机接口应用中的新方向。

\begin{takeaway}
人脸重建与动画是 3DGS 化身技术最成熟的领域：FLAME 参数模型与高斯基元的绑定实现表情驱动，神经纹理和潜码解耦身份与表情，而 Mind-to-Face 等工作将驱动信号源从视觉扩展到脑电，开辟了脑机接口新方向。
\end{takeaway}

%=====================================================================

\subsection{近期方法脉络：头部Avatar重建、人脸头像生成、Talking head与表情动画等}

这一组近期工作延续 3DGS avatar、3DGS、3DGS talking head，主要把改动集中在扩展、应用、编辑等环节。下面只保留相对既有路线的变化，便于和前文方法直接对照。

\paragraph{SpatialAvatar-0}
SpatialAvatar-0（2026.06）~\cite{2606_15659} 延续 3DGS avatar 在“头部Avatar重建”上的基本设定，主要改变是：SpatialAvatar-0：多阶段4D head avatar重建，先生成可泛化高斯布局再进行保持布局的精修。高斯基元仍承担可渲染动态场景表示。
其设计动机是：提出多阶段重建流程，缓解前馈模型数据集偏置和精修阶段破坏上游高斯布局的问题。
具体而言，方法以3DGS头部表示为核心，先从少量源图得到可泛化初始化，再通过后续精修提升个体质量，同时避免自适应致密化完全打散初始布局。
适用边界是：仍依赖人脸/头部先验；极端表情、遮挡和发丝细节可能困难。

\paragraph{Splatshot}
Splatshot（2026.05）~\cite{2606_01493} 延续 3DGS 在“人脸头像生成”上的基本设定，主要改变是：SplatShot：训练自由框架，在扩散去噪过程中反复重拟合3DGS并用3D反馈约束多视角一致性。高斯基元仍承担可渲染动态场景表示。
其设计动机是：提出每个去噪步的3D反馈循环：从多视角预测图像重拟合3DGS，再把光度差异反传到扩散过程，以结合2D先验和显式3D一致性。
具体而言，给定基础3DGS人脸模型和参考照片，方法联合去噪多个目标视角；每步估计干净图像、拟合3DGS、计算跨视角3D反馈，从而减少扩散多视角不一致。
适用边界是：依赖基础3DGS人脸模型和扩散先验；极端表情、遮挡或非人脸输入可能失败；训练自由但迭代重拟合成本较高。

\paragraph{EmoZone-Talker}
EmoZone-Talker（2026.06）~\cite{2606_15848} 延续 3DGS talking head 在“Talking head与表情动画”上的基本设定，主要改变是：EmoZone-Talker：用Facial Action Units进行区域语义控制，协调音频动态和显式表情信号。高斯基元仍承担可渲染动态场景表示。
其设计动机是：提出显式空间区域和时间协调机制，缓解语音驱动和表情控制信号之间的空间纠缠与时间不稳定。
具体而言，方法以AU为可解释控制变量，将面部区域划分后分别协调音频和表情条件，使3DGS talking head支持局部可编辑表情。
适用边界是：依赖AU标注或估计质量；夸张表情和跨身份泛化仍具挑战。

\paragraph{GaussianEmoTalker}
GaussianEmoTalker（2026.07）~\cite{2607_00959} 延续 3DGS avatar 在“Emotional Gaussian talking heads”上的基本设定，主要改变是：GaussianBlendshapes 和 neutral-to-emotional residual deformation 的音频驱动 3DGS 框架。高斯基元仍承担可渲染动态场景表示。
其设计动机是：将情绪动画表述为从 neutral talking space 到 emotional avatar 的 residual deformation，兼顾唇同步、身份保真和情绪控制。
具体而言，先构建 identity-specific neutral talking space with GaussianBlendshapes，再由音频和情绪条件预测残差变形，实现高保真可控表情。
适用边界是：依赖身份特定中性空间和 blendshape 质量；夸张表情、跨语言和强头部运动仍可能困难。

\paragraph{Edit3DGS}
Edit3DGS（2026.06）~\cite{2606_17432} 延续 3DGS head avatar 在“动态头像编辑”上的基本设定，主要改变是：Edit3DGS：2D instruction-guided diffusion编辑视频帧，再聚合到3D Gaussian Splatting动态头部表示。高斯基元仍承担可渲染动态场景表示。
其设计动机是：提出统一框架，将图像域语义可控性和3DGS的时序一致表示结合，支持表情、属性和外观细化。
具体而言，给定输入视频，方法mask可编辑面部区域，用文本条件扩散生成编辑帧，再通过3DGS聚合形成时间一致的动态头部资产。
适用边界是：扩散编辑可能引入跨帧不一致；mask和文本指令质量影响最终3D结果。
\section{身体与运动}
\label{sec:16_2}

身体与运动的建模是人体化身中最为复杂的方向，需要处理多关节联动、衣物非刚性变形以及各部位异质外观。此外，铰接物体重建、人体-物体交互和人体新视角合成也为身体级 3DGS 重建提供了丰富的研究问题。本节从全身化身、铰接物体与交互、场景中的人体动画、人体新视角合成等角度展开讨论。如果说第~\ref{sec:16_1}~节聚焦的"头部"已经足够复杂，那么"全身"则带来了数量级的挑战：关节数从数十增长到上百，衣物类型从头发扩展到上衣、裤子、裙子和鞋子，交互对象从无到有。

\subsection{化身表示与动画驱动}

本节讨论~2~条技术路线：（1）基于网格-高斯混合表示的化身；（2）变形与动画驱动。

3DGS 应用于全身化身的开创性工作出现在 2023 年末。

\paragraph{DGS-Avatar}

DGS-Avatar（2023.12）~\cite{2312_09228} 首次展示了从单目视频使用 3DGS 创建可动画化身穿身的可行性。

\textbf{Baseline 与问题}：此前的 HumanNeRF 等隐式方法虽然能从单目视频重建人体，但其体渲染管线导致训练和推理速度极慢（每帧数秒），完全无法满足实时交互需求。

\textbf{改了什么}：3DGS-Avatar 将人体建模为规范空间中的 3D 高斯集合，通过可学习的非刚性变形网络将高斯从规范空间映射到观测空间，利用 3DGS 的可微分光栅化器实现实时渲染。

\textbf{实现原理}：核心架构包含三个模块。(1)~非刚性变形模块：分解为姿态依赖的非刚性变形 $F_{nr}$ 和刚性变换 $F_r$。$F_{nr}$ 使用多级哈希网格编码规范位置 $x_c$，与姿态潜编码 $Z_p$（由轻量级层次姿态编码器从 SMPL 姿态参数提取）拼接后送入浅层 MLP（2 个隐藏层，宽度 128），输出高斯位置、缩放、旋转的偏移量以及特征向量 $z$。变形规则为：$x_d = x_c + \Delta x$、$s_d = s_c \cdot \Delta s$、$q_d = q_c \cdot \Delta q$（四元数乘法等价于旋转矩阵相乘）。

(2)~刚性变换模块：学习一个蒙皮 MLP 从变形位置 $x_d$ 预测蒙皮权重，通过前向 LBS 变换高斯位置和旋转到观测空间。(3)~颜色 MLP：将变形网络输出的特征 $z$ 通过球谐基函数解码为视角依赖的颜色。初始化时在 SMPL 规范网格表面随机采样 $N=50k$ 个点作为高斯初始位置。

\textbf{设计动机}：将复杂的人体变形分解为非刚性（衣物褶皱等姿态依赖变形）和刚性（骨骼驱动的整体变换）两个层级，借鉴了 HumanNeRF 的设计思想但用 3DGS 替代 NeRF 实现了数量级的加速。

\textbf{公式与参数}：引入尽可能等距正则化（as-isometric-as-possible），约束规范空间中相邻高斯在变形后保持距离，提升对未见姿态的泛化。在 30 分钟内完成训练并以超过 50 FPS 的帧率渲染。

\paragraph{HUGS}

HUGS（2023.11）~\cite{2311_17910}（Human Gaussian Splats）从仅 50--100 帧的单目视频中自动分离静态场景和可动画人体化身。

\textbf{Baseline 与问题}：3DGS-Avatar 需要预拟合的 SMPL 参数作为输入，且人体和场景无法自动分离——限制了其在真实视频中的适用性。

\textbf{改了什么}：HUGS 实现了人体-场景的自动分离，将整个场景分解为人体高斯（由 SMPL 初始化并参数化）和场景高斯（由 COLMAP 点云初始化），两个子系统在统一框架中联合优化。

\textbf{实现原理}：利用 SMPL 体模型参数化人体高斯：每个高斯在 SMPL 规范空间中定义，通过学习到的 LBS 权重和 SMPL 骨架变换到观测空间。关键设计是联合优化 LBS 权重——预定义的 SMPL 蒙皮权重仅基于关节位置，无法区分同一关节影响范围内的衣物和皮肤，学习到的权重可以捕捉这种差异。场景高斯则遵循标准 3DGS 流程，通过 COLMAP 稀疏点云初始化。两套高斯共享同一光栅化器进行渲染，通过前景/背景分离损失和 SMPL 前景掩码实现自动分割。

\textbf{设计动机}：人体-场景分离是实现化身重用的关键——只有将人体从场景中分离出来，才能将化身放到新场景中。利用 SMPL 的参数化能力将高斯绑定到人体骨架上，使得动画驱动仅需改变姿态参数即可。

\paragraph{SplatArmor}

SplatArmor（2023.11）~\cite{2311_10812} 通过 SE(3) 场捕获姿态依赖的高斯位置和各向异性变化，并引入神经颜色场（neural color field）替代逐高斯独立颜色优化——该颜色场将规范空间中的 3D 坐标映射到颜色，隐式约束邻近高斯的颜色连续性，同时为高斯中心提供 3D 几何监督。

\subsubsection{基于网格-高斯混合表示的化身}

本路线介绍将参数化体模型网格拓扑与~3DGS~高保真渲染能力结合的混合表示方法。

上述先驱工作虽然验证了 3DGS 用于全身化身的可行性，但它们要么将高斯直接放在规范空间中自由优化（如 3DGS-Avatar），要么依赖 SMPL 骨架做松耦合（如 HUGS），在复杂姿态和宽松衣物场景下容易产生伪影。后续工作发现，将参数化体模型（SMPL、SMPL-X）的网格拓扑与 3DGS 的高保真渲染能力结合，可以兼顾几何约束与外观质量——这一思路成为全身化身最主流的架构。

\paragraph{GoMAvatar}

GoMAvatar（2024.04）~\cite{2404_07991} 提出了"网格上的高斯"（Gaussians-on-Mesh, GoM）混合表示。

\textbf{Baseline 与问题}：纯 3DGS 化身缺乏表面约束，高斯可以漂移到任意位置导致几何退化；纯网格化身虽然拓扑明确但渲染质量受限于纹理分辨率。

\textbf{改了什么}：GoM 将每个 3D 高斯附着到可变形网格的三角形面片上，结合高斯泼溅的渲染质量和网格的几何正则化能力。

\textbf{实现原理}：规范空间中的 GoM 表示为 $\Theta^c = \{v_{\Theta,i}^c\}_{i=1}^{V}, \{f_{\Theta,j}\}_{j=1}^{F}\}$，其中每个面片 $f_{\Theta,j} = (r_{\Theta,j}, s_{\Theta,j}, c_{\Theta,j}, \Delta_{j,k})_{k=1}^{3}$ 定义了局部高斯的旋转、缩放、颜色和三个顶点索引。关键设计是局部-世界坐标变换：每个面片的局部高斯 $N(0, \Sigma_j^{\text{local}})$ 通过仿射变换 $f(x) = A_j x + b_j$ 映射到世界坐标。均值 $b_j$ 设置为三角形重心；变换矩阵 $A_j$ 的前两列定义在 Steiner 椭圆的两个半轴方向上（Steiner 椭圆是三角形的最大内接椭圆），第三列沿面片法线方向，法线方向的缩放设为极小值（$\epsilon = 10^{-3}$），确保高斯呈薄片状贴合表面。当局部旋转为零、局部缩放为一时，椭球在三角形上的投影恰好恢复 Steiner 椭圆。

\textbf{设计动机}：通过将高斯"钉"在网格面片上，几何变形完全由网格的 LBS 驱动——无需额外学习高斯的运动，只需学习其外观属性。网格提供了拓扑一致性和运动正则化，高斯则提供了超越纹理映射的渲染质量。

\textbf{公式与参数}：实现了实时、内存高效的高质量可动画人体建模，每个化身仅需 3.63 MB 存储，渲染速度达 43 FPS。

\paragraph{HAHA}

HAHA（2024.04）~\cite{2404_01053} 学习在 SMPL-X 网格的哪些区域使用高斯泼溅（如头发和超出网格的衣物），通过三阶段管线实现：先学习完整高斯化身并精修 SMPL-X 参数，再从精修网格学习 RGB 纹理映射，最后合并两个化身并通过可微渲染优化不透明度，自动删除可被纹理替代的冗余高斯，将所需高斯数量减少到不到三分之一（从约 20K 降至约 6K）。

\paragraph{ExAvatar}

ExAvatar（2024.07）~\cite{2407_21686} 继承了 GoMAvatar 的"网格+高斯"混合思路，但在绑定策略和表情驱动两个维度做出了关键改进。

\textbf{Baseline 与问题}：GoMAvatar 将高斯绑定到面片中心，高斯的运动完全由面片的 LBS 变换决定，面片之间的相邻关系未被显式建模。这导致在新颖表情和姿态下，相邻高斯可能产生不合理的空间关系（如面部高斯在夸张表情时相互穿越）。此外，大多数全身化身方法不支持面部表情驱动——仅用身体姿态参数化变形空间。

\textbf{改了什么}：ExAvatar 将每个 3D 高斯视为 SMPL-X 网格表面上的顶点（而非面片中心的附属物），通过预定义的三角形面片连通性信息建立高斯之间的邻域关系。这一表示使得：(1)~连通性正则化器可约束相邻高斯在变形后保持合理的相对位置，减少新颖表情和姿态下的伪影；(2)~化身直接继承了 SMPL-X 的表情参数空间，可通过驱动面部表情系数生成训练视频中未见过的表情。

\textbf{实现原理}：训练从短视频（约 30 秒单目视频）进行。初始化时在每个 SMPL-X 三角形面片的中心和顶点处放置高斯。

优化阶段，位置和旋转偏移通过 MLP 从姿态和表情参数预测，连通性正则化器惩罚相邻高斯之间的距离偏离其在规范空间中的初始比例——即若两个高斯在规范空间中距离为 $d_0$，变形后距离 $d$ 应满足 $d \approx s \cdot d_0$（$s$ 为局部缩放因子），防止高斯在极端表情下相互穿越或产生空洞。表情驱动通过 SMPL-X 的表情基（约 10 维表情系数）直接参数化，无需额外的表情学习模块。

\textbf{设计动机}：将高斯视为顶点而非面片附属物的核心洞察是——顶点天然具有连通性信息（每个顶点连接若干三角形），而面片中心之间没有显式关系。这种结构化表示在几何约束上更强，但也更安全：即使在新颖表情下，连通性正则化器也能防止灾难性的几何退化。

\paragraph{DEGAS}

DEGAS（2024.08）~\cite{2408_10588} 首次在全身化身中整合丰富的面部表情，与 ExAvatar 直接利用 SMPL-X 表情空间不同，DEGAS 学习独立的表情生成模型。

\textbf{Baseline 与问题}：ExAvatar 依赖 SMPL-X 的固定表情基（约 10 维），其表达能力受限于参数化模型的线性假设——复杂表情（如皱眉+撇嘴的组合）难以被有限维的表情系数精确表示。此外，全身化身训练中的面部表情数据通常非常有限（短视频中表演者的表情多样性不足），限制了表情驱动的质量。

\textbf{改了什么}：DEGAS 学习条件变分自编码器（cVAE，Conditional Variational Auto-Encoder，一种结合条件信息和变分推断的生成模型），从身体运动和面部表情的联合信号生成 SMPL-X UV 布局中的高斯图——每个 UV 像素参数化一个 3D 高斯基元（位置、旋转、缩放、不透明度、球谐系数）。关键创新是使用仅在 2D 肖像图像上训练的表情潜在空间来驱动面部表情：cVAE 的编码器从 2D 肖像图像提取表情嵌入，解码器将该嵌入与身体运动信号一起输入，生成 UV 空间中的完整高斯图。

\textbf{设计动机}：2D 肖像数据（如 VoxCeleb 等大规模人脸视频数据集）远比全身多视角面部表情数据容易获取。通过在 2D 肖像上学习表情潜在空间，DEGAS 绕过了全身多视角表情数据的稀缺瓶颈，使表情驱动能力不再受限于训练视频中表演者的表情多样性。

\paragraph{SAGA}

SAGA（2024.12）~\cite{2412_00845}（Surface-Aligned Gaussian Avatar）针对上述 GoMAvatar 和 ExAvatar 的"严格绑定"策略做出反思。

\textbf{Baseline 与问题}：GoMAvatar 和 ExAvatar 将高斯严格绑定在网格面片或顶点上，提供了良好的几何约束，但也限制了表达力——高斯无法偏离网格表面来表示超出网格几何的细节（如蓬松的头发、宽松衣物的褶皱），导致渲染结果"过于平滑"。另一方面，纯自由 3DGS（如 3DGS-Avatar）在高斯脱离几何约束后容易在单目重建的严重病态条件下过拟合训练视图中的时变区域（衣物褶皱、阴影），导致几何噪声和新视角/新姿态泛化失败。

\textbf{改了什么}：SAGA 提出"先紧后松"的两阶段对齐策略。附着阶段（Adhered Stage）将高斯附着在网格上，继承几何约束；但与 GoMAvatar 的固定绑定不同，SAGA 允许高斯在网格表面"流动"（通过优化重心坐标实现位置更新），增加了灵活性。

分离阶段（Detached Stage）将高斯从网格表面释放，允许其偏离到网格之外以表达更多细节；但引入高斯-网格对齐正则化（最小化高斯与原始绑定三角面片之间的位置和朝向偏差），在释放表达力的同时维持几何对齐。

\textbf{实现原理}：附着阶段中，每个高斯通过重心坐标 $(b_1, b_2, b_3)$ 绑定到一个三角面片，优化过程中高斯可沿网格表面移动（"行走于网格"策略动态更新绑定关系——当高斯漂移到相邻面片时自动重新绑定）。分离阶段中，高斯位置和旋转被释放为独立优化变量，正则化项 $\mathcal{L}_{\text{align}} = \lambda_p \|\boldsymbol{\mu}_i - \hat{\boldsymbol{\mu}}_i\|^2 + \lambda_r \|\boldsymbol{q}_i - \hat{\boldsymbol{q}}_i\|^2$ 约束高斯位置 $\boldsymbol{\mu}_i$ 和旋转 $\boldsymbol{q}_i$ 不偏离其网格绑定点的投影位置 $\hat{\boldsymbol{\mu}}_i$ 和面片朝向 $\hat{\boldsymbol{q}}_i$ 过远。

\textbf{设计动机}：两阶段策略的直觉是"先学好走路再学跑步"——先用严格几何约束学习基本的化身形状和运动模式，确保几何正确；再在可靠的基础上释放表达力，让高斯自由优化外观细节。行走于网格策略解决了固定绑定的僵化问题：当高斯的视觉最优位置不在其初始绑定面片上时，它可以沿表面迁移到更合适的面片。

\paragraph{iHuman}

Human（2024.07）~\cite{2407_11174} 从另一个角度解决高斯与网格的对齐问题。

\textbf{Baseline 与问题}：GoMAvatar 和 ExAvatar 将高斯的几何属性（位置、旋转、缩放）完全由所在面片的几何决定，高斯无独立自由度。这意味着外观细节的表达完全受限于网格的分辨率和变形精度——当面片变形不准确时，高斯的渲染质量也会下降。

\textbf{改了什么}：iHuman 通过隐式 MLP 建模表面位移和颜色的球谐系数变化——3D 高斯仍绑定到体模板的三角形面上作为基础位置，但 MLP 为每个高斯预测额外的位置偏移和 SH 系数调整，使外观细节不完全受限于网格几何。辅助监督方面，iHuman 利用显式 2D 高斯（surfels）嵌入 3D 空间来编码法线信息，通过法线渲染的辅助损失约束几何质量。

\textbf{设计动机}：显式绑定提供稳定的几何基础，隐式偏移补充网格精度不足的细节——这种"显式锚定+隐式精修"的混合策略在训练速度和渲染质量之间取得了平衡，实现了秒级训练和快速的人体数字化。

\subsubsection{变形与动画驱动}

本路线介绍高斯化身的蒙皮变形、姿态驱动和运动优化方法。

精确的变形建模是实现高质量动画的关键。线性混合蒙皮（LBS，将顶点位置通过加权平均多个骨骼变换来得到变形后位置的方法）是最常用的变形方法，但其对高斯旋转等非线性属性的处理存在固有缺陷：旋转的线性混合不保证结果仍然是有效旋转。

\paragraph{On the Skinning of Gaussian Avatars}

On the Skinning of Gaussian Avatars（2025.09）~\cite{2509_11411} 深入分析了这一问题，提出基于四元数平均的完整高斯蒙皮方案，可同时正确处理位置、旋转和球谐系数三类属性的蒙皮变形，且无需修改高斯光栅化器本身。

\textbf{Baseline 与问题}。3DGS 化身的每个高斯包含五类属性：位置 $\boldsymbol{\mu} \in \mathbb{R}^3$、旋转 $\boldsymbol{q} \in \mathbb{H}$（四元数）、缩放 $\boldsymbol{\sigma} \in \mathbb{R}^3$、不透明度 $\alpha \in (0,1)$、以及三阶球谐系数 $\boldsymbol{s} \in \mathbb{R}^{16 \times 3}$。其中缩放和不透明度是旋转不变量（无论怎么旋转，椭球的轴长和透明度不变），无需特殊处理；但位置、旋转和球谐系数都依赖于骨骼变换中的旋转分量。在 LBS 框架下，每个高斯附带一组蒙皮权重 $\boldsymbol{w}_i \in \mathbb{R}^{J}$（$J$ 为关节数），位置通过加权混合骨骼变换来变形：

\[
\boldsymbol{\mu}_{i}^{t} = \mathcal{T}_{i}^{t}\boldsymbol{\mu}_{i}, \quad \mathcal{T}_{i}^{t} = \sum_{j=0}^{J} \boldsymbol{T}_{j}^{t} w_{i}^{j}
\]

其中 $\boldsymbol{T}_{j}^{t} \in \mathbb{SE}(3)$ 是骨骼变换矩阵。这一线性混合对位置变形是标准做法，问题在于：骨骼变换中的旋转分量经线性混合后，在很大比例的典型蒙皮权重下并不产生有效旋转（2007）~\cite{kavan2007skinning}。对位置而言这仅导致体积收缩（volume loss），但对旋转属性和球谐系数而言后果更为严重——无效旋转会直接破坏高斯的朝向和视角相关外观。

此前方法采用了三类绕行策略：(1)~用神经网络直接预测高斯旋转（如 GaussianAvatar、HUGS），但这增加了模型复杂度且旋转质量受限于网络容量；(2)~从网格面片（法线-切线-副切线）推导旋转（如 SplattingAvatar、iHuman），但这会丢失旋转自由度，且旋转更新受位置更新干扰；(3)~将视角方向反变换到规范空间来计算 SH（如 Animatable 3DGS），但这需要修改光栅化器且旋转仍然不一致。此外，几乎所有方法要么只使用零阶 SH 系数（放弃视角依赖效果），要么用神经网络预测高阶系数（增加模型负担），无法直接优化完整的球谐表示。

\textbf{改了什么}。On the Skinning of Gaussian Avatars 放弃上述绕行策略，直接从数学上解决问题：用加权四元数平均（weighted quaternion averaging）替代线性旋转混合来计算蒙皮旋转，然后用这个正确的蒙皮旋转同时处理高斯旋转和球谐系数的变形。这使方法成为纯顶点级别的动画技术，高度并行，且完全独立于具体的光栅化器实现。

\textbf{实现原理}。直觉上，问题出在"对旋转做线性平均"这一步——多个骨骼旋转的加权平均不应该在矩阵空间做线性插值，而应该在旋转空间做球面插值。四元数平均正是解决这一问题的经典方法。

具体地，对于时间步 $t$ 的姿态，通过正向运动学得到骨骼变换集合 $\{\boldsymbol{T}_{j}^{t}\}$。首先从每个骨骼变换中提取旋转分量，转换为四元数：

\[
\mathbf{b}_{j}^{t} = \mathcal{R}^{-1}(\boldsymbol{R}_{j}^{t})
\]

其中 $\mathcal{R}^{-1}$ 是旋转矩阵到四元数的转换函数。然后，利用高斯的蒙皮权重 $\boldsymbol{w}_i$ 构造加权外积矩阵：

\[
\mathbf{A}_{i}^{t} = \sum_{j=0}^{J} w_{i}^{j} (\mathbf{b}_{j}^{t})^{T}\mathbf{b}_{j}^{t}
\]

该矩阵最大特征值对应的特征向量即为加权平均四元数 $\bar{\boldsymbol{q}}_{i}^{t}$（直观理解：四元数的外积矩阵在加权平均意义下将多个旋转"汇聚"到最近的公共旋转）。随后用这个蒙皮旋转变形高斯的旋转分量：

\[
\boldsymbol{q}_{i}^{t} = \bar{\boldsymbol{q}}_{i}^{t} \cdot \boldsymbol{q}_{i}
\]

四元数乘法等价于旋转复合——先将规范空间的高斯旋转到某个朝向，再叠加骨骼驱动的蒙皮旋转。

球谐系数的旋转则需要另一组工具：Wigner D-矩阵 $\mathcal{W}^{l}(\boldsymbol{R})$，它将 $\mathbb{SO}(3)$ 中的旋转映射为 $(2l+1) \times (2l+1)$ 矩阵，可以正确旋转每个阶次 $l$ 的 SH 系数：

\[
\boldsymbol{s}_{i,0:1}^{t} = \boldsymbol{s}_{i,0:1}^{t}, \quad \boldsymbol{s}_{i,1:4}^{t} = \mathcal{W}^{1}(\bar{\boldsymbol{R}}_{i}^{t}) \boldsymbol{s}_{i,1:4}^{t}
\]
\[
\boldsymbol{s}_{i,4:9}^{t} = \mathcal{W}^{2}(\bar{\boldsymbol{R}}_{i}^{t}) \boldsymbol{s}_{i,4:9}^{t}, \quad \boldsymbol{s}_{i,9:16}^{t} = \mathcal{W}^{3}(\bar{\boldsymbol{R}}_{i}^{t}) \boldsymbol{s}_{i,9:16}^{t}
\]

其中零阶系数（常量颜色）无需旋转，而 1-3 阶系数分别用对应阶次的 Wigner D-矩阵进行旋转。这里 $\bar{\boldsymbol{R}} = \mathcal{R}(\bar{\boldsymbol{q}})$ 是将蒙皮四元数转回旋转矩阵。至此，位置（标准 LBS）、旋转（四元数平均）和球谐系数（Wigner D-矩阵）三类属性分别得到正确处理，缩放和不透明度因旋转不变性而无需变形。

\textbf{设计动机}。核心洞察是：此前方法的各种绕行策略（网络预测旋转、网格法线推导旋转、视角反变换）本质上都在回避"LBS 不能正确混合旋转"这一数学问题。On the Skinning of Gaussian Avatars 选择正面解决——四元数平均是旋转空间中加权平均的理论正确方法，且计算开销很小（仅需矩阵 $\mathbf{A}$ 的特征分解，实践中用 15 次幂迭代即可收敛）。这一设计使方法具有三个独特优势：(1)~完全独立于光栅化器——任何标准 GS 光栅化器都可以直接使用；(2)~高度并行——纯顶点级计算，无需神经网络推理；(3)~可即插即用地集成到现有化身方法中（实验中替换了 HUGS、iHuman、SplattingAvatar、Animatable 3DGS 等方法的蒙皮模块，均获得一致性提升）。

\textbf{公式与参数}。四元数平均矩阵 $\mathbf{A} \in \mathbb{R}^{4 \times 4}$，其最大特征向量通过 15 次幂迭代求得。Wigner D-矩阵使用 e3nn 库实现。在 Unity3D 引擎中，6 个同时动画的化身可在笔记本 GPU（3080 Ti Mobile）上达到 200+ FPS 的渲染速度。方法在 NeuMan、PeopleSnapshot、GalaBasketball 和 XHuman 四个数据集上均有验证，集成到 HUGS 后 PSNR 一致性提升，集成到 SplattingAvatar 后不仅质量提升还减少了所需高斯数量（表明优化过程本身也得到改善）。

\paragraph{SeqAvatar}

SeqAvatar（2024.11）~\cite{2411_16768} 修改了 3DGS 化身的非刚性变形环节——用层次化时序运动上下文替代单帧姿态条件——LBS 刚性变换和可微分光栅化渲染管线保持不变。

\textbf{Baseline 与问题}。此前 GauHuman、3DGS-Avatar 等方法仅使用当前帧的 SMPL 姿态参数作为非刚性变形的条件。这一设计存在两个不足：(1)~时序盲——同一姿态在不同运动阶段可能对应不同外观（例如手臂上举静止时和快速上举过程中的衣物褶皱完全不同），但单帧姿态无法区分这种"同一姿态、不同外观"的一对多映射；(2)~局部盲——骨骼姿态仅提供全局关节信息，无法捕获远离骨架的衣物（裙摆、袖口）的局部运动细节。此前 NeRF-based 的 Dyco 尝试用姿态序列建模时序信息，但其全局性的关节运动无法捕获逐点的精细运动，且未利用 3DGS 的显式点表示特性。

\textbf{改了什么}。SeqAvatar 引入从粗到细的层次化运动上下文（coarse-to-fine motion context），将非刚性变形的条件从"当前帧姿态"扩展为"粗粒度骨骼运动 $+$ 细粒度逐点运动"，并通过时空多尺度采样策略增强运动条件的鲁棒性。

\textbf{实现原理}。整体流程分为三级：粗粒度骨骼运动编码、细粒度逐点运动编码、以及时空多尺度采样整合。

(1)~\textit{粗粒度骨骼运动}。直觉上，人体运动是一个连续过程而非离散快照——手臂是"正在加速上举"还是"缓慢放下"会显著影响衣物形态。因此，对于时间步 $t$，SeqAvatar 考虑一系列等间隔采样的历史帧 $T = \{t-s, t-2s, \ldots, t-Ls\}$（$L$ 为序列长度，$s$ 为时间间隔），计算相邻帧之间的骨骼姿态差分作为运动描述符：

\[
\Delta P^t = P^t - P^{t-s}, \quad t \in T
\]

其中 $P^t \in \mathbb{R}^{K \times 3}$ 为 $K$ 个关节的轴角表示。沿时间维度展平后，通过 MLP 编码为骨骼运动嵌入：

\[
f_{\Delta P} = E_{\Delta P}(\Delta P) \in \mathbb{R}^{32}
\]

这一嵌入捕获了全身骨骼的运动趋势（加速、减速、方向变化），所有高斯共享同一个 $f_{\Delta P}$。

(2)~\textit{细粒度逐点运动}。直觉上，裙摆的运动无法仅从骨架关节推断——它取决于裙摆本身的局部位置和速度。利用 3DGS 的显式点表示，SeqAvatar 为每个高斯计算逐点速度向量 $v_i^t$，建模局部化的身体运动。然而，直接在优化过程中使用高斯位置计算速度存在循环依赖问题（计算速度需要观测空间坐标，而获得观测空间坐标又需要先完成非刚性变形）。为此，SeqAvatar 引入运动模板场（Motion Template Field）$\mathcal{V} = \{v_i\}_{i=1}^{N}$，用 SMPL 模板顶点的速度作为稳定参考。具体地，利用模板 LBS 权重 $W$ 将规范空间的 SMPL 顶点 $\mathcal{T}$ 变换到观测空间：

\[
\mathcal{T}_o^t = \mathrm{LBS}(\mathcal{T}, B^t, W)
\]

然后计算每个模板顶点的速度 $V^t = (\mathcal{T}_o^t - \mathcal{T}_o^{t-s})/s$。每个高斯的速度从运动模板场中按距离采样获得。

(3)~\textit{时空多尺度采样（STMS）}。为增强运动条件的鲁棒性、缓解过拟合，SeqAvatar 从空间和时间两个维度扩展感受野。空间维度上，对每个高斯采样 $K$ 个最近邻模板顶点的速度，通过 MLP $E_{\text{knn}}$ 学习局部运动嵌入：

\[
e_i^t = E_{\text{knn}}(v_j^t), \quad j \in \mathrm{KNN}(\mathcal{T}, x_i)
\]

时间维度上，采用多尺度序列采样策略，以递增间隔 $S = \{s_0 + i \cdot \Delta s\}_{i=0}^{m}$ 采样多个时间尺度的运动序列，同时捕获长期运动趋势和短期帧间细节。最终逐点运动嵌入维度为 $f_V \in \mathbb{R}^{96}$，与骨骼运动嵌入 $f_{\Delta P} \in \mathbb{R}^{32}$ 拼接为 $f_i' \in \mathbb{R}^{128}$ 的完整运动条件。

(4)~\textit{非刚性变形预测}。将层次化运动条件 $f_i'$ 与规范位置 $x$ 和当前姿态 $P$ 一起送入非刚性变形 MLP（3 个隐藏层，ReLU 激活），预测每个高斯的位置偏移、缩放偏移和旋转偏移：

\[
\Delta x, \Delta s, \Delta r = E_{\text{non-rigid}}(x, P, f_{\Delta P}, f_V)
\]

变形后的规范高斯为 $x' = x + \Delta x$、$s' = s + \Delta s$、$r' = r \cdot \Delta r$（四元数乘法等价于旋转复合）。随后通过标准 LBS 变换到观测空间，可微分光栅化渲染图像。

\textbf{设计动机}。核心洞察是：3DGS 的显式点表示为运动建模提供了独特优势——每个高斯都有明确的 3D 位置，可以直接计算逐点速度，这是 NeRF 等隐式表示无法做到的。因此，SeqAvatar 利用这一特性设计逐点运动条件，突破此前方法仅依赖全局骨骼姿态的局限。运动模板场的引入则解决了优化过程中高斯位置不断更新导致的速度计算不稳定问题。时空多尺度采样进一步扩展了运动上下文的感受野，使模型能够鲁棒地处理复杂运动。

\textbf{公式与参数}。骨骼运动嵌入维度 32，逐点运动嵌入维度 96，拼接后运动条件共 128 维。非刚性变形 MLP 使用 3 个隐藏层，每层后接 ReLU 激活。在 I3D-Human 数据集上，SeqAvatar 在大多数指标上超越此前最优方法，同时保持实时渲染（$\geq 45$ FPS）。

\paragraph{SFGS}

SFGS（2026.04）~\cite{2604_09324} 在 SeqAvatar 的层次化运动条件基础上，进一步改进了特征编码方式。

\textbf{Baseline 与问题}：SeqAvatar 使用模板顶点速度作为运动条件，结合多分辨率哈希网格编码规范位置特征。然而哈希网格的空间分辨率固定，无法自适应地为高动态面部区域和低动态身体区域分配不同的特征精度。此外，SeqAvatar 的时序建模依赖于速度向量，缺乏对长期时序连续性的显式编码。

\textbf{改了什么}：SFGS 将特征编码从哈希网格替换为三平面-六平面混合架构——静态三平面 $f_{\text{tri}}(v_i) \in \mathbb{R}^{3d}$ 捕获姿态依赖的空间细节，动态六平面 $f_{\text{hex}}(v_i, t_i) \in \mathbb{R}^{6d}$ 引入时间维度建模时序连续性。此外，面部区域使用更高分辨率的专用平面，并通过基于精细手部重建的残差精修模块改进手部变形。

\textbf{设计动机}：三平面在空间效率上优于哈希网格（三个 2D 平面 vs. 多层哈希表），且面部分辨率自适应的设计将更多参数预算分配给视觉上最敏感的区域。六平面在 HexPlane 的基础上增加了时间维度，使特征场能同时编码空间位置和时间变化。

\paragraph{HumanSplatHMR}

HumanSplatHMR（2026.05）~\cite{2605_02784} 提出了人体网格恢复与高斯化身之间的闭环联合优化框架。

\textbf{Baseline 与问题}：上述全身化身方法（包括 3DGS-Avatar、GoMAvatar 及各种网格-高斯混合架构）均将人体姿态估计和化身重建视为两个独立步骤：先用外部工具（运动捕捉系统或离线 SMPL 拟合器）获取姿态参数，再基于固定姿态训练化身。这一流水线存在两个问题：ViT-based 姿态估计器并不总是可靠，可能过拟合到 2D 视角；而 NeRF/3DGS 化身方法将姿态和外观分开处理，限制了化身对新姿态的泛化——不准确的姿态输入会直接导致化身渲染伪影，且错误无法在训练中被纠正。

\textbf{改了什么}：HumanSplatHMR 将姿态估计和化身重建统一为联合优化问题——不再将 SMPL 姿态仅作为变形先验使用，而是通过可微分渲染器将光度损失、分割损失和深度损失反向传播到 SMPL 姿态参数和全局位置，使姿态精修与化身学习相互促进。

\textbf{实现原理}：框架核心是"闭环"设计。

首先用现有人体姿态估计器（如 HMR 2.0）获取初始 SMPL 参数，而非依赖昂贵的运动捕捉或离线精修，以模拟真实世界条件。然后在可微分 3DGS 渲染管线中引入松弛耦合损失（loose coupling loss），约束 SMPL 网格与高斯点云之间的一致性——高斯点云提供的视觉监督信号通过梯度回传同时更新化身参数和 SMPL 姿态参数。随着训练的推进，姿态估计的精度持续提升（因为渲染损失提供了像素级的监督），而更准确的姿态反过来又改善了化身的几何质量，形成正反馈循环。

\textbf{设计动机}：将姿态和化身解耦处理是当前化身方法泛化性不足的根源之一。当输入姿态有误差时，化身被迫在不正确的几何约束下学习外观，导致新姿态合成质量下降。HumanSplatHMR 的闭环设计使姿态误差可以在化身训练过程中被自动修正。

\textbf{公式与参数}：实验表明，该方法在姿态恢复精度和新视角/新姿态合成质量上均持续优于仅做图像级精修的姿态基线和将姿态估计与化身重建解耦的化身基线。

\subsection{单图像与少输入创建}
\label{ss:16_2_单图像与}

本节介绍从单张图像或少量输入创建全身化身的方法。

\subsubsection{从单图像/少输入创建化身}

本路线介绍从单张图像或稀疏输入快速生成全身~3DGS~化身的方法。

从单张图像创建全身化身是最具挑战性也最实用的方向之一。

\paragraph{LHM}

LHM（2025.03）~\cite{2503_10625}（Large Animatable Human Reconstruction Model）修改了 3DGS 的化身创建管线——用多模态 Transformer 前向推理替代多视角优化循环——3D 高斯的参数化表示和可微分光栅化渲染管线保持不变。

\textbf{Baseline 与问题}。上文介绍的 3DGS-Avatar 和 GoMAvatar 等先驱方法都依赖单目视频序列进行逐帧优化（3DGS-Avatar 需 30 分钟训练），且需要预拟合的 SMPL 参数。核心瓶颈在于：标准 3DGS 的化身创建本质是一个优化问题（从 COLMAP 点云或 SMPL 网格初始化高斯，然后迭代优化参数），而非推理问题——给定一张图像，系统无法"一步"生成化身。

\textbf{改了什么}。LHM 将化身创建从优化问题重新定义为回归问题：设计多模态 Transformer 架构，以单张 RGB 图像为输入，在单次前向传播中直接预测规范空间中每个 3D 高斯的全部参数（位置偏移、旋转、缩放、球谐系数）。旧机制——多视角输入 + COLMAP 初始化 + 迭代优化——被新机制——单图像输入 + SMPL-X 采样点初始化 + Transformer 回归——替代。

\textbf{实现原理}。LHM 的核心是多模态体-头 Transformer（Multimodal Body-Head Transformer, MBHT），它融合三类 token 来回归高斯参数。
(1)~\textit{几何 token}：从 SMPL-X 规范网格表面采样 3D 点 $\{X_i\}_{i=1}^{N_{\text{points}}}$，经正弦位置编码和 MLP 投影得到几何 query token：

\[
T_{3D} = \phi_{\text{proj}}(\gamma(X)) \in \mathbb{R}^{N_{\text{points}} \times C}
\]

直觉上，SMPL-X 网格上的采样点充当"问题点"——每个点询问图像"我在这个位置应该长什么样"。$C$ 是 token 维度。
(2)~\textit{身体图像 token}：使用在千万级人体图像上预训练的 Sapiens-1B 编码器提取语义身体特征。由于 Sapiens 是专门针对人体预训练的视觉 Transformer，它能捕获比通用编码器更丰富的衣物纹理和身体形状信息。
(3)~\textit{头部特征金字塔 token}：人的头部在输入图像中只占很小的区域，经过编码器的空间下采样后关键面部细节（皱纹、瞳孔）几乎完全丢失。为解决这一问题，LHM 设计了头部特征金字塔（Head Feature Pyramid），从 DINOv2（自监督视觉Transformer，在大型数据集上通过自蒸馏训练的视觉基础模型）的第 4、11、17、23 层提取多尺度特征并融合：

\[
T_{\text{head}} = F_{\text{fusion}}(E_{\text{dino}}^4, E_{\text{dino}}^{11}, E_{\text{dino}}^{17}, E_{\text{dino}}^{23}) \in \mathbb{R}^{N_{\text{head}} \times C}
\]

直觉上，浅层保留高频纹理细节（毛孔、细纹），深层编码稳健的头部几何先验（整体形状），两者融合使面部身份和细节同时得到保持。

三类 token 在 MBHT 块中通过注意力机制交互融合。此外，LHM 还从身体图像 token 中提取全局上下文特征 $F_{\text{global}} = \text{MLP}_{\text{global}}(\text{MaxPool}(T_{\text{body}}))$，用于调制各注意力块的计算，为局部 token 交互提供整体人体结构的全局引导。MBHT 的关键设计是分步融合：头部 3D 点 token 先与头部图像 token 交互（增强面部细节），然后与身体 3D 点 token 拼接，再与身体图像 token 交互——这种"先头后身"的顺序确保面部细节不会在大量身体 token 的注意力竞争中被淹没。

经过 $N_{\text{layer}}$ 层 MBHT 块后，MLP 回归头从每个几何 token 预测 3DGS 参数：

\[
(\Delta p_i, r_i, f_i, s_i, o_i) = \phi_{\text{regress}}(T_{\text{points}}^{(i)}), \quad p_i = x_i + \Delta p_i
\]

其中 $\Delta p_i \in \mathbb{R}^3$ 是相对于 SMPL-X 初始位置 $x_i$ 的残差偏移，$r_i, f_i, s_i, o_i$ 分别为旋转、球谐系数、缩放和不透明度。训练时通过 LBS 将规范化身变换到目标视角（为更好地建模衣物变形，LHM 采用扩散体素蒙皮方法而非标准 LBS 来计算蒙皮权重），使用光度损失和规范空间正则化进行端到端训练。

\textbf{设计动机}。(1)~将 3D 点云而非 2D patch 作为 query 是因为 3DGS 的输出天然是 3D 高斯点集，点云 query 可直接回归每个高斯的位置和属性。(2)~头部特征金字塔的必要性源于尺度失衡：头部面积小但细节密度高（眼睛、嘴唇的微小变化对身份识别至关重要），通用编码器的下采样会不可逆地丢失这些信息，因此需要用专门的浅层特征补偿。(3)~全局上下文特征的引入解决了局部注意力的视野局限：单个 3D 点 token 仅关注自身邻域的图像特征，缺乏对身体整体姿态和比例的感知，max pooling + MLP 从全局身体 token 中提取的上下文嵌入为所有注意力计算提供了人体整体结构的先验。

(4)~头部 token 收缩正则化（Head Token Shrinkage Regularization）：实验发现注意力机制过度依赖头部特征，削弱了身体部位的学习。受 MAE 启发，训练时以 0\%--50\% 的比例随机遮蔽头部图像 crop，迫使模型通过身体上下文补偿，提升身体部位重建质量。

\textbf{公式与参数}。规范空间正则化由两项组成：(1)~球形化损失（as-spherical-as-possible），惩罚高斯的过度各向异性，避免产生针状椭球：$L_{\text{ASAP}} = \frac{1}{N}\sum_{i=1}^{N} \|S_i - \text{diag}(1)\|_F^2$，其中 $S_i$ 为协方差矩阵，$\text{diag}(1)$ 为单位矩阵——该项约束协方差矩阵接近各向同性，使高斯趋向球形；(2)~位置锚定损失（as-close-as-possible），约束高斯位置不偏离 SMPL-X 初始位置过远：$L_{\text{ACAP}} = \frac{1}{N_{\text{points}}}\sum_{i=1}^{N_{\text{points}}} \max(\|\Delta p_i\|_2 - d, 0)$，铰接距离阈值 $d = 5.25\text{cm}$，允许局部调整但防止大幅漂移。综合正则化 $L_{\text{reg}} = 50 L_{\text{ASAP}} + 10 L_{\text{ACAP}}$。光度损失权重 $\lambda_{\text{rgb}} = 1.0$、$\lambda_{\text{mask}} = 0.5$、$\lambda_{\text{lpips}} = 1.0$。整个推理过程在秒级完成。

\paragraph{LHM++}

LHM++（2025.06）~\cite{2506_13766} 进一步改进 LHM 的效率和输入灵活性。

\textbf{Baseline 与问题}：LHM 虽然实现了单图像前馈化身重建，但其多模态体-头 Transformer 中的全局注意力机制计算复杂度随点 token 数量二次增长（规范网格上采样点越多，注意力矩阵越大），限制了可处理的点密度和推理速度。此外，LHM 要求输入图像附带准确的相机参数和 SMPL-X 姿态标注，在用户随意拍摄的照片中这些信息通常不可用。

\textbf{改了什么}：LHM++ 在架构层面引入编码器-解码器点-图像 Transformer（Encoder-Decoder PIT），将 LHM 的单层全局注意力替换为渐进式编解码流程；在输入层面去除了对相机参数和姿态标注的依赖，支持从一张或多张无标注的随意拍摄图像直接生成可动画化身。

\textbf{实现原理}：核心架构 PIT 的工作方式如下。

编码器阶段采用 Point Transformer v3 的分块自注意力（partitioned self-attention），将规范空间中的 3D 几何点 token 按空间邻近性分块处理，通过网格池化（grid pooling）逐层压缩点 token 数量——浅层处理密集的局部点特征，深层处理稀疏的全局语义特征。与 LHM 对所有点 token 施加全局注意力（$O(N^2)$ 复杂度）不同，分块注意力的复杂度近似线性。解码器阶段通过多模态注意力将分层的 3D 点特征与图像特征融合，解码后的特征用于预测 3D 高斯参数。此外，LHM++ 引入二分图软匹配（Bipartite Soft Matching）策略合并冗余图像 token——当多张输入图像包含大量重叠区域时，图像 token 之间存在高度冗余，软匹配通过二分图匹配保留信息量最大的 token 子集，进一步降低计算开销。

LHM++ 还配备轻量级 3D 感知神经动画渲染器（3D-aware Neural Animation Renderer），在实时推理中对重建化身的渲染质量进行精修。

\textbf{设计动机}：LHM 的全局注意力在点 token 数量较大时成为效率瓶颈，而实际应用中高质量的化身需要密集的采样点。渐进式编解码将"理解全局结构"和"生成局部细节"分配到不同层级处理，既保持了多尺度特征的完整性又大幅降低了计算量。去除相机和姿态标注的要求则使方法在真实世界场景中更加实用。

\textbf{公式与参数}：在无需相机或姿态标注的条件下，LHM++ 可在数秒内从一张或多张随意拍摄的图像生成高质量可动画化身。

\paragraph{PERSONA}

PERSONA（2025.08）~\cite{2508_09973} 与 LHM++ 的前馈回归路线互补，采用"扩散模型生成训练数据 + 3DGS 优化化身"的策略。

\textbf{Baseline 与问题}：LHM/LHM++ 的前馈模型在训练分布内表现出色，但单图像输入的信息量有限——无法直接观察到人体的背面和侧面，前馈模型只能依靠训练数据中的统计先验"猜测"未观察区域，导致几何和纹理的精细度受限。多视角优化方法（如 GoMAvatar）虽然质量更高，但需要多视角视频输入。

\textbf{改了什么}：PERSONA 结合了两种路线的优势——利用扩散模型从单图像生成姿态丰富的合成视频作为训练数据，然后用 3DGS 优化方法从这些合成数据中创建高质量化身。关键设计是平衡采样策略（在多样性——覆盖更多姿态——和质量——保证渲染精度——之间取得平衡），以及几何加权优化（对不同身体部位施加差异化权重，确保跨多样姿态的渲染质量）。

\textbf{设计动机}：扩散模型提供了"虚拟多视角拍摄"能力——从单张图像生成覆盖 360 度的合成训练数据，使 3DGS 优化方法能获得近似多视角的监督信号。

\paragraph{AdaHuman}

AdaHuman（2025.05）~\cite{2505_24877} 进一步提升了扩散辅助的化身创建质量。

\textbf{Baseline 与问题}：PERSONA 的扩散模型生成整体人体图像，但不同身体部位（头部、躯干、手部）在渲染质量上的重要性差异很大——面部和手部是视觉注意力集中的区域，质量要求远高于躯干。

\textbf{改了什么}：AdaHuman 引入姿态条件化 3D 联合扩散模型，在任意姿态下合成一致的多视角图像。同时通过组合式 3DGS 精修模块分别处理头部、躯干、手部等局部身体部位的细节增强——为不同部位分配独立的精修网络，使计算资源集中在视觉上最敏感的区域。

在单目视频设定下，几何信息的匮乏使得化身重建更加困难。以下方法通过引入基础模型先验和修复策略来弥补这一缺陷。

\paragraph{FMGS-Avatar}

FMGS-Avatar（2025.09）~\cite{2509_14739} 针对单目视频几何信息不足的问题引入基础模型先验。

\textbf{Baseline 与问题}：多视角 3DGS 方法依赖 COLMAP 提供的稀疏点云或 SMPL 网格初始化高斯位置，但单目视频无法提供可靠的多视角几何——这导致高斯初始化不准确，优化容易陷入局部最优。

\textbf{改了什么}：FMGS-Avatar 利用 Sapiens 等大规模预训练模型提取前景掩码、伪深度、法线和语义分割等多模态 2D 先验，补充单目视频的有限视觉线索。在表示上选择 2D 高斯（surfels）而非 3D 高斯——前者作为平面基元更适合表示衣物和皮肤等薄表面，避免了 3D 椭球必须极度扁平化导致的训练不稳定。

\textbf{设计动机}：基础模型（如 Sapiens 在千万级人体图像上预训练）提供的 2D 先验虽然是近似的，但足以替代缺失的多视角几何信息，为 3DGS 优化提供合理的初始化和正则化约束。

\paragraph{InpaintHuman}

InpaintHuman（2026.01）~\cite{2601_02098} 专门解决遮挡场景下的化身重建，与 AHOY 从野外视频处理遮挡的路线互补。

\textbf{Baseline 与问题}：AHOY 用视频扩散模型为未观察区域"幻觉"出监督，但视频扩散的逐帧随机性可能引入时序不一致。InpaintHuman 从图像修复而非视频生成的角度解决遮挡问题。

\textbf{改了什么}：InpaintHuman 通过多尺度 UV 参数化表示和身份保持扩散修复模块实现完整化身的生成。该模块对个性化扩散模型进行身份微调后，为未观察到的身体区域生成像素级监督信号——即使人体的一半被遮挡物覆盖，也能推断出完整的 3D 外观。身份微调确保修复结果与已观察区域保持身份一致性，不会生成属于其他人的外观。

\subsection{物理与衣物动态}
\label{ss:16_2_物理与衣}

本节讨论~2~条技术路线：（1）物理模拟与衣物动态；（2）遮挡、模糊和野外场景的处理。

\subsubsection{物理模拟与衣物动态}

本路线介绍基于物理模拟的衣物动态和人体运动建模方法。

宽松衣物的物理可信动画是全身化身的重要挑战。前文讨论的化身方法主要采用数据驱动的变形策略，即从姿态参数直接回归衣物变形。然而衣物可以在同一姿态下处于不同状态（如袖子贴臂或滑落），这种一对多的映射关系难以用数据驱动方法准确建模。物理模拟提供了另一条路线：从物理定律出发推演衣物运动，不依赖训练数据的覆盖范围。

\paragraph{PICA}

PICA（2024.07）~\cite{2407_05324}（Physics-Integrated Clothed Avatar）修改了 3DGS 化身的表示与驱动架构——用双层独立 3DGS（body + clothing）替代单层统一 3DGS，并用基于图神经网络的物理模拟替代骨架 LBS 驱动衣物——多视角光栅化渲染管线保持不变。

\textbf{Baseline 与问题}。前述 GoMAvatar、ExAvatar 等方法均将衣物和人体统一表示为基于 SMPL(-X) 网格的单层 3DGS，通过 LBS 驱动整体变形，再用 MLP 学习姿态依赖的非刚性偏移来近似衣物动态。这种数据驱动的策略存在两个根本问题：第一，衣物可以在同一人体姿态下处于不同状态（例如手臂抬起时袖子可能贴臂也可能滑落），因此将衣物动态建模为单帧姿态的函数 $f(\theta)$ 本身就是欠定的；第二，数据驱动方法在训练分布之外的姿态上容易产生不合理的衣物运动伪影，而覆盖足够多姿态又需要极长的训练序列。

\textbf{改了什么}。PICA 将单层化身拆分为双层独立 3DGS 表示：一套高斯锚定在体网格 $M^{\text{body}}$（直接使用 SMPL-X 模板）上表示人体，另一套高斯锚定在衣物网格 $M^{\text{clothing}}$ 上表示衣物。衣物网格的获取方式因衣物类型而异：对于紧身衣物，直接复用 SMPL-X 网格；对于宽松衣物或裙装，则通过 Animatable Gaussians 提供的方法从单帧重建模板网格，再将 2D 分割掩码（由 PerSAM 和 Grounded SAM 自动生成）投影到全身网格上裁剪出衣物区域。对于长发表演者，还在头部后方附加额外三角面片以建模头发。两套高斯的变形特性完全独立——体网格仍由 LBS 驱动，衣物网格则由基于图神经网络（GNN）的物理模拟器逐帧预测加速度来驱动。两套网格各配备独立的非刚性变形 MLP（以帧索引 $t$ 为条件，预测逐顶点偏移 $\Delta \mathbf{p} = f_{\text{part}}(\mathbf{p}, t)$），衣物网格的 LBS blend weights 初始化为体网格最近邻顶点的权重，并在训练中联合优化。

\textbf{实现原理}。核心架构包含三个紧密耦合的模块。

(1)~双层网格-高斯表示。每个三角面片上锚定 13 个 3D 高斯。高斯的几何属性（位置、旋转、缩放）完全由所在面片的几何决定：位置通过重心坐标加法线偏移定义，$x = \sum_{j=1}^{3} b_j p_{ij} + \epsilon n_i$；旋转矩阵的列向量分别沿面片法线 $n_i$、边方向 $p_{i1} - p_{i0}$ 及其叉积排列，确保高斯与面片对齐；法线方向缩放设为极小常数，使高斯呈薄片状贴合表面。这种设计使得高斯随面片运动而自动更新，无需独立学习高斯的运动场。颜色方面，PICA 用姿态条件化颜色 MLP 替代标准球谐函数：$c_i = f_{\text{color}}(f_i, \theta, h_t, \hat{d})$，其中 $f_i$ 为逐高斯特征向量，$\theta$ 为 SMPL-X 姿态参数，$h_t$ 为帧索引的时间嵌入（编码时变因素如衣物褶皱的展开），$\hat{d} = R_i^T d$ 为规范化的视角方向。这一替换的直觉是：球谐系数只能表示视角依赖的颜色变化，无法捕捉衣物自阴影和褶皱引起的姿态依赖外观。

(2)~基于层次化图神经网络的物理模拟。PICA 借鉴 HOOD 框架，从 $M^{\text{body}}$ 和 $M^{\text{clothing}}$ 构建层次化图 $\mathcal{G} = (V^{\text{body}} \cup V^{\text{clothing}},\; E^{\text{clothing}} \cup E^{\text{coarse}} \cup E^{\text{clothing2body}})$，其中 $E^{\text{clothing}}$ 是衣物网格的边，$E^{\text{coarse}}$ 是简化衣物网格上的边（提供远程连接），$E^{\text{clothing2body}}$ 是从衣物顶点到最近体顶点的跨层边（动态更新，阈值内才连接）。图的节点特征包含当前位置 $p_t$、速度 $v_t$、法线 $n_t$ 和物理属性 $\phi$（质量、弯曲系数等），边特征存储两节点的相对位置。GNN 基于 MeshGraphNets 架构，逐帧预测衣物顶点的加速度：

\begin{equation}
    a_t = \text{GNN}(\mathcal{G}_{t-1}, \phi), \quad v_t = v_{t-1} + a_t, \quad p_t = p_{t-1} + v_t,
\end{equation}

即"给定上一帧图状态和物理属性，预测当前帧加速度，然后通过半隐式欧拉积分更新速度和位置"。物理属性 $\phi \in \mathbb{R}^4$（包含质量、弯曲刚度等四维参数）是可优化变量，通过自监督物理损失从重建数据中学习：

\begin{equation}
    \mathcal{L}_{\text{phys}}(\phi) = \sum_{t=0}^{T} \| p_{t+1}^{r} - p_{t+1}^{s}(p_t^{r}, p_{t-1}^{r}, \phi) \|_2^2,
\end{equation}

其中 $p_{t+1}^{r}$ 是第 $t+1$ 帧的重建位置（来自多视角优化），$p_{t+1}^{s}$ 是模拟器从第 $t$ 和 $t-1$ 帧重建位置出发、以物理属性 $\phi$ 为条件预测的下一帧位置。直觉上，这个损失迫使物理属性 $\phi$ 的取值使得模拟器的轨迹尽可能贴近真实观测——不需要真实物理参数的标注，只需渲染质量的间接监督。

(3)~双层分割与碰撞处理。训练时引入分割损失 $\mathcal{L}_{\text{seg}}$ 确保衣物高斯只在衣物网格上渲染、体高斯只在体网格上渲染（通过 BCE 损失约束每个高斯的渲染贡献是否应出现在衣物区域）。碰撞损失 $\mathcal{L}_{\text{collision}} = \frac{\lambda}{n} \sum_{i=1}^{n} \max(\epsilon - (v_i^c - v_j^b) \cdot n_j^b,\; 0)^3$ 惩罚衣物顶点 $v_i^c$ 穿入体网格，其中 $v_j^b$ 是衣物顶点在体网格上的最近邻顶点，$n_j^b$ 为对应的法线向量，$\epsilon$ 为碰撞阈值小量——当衣物到体顶点的向量与体法线点积小于 $\epsilon$ 时施加惩罚，确保衣物网格始终处于体网格的外法线方向。此外还引入不透明度稀疏损失 $\mathcal{L}_{\text{opac}}$ 鼓励高斯不透明度趋向 0 或 1（避免半透明伪影），以及 Laplacian 平滑损失和法线一致性损失来保持网格几何规整，距离损失防止体网格偏离 SMPL-X 初始化过远。

\textbf{设计动机}。双层表示的核心动机是解耦衣物和人体的运动特性——人体运动由骨架刚性地驱动，而宽松衣物的运动由惯性、重力和接触力决定，两者遵循完全不同的物理规律。用 GNN 物理模拟替代数据驱动的姿态依赖变形，本质上是将"从姿态查表得变形"的范式替换为"从物理定律推演变形"的范式，使系统对训练分布外的姿态也能产生物理合理的衣物运动。物理属性 $\phi$ 的可优化设计则弥合了通用预训练 GNN 与特定衣物材质之间的差距。

\textbf{公式与参数}。在 ActorsHQ 数据集（25 个均匀分布相机、$1494 \times 1022$ 分辨率、200 帧）上训练约 40K 次迭代（单卡 A100 约 100 分钟）。每个面片锚定 13 个高斯。HOOD 的 GNN 模型在大型布料模拟数据集上预训练，可泛化到训练时未见过的衣物类型和形状，PICA 直接使用该预训练模型而不再重新训练。训练时对颜色网络的姿态输入和隐编码采用退火策略（annealing），逐步降低帧索引编码的影响以避免过拟合。推理时，体网格通过 LBS 驱动，衣物网格通过 HOOD 物理模拟驱动，整个动画系统以约 12 FPS 运行；仅渲染（不含物理模拟）可达约 20 FPS。得益于双层表示的独立性，PICA 还支持虚拟试穿（virtual try-on）——将一个表演者的衣物网格和高斯替换为另一个表演者的，即可在同一人体上驱动全新衣物的动画。

\paragraph{PBDyG}

PBDyG（2024.12）~\cite{2412_04433}（Position Based Dynamic Gaussians）沿 PICA 的"物理模拟替代数据驱动"路线，但选择了不同的物理引擎。

\textbf{Baseline 与问题}：PICA 的 GNN 物理模拟需要预训练的 GNN 模型（HOOD），且 GNN 的推理开销使其难以实时运行（整体约 12 FPS）。此外，PICA 将高斯锚定在面片上，高斯的位置完全由网格决定，无法建模超出网格表面的大幅衣物偏移。

\textbf{改了什么}：PBDyG 用基于位置的动力学（PBD，Position Based Dynamics）替代 GNN 物理模拟。PBD 是游戏引擎中广泛使用的实时物理方法，通过迭代约束投影求解位置，无需神经网络推理，天然支持实时计算。高斯直接在 3D 空间中自由放置（不绑定到网格面片），用 Delaunay 三角剖分在高斯中心之间建立四面体连通结构（保留 $k=30$ 邻域内的有效四面体），然后通过 PBD 的距离约束（维持相邻高斯的间距）和碰撞约束（防止穿入身体）迭代求解高斯位置。

\textbf{实现原理}：物理属性（质量、材料刚度）从多视角 RGB 视频中估计：通过光度损失反向传播优化物理参数，使模拟轨迹逼近真实观测，与 PICA 的自监督物理损失思想类似。PBD 的约束求解在 Gauss-Seidel 迭代中逐个满足约束，每帧仅需数次迭代即可收敛。

\textbf{设计动机}：PBD 的核心优势在于简单高效：无需预训练数据、无需神经网络推理，纯几何约束即可产生物理合理的运动，适合实时应用。四面体连通结构保持了高斯之间的拓扑关系，使衣物不会断裂或自穿。

\paragraph{MPMAvatar}

MPMAvatar（2025.10）~\cite{2510_01619} 在物理模拟的逼真度上更进一步。

\textbf{Baseline 与问题}：PICA 的 GNN 和 PBDyG 的 PBD 都是位置/速度层面的简化物理模型，无法精确模拟连续介质的大变形（如裙摆的剧烈翻转）和拓扑变化（如衣物的自接触和折叠）。PBD 的距离约束是硬约束，对大幅变形的处理容易产生锁死（locking）现象。

\textbf{改了什么}：MPMAvatar 使用材料点方法（MPM，Material Point Method，一种在连续介质力学中同时使用拉格朗日粒子和欧拉网格的混合模拟方法）模拟器。MPM 通过拉格朗日粒子（携带质量、动量）和欧拉网格（求解动量守恒方程）的双向耦合求解：粒子在网格之间传输物理量，网格求解 Navier-Stokes 方程后再将结果插值回粒子，兼具拉格朗日方法的大变形追踪能力和欧拉方法的数值稳定性。

\textbf{设计动机}：MPM 的核心优势是首次实现了对未见交互的零样本泛化：训练时用 A 姿态训练衣物模拟，推理时衣物对全新的身体运动和外部交互（如被风吹、被手拉扯）仍能产生物理合理的响应。这是因为 MPM 从连续介质力学方程出发（而非从训练数据学习运动模式），物理定律对训练分布外的场景同样适用。

\paragraph{RealityAvatar}

RealityAvatar（2025.04）~\cite{2504_01559} 从运动分解的角度处理衣物动态。

\textbf{Baseline 与问题}：PICA、PBDyG、MPMAvatar 用物理模拟驱动衣物，需要显式的碰撞检测和物理参数估计。对于运动趋势明显的衣物（如裙摆的摆动主要由身体运动方向决定而非复杂碰撞），物理模拟的计算开销可能是不必要的。

\textbf{改了什么}：RealityAvatar 引入运动趋势模块（motion trend module）和潜在骨骼编码器（latent bone encoder），将规范空间到观测空间的变换分解为运动趋势（捕获整体运动方向）和人体变换（刻画姿态驱动的形状变化）两个组件。潜在骨骼编码器借鉴"运动后稳态"假设——同一姿态在运动后稳态下的衣物变形主要取决于当前姿态参数——通过学习潜在骨骼结构来建模姿态驱动的衣物变形，显式处理服装行为的时间变化。

\paragraph{STG-Avatar}

STG-Avatar（2025.10）~\cite{2510_22140} 从高斯表示本身的可扩展性出发。

\textbf{Baseline 与问题}：上述方法均使用静态高斯（规范空间中的固定点集），动态效果通过变形网络或物理模拟驱动。这种间接建模方式在处理高频时变细节（如衣物褶皱的快速出现和消失）时容易产生模糊。

\textbf{改了什么}：STG-Avatar 将时空高斯（Spatio-Temporal Gaussian, STG）与 LBS 协同集成——为每个高斯引入时间属性使其能直接表示动态场景。利用光流约束识别高动态区域并引导自适应致密化（在高动态区域增加更多时空高斯），配合紧凑 MLP 解码器高效预测颜色。

\textbf{设计动机}：直接在时空中建模高斯使每个高斯可以在不同时间步具有不同的属性，而非仅通过变形间接产生时间变化。光流引导的致密化确保了参数预算被分配到最需要时变细节的区域。

\subsubsection{遮挡、模糊和野外场景的处理}

本路线介绍从遮挡、运动模糊等恶劣条件下重建~3DGS~人体化身的方法。前文讨论的方法大多假设输入视频清晰完整，但真实世界的视频往往充满各种退化因素。

真实世界场景中人体经常被遮挡、运动模糊或处于复杂的野外环境。前文讨论的方法大多假设输入视频质量良好（清晰、无遮挡、相机标定准确），但真实世界的数据远不理想。本节介绍处理这些恶劣条件的方法。

AHOY~\cite{2603_17975}（全称 Animatable Humans under Occlusion from YouTube Videos，2026.03 发布）引入了"幻觉即监督"（hallucination-as-supervision）这一全新的学习范式，利用身份微调的视频扩散模型为未观察到的身体区域生成密集的像素级监督信号，从而实现从严重遮挡的野外单目视频重建完整可动画 3DGS 化身——多视角光栅化渲染管线保持不变。

\textbf{Baseline 与问题}。现有 3DGS 化身方法（3DGS-Avatar、HUGS、GaussianAvatar 等）均假设输入视频中人体完全可见、无遮挡，且通常需要规范姿态的多视角覆盖。这一假设排除了绝大多数真实世界视频——YouTube 上的日常活动视频（烹饪、运动、对话）中人体频繁被家具、物体或其他人物遮挡，大量身体区域从未被观察到，且每个姿态下不存在多视角监督。从这种严重遮挡的单目视频重建完整化身面临两个根本挑战：第一，大范围身体区域无任何直接观测数据可供监督；第二，缺少多视角覆盖使得姿态依赖的外观变化（如衣物褶皱随姿态的变化）无法被学习。

\textbf{改了什么}。AHOY 提出四阶段流水线，通过扩散模型"幻觉"出的合成数据弥补缺失监督。第一阶段从部分观察的视频帧构建粗化身：利用 DensePose 将可见纹理映射到规范姿态 UV 图集，用 FLUX 修复未观察区域，再用多视角扩散模型生成 $V{=}4$ 个正交视角（前、后、左、右）的规范姿态图像，通过 StyleUNet 预测规范高斯图集（canonical Gaussian maps）并经 LBS 变形至观测帧，在可见区域掩码监督下训练。第二阶段用 LoRA 在 Wan 2.2 视频扩散模型上微调主体身份（将视频帧视为无序集合以解耦外观与运动），然后让粗化身执行精心设计的结构化运动序列（$360^{\circ}$ 旋转配对目标姿态），将粗渲染结果通过 RF-Inversion 经由身份微调后的扩散潜空间进行前向-反向传播，生成覆盖未观察区域的"幻觉"监督视频。第三阶段利用幻觉视频升级为姿态依赖高斯图集（pose-dependent Gaussian maps），引入 map-pose/LBS-pose 解耦和联合姿态-相机优化来吸收扩散数据的多视角不一致性，并用 FLAME 头部路径保持面部身份。第四阶段支持任意姿态驱动和 3DGS 场景合成。

\textbf{实现原理}。核心架构包含四个紧密耦合的模块。

(1)~纹理映射与粗化身构建。给定单目视频帧 $\{\mathbf{I}_t\}_{t=1}^{T}$，先用 NLF 估计逐帧 SMPL 参数 $\{(\bm{\theta}_t, \bm{\beta})\}$，用 SAM 3 获取可见性掩码 $\{\mathbf{M}_t\}$，再通过 DensePose 将可见像素映射到规范纹理图集（first-write-wins 规则避免时序混合伪影）。修复后的规范图像经多视角扩散模型生成 $V{=}4$ 视角。StyleUNet $S$ 处理规范高斯图集，预测逐像素高斯属性偏移和学习皮肤权重，训练目标为 $\mathcal{L}_{\text{coarse}} = \sum_{v=1}^{V} \|\mathbf{I}_*^{v} - \mathcal{R}(\mathcal{G}^{H}; \bm{\gamma}_{v})\|_1 + \sum_{t=1}^{T} \|(\mathbf{I}_t - \hat{\mathbf{I}}_t) \odot \mathbf{M}_t\|_1$，其中 $\hat{\mathbf{I}}_t = \mathcal{R}(\mathcal{T}(\mathcal{G}^{H}; \bm{\theta}_t); \bm{\gamma}^{\text{obs}})$ 为渲染图像。

(2)~幻觉监督生成。用 LoRA 微调 Wan 2.2，学习身份绑定文本 token $[\text{v}]$ 和低秩权重更新 $\mathbf{W}' = \mathbf{W} + \mathbf{A}\mathbf{B}$，训练目标为流匹配损失 $\mathcal{L}_{\text{id}} = \mathbb{E}_{\tau, \mathbf{z}_0}[\|(\mathbf{z}_1 - \mathbf{z}_0) - \mathbf{v}_\phi(\mathbf{z}_\tau, \tau; \text{``a }[\text{v}]\text{ person''})\|^2]$。微调后的模型通过 RF-Inversion 将粗渲染视频编码到噪声潜空间再反向解码：$\mathbf{z}_1 = \mathbf{z}_0 + \int_0^1 \mathbf{v}_\phi(\mathbf{z}_\tau, \tau)\,d\tau$，$\tilde{\mathbf{z}}_0 = \mathbf{z}_1 + \int_1^0 \mathbf{v}_\phi(\mathbf{z}_\tau, \tau)\,d\tau$，$\tilde{\mathbf{V}}_n = \mathsf{Dec}(\tilde{\mathbf{z}}_0)$。前向过程编码粗渲染的空间结构（身体姿态和轮廓），反向过程通过 LoRA 将外观导向主体身份，用合理的纹理填充先前缺失的区域。结构化运动序列分为 turn 类型（$360^{\circ}$ 旋转配对目标动作，提供单姿态多视角监督）和 static 类型（坐姿等多角度视频，跨视频同步提供多视角覆盖）。

(3)~map-pose/LBS-pose 解耦。这是 AHOY 应对扩散数据多视角不一致性的关键设计。每个幻觉视频序列 $n$ 引入两种姿态：map pose $\bm{\theta}_n^{\text{map}}$ 作为 StyleUNet 的输入，在整个序列中保持不变，确保同一姿态下高斯属性预测的一致性；LBS pose $\bm{\theta}_{n,t}^{\text{lbs}}$ 逐帧优化（加上修正 $\Delta\bm{\theta}_{n,t}$），吸收扩散模型产生的帧间几何不一致性。身体监督损失为 $\mathcal{L}_{\text{body}} = \sum_{n,t} \|(\tilde{\mathbf{I}}_{n,t} - \hat{\mathbf{I}}_{n,t}) \odot (1 - \mathbf{M}_{n,t}^{\text{head}})\|_1$，其中头部像素被排除（扩散生成的面部身份不可靠）。

(4)~头部身份保持。通过 SMPL-to-FLAME 顶点对应将规范高斯分为头部和身体区域。头部区域使用 CAP4D 的多视角面部扩散模型生成身份一致的面部图像 $\{\mathbf{I}_p^{\text{face}}\}_{p=1}^{P}$，以 FLAME 表情参数 $\bm{\psi}_p$ 驱动变形：$\mathcal{L}_{\text{head}} = \sum_{p=1}^{P} \|\mathbf{I}_p^{\text{face}} - \mathcal{R}(\mathcal{T}_{\text{FLAME}}(\mathcal{G}^{\text{head}}; \bm{\psi}_p); \bm{\gamma}_p^{\text{face}})\|_1$。总目标函数为 $\mathcal{L} = \lambda_{\text{body}}\mathcal{L}_{\text{body}} + \lambda_{\text{obs}}\mathcal{L}_{\text{obs}} + \lambda_{\text{head}}\mathcal{L}_{\text{head}} + \lambda_{\text{per}}\mathcal{L}_{\text{per}} + \lambda_{\text{reg}}\mathcal{L}_{\text{reg}}$。

\textbf{设计动机}。幻觉即监督的核心动机是：对于大范围未观察的身体区域，"有监督但不完美"远优于"完全无监督"——扩散模型生成的内容虽然不一定完全正确，但提供了合理的几何和外观先验。两阶段规范-姿态依赖架构的动机在于：粗化身阶段仅有 $V{=}4$ 个规范视角的监督，不足以学习姿态依赖变形；幻觉视频提供了密集的多姿态多视角覆盖，才使升级到姿态依赖高斯图集成为可能。map-pose/LBS-pose 解耦的动机是：每个幻觉视频独立采样自扩散模型，跨视频甚至同一视频内不同帧之间不具备几何一致性；通过将"决定高斯属性"的 map pose 固定、"决定空间变形"的 LBS pose 逐帧调节，使得规范高斯保持一致的同时变形能够自适应每帧的几何偏差。

\textbf{公式与参数}。AHOY 在 YouTube 视频（含 Obama、Curry 等公众人物）和 BEHAVE 多视角数据集上评估，在动画质量和静态重建质量上均优于 LHM 和 IDOL 等前馈方法以及 SyncHuman 和 PSHuman 等扩散方法。推理时通过 PCA 姿态投影将新驱动姿态约束在训练分布内，防止分布外伪影。化身可驱动任意 SMPL 姿态并合成到由手机视频重建的 3DGS 场景中，实现端到端的从消费级 RGB 输入到照片级真实感人体-场景渲染。

上述方法均假设相机标定和人体姿态估计的输出是精确的——但这一假设在野外视频中常常不成立。

\paragraph{JOintGS}

JOintGS（2026.02）~\cite{2602_04317}（Joint Optimization of Cameras, Bodies and 3D Gaussians）正是针对这一前提的脆弱性，提出相机外参、人体姿态和 3D 高斯表示的统一联合优化框架，在规范空间中用三平面表示解码人体高斯属性，通过前景-背景分离实现协同精修——多视角光栅化渲染管线保持不变。

\textbf{Baseline 与问题}。现有 3DGS 化身方法（3DGS-Avatar、HUGS、GaussianAvatar、ExAvatar 等）严重依赖精确的相机标定和姿态标注作为固定输入。然而在野外单目视频场景中，COLMAP 等运动恢复结构工具因动态人体运动导致静态特征对应不足，产生噪声相机轨迹；HMR2.0 等单目姿态估计器因深度歧义和遮挡产生不准确的 3D SMPL 参数。即使是微小的初始误差也会级联放大为对齐不准、时序不一致和人体-场景穿透等严重伪影。

\textbf{改了什么}。JOintGS 将动态人体重建建模为统一优化问题，通过显式前景-背景分离实现三个互补路径的闭环协同精修：(1) 静态背景高斯 $\mathcal{G}_B$ 跨帧一致，为相机估计提供多视角几何约束；(2) 精修后的相机通过准确的时间对应改善人体姿态对齐；(3) 优化的人体姿态通过消除动态伪影增强场景重建。人体表示在规范空间中定义，通过三平面解码高斯属性（位置、旋转、缩放、不透明度、球谐系数和学习 LBS 权重），并引入时间偏移模块和残差颜色场分别捕捉非刚性几何变形和光照变化。

\textbf{实现原理}。核心架构包含四个关键组件。

(1)~规范高斯场与姿态驱动变形。在 SMPL 静止姿态（A-pose）的规范空间中均匀采样 $N_H{=}110\text{K}$ 个高斯，每个高斯的属性（位置 $\boldsymbol{\mu}_i^c$、旋转 $\mathbf{R}_i^c$、缩放 $\mathbf{S}_i^c$、不透明度 $\alpha_i$、球谐系数 $\mathbf{c}_i$ 和 LBS 权重 $\mathbf{w}_i \in \mathbb{R}^{K}$，$K{=}24$）从三平面表示中解码。LBS 权重从最近 SMPL 顶点通过重心插值初始化。给定帧 $t$ 的姿态 $\boldsymbol{\theta}_t$，通过 LBS 变形将规范高斯变换到姿态空间：$\boldsymbol{\mu}_i^t = \sum_{k=1}^{K} w_{i,k} \mathcal{T}_k(\boldsymbol{\theta}_t) [\boldsymbol{\mu}_i^c; 1]$，$\mathbf{R}_i^t = (\sum_{k=1}^{K} w_{i,k} \mathbf{R}_k(\boldsymbol{\theta}_t)) \mathbf{R}_i^c$。

(2)~时间动态建模。LBS 变形仅处理骨架驱动的刚体运动，无法捕捉衣物褶皱等非刚性动态和光照变化。时间偏移模块用浅层 MLP $\mathcal{F}_{\text{offset}}$ 以规范位置的多分辨率哈希编码和帧索引的位置编码为条件，预测逐帧位置和旋转偏移：$[\Delta\boldsymbol{\mu}_i^t, \Delta\mathbf{R}_i^t] = \mathcal{F}_{\text{offset}}(\mathcal{E}(\boldsymbol{\mu}_i^c), \gamma(t))$，最终姿态为 $\boldsymbol{\mu}_i^t \leftarrow \boldsymbol{\mu}_i^t + \Delta\boldsymbol{\mu}_i^t$，$\mathbf{R}_i^t \leftarrow \mathbf{R}_i^t \cdot \Delta\mathbf{R}_i^t$。残差颜色模块用另一个 MLP $\mathcal{F}_{\text{color}}$ 预测颜色残差：$\mathbf{c}_i^t(\boldsymbol{\mu}_i^c) = \mathbf{c}_i^{\text{SH}}(\boldsymbol{\mu}_i^c) + \mathcal{F}_{\text{color}}(\mathcal{E}(\boldsymbol{\mu}_i^c), \gamma(t))$。两个模块均通过正则化约束为小幅修正，确保仅捕捉骨架运动之外的残差动态。

(3)~协同精修机制。相机修正参数化为可学习的 $\Delta\mathbf{T}_t \in SE(3)$，作用于 COLMAP 初始值：$\mathbf{T}_t = \Delta\mathbf{T}_t \circ \hat{\mathbf{T}}_t$。三个互补路径形成闭环：背景锚定相机精修——用 SAM 生成人体掩码 $\mathbf{M}_t$，仅在静态区域最小化光度误差 $\mathcal{L}_B = \|(1-\mathbf{M}_t) \odot (\mathbf{I}_t - \hat{\mathbf{I}}_t^B)\|_1$，隔离相机优化与动态人体运动的影响；相机引导人体姿态优化——精修相机建立准确的时间对应后，通过光度损失和轮廓损失 $\mathcal{L}_H = \|\mathbf{M}_t \odot (\mathbf{I}_t - \hat{\mathbf{I}}_t^H)\|_1 + \|\mathbf{M}_t - \hat{\mathbf{M}}_t^H\|_1$ 优化 SMPL 参数，梯度通过可微分 LBS 变换直接流到姿态参数；姿态感知高斯优化——精修的相机和姿态改善前景-背景分离，减少前景泄漏到 $\mathcal{G}_B$ 和背景伪影出现在 $\mathcal{G}_H$，用 RGB、SSIM 和 LPIPS 组合损失优化高斯参数，更干净的场景分离反过来稳定相机估计，完成反馈闭环。

(4)~三阶段训练策略与正则化。Warm-up 阶段仅优化高斯参数以建立可靠几何先验；独立优化阶段同时开启相机和 SMPL 参数优化，但两者独立更新以避免梯度干扰；联合优化阶段使用完整损失进行协同精修。额外正则化包括：LBS 权重正则化 $\mathcal{L}_{\text{lbs}} = \sum_{i} \|\mathbf{w}_i - \mathbf{w}_i^{\text{SMPL}}\|_2^2$ 防止皮肤权重过拟合训练姿态；偏移正则化 $\mathcal{L}_{\text{offset}} = \sum_{i,t}(\|\Delta\boldsymbol{\mu}_i^t\|_2^2 + \|\Delta\mathbf{R}_i^t - \mathbf{I}\|_F^2 + \|\Delta\mathbf{c}_i^t\|_2^2)$ 惩罚大变形；规范正则化 $\mathcal{L}_{\text{canonical}} = \sum_i \min_j \|\boldsymbol{\mu}_i^c - \bar{\mathbf{v}}_j\|_2^2$ 将人体高斯软约束在 SMPL 网格表面附近。

\textbf{设计动机}。协同精修的核心动机是：相机、姿态和高斯三者的误差是相互耦合的——相机不准导致人体投影错误，姿态不准导致变形不合理，高斯不准导致渲染质量差——但它们之间的依赖关系也意味着精修其中一个可以改善其余两个。显式前景-背景分离的关键在于：背景高斯天然跨帧一致（静态场景），为相机估计提供纯净的多视角约束，不受动态人体运动干扰；而精修后的相机又为人体提供了更准确的时间对应。三阶段训练策略的动机是防止退化解：如果从一开始就联合优化所有参数，梯度干扰可能导致相机和姿态同时偏离正确方向。

\textbf{公式与参数}。在 NeuMan 数据集（6 个野外序列，手持手机拍摄 10--20 秒）上，JOintGS 在 NeuMan 上取得 2.1 dB PSNR 提升超过最强基线 HUGS，同时保持实时渲染性能。在 EMDB 数据集（81 个野外序列，58 分钟运动数据）上同样取得最优重建质量。鲁棒性方面，当噪声水平 $\sigma{=}0.01$ 时 JOintGS 仅下降 0.9 dB PSNR，而 HUGS 下降 3.7 dB，展现出对噪声初始化的显著鲁棒性。消融实验确认了联合优化策略和各组件的必要性。

\paragraph{Deblur-Avatar}

Deblur-Avatar（2025.01）~\cite{2501_13335} 处理运动模糊这一常见但被忽视的退化因素。

\textbf{Baseline 与问题}：上述方法（包括 AHOY、JOintGS）均假设输入视频帧是清晰的。然而在真实拍摄中，快速运动的人体（如舞蹈、体育场景）在曝光时间内会产生运动模糊——模糊方向取决于运动轨迹，模糊程度取决于运动速度和曝光时间。直接用模糊图像作为光度监督训练 3DGS 会导致化身吸收模糊模式，渲染结果永久模糊。

\textbf{改了什么}：Deblur-Avatar 将基于人体运动的运动模糊模型集成到 3DGS 优化循环中。核心思想是"正向生成模糊"：在曝光窗口内采样多个时间步的骨架姿态，对每个时间步进行 LBS 变换和光栅化，然后将多帧渲染结果平均合成模糊图像。生成的模糊图像与真实模糊输入比对，光度损失反向传播同时优化运动轨迹和 3D 高斯参数。

\textbf{设计动机}：与外部去模糊预处理不同，将模糊模型嵌入优化循环使化身直接学习清晰外观。因为模糊是可微渲染管线的一部分，梯度信号会引导高斯参数收敛到清晰状态。

\subsection{压缩流式与扩展应用}
\label{ss:16_2_压缩流式}

本节讨论~8~条技术路线：（1）压缩、流式传输和~VR~集成；（2）铰接物体重建与运动分析；（3）人体-物体交互；（4）场景中的人体动画；（5）数据增强与隐私保护；（6）泛化式方法；（7）基于先验和扩散的方法；（8）特殊场景的人体新视角合成。

\paragraph{CrowdGaussian}

CrowdGaussian（2026.03）~\cite{2603_17779} 从单图像扩展到多人场景。

\textbf{Baseline 与问题}：AHOY 处理单人的严重遮挡，但真实场景中（如演唱会、体育赛事）图像通常包含多个被互相遮挡的人。将单人方法逐个应用不可行，因为不同人之间的遮挡关系无法从单张图像中分离。

\textbf{改了什么}：CrowdGaussian 分两阶段进行多人 3D 重建。粗略群体完成阶段用多人 HMR 估计姿态并分割个体，再通过 LORM（Learned Object Reconstruction Module）为每个被遮挡的个体生成粗 3DGS 表示。扩散精修阶段用单步扩散模型配合自校准学习策略（Self-Calibrated Learning, SCL）自适应保留已恢复区域同时增强细节不足区域，从严重遮挡的输入中重建完整的多人 3DGS 表示。

\textbf{设计动机}：两阶段设计将"粗略形状恢复"和"细节增强"解耦——粗阶段利用姿态先验快速建立每个人的基本几何，精修阶段用扩散模型的生成能力为每个个体补全被遮挡区域。SCL 的自适应机制避免了扩散模型过度修改已有良好恢复的区域。

\subsubsection{压缩、流式传输和VR集成}

本路线介绍~3DGS~化身的模型压缩、流式传输和~VR~实时部署方法。

化身在实用化部署中面临存储和传输的挑战。前文讨论的方法大多关注渲染质量和动画能力，但在实际部署中，模型大小、推理速度和传输带宽同样是关键约束。本节从模型压缩（GEM）、计算调度（LoDAvatar）到实时 VR/远程呈现部署，逐步推进实用化链条。

Gaussian Eigen Models（GEM）（2024.07）~\cite{2407_04545} 从模型压缩的角度重新审视化身表示。

\textbf{Baseline 与问题}：上述所有化身方法在推理时仍需 CNN 解码器（如 StyleUNet）或 MLP 网络来预测高斯属性，这些网络的参数量构成了存储和推理的主要开销。对于移动端或 VR 应用，部署完整的 CNN 解码器是不可接受的。

\textbf{改了什么}：GEM 将化身外观表示为一组线性特征基（类似 PCA 降维的思想），通过低维参数向量的线性组合高效生成高斯基元。具体地，先用高质量 CNN 模型（如 Animatable Gaussians 的改编版本）为每帧生成 3D 高斯参数，然后对帧序列中的高斯属性分别执行 PCA，得到位置、旋转、缩放和不透明度的特征基 $\mathbf{B}_\theta, \mathbf{B}_\phi, \mathbf{B}_\alpha, \mathbf{B}_\sigma$ 和对应的均值 $\boldsymbol{\mu}$。推理时，人脸实例 $\mathbf{G} = \{\boldsymbol{\mu}_i + \mathbf{B}_i \mathbf{k}_i \mid i \in \{\theta, \phi, \alpha, \sigma\}, \mathbf{c}\}$ 由低维系数 $\mathbf{k}$ 的线性组合生成，颜色 $\mathbf{c}$ 作为全局纹理直接优化。

\textbf{设计动机}：GEM 在测试时完全不需要 CNN 或 3DMM——仅通过 PCA 特征基的线性组合和轻量级图像回归网络（ResNet，约 20--50 个系数）即可从单张图像驱动化身。这使高质量化身被压缩为轻量级模型，适合实时移动端部署。特征基通过 QR 分解定期正交化来精修，确保线性组合的数值稳定性。

在 VR 场景中同时渲染多个化身时，化身高斯数量直接决定帧率——高斯越多视觉质量越高，但计算成本也越大。

\paragraph{LoDAvatar}

LoDAvatar（2024.10）~\cite{2410_20789} 首次将多层次细节（Level of Detail, LoD）概念引入可驱动高斯化身，通过控制高斯数量来平衡视觉质量与计算成本。其核心流程分为四步：数据准备、高斯嵌入、高斯优化和选择性细节增强。在高斯嵌入阶段，方法在 SMPL 网格的每个三角形面片上建立局部坐标系（以面片中心为原点，法线方向为 $z$ 轴），并在三个顶点和面片中心各初始化一个高斯，将高斯参数从世界坐标系变换到局部坐标系，使其随网格变形而同步运动。优化阶段固定顶点处高斯的位置参数，仅优化面片中心高斯及其他参数（旋转、缩放、颜色、不透明度），且不使用 3DGS 原始的自适应密度控制——既不分裂也不剪枝，保持高斯数量恒定。优化完成后，将面片中心高斯的优化后位置与三角形顶点连接，形成三个新三角形面片，并在此之上继续初始化和优化新的高斯。通过 3--4 次迭代，即可从同一 SMPL 网格生成从 75k 到 1.17M 个高斯的化身系列（LoD1 到 LoD3）。

此外，LoDAvatar 引入选择性细节增强机制：利用 SAM 等图像分割算法生成面部和手部的掩码，将网格顶点投影到图像平面以识别对应三角形面片，仅在这些区域执行额外的三角形细分和高斯嵌入。这一策略使 LoD2+选择性细节增强版（约 379k 高斯）在 VR 近距离交互中达到与 Gaussian Avatars（182k 高斯）和 Splatting Avatars（412k 高斯）相当的视觉质量，而总优化时间（22.5 小时）少于后两者（分别为 23.4 和 32.0 小时），且低 LoD 版本在远距离观察时帧率显著提升。

在 VR 应用方面，3DGS 化身需要满足更严格的实时性和交互性要求。

\paragraph{VRGaussianAvatar}

VRGaussianAvatar（2026.02）~\cite{2602_01674} 面向 VR 头显的实时全身化身需求。

\textbf{Baseline 与问题}：VR 场景的特殊约束是：(1)~输入仅来自 HMD 的 IMU 追踪信号（6 自由度头部姿态+手柄位置），缺乏全身骨骼的完整观测；(2)~立体渲染要求同时生成左右眼视角，计算量近似翻倍；(3)~延迟要求极严格（$<$20ms 以避免晕动症）。

\textbf{改了什么}：VRGaussianAvatar 实现了仅使用 HMD 追踪信号的实时全身 3DGS 化身。系统采用 VR 前端和 GA 后端的并行架构：前端从 HMD 传感器估计全身姿态（从稀疏的头部+手部信号推断全身关节位置），后端通过双目批处理技术（Binocular Batching）将左右眼视角联合处理以减少冗余计算——避免对同一组高斯进行两次独立的光栅化，而是复用中间结果。

\paragraph{Mon3tr}

Mon3tr（2026.01）~\cite{2601_07518} 面向远程呈现的带宽效率问题。

\textbf{Baseline 与问题}：现有 3DGS 远程呈现系统需要传输完整的高斯参数集（位置、颜色、形状等），带宽需求达数十 Mbps，无法在普通网络连接上实时传输。

\textbf{改了什么}：Mon3tr 提出基于 3DGS 参数化人体建模的单目 3D 远程呈现框架，通过 WebRTC 数据通道以低于 0.2 Mbps 的带宽传输运动和外观特征（仅传输低维的运动参数和轻量级外观编码，而非完整高斯参数），在接收端以约 60 FPS 合成逼真化身。

\paragraph{Animated 3DGS Avatars}

Animated 3DGS Avatars（2026.01）~\cite{2601_01660} 提出了深度高斯阴影图（Deep Gaussian Shadow Maps, DGSM）技术，解决了可动画化身插入静态 3DGS 场景时的光照一致性问题。

\textbf{Baseline 与问题}：当可动画的 3DGS 化身被放入静态 3DGS 场景中时（例如将人体化身放入扫描重建的室内场景），光照不一致的问题变得突出——化身看起来像是"漂浮"在场景中，既不接收场景中的阴影，也不向场景投射阴影。传统解决方案需要将 3DGS 转换为网格表示后使用经典阴影映射算法，但这会丢失 3DGS 体积表示的细节优势。

\textbf{改了什么}：DGSM 将经典深度阴影映射（deep shadow mapping）的思想迁移到体积 3DGS 表示中，利用 3DGS 沿光线方向累积透射率的闭式解实现体积阴影计算，完全在 3DGS 表示内操作而无需网格化。

\textbf{实现原理}：方法包含两个核心机制。(1)~深度高斯阴影图：首先从场景高斯的球谐系数中估计 $k$ 个主要光源位置，然后以每个光源为视点将场景高斯光栅化。

与经典阴影映射记录最近深度不同，DGSM 沿光线方向在每个同心径向壳层（concentric radial shell）上列表化累积透射率，并将结果存储在八面体图集（octahedral atlas）中——现代 GPU 可实时查询这些图集来衰减受阴影影响的高斯，从而实现投射和接收阴影的一致性。3DGS 的体积表示恰好支持沿光线的透射率累积的闭式计算，无需体素化或网格化。(2)~SH 环境光照重光照：为使移动的化身与场景的环境光照一致，方法用 HDRI 探针表示局部环境光照（以球谐基函数展开），并对每个化身高斯应用快速逐高斯辐射传输（per-Gaussian radiance transfer），避免了显式 BRDF 估计或离线优化。

\textbf{设计动机}：3DGS 的体积表示天然适合计算沿光线的累积不透明度——这正是阴影映射所需的核心量。经典深度阴影映射需要处理半透明物体的阴影（如头发、烟雾），在网格表示中需要额外数据结构，而在 3DGS 中只需查询高斯的 $\alpha$ 值沿光线的加权累积即可。

\textbf{公式与参数}：方法在 AvatarX 和 ActorsHQ 化身与 ScanNet++、DL3DV、SuperSplat 场景的合成实验中验证了 DGSM 和 SH 重光照的有效性，支持单化身和多化身设置下的协调阴影与重光照，全程在体积 3DGS 表示中完成。

除全身化身外，3DGS 还被广泛应用于铰接物体重建、人体-物体交互和场景中的人体动画等方向。这些方向虽然不以"化身驱动"为核心目标，但它们与化身技术共享底层表示和方法论，形成了紧密的研究生态。

\subsubsection{铰接物体重建与运动分析}

本路线介绍基于~3DGS~的铰接物体重建与运动参数估计方法。

铰接物体（如笔记本电脑、抽屉、橱柜等具有可移动部件的物体）的 3D 重建在机器人操作和 AR/VR 中具有重要价值。与人体化身不同，铰接物体的挑战在于：运动模式更简单（通常只有 1--2 个自由度的旋转或平移），但缺少成熟的参数化先验模型（如人体有 SMPL），需要从观测数据中自动发现运动结构。

\paragraph{SplArt}

SplArt（2025.06）~\cite{2506_03594} 提出了自监督、类别无关的铰接物体重建框架。

\textbf{Baseline 与问题}：此前的铰接物体重建方法需要部件级标注或已知铰接类型，且联合优化重建和铰接估计容易陷入局部最优。

\textbf{改了什么}：SplArt 为每个高斯增加一个可微移动性参数（mobility parameter），实现精细的部件分割，并通过解耦的多阶段优化策略将重建、部件分割和铰接估计逐步处理。

\textbf{实现原理}：方法假设物体由静态父部件和可移动子部件组成，子部件通过旋转或平移铰接相对于父部件运动。核心流程分三个阶段：阶段 1 对铰接的两个端状态分别独立优化标准 3DGS 模型；阶段 2 引入跨静态公式（cross-static formulation），通过比较两个独立重建的高斯属性差异，估计每个高斯的移动性参数（0 表示静态、1 表示可动），实现部件级的软分割；阶段 3 使用跨移动公式（cross-mobile formulation），利用移动性参数将可动部件的高斯从一个状态变换到另一个状态，同时优化铰接参数（旋转轴 $a$、枢轴点 $p$、旋转角度 $\theta$），并通过渲染损失精修移动性。

\textbf{设计动机}：将重建和铰接估计解耦到不同阶段是避免局部最优的关键——先确保每个状态的重建质量，再比较差异识别可动部件，最后估计运动参数。

\textbf{公式与参数}：旋转铰接建模为 $^1T_0 x = R_{a,\theta}(x-p) + p$，平移铰接建模为沿轴 $a$ 的平移 $d$。

\paragraph{ArtGS}

ArtGS（2025.02）~\cite{2502_19459} 利用规范高斯和粗到细初始化对齐不同物体状态的铰接部件信息，通过蒙皮启发的部件动态建模模块改进部件网格重建和铰接学习。

\paragraph{FreeGaussian}

FreeGaussian（2024.10）~\cite{2410_22070} 是首个免标注的铰接物体可控重建方法，其核心创新在于通过微分分析解析推导动态高斯流（dynamic Gaussian flow，即每个高斯基元在连续帧间的 3D 运动向量场）。

\textbf{Baseline 与问题}：此前方法需要 mask 标注才能识别动态区域。

\textbf{改了什么}：FreeGaussian 建立了光流、相机运动与高斯流之间的数学关系（Lemma 1），从而将相机自运动（平移速度 $\mathbf{v}$ 和旋转速度 $\boldsymbol{\omega}$）与铰接运动在数学上完全解耦，无需 mask 即可自动发现可交互结构。具体流程为：先从连续帧间的光流和相机自运动推导动态高斯流，再通过 HDBSCAN 聚类获取运动轨迹实现部件分割。在此基础上，FreeGaussian 引入 3D 球面向量控制方案（3D spherical vector control），利用高斯轨迹作为状态表示直接在 3D 空间定义控制参数，与 splatting 光栅化管线天然对齐，替代了传统需要手动定义的 1D 控制信号。

\paragraph{GEAR}

GEAR（2026.04）~\cite{2604_07728} 从另一个角度解决 SplArt 遗留的优化问题。

\textbf{Baseline 与问题}：直接联合优化部件分割和运动参数会导致误差归因模糊——优化器倾向于扭曲几何来补偿分割错误。

\textbf{改了什么}：GEAR 将 EM 算法思想引入铰接建模，采用交替优化策略：E 步固定运动参数，优化每个高斯的可学习掩码向量 $\mathbf{m}_i$（即部件分割，视为隐变量）；M 步固定分割结果，优化各部件的 SE(3) 运动参数（视为显变量）。为规范态构建两组高斯 $G_0$ 和 $G_1$，变换过程 $\mathbf{g}_i' = T(\sum_k m_{i,k} \cdot T_k)\mathbf{g}_i$ 将规范态高斯变换到目标状态进行渲染验证。GEAR 还利用基础 2D 分割模型提供多视角部件先验作为弱监督约束，无需在分布内数据上微调即可泛化到复杂多关节物体。

\paragraph{FreeArtGS}

FreeArtGS（2026.03）~\cite{2603_22102} 是首个在自由运动场景下重建铰接物体的方法，仅需单目 RGB-D 视频输入，无需假设物体具有静态基座或预定义运动模式。

\textbf{Baseline 与问题}：此前方法（包括 SplArt、FreeGaussian）假设物体至少有一个部件静止或运动模式已知，无法处理完全自由运动的物体。

\textbf{改了什么}：FreeArtGS 将重建分解为三个模块：（1）自由运动部分分割——利用 AllTracker 的密集 2D 轨迹和 DINOv3（DINO系列的第三代自监督视觉Transformer）特征进行优化，在滑动窗口内对每对帧求解两个刚体变换和软部件权重，通过 Huber 损失优化刚体变换，并用熵惩罚和空间相干正则化确保分割质量；（2）关节估计——为每个部分独立重建初始高斯，通过部分间相对变换确定关节类型（旋转或平移）和关节轴；（3）基于 3DGS 的端到端优化——使用混合渲染联合微调外观、几何、相机和关节参数。

\paragraph{ArtSplat}

ArtSplat（2026.05）~\cite{2605_24304} 是首个前馈式铰接 3DGS 框架，从稀疏多视角多状态图像在单次前向传播中同时重建几何和关节参数，速度比基线快 400 倍以上。

\textbf{改了什么}：ArtSplat 基于 VGGT 架构扩展，将关节参数建模为 11 维逐像素密集表示（关节类型 logits、轴方向、枢轴位置、旋转角度、平移位移），完全可微且端到端训练。输入为覆盖两个铰接状态的稀疏多视角图像，经 DINOv2（自监督视觉Transformer）patch 化编码后加入可学习的相机和状态 token，通过交替帧注意力和全局注意力处理。关键设计是跨状态注意力（cross-state attention），带状态 token 的注意力块捕获两个铰接状态间的离散运动差异。每个像素同时预测深度、高斯参数和关节参数，形成状态条件化高斯集。

\subsubsection{人体-物体交互}

本路线介绍基于~3DGS~的人体-物体交互重建与接触建模方法。

人体与物体的交互重建需要同时处理人体运动、物体运动和接触关系。这比单独重建人体或物体更困难，因为两个动态实体之间的接触区域会产生复杂的光遮挡和几何穿插。

\paragraph{GASPACHO}

GASPACHO（2025.03）~\cite{2503_09342} 修改了 3DGS 的物体表示与训练策略——用 2D 表面流形上的高斯图（Gaussian Maps）替代 3D 体积中的自由高斯学习，用遮挡感知的组合损失替代标准光度损失——基于 SMPL 的人体 LBS 蒙皮变形和基于 6D 刚体变换的物体运动建模保持不变。

\textbf{Baseline 与问题}。读者已经从前文 SplArt 了解到铰接物体重建需要参数化运动模型。GASPACHO 面临的核心挑战更进一步：在人体-物体交互中，物体通常被人体严重遮挡（即使多相机设置下，大量视角中的物体仍不可见），直接使用标准 3DGS 优化物体时无法收敛；此外，先前方法要么将物体当作背景处理（无法动画化），要么假设物体静态不动，均无法实现可控制的人-物交互合成。

\textbf{改了什么}。GASPACHO 引入三项关键修改：（1）用基于三平面特征（triplane feature）的参数化替代直接优化 3DGS 参数来学习物体模板——特征充当正则化和平滑先验，使物体在严重遮挡下仍可重建；（2）用 2D 高斯图（Gaussian Maps）替代 3D 体积中的自由高斯学习——将物体高斯属性（颜色、协方差、尺度、不透明度）参数化为 2D UV 图像上的像素值，由 StyleUNet 网络预测与规范模板的偏移量；（3）用遮挡引导的组合损失替代标准光度损失——分别渲染人体和物体高斯，并用分割掩码屏蔽被另一实体遮挡的像素区域。

\textbf{实现原理}。整个流程分为粗到细的三个阶段。

第一阶段学习物体规范模板：在所有时间帧中选择物体遮挡最少的帧 $t^*$，在该帧上通过三平面特征参数化优化物体高斯。具体而言，学习三个 $M \times W$ 维的特征平面 $\mathbf{F} \in \mathbb{R}^{3 \times M \times W \times l}$（$l$ 为特征维度），每个物体高斯位置 $\mathbf{x}_p$ 通过正交投影采样三个平面获取拼接特征 $\mathbf{f}_p = \pi(\mathbf{x}_p; \mathbf{F}) \in \mathbb{R}^{3l}$，再通过小型 MLP 映射为高斯参数。优化目标为 $\mathcal{G}_O^* = \arg\min_{\mathcal{G}_O} \sum_{c=1}^{N} \|\mathcal{R}(\mathcal{G}_O; \gamma_c) - \mathbf{O}_{t^*}^c \cdot \mathbf{I}_{t^*}^c\|_2$，即仅使用该帧中物体可见像素进行监督。简单来说，特征平面为被遮挡的高斯提供了跨空间的信息插值，避免因观测不足而退化。

第二阶段进行物体追踪：利用学到的模板 $\mathcal{G}_O^*$，通过可优化的 6D 刚体变换 $\boldsymbol{\phi}_t$（旋转 $\mathbf{R}_{\boldsymbol{\phi}}$ 和平移 $\mathbf{t}_{\boldsymbol{\phi}}$）将模板变换到每一帧的姿态，并最小化渲染误差来估计每帧的物体位姿。变换公式为 $\mathbf{x}' = \mathbf{R}_{\boldsymbol{\phi}} \mathbf{x} + \mathbf{t}_{\boldsymbol{\phi}}$，$\boldsymbol{\Sigma}' = \mathbf{R}_{\boldsymbol{\phi}} \boldsymbol{\Sigma} \mathbf{R}_{\boldsymbol{\phi}}^\top$——与 SplArt 的部件变换公式形式类似，但这里作用于整个物体而非单个部件。为确保收敛，每帧位姿 $\boldsymbol{\phi}_t$ 从上一帧 $\boldsymbol{\phi}_{t-1}$ 初始化，形成时序上的连续优化。

第三阶段使用高斯图精修：给定人体规范模板（基于 SMPL 网格密集采样）和物体规范模板，将两者的规范高斯分别从前、后两个视角投影到 2D UV 图上，建立像素与 3D 高斯的一一对应。人体使用姿态依赖的位置图（通过 LBS 变形 SMPL 顶点后存储到 UV 图），物体使用姿态无关的固定位置图（直接存储规范空间中的 3D 坐标）。两个 StyleUNet $S_H$、$S_O$ 分别预测高斯属性（位置偏移、颜色、协方差、不透明度）与规范模板的偏差；人体 StyleUNet $S_H$ 还额外预测逐高斯蒙皮权重 $\mathbf{w}_k^H$，用于后续 LBS 变形。这里的关键洞察是：物体的姿态无关图（pose-independent map）使 6D 位姿在训练中可被持续精修，而不会与外观学习相互干扰——论文实验证实这一设计对物体纹理质量至关重要。

\textbf{遮挡引导的组合损失}。对于每个相机和时间步，分别渲染人体 $\mathbf{I}_H^{ct}$、物体 $\mathbf{I}_O^{ct}$ 和组合 $\mathbf{I}_A^{ct}$，并使用分割掩码屏蔽遮挡区域：人体损失仅在物体掩码外的像素上计算（$\mathcal{L}_H^{ct} = \|\mathbf{I}_t^c \mathbf{H}_t^c (1 - \mathbf{O}_t^c) - \mathbf{I}_H^{ct} (1 - \mathbf{O}_t^c)\|_1$），物体损失仅在人体掩码外计算（$\mathcal{L}_O^{ct}$ 形式对称），组合损失则在人体加物体掩码内计算。这避免了在人体遮挡物体的像素上错误地惩罚物体渲染。除掩码 $L_1$ 损失外，最终目标还包含感知损失 $\mathcal{L}_{per}$（使用相同掩码策略）和蒙皮权重正则化 $\mathcal{L}_{reg}$（鼓励预测的蒙皮权重接近 SMPL 推导值），即 $\mathcal{L} = \lambda_{L1}\mathcal{L}_1 + \lambda_{per}\mathcal{L}_{per} + \lambda_{reg}\mathcal{L}_{reg}$。多视角一致性确保被忽略的高斯在其他帧或视角中仍能得到监督。

\textbf{高斯空间接触约束}。在推理阶段合成新人-物交互时，朴素地组合两套高斯会导致穿模或缺失接触。GASPACHO 定义了 SMPL 接触顶点集合 $V_c^{\text{SMPL}}$（脚、臀、手），通过在规范空间中的最近邻搜索识别对应的人体接触高斯索引集 $\mathcal{I}_c$，并优化接触高斯的位移 $\boldsymbol{\Delta}$。接触帧通过 SMPL 运动线索（低垂直速度和向下加速度）自动检测：

\begin{equation}
\mathcal{L} = \lambda_c \sum_{k \in \mathcal{I}_c} \|\mathbf{x}_{k,t}^H - \mathbf{x}_{p(k,t)}^O\|_2^2 + \lambda_p \sum_{k \notin \mathcal{I}_c} h_r(d_\beta(\mathbf{x}_{k,t}^H, \mathcal{G}_O^t))^2 + \lambda_r \|\boldsymbol{\Delta}\|_2^2 + \lambda_t \|\boldsymbol{\Delta}_t - \boldsymbol{\Delta}_{t-1}\|_2^2
\end{equation}

直觉上：第一项将接触高斯拉向最近的物体表面（使手"抓住"物体），第二项通过铰链间隔 $h_r(d) = \max(0, r - d)$ 将非接触高斯推离物体至少距离 $r$，第三、四项分别惩罚过大的位移和时序不连续。软最近邻距离 $d_\beta(\mathbf{x}) = -\frac{1}{\beta} \log \sum_{p} \exp(-\beta \|\mathbf{x} - \mathbf{x}_p^O\|)$ 提供了对物体高斯集的可微距离度量。

\textbf{设计动机}。三平面特征参数化的动机是直接 3DGS 在遮挡下不收敛（论文实验证实），而特征平面提供跨空间的信息平滑和正则化。高斯图的动机是利用 CNN 的强归纳偏置并建立像素与高斯的一一对应，使物体位姿优化更稳定。接触约束的动机是：与 SMPL 网格不同，高斯化身没有固定拓扑对应关系，无法手动指定接触点，因此需要通过最近邻搜索自动识别接触高斯。该方法在 BEHAVE（4 相机稀疏设置）、NeuralDome 和 DNA-Rendering（60 相机密集设置）三个基准上均取得了最佳重建质量，并首次实现了从不同数据集重建的人体和物体在 3D 场景中的可控制组合。

\paragraph{HOIGS}

HOIGS（2026.04）~\cite{2604_04016} 是首个显式建模人-物交互的 4DGS 框架。

\textbf{Baseline 与问题}：此前的人-物交互方法将人体和物体分别重建后简单组合，忽略了两者之间因接触和操作引起的相互变形影响。

\textbf{改了什么}：HOIGS 为人体和物体设计不同的变形基线以适配各自的动态特性——物体使用三次 Hermite 样条（CHS）建模连续平滑运动轨迹，确保位置和速度连续性，可学习控制点表示关键帧平均位置；人体使用 HexPlane 特征，通过时间不变空间特征学习规范 T-pose 纹理，再经 LBS 变形到世界空间。两者的运动特征随后输入交叉注意力 HOI 模块，通过交叉注意力回归交互感知的几何修正，捕捉双向运动依赖关系（如人拿起杯子时手指对杯子的推动、杯子对手部皮肤的压迫），最终通过集成目标函数联合优化整个场景。

\paragraph{PhyGenHOI}

PhyGenHOI（2026.05）~\cite{2605_30268} 将生成式人体运动与显式物理物体模拟统一在 3DGS 表示下。

\textbf{Baseline 与问题}：纯生成式方法产生的人-物交互缺乏物理约束，容易出现鬼影和穿透伪影，物体运动轨迹不符合物理定律。

\textbf{改了什么}：PhyGenHOI 将场景中的每个实体建模为不同类型的智能体——人体是由运动扩散模型（MDM）驱动的语义智能体，物体是通过材料点方法（MPM）模拟的物理智能体。两者之间的协调通过三个耦合机制实现：（1）窗口吸引损失（windowed attraction loss）在时间窗口内将人体运动轨迹吸引到物体位置，使人适时接触物体；（2）接触驱动再模拟（contact-driven re-simulation）在检测到接触时触发碰撞动量转移，使物体被推动后的运动符合物理定律；（3）掩码视频-SDS 目标通过 2D 视频先验蒸馏优化整体视觉质量。这种"生成驱动人、物理驱动物"的混合策略消除了纯生成方法的鬼影和穿透伪影。

\subsubsection{场景中的人体动画}

本路线介绍将~3DGS~人体化身融入~3D~场景并生成自然交互动画的方法。

将人体化身放置在 3D 场景中并实现自然的交互动画是实际应用的关键需求。一个"悬空"在场景中或与家具穿模的化身会立即破坏沉浸感。

\paragraph{AHA}

AHA（2025.11）~\cite{2511_09827}（Animating Human Avatars）解决了将高斯人体化身直接放置到 3DGS 重建场景中并生成自然交互动画的问题。

\textbf{Baseline 与问题}。传统人体动画流水线通常将运动合成与场景渲染分开处理，缺乏对场景几何的感知：合成的步态可能与地形不匹配，手可能穿过桌面，脚可能悬空在地板之上。当场景和人体都用 3D 高斯表示时，问题更加棘手——3DGS 没有显式网格表面，无法直接使用碰撞检测或接触约束，必须设计新的高斯-高斯交互机制。

\textbf{改了什么}。AHA 提出将渲染与运动合成解耦的框架，包含三项关键设计：（1）从多视角视频学习可变形的人体高斯表示，通过 SMPL 姿态驱动动画；（2）高斯对齐的运动合成模块，利用强化学习和潜变量优化在 3DGS 场景中生成与场景几何对齐的运动参数；（3）可微分接触精修算法，自动检测接触帧并优化平移向量，使人-场景高斯之间产生正确的接触而避免穿模。

\textbf{实现原理}。整个流程分为三个阶段。

第一阶段进行高斯重建。场景高斯通过标准 3DGS 从 RGB 图像重建，得到 $N_S$ 个各向异性高斯 $G^S = \{(\boldsymbol{\mu}_i, \boldsymbol{\Sigma}_i, \alpha_i, c_i)\}_{i=1}^{N_S}$。人体高斯则采用基于 SMPL 的规范参数化：在 SMPL 网格的前后 UV 图集上放置高斯，通过姿态条件图 $P_t \in \mathbb{R}^{H \times W \times C}$ 和 StyleUNet 预测规范高斯 $G^C_t = f_\theta(P_t)$，再经线性混合蒙皮（LBS）变换到姿态空间，得到 $G^P_t = \text{LBS}(G^C_t)$。推理时，新姿态通过 PCA 子空间映射和同样的 StyleUNet+LBS 流程生成姿态化高斯。

第二阶段进行高斯对齐的运动合成。该模块以场景高斯为环境，在潜变量空间中使用强化学习策略合成运动参数。核心思路是将运动参数的优化转化为在高斯场景中的导航问题：运动模块的奖励信号来自渲染质量反馈，使合成的步态自动适应场景中的地面高度、障碍物位置等几何特征。这避免了传统方法需要显式场景几何（如网格或深度图）来约束运动的限制。

第三阶段进行可微分接触精修。运动合成给出的人体高斯序列在接触帧处可能存在穿模或悬空。AHA 检测出接触帧，并为接触帧中的人体高斯优化一组平移向量 $\boldsymbol{\Delta}$，目标函数平衡三个因素：接触高斯与场景表面的距离最小化（使脚"踩实"地面）、非接触高斯与场景的穿模惩罚（避免身体穿入家具）、以及时间连续性约束（避免逐帧抖动）。由于 3DGS 没有显式表面，接触检测通过高斯间距离和密度阈值实现，穿模惩罚则基于软最近邻距离 $d_\beta(\mathbf{x}) = -\frac{1}{\beta} \log \sum_p \exp(-\beta \|\mathbf{x} - \mathbf{x}_p^S\|)$ 度量人体高斯到场景高斯集的可微距离。

\textbf{设计动机}。将渲染与运动合成解耦的动机是：场景重建和人体动画本身已各自成熟，但两者的结合需要在运动层面（而非渲染层面）对齐。高斯对齐运动模块的动机是利用 3DGS 的可微分渲染作为运动优化的监督信号，天然适配高斯-高斯的场景设定。接触精修作为后处理的动机是：端到端的运动合成难以保证每帧的物理正确性，而专门的精修步骤可以聚焦处理接触帧，计算量远小于全程约束。

\paragraph{ODHSR}

ODHSR（2025.04）~\cite{2504_13167} 提出首个仅需单目 RGB 视频即可同时执行相机追踪、人体姿态估计和密集光真实人体-场景重建的在线框架，比先前工作快 75 倍。

\textbf{Baseline 与问题}：此前的人体-场景重建需要多视角输入或分步处理（先重建场景再植入人体），无法从单目视频在线完成。

\textbf{改了什么}：ODHSR 基于 3DGS 设计统一 SLAM 管道，将人体化身在规范空间中表示为 3D 高斯集合，通过 SMPL 驱动的 LBS 变形到观测空间，并将变形分解为刚性变换和非刚性位移两部分以处理动态服装。非刚性变形采用多分辨率哈希编码网络建模，同时建模姿态相关的环境光遮蔽。场景使用独立的 3DGS 表示。系统结合 Depth Anything V2 单目几何先验和 SMPL 人体先验进行初始化，引入遮挡感知的人体轮廓项辅助人体和场景重建的分解，仅维护少量关键帧以保证高效在线映射。

\paragraph{Make-It-Animatable}

Make-It-Animatable（2024.11）~\cite{2411_18197} 是数据驱动方法，在一秒内使任何 3D 人形模型准备好角色动画，统一生成高质量的混合权重、骨骼和姿态变换。

\textbf{Baseline 与问题}：传统角色绑定（rigging）是高度手工的过程，需要艺术家为每个 3D 模型手动创建骨骼结构和蒙皮权重，且绑定结果与输入表示强耦合（网格和 3DGS 各需不同的绑定管线）。

\textbf{改了什么}：Make-It-Animatable 的核心是粒子形状自编码器（particle shape autoencoder），将任意 3D 表示（网格或 3DGS）统一为点云表示——通过 FPS 下采样和局部自注意力编码为紧凑潜在特征，解码器使用 cross-attention 从潜在特征预测每个粒子的动画属性。方法采用从粗到精策略：先预测粗粒度骨骼属性（关节位置和变换），再预测细粒度蒙皮权重，并通过骨骼感知注意力机制确保相邻骨骼的一致性。其理论框架将动画视为低秩分解——时间基为关节变换矩阵 $\Phi_t \in \mathbb{R}^{T \times K \times d}$，空间权重为蒙皮权重矩阵 $\Phi_s \in \mathbb{R}^{K \times N}$，其中 $K$ 为骨骼数量，$d$ 为自由度（刚体变换为 6），$N$ 为粒子数。该方法支持非标准骨骼结构（如手、耳朵、尾巴等额外骨骼），在大规模 3D 角色数据集上训练。

\subsubsection{数据增强与隐私保护}

本路线介绍利用~3DGS~进行训练数据增强和隐私保护的方法。

在数据层面，3DGS 不仅可以用于重建和渲染，还可以通过可微分渲染生成高质量的训练数据或保护化身的身份隐私。

\paragraph{ISExplore}

ISExplore（2025.11）~\cite{2511_07940}（Informative Segment Explore）揭示了个性化说话人头像训练中一个反直觉的发现：精心选择的几秒参考片段即可达到与完整视频相当的合成质量。

\textbf{Baseline 与问题}。主流 3DGS 和 NeRF 说话人头像方法通常需要数分钟的参考视频从头训练个性化模型，训练耗时长且数据处理成本高。一个自然的疑问是：参考视频中的每一秒对模型学习同等重要吗？如果不是，是否存在少量"高信息量"片段足以支撑高质量训练？

\textbf{改了什么}。ISExplore 将参考视频选择重新表述为数据选择问题，提出三阶段自动片段选择策略：先用音频特征多样性对所有候选片段排序，再对排名靠前的候选评估运动复杂度，最后通过信息量评分选择最优片段。该方法无需修改底层 TFG 模型本身，可作为即插即用的数据预处理模块。

\textbf{实现原理}。整个流程分为三步。

第一步进行音频多样性排序。将参考视频以 1 秒步长切分为 $n$ 秒等长片段，构成候选集 $X = \{x_1, x_2, \ldots, x_N\}$。对每个片段，用 AVE 编码器提取帧级音频嵌入，计算所有特征向量之间的平均欧氏距离作为多样性得分 $D$。论文通过预备实验比较了四种候选度量（最大主成分方差、累积方差、语义熵、欧氏距离），发现欧氏距离与最终唇形同步质量的相关性最强（$p = 0.0019$，$R^2 = 0.4653$）。按 $D$ 降序排列后保留前 $m$ 个候选。

第二步进行运动复杂度估计。对每个保留的候选片段，提取 3DMM 唇部特征点（第 48--68 号点），分别建模两类运动：局部唇形变形（嘴部开合）和全局特征点位移（头部运动与视角变化）。将特征点轨迹通过离散傅里叶变换（DFT）转换到频域，定义高频比（HFR）为高频分量幅度之和与全频段幅度之和的比值，作为运动复杂度得分 $MC$：当阈值 $T = 0.25$ 时，样本分布最均匀，区分最稳定。

第三步计算信息量评分并选择。核心公式为 $I = D / (MC + 10^{-8})$——信息量定义为音频多样性与运动复杂度的比值。直觉上，理想片段应包含丰富的音频内容（提供充足的唇形监督信号），同时避免过于复杂的运动模式（降低模型拟合难度）。选择信息量最高的片段 $x_{\text{ISExplore}} = \arg\max_{x_i} I(x_i)$ 作为训练数据。

\textbf{设计动机}。三个维度分别解决不同问题。音频多样性解决"学什么"的问题——丰富的音素分布让模型学习更全面的唇形-音频对应关系。唇形幅度解决"能否学得会"的问题——神经网络的频谱偏置（spectral bias）使其倾向于优先拟合低频信号，低运动复杂度的片段更容易被模型正确学习。视角多样性解决"姿态与语音的耦合问题"——当参考片段中头部运动过于剧烈时，音频驱动的唇形运动与姿态驱动的外观变化相互纠缠，模型难以将两者正确解耦。

\textbf{公式与参数}。在实际实验中，片段长度 $n$ 通常设为 3--5 秒，保留候选数 $m = 5$ 或 $10$，高频阈值 $T = 0.25$。论文表明，使用 ISExplore 选择的 3 秒片段训练，其唇形同步质量与使用完整参考视频训练相当，数据处理和训练时间减少 5 倍以上。

\paragraph{AEGIS}

AEGIS（2025.11）~\cite{2511_17747} 提出了首个面向 3D 高斯化身的隐私保护身份遮蔽框架。

\textbf{Baseline 与问题}：基于 3DGS 的光真实 3D 面部化身日益普及，但这也带来了身份盗用风险——恶意方可以从在线共享的化身模型中提取面部特征，绕过生物识别认证系统。2D 图像领域已有成熟的对抗性遮蔽方法，但 3D 化身的动态性和多视角一致性要求使得直接迁移不可行：2D 扰动仅对特定视角有效，而 3D 高斯从任意角度渲染时扰动模式可能失效或产生可见伪影。

\textbf{改了什么}：AEGIS 通过对抗性扰动 3D 高斯的颜色系数（球谐函数参数），在保持化身感知质量的同时实现跨视角一致的去识别。核心设计利用预训练的人脸验证网络作为对抗性损失的身份判别器，引导颜色扰动方向使验证网络无法匹配原始身份，但不修改化身的几何（位置、缩放、旋转、不透明度），从而保持面部结构和表情完整性。

\textbf{实现原理}：优化目标为两部分的加权和：对抗性身份损失最大化预训练人脸验证网络的身份嵌入距离（迫使其无法识别原始身份），感知保真损失最小化原始渲染与扰动渲染之间的 SSIM 差异（保持视觉质量）。优化仅作用于高斯颜色系数，几何属性完全冻结。扰动后的高斯化身从任意视角渲染时，身份信息被有效遮蔽，但年龄、种族、性别、表情等非身份属性得以保留。

\textbf{公式与参数}：AEGIS 达到 100\% 去识别率（人脸检索和验证准确率均降至 0\%），同时保持高感知质量（SSIM = 0.9555，PSNR = 35.52~dB）。无需重新训练或修改化身几何结构，适用于任何已有 3D 高斯化身。

人体新视角合成是 3DGS 在人体领域最基础的应用之一，重点在于渲染质量而非动画驱动能力。与前文讨论的"化身"不同，新视角合成不需要支持任意姿态驱动，而是专注于从已有观测中合成新的相机视角。

\subsubsection{泛化式方法}

本路线介绍无需逐人优化的前馈式人体新视角合成方法。

泛化式方法旨在训练一个通用模型，对未见身份直接进行新视角合成，无需逐人优化。

\paragraph{GPS-Gaussian}

GPS-Gaussian（2023.12）~\cite{2312_02155} 是该方向的先驱工作。

\textbf{Baseline 与问题}：传统基于优化的 3DGS 方法需要对每个新身份从头优化高斯参数，训练耗时数分钟到数十分钟，完全无法满足实时应用需求。

\textbf{改了什么}：GPS-Gaussian 提出了无需优化的即时新视角合成框架：给定稀疏源视图，直接通过前馈网络预测 3D 高斯属性，无需任何逐场景优化。

\textbf{实现原理}：核心流程分为三步。(1)~视角选择与深度估计：从 $N$ 个输入图像中选择与目标视角最邻近的两个源视图（左视图 $v_l$ 和右视图 $v_r$），通过共享图像编码器提取多尺度特征图，然后使用受 RAFT-Stereo 启发的迭代式双目深度估计器预测视差图（等价于深度图）。选择迭代式估计的原因是它避免了构建高计算成本的代价体积（cost volume），更适合高分辨率输入。(2)~高斯参数预测：对于源视图中的每个前景像素，其深度值和 RGB 值分别确定对应 3D 高斯的位置和颜色，其余参数（缩放、旋转、不透明度）由像素级回归网络预测。这些 2D 参数图通过相机参数反投影到 3D 空间，形成 3D 高斯表示。(3)~跨视图聚合与可微分渲染：来自两个源视图的 3D 高斯在 3D 空间中聚合，通过可微分光栅化渲染到目标视角，整个流程端到端可训练。

\textbf{设计动机}：将 3D 高斯的属性预测分解为"从深度+RGB确定位置和颜色"和"从回归网络预测其余参数"两个步骤，充分利用了双目立体视觉提供的几何约束，同时让回归网络仅需学习高斯的形状和透明度属性，大幅降低了学习难度。

\textbf{公式与参数}：通过在大规模人体扫描数据上训练的回归模块，支持 2K 分辨率的实时渲染。

Generalizable Human Gaussians (GHG)（2024.07）~\cite{2407_12777} 与 GPS-Gaussian 的核心区别在于高斯参数的回归空间：GPS-Gaussian 在像素级 2D 图像空间预测高斯属性，而 GHG 将回归过程转移到人体模板的 2D UV 空间。

\textbf{Baseline 与问题}：GPS-Gaussian 等前馈方法在像素空间直接预测高斯参数，缺乏对人体几何结构（关节、自遮挡、复杂表面如头发）的显式建模，导致稀疏视角（如仅 3 视角）下几何和纹理重建严重失真。另一方面，NeRF-based 泛化方法虽然利用了人体先验，但逐像素体渲染的计算瓶颈使其无法实时推理。

\textbf{改了什么}：GHG 将 3D 高斯参数的学习重新表述为人体模板 2D UV 空间上的回归过程，利用 SMPL 模板提供的强几何先验和 2D CNN 的邻域信息聚合能力，并引入多脚手架（multi-scaffold）表示弥补模板与真实人体之间的几何差距。

\textbf{实现原理}：核心流程分为三步。

(1)~UV 空间高斯参数化：将 SMPL 模板网格展开为 2D UV 图，每个 UV 像素对应模板上的一个 3D 点，该点的位置即为高斯中心位置的锚定。通过将稀疏输入视图的特征反投影到 UV 空间，获得与人体几何对齐的特征图。(2)~高斯参数图生成：在 UV 特征图上使用 2D CNN（而非 MLP）预测每个 UV 像素处的高斯参数（缩放、旋转、不透明度、球谐系数）。相比 GPS-Gaussian 的像素级 MLP 回归，2D CNN 能利用 UV 图中相邻像素的局部相关性——相邻身体部位的纹理和几何高度相似，CNN 的感受野能有效聚合这些信息。(3)~多脚手架表示：SMPL 模板仅提供裸体几何，无法表示衣服褶皱和头发等偏移细节。GHG 通过膨胀模板网格生成多层偏移网格（scaffold），在不同距离层次上放置高斯，每层脚手架独立预测偏移量，从而捕获超出单层模板表面的位移细节。此外，UV 空间表示使模型与修复（inpainting）模型天然兼容——未被观测到的 UV 区域可直接由预训练修复模型填充。

\textbf{设计动机}：将高斯学习从 3D 空间转移到 2D UV 空间的动机是双重的：一是利用 SMPL 提供的人体几何先验将问题从无约束的 3D 回归简化为有结构先验的 2D 回归；二是 2D CNN 的局部感受野天然适合人体纹理的空间连续性，且推理效率远高于 NeRF 的逐射线采样。多脚手架的动机则是弥补参数化模板的几何局限——真实人体的衣服和头发在模板表面之外，单一锚定面无法表示这些偏移。

\textbf{公式与参数}：GHG 在 THuman~2.0 和 RenderPeople 数据集上评估，在数据集内和数据集间的泛化设定中均优于当时方法。推理时无需任何测试时优化或微调，仅从 3 个稀疏视角即可生成光真实的新视角渲染。

\paragraph{GRGS}

GRGS（2025.05）~\cite{2505_21502} 将泛化式人体新视角合成从固定光照条件扩展到可编辑重光照场景，是前述 GPS-Gaussian 和 GHG 的功能升级——不仅泛化到新身份，还泛化到新光照环境。

\textbf{Baseline 与问题}：现有可泛化人体 NVS 方法（如 GPS-Gaussian、GHG）将光照烘焙到高斯颜色中，渲染结果无法适应新光照条件。而基于逆渲染的 3DGS 重光照方法需要逐人优化，且几何和材质估计的不准确性会级联影响重光照质量。2D 图像重光照方法虽能泛化，但缺乏 3D 一致性，多视角渲染时产生闪烁伪影。

\textbf{改了什么}：GRGS 提出全监督的前馈框架，将几何、材质和光照信息从多视角 2D 观测投影到 3D 高斯属性中，实现一次性前向推理即可获得可重光照的高斯表示。核心包含两个模块：光照鲁棒几何精修（LGR）和物理基础神经渲染（PGNR）。

\textbf{实现原理}：(1)~LGR 模块：首先用立体匹配方法估计粗深度，然后在合成重光照数据上训练深度感知法线精修网络，预测逐高斯表面法线。

关键设计是在多种光照条件下训练同一几何，迫使网络学习光照不变的几何特征——无论光照如何变化，同一个人体的深度和法线预测保持一致。(2)~PGNR 模块：在高质量几何的基础上，为每个高斯回归几何感知的材质参数（反照率、粗糙度、金属度），结合环境光照参数化进行基于物理的着色。着色过程整合了直接光照、阴影和间接光照，支持可编辑的环境光照切换。为避免昂贵的光线追踪，GRGS 设计了 2D-to-3D 投影训练方案：利用可微的环境光遮蔽、直接光照和间接光照图作为监督信号，将 2D 光照估计投影到 3D 高斯域，大幅降低计算成本。

\textbf{设计动机}：采用全监督而非逆渲染策略的动机是：逆渲染本质上是病态问题，几何和材质估计误差会相互放大；全监督方法通过合成数据直接提供几何和材质真值，绕过了逆渲染的不确定性。在合成重光照数据上训练 LGR 的动机是：真实数据无法获得同一姿态在不同光照下的多组观测，合成数据则可以精确控制光照变量，使网络学到光照不变的几何先验。

\textbf{公式与参数}：GRGS 在多种角色和光照条件下实现了优越的视觉质量、几何一致性和泛化能力，支持带有阴影和间接光照的实时可编辑重光照。

\subsubsection{基于先验和扩散的方法}

本路线介绍利用扩散模型和先验知识辅助稀疏视角人体新视角合成的方法。

从稀疏视角甚至单张图像进行人体新视角合成需要强先验的辅助。前文讨论的泛化式方法虽不依赖逐人优化，但在极端稀疏输入下仍面临挑战。扩散模型的强大生成能力为此提供了新的解决方案。

\paragraph{HuGDiffusion}

HuGDiffusion（2025.01）~\cite{2501_15008} 是首个将条件扩散模型引入单视角人体 3DGS 新视角合成的工作。与前面 GPS-Gaussian 等前馈回归方法不同，HuGDiffusion 不再将高斯属性预测视为回归问题，而是视为条件生成问题，利用扩散模型在不可见区域生成合理外观的能力来弥补回归模型的固有缺陷。

\textbf{Baseline 与问题}。现有可泛化 3DGS 框架（如 GPS-Gaussian）遵循统一的训练范式：参数化学习模型 $\mathcal{M}(\cdot;\theta)$ 以输入图像为条件，直接回归 3D 高斯属性集 $\mathcal{A}$，再通过可微分渲染施加像素级监督 $\min_\theta \sum_{k=1}^{K} \mathcal{L}_\text{pmet}(\mathcal{R}(\mathcal{M}(\cdot;\theta); \mathbf{V}^{(k)}), \mathbf{I}^{(k)})$。这种回归式监督存在根本性局限：单视角输入中大量人体区域（背面、遮挡部位）不可见，回归模型缺乏对这些区域生成合理外观的能力——即使数值指标尚可，渲染出的不可见区域也经常出现大片黑色伪影和不自然纹理。

\textbf{改了什么}。HuGDiffusion 对上述训练范式进行了两层根本性改变。（1）监督信号从像素级变为属性级：不再仅通过渲染图像与真实图像的像素差异来优化，而是直接在 3D 高斯属性空间上构建代理真值（proxy ground-truth），实施属性级监督 $\min_\theta \mathcal{L}_\text{setdiff}(\mathcal{A}, \hat{\mathcal{A}})$。（2）生成模型从回归网络变为条件扩散模型：用扩散过程的逐步去噪来生成高斯属性，而非一次性回归输出。

\textbf{实现原理}。整体框架包含三个核心阶段。

\underline{阶段一：代理真值高斯属性的构建}。为了训练扩散模型，需要属性级的监督信号——即每条人体数据对应的"标准答案"高斯属性集。然而，直接对每个场景独立运行 vanilla 3DGS 优化会产生分布不一致的高斯属性集（即使同一场景的两次运行，属性值和基元排列也可能截然不同），这样的数据无法供扩散模型有效学习。为此，HuGDiffusion 设计了两阶段构建流程。第一阶段（逐场景过拟合）：对每个人体捕获，先从真实人体表面均匀采样 $N = 20{,}000$ 个 3D 点作为固定的高斯位置，其余四类属性（不透明度 $\alpha^{(n)}$、缩放 $s^{(n)}$、旋转 $q^{(n)}$、球谐系数 $c^{(n)}$）由 Point Transformer 网络 $\mathcal{F}_1$ 输出，即 $\{\alpha^{(n)}, s^{(n)}, q^{(n)}, c^{(n)}\}_{n=1}^{N} = \mathcal{F}_1(\{p^{(n)}\}_{n=1}^{N}; \theta_1)$。关键设计在于将高斯属性作为网络输出而非直接优化变量，利用神经网络的内在平滑性来抑制属性的剧烈波动。同时引入辅助正则约束以抑制缩放和不透明度的高度非均匀分布：

\begin{equation}
\mathcal{L}_\text{aux} = \sum_{n=1}^{N}\|\mathtt{radius}(\bar{\mathbf{s}}^{(n)}) - \mathtt{kdist}(\hat{\mathbf{p}}^{(n)})\|^2 + \sum_{n=1}^{N}\|\bar{\alpha}^{(n)} - 1\|^2
\end{equation}

其中 $\mathtt{radius}(\cdot)$ 根据缩放属性 $\mathbf{s}^{(n)}$ 计算高斯半径，$\mathtt{kdist}(\cdot)$ 计算高斯位置 $\hat{\mathbf{p}}^{(n)}$ 到其最近邻高斯的平均距离。直觉上，第一项确保相邻高斯的缩放与其空间间距保持一致（避免某些高斯过大或过小），第二项鼓励不透明度趋向 $1$ 以避免半透明碎片的产生。论文通过实验验证了这一设计的必要性：固定所有随机种子后对同一场景运行 $200$ 次 vanilla-3DGS 优化，球谐系数求和值呈现巨大波动，局部区域（如同一颜色衣物表面）的球谐值范围从 $-12.5$ 到 $6$，缩放和不透明度属性同样高度混乱——这样的数据完全无法作为扩散模型的监督信号。第二阶段（分布统一）：用另一个 Point Transformer $\mathcal{F}_2$ 对全部 $J = 480$ 条人体数据联合过拟合（batch size 为 $4$），以批量方式逐步优化，对齐不同场景间的属性分布。由于无法一次性输入全部 $480$ 条数据，采用分批迭代策略逐步统一属性分布。实验表明，两阶段构建使属性最大值的方差从 $0.1114$ 降至 $0.1094$，最小值的方差从 $0.0707$ 降至 $0.0329$，显著收窄了解空间。论文还进一步验证了直接用 vanilla-3DGS 做第一阶段再由 Point Transformer 做第二阶段统一也无法获得足够一致的分布，证明两阶段中均使用 Point Transformer 的必要性。

\underline{阶段二：高斯位置的生成}。训练时直接使用代理真值构建流程中已准备好的高斯位置 $\{\hat{\mathbf{p}}_j^{(n)}\}_{n=1}^N$。推理时需要从输入图像估计高斯位置，位置生成器的流程如下：先通过 EcoDepth 单目深度估计网络将输入图像 $\mathbf{I}_s$ 转为深度图并反投影为部分 3D 点云；同时从 $\mathbf{I}_s$ 估计 SMPL 人体模型参数，并用部分点云修正 SMPL 姿态以获得更准确的 3D 空间姿态（论文指出，SMPL 的 2D 投影姿态可能看似正确但 3D 空间中姿态错误，需要 3D 修正）。修正后的 SMPL 模型和部分点云一同送入基于 HaP（Human-as-Points）的点云生成网络，输出完整的 3D 高斯位置，再经点云上采样和最远点采样（FPS）确保点云的均匀性和密度一致性。

\underline{阶段三：条件扩散生成其余属性}。在获得高斯位置后，HuGDiffusion 提取两类人体中心化特征作为扩散条件：（a）像素对齐特征 $\bm{\beta}^{(n)}$——将 3D 高斯位置分为可见和不可见两组，可见位置投影到输入图像 $\mathbf{I}_s$ 上提取 2D CNN 特征，不可见位置投影到由 fine-tuned Stable Diffusion 预测的背面图像 $\mathbf{I}_\text{back}$ 上提取特征，两组特征拼接形成完整的像素对齐特征（特征提取使用预训练 ResNet18）；（b）SMPL 语义特征 $\bm{\gamma}^{(n)}$——为每个高斯检索最近 SMPL 表面点的语义标签、距离和索引，通过 MLP 编码为结构先验，注入显式的空间身体结构信息以减少噪声并产生更清晰的相邻部位边界。属性扩散过程将四类属性分成两组独立处理（实验发现同时扩散全部四类属性会导致训练崩溃）。扩散模型的主干网络采用 PointNet++（论文发现 Point Transformer 的体素大小会严重影响性能，因此不适合作为扩散模型 backbone）。第一组为球谐系数，通过属性扩散器 $\mathtt{GSDIFF}_{\psi_1}$ 在时间步 $t$ 预测噪声：

\begin{equation}
\hat{\epsilon}^{(n)} = \mathtt{GSDIFF}_{\psi_1}(\tilde{\mathbf{c}}_t^{(n)}, \mathbf{p}^{(n)}, \bm{\beta}^{(n)}, \bm{\gamma}^{(n)}, t), \quad \mathcal{L}_1 = \mathbb{E}_\epsilon \|\hat{\epsilon}^{(n)} - \epsilon^{(n)}\|^2
\end{equation}

其中 $\tilde{\mathbf{c}}_t^{(n)}$ 是添加了噪声的球谐系数，$\bm{\beta}^{(n)}$ 和 $\bm{\gamma}^{(n)}$ 分别为像素对齐特征和 SMPL 语义特征。推理时从高斯分布采样 $\mathbf{c}_T^{(n)}$，逐步去噪得到 $\mathbf{c}_0^{(n)}$。然而，由于在 3D 域中证据下界（ELBO）的归纳偏置不清晰，去噪后的 $\mathbf{c}_0^{(n)}$ 仍含有残余噪声。因此，HuGDiffusion 设计了一个额外的精修步骤（extra-step）：用第二个扩散器 $\mathtt{GSDIFF}_{\psi_2}$ 在去噪后的球谐系数 $\mathbf{c}_0^{(n)}$ 上同时预测球谐残差 $\epsilon^{(n)}$ 以及其余三类属性（不透明度 $\alpha^{(n)}$、缩放 $\mathbf{s}^{(n)}$、旋转 $\mathbf{q}^{(n)}$），最终球谐系数修正为 $\mathbf{c}^{(n)} = \mathbf{c}_0^{(n)} - \epsilon^{(n)}$，施加四项 $L_2$ 监督：

\begin{equation}
\mathcal{L}_2 = \|\mathbf{c}^{(n)} - \hat{\mathbf{c}}^{(n)}\|^2 + \|\alpha^{(n)} - \hat{\alpha}^{(n)}\|^2 + \|\mathbf{s}^{(n)} - \hat{\mathbf{s}}^{(n)}\|^2 + \|\mathbf{q}^{(n)} - \hat{\mathbf{q}}^{(n)}\|^2
\end{equation}

这一 extra-step 设计同时解决了球谐去噪不彻底和其他属性生成两个问题，避免了单独为不透明度、缩放、旋转再设计一组完整扩散流程的复杂度。

\textbf{设计动机}。将生成范式从回归改为扩散的动机是：单视角 NVS 本质上是条件生成任务——不可见区域没有回归目标，需要模型"想象"合理外观，而扩散模型天然擅长在数据分布中采样出合理的补全结果。构建代理真值属性的动机是：扩散模型需要属性级监督信号（而非像素级），但 vanilla 3DGS 优化结果的随机性使属性分布混乱不可学，因此必须通过 Point Transformer 的网络平滑性和两阶段分布统一来创造一致的监督信号。将球谐系数与其余属性分开扩散的动机是：四类属性的数值范围和分布特性差异巨大（如球谐系数可跨越正负数十，而不透明度在 $[0,1]$），同时扩散会导致训练崩溃。

\textbf{公式与参数}。HuGDiffusion 在 Thuman2 数据集上使用 $J = 480$ 个人体模型构建代理真值，每人渲染 $360$ 个视角、均匀采样 $N = 20{,}000$ 个高斯位置，图像分辨率为 $512 \times 512$。逐场景过拟合 $4{,}000$ epochs，分布统一阶段 $1{,}300$ epochs（batch size $4$），属性扩散模型训练 $300$ epochs（batch size $4$），学习率均为 $0.0002$，均使用 Adam 优化器。最终在 Thuman2（$20$ 人测试）、CityuHuman（$20$ 人）、2K2K（$25$ 人）、CustomHuman（$40$ 人）四个数据集上均取得 PSNR、SSIM、LPIPS 全部最优，Thuman2 上 PSNR 达 $30.03$~dB（SSIM $0.953$，LPIPS $0.065$），CityuHuman 上 PSNR $32.47$~dB，2K2K 上 PSNR $30.64$~dB，CustomHuman 上 PSNR $34.82$~dB。论文的消融实验还表明：（1）用回归模型（无论像素级还是属性级监督）替代扩散模型时，不可见区域出现大面积黑色伪影（如黑色头部和手臂），扩散模型能生成合理外观；（2）移除背面图像预测模块时，面部身份信息丢失，PSNR 从 $30.03$ 降至 $29.63$；（3）移除 SMPL 语义特征时，颈部等部位边界模糊且出现红色噪声像素。

\paragraph{PoseGaussian}

PoseGaussian（2026.02）~\cite{2602_05190} 与 GPS-Gaussian 的关键区别在于对人体姿态信息的利用方式——GPS-Gaussian 仅将姿态用于背景分离和视角选择，而 PoseGaussian 将姿态信号深度嵌入几何估计和时序稳定的两个阶段。

\textbf{Baseline 与问题}：现有 3DGS 人体方法要么完全忽略人体语义（如 GPS-Gaussian 不使用姿态或身体结构信息来指导高斯参数预测），要么过度依赖精确的 SMPL 对齐（如 GauHuman 在姿态噪声或体型偏离训练分布时泛化能力显著下降）。此外，在快速运动和严重自遮挡场景下，现有方法常出现运动模糊和背景泄漏。

\textbf{改了什么}：PoseGaussian 将人体姿态作为双重信号引入渲染管线：(1) 作为结构先验，与颜色编码器融合来精修深度估计；(2) 作为时间线索，通过专用姿态编码器增强帧间时序一致性。整体流程为端到端可微、可训练的统一架构。

\textbf{实现原理}：管线包含三个核心组件。

(1)~姿态引导特征融合用于深度推理：从输入图像中提取颜色特征的同时，从 SMPL 姿态参数中提取姿态特征，通过交叉注意力机制将两者融合。融合后的特征用于精修深度估计，使深度预测不仅依赖外观匹配（在自遮挡和纹理缺失区域容易失败），还能利用身体结构的几何先验（例如，知道手臂的关节位置有助于推断被遮挡的前臂深度）。(2)~时序姿态稳定器（Temporal Pose Stabilizer, TPS）：专用姿态编码器处理连续帧的姿态序列，提取时序姿态变化特征作为高斯参数预测的额外条件。该模块通过约束相邻帧的高斯参数变化与姿态变化一致，减少快速运动时的渲染抖动。(3)~姿态条件化损失：引入姿态感知的正则化项，当姿态变化剧烈时增强时序一致性约束，当姿态稳定时放松约束以保留细节。

\textbf{设计动机}：将姿态嵌入几何阶段而非仅作为条件的动机是：姿态提供了人体结构的强语义锚点，直接参与深度推理比仅在渲染阶段使用姿态约束更为有效。

时序姿态稳定器的设计动机是：纯外观的时序约束在快速运动时失效（相邻帧外观变化太大），而姿态提供了比外观更稳定的时序线索。

\textbf{公式与参数}：PoseGaussian 以 100 FPS 的实时渲染速度在 ZJU-MoCap 和 THuman2.0 数据集上实现了最优的感知质量和结构精度（ZJU-MoCap 上 PSNR 30.86，SSIM 0.979，LPIPS 0.028），同时在未见姿态和未见身份上展现良好泛化能力。

Double Unprojected Textures (DUT)（2024.12）~\cite{2412_13183} 与前述 GHG 类似地在 UV 纹理空间操作，但核心创新在于"双重纹理反投影"策略——利用几何精修后的第二次反投影来改善外观学习质量。

\textbf{Baseline 与问题}：基于纹理空间的 2D CNN 方法（如 GHG 的前身 ANR 等）依赖纹理反投影（texture unprojection）将多视角图像信息对齐到 2D UV 空间。然而，反投影的精度取决于模板几何与真实人体几何的对齐程度。现有方法要么将几何和外观联合学习（从不完美的纹理反投影中同时学两者，质量下降），要么在几何估计时完全忽略稀疏图像信息（仅依赖运动学驱动的模板变形，导致衣服褶皱等细节丢失）。

\textbf{改了什么}：DUT 将粗略几何变形估计与外观合成显式解耦，通过两次纹理反投影实现级联精修：第一次反投影基于粗模板几何，用于估计精修后的几何；第二次反投影基于精修几何，产生更准确的纹理图供外观学习使用。

\textbf{实现原理}：管线分为两个阶段。

(1)~图像条件化模板变形：先将人体模板网格通过 LBS 变换到当前姿态，使其大致对齐稀疏输入图像。然后进行第一次纹理反投影：将多视角图像信息投影到 2D UV 纹理空间，获得粗纹理图。该粗纹理图包含了图像中的几何线索（如轮廓、褶皱阴影等），输入 UNet 网络预测模板顶点的偏移量，生成精修后的几何。关键区别在于这是图像条件化而非运动条件化的变形——网络直接从局部图像特征学习变形，而非仅依赖运动参数，使其对分布外动作（如训练中未见过的跳跃动作）具有更强的鲁棒性。(2)~高斯外观预测：将精修后的几何再次通过 LBS 变换，进行第二次纹理反投影。由于几何更接近真实表面，第二次反投影产生的纹理图伪影更少、与输入视图对齐更好。该精修纹理图作为条件输入第二个 UNet，预测高斯泼溅参数（位置、颜色、缩放、旋转、不透明度）。此外，DUT 引入高斯缩放精修（Gaussian scale refinement），利用模板变形信息指导高斯缩放参数，缓解几何变形对高斯外观的负面影响。

\textbf{设计动机}：双重反投影的动机是：第一次反投影的几何不准确性会导致纹理错位和伪影，但如果先用粗纹理改善几何再反投影一次，纹理质量显著提升——这是一个正反馈循环。将几何和外观解耦的动机是：联合学习时，不准确的几何会污染外观学习，而外观噪声又反过来干扰几何优化；显式分阶段处理打破了这种恶性耦合。

\textbf{公式与参数}：DUT 支持 4K 分辨率的实时光真实渲染（CVPR 2025），在定量和定性实验中显著超越先前最优方法，同时对分布外动作保持鲁棒性。

\subsubsection{特殊场景的人体新视角合成}

本路线介绍面向遮挡、稀疏视角等特殊场景的人体新视角合成方法。

\paragraph{U-4DGS}

U-4DGS（2026.02）~\cite{2602_06343} 专门处理单目遮挡人体的渲染，核心思想是将任务重新表述为异方差观测噪声下的最大后验（MAP）估计问题，用可学习的不确定性取代脆弱的硬掩码或外部生成先验。

\textbf{Baseline 与问题}。单目遮挡人体渲染面临一个两难困境：若将所有像素一视同仁地作为监督信号（同方差 MLE），变形场被迫建立错误对应，导致遮挡物纹理被烘焙到人体几何中；若直接施加硬可见性掩码排除遮挡区域，则这些区域的几何完全失去梯度约束，无法致密化或保持连贯，产生空洞。现有两类解决路线各有局限：几何启发式方法（如左右对称约束）无法处理不对称服装，生成先验方法（如扩散模型修复）的逐帧随机性会引入严重的时序闪烁。U-4DGS 的关键洞察是：视频的时序冗余性意味着遮挡是瞬态的——某帧被遮挡的部位在其他帧中大概率可见，因此不需要"想象"缺失内容，而是需要一个动态判别机制来评估每次观测的可靠程度。

\textbf{改了什么}。U-4DGS 对标准 4DGS 管线进行了三处修改。（1）概率变形网络：在标准变形 MLP 输出几何偏移量（$\Delta\boldsymbol{\mu}, \Delta\mathbf{r}, \Delta\mathbf{s}$）之外，额外预测逐高斯的认知不确定性标量 $\sigma \in \mathbb{R}^+$（Softplus 激活），为每个高斯基元的时序对应关系提供置信度度量。（2）联合光栅化：在颜色渲染之外并行渲染像素对齐的不确定性图 $\hat{U}(\mathbf{u}) = \sum_{i \in \mathcal{N}} T_i \alpha_i \sigma_i$，使用与颜色相同的 $\alpha$-混合 权重累积不确定性，确保渲染外观与置信度的几何一致性。（3）置信度感知正则化：利用学到的 $\sigma$ 选择性地在不可靠区域施加空间一致性和时序惯性约束，防止几何漂移。

\textbf{实现原理}。MAP 目标分解为似然项 $P(\mathcal{D}|\Theta)$ 和先验项 $P(\Theta)$。似然项假设像素残差服从 Laplacian 分布（位置参数为预测颜色，尺度参数为不确定性 $\hat{U}$），取负对数似然得到不确定性加权光度损失：

\begin{equation}
\mathcal{L}_{NLL} = \sum_{\mathbf{u} \in \Omega} \left( \frac{\|\mathbf{C}_{gt}(\mathbf{u}) - \hat{\mathbf{C}}(\mathbf{u})\|_1}{\hat{U}(\mathbf{u}) + \epsilon} + \lambda_{reg} \log(\hat{U}(\mathbf{u}) + \epsilon) \right)
\end{equation}

该目标充当自适应梯度调节器：可见区域中优化器降低 $\hat{U}$ 让 RGB 监督主导；遮挡区域中优化器升高 $\hat{U}$ 自动衰减梯度幅度，从而保护规范几何免受误导信号污染。先验项包含两组置信度感知正则化：空间一致性损失 $\mathcal{L}_{spa}$ 在规范空间构建 KNN 图，以 $\text{sg}(\sigma_i)$ 为注意力权重约束相邻高斯的变形偏移一致——当某高斯被遮挡（$\sigma_i$ 高）时，光度梯度消失，空间正则项接管，将其变形强制同步到由光度损失牢牢锚定的可见邻居上，实现从可靠边界向遮挡区域内部的几何有效性传播；时序惯性损失 $\mathcal{L}_{temp}$ 对不确定高斯的变形轨迹施加二阶差分平滑，防止无视觉线索时的轨迹抖动。

\textbf{设计动机}。选择不确定性建模而非生成先验的动机是：生成模型的逐帧随机性不可避免地导致身份漂移和时序闪烁（例如衬衫颜色在帧间不一致变化），而不确定性建模通过严格聚合有效的时序信息来保证物理一致性。选择 MAP 而非硬掩码的动机是：硬掩码留下数学上完全无约束的盲区，而 MAP 框架通过不确定性 $\sigma$ 实现从"完全信任"到"完全忽略"的连续过渡，并通过正则化先验在无约束区域注入物理约束。将 $\sigma$ 定义在规范空间而非观测空间的动机是：遮挡是视角相关的瞬态现象，将不确定性绑定为高斯基元的内在属性（LBS 变换期间保持不变），使网络能学习空间-时序失效模式（如频繁自遮挡的关节区域），而不受全局视角的影响。

\textbf{公式与参数}。U-4DGS 在 ZJU-MoCap 数据集上达到 PSNR $24.62$~dB、SSIM $0.9606$、LPIPS $31.72$，在 OcMotion 数据集上达到 PSNR $20.11$、SSIM $0.9030$、LPIPS $79.17$，两项指标均取得全部最优。训练约 $15$ 分钟。

\paragraph{OccFusion}

OccFusion（2024.07）~\cite{2407_00316} 与 U-4DGS 解决相同的遮挡人体渲染问题，但采用了截然不同的技术路线：U-4DGS 用可学习的不确定性建模在时序上聚合有效观测，OccFusion 则用预训练 2D 扩散模型作为生成先验来"想象"被遮挡区域的内容。

\textbf{Baseline 与问题}：大多数人体渲染方法假设输入视频中人体始终完全可见，但真实场景中遮挡无处不在（如医院、体育场、建筑工地中的物体遮挡）。现有遮挡处理方法（如 OccNeRF）计算成本高、训练时间长，而直接在遮挡区域忽略监督会导致浮点伪影、不完整的人体部位和可见与不可见区域的明显接缝。

\textbf{改了什么}：OccFusion 提出三阶段管线，利用 3DGS 的高效渲染和预训练 2D 扩散模型的生成能力来处理遮挡：(1) 初始化阶段从部分可见掩码恢复完整人体掩码；(2) 优化阶段用分数蒸馏采样（SDS）正则化确保人体几何完整性；(3) 精修阶段用上下文修复（in-context inpainting）改善少观测区域的外观质量。

\textbf{实现原理}：(1)~初始化阶段：利用分割和姿态先验，将每帧的部分可见性掩码修复为完整的人体占据掩码。这为后续阶段提供了完整的人体区域监督信号。(2)~优化阶段：初始化一组 3D 人体高斯，在观测到的人体区域进行正常光度监督优化。对于未观测区域，引入姿态条件化的 SDS 损失，利用预训练扩散模型（如 Stable Diffusion）的 2D 先验来约束遮挡区域的渲染结果在语义上合理。SDS 同时在姿态空间和规范空间施加约束：姿态空间确保每帧渲染的遮挡区域看起来像完整人体，规范空间确保不同帧的遮挡区域在高斯基元层面保持 3D 一致性。(3)~精修阶段：对于少观测的人体部位（如背面），利用预训练修复模型结合部分观测和前阶段渲染结果进行上下文修复，进一步提升这些区域的渲染质量。

\textbf{设计动机}：选择扩散先验而非不确定性建模（U-4DGS 路线）的动机是：当遮挡持续时间长或遮挡面积大时，时序聚合可能无法获得足够的有效观测；扩散模型能生成语义合理的内容来填补这些缺失。

但代价是扩散先验的逐帧随机性可能引入时序不一致（U-4DGS 在其问题陈述中明确指出了这一局限）。OccFusion 的 SDS 在规范空间施加约束正是为了缓解这一问题。

\textbf{公式与参数}：OccFusion 在 ZJU-MoCap（模拟遮挡）和 OcMotion（真实遮挡）数据集上取得遮挡人体渲染的最优性能。训练仅需约 10 分钟，显著快于基于 NeRF 的遮挡处理方法。

\paragraph{Tele-Aloha}

Tele-Aloha（2024.05）~\cite{2405_14866} 是一个将 3DGS 新视角合成技术工程化为完整远程呈现系统的代表性工作，其新视角合成算法可视为 GPS-Gaussian 在宽基线稀疏相机设定下的工程优化版本。

\textbf{Baseline 与问题}：现有沉浸式远程呈现系统（如 Google Starline）需要密集相机阵列和定制硬件，成本高达数万美元，难以推广到普通消费者。基于 NeRF 的方法推理速度太慢，无法满足实时通信要求（端到端延迟需低于 150~ms）。GPS-Gaussian 等前馈方法虽然快速，但它们假设窄基线的密集视角，在 Tele-Aloha 的宽基线四相机设定下，立体匹配和特征聚合的准确性严重下降。

\textbf{改了什么}：Tele-Aloha 提出面向上半身通信的低预算远程呈现系统，仅需 4 个消费级 RGB 相机（无深度传感器）、1 块消费级 GPU 和 1 块自动立体显示屏。核心算法包含三项关键设计：级联视差估计器处理宽基线索件、隐空间高斯泼溅保证渲染完整性、加权混合机制恢复 2K 分辨率。

\textbf{实现原理}：(1)~级联视差估计：4 个相机的基线较宽，传统立体匹配方法难以建立可靠对应。Tele-Aloha 设计级联架构，从粗到精逐步估计视差图，在宽基线下仍能获得鲁棒的几何线索。(2)~隐空间高斯泼溅（Latent-based Gaussian Splatting）：与 GPS-Gaussian 直接预测 RGB 颜色的高斯不同，Tele-Aloha 在低分辨率隐空间中进行高斯泼溅——将像素对齐的隐特征而非球谐系数投影到目标视角，再通过神经图像解码器恢复到较低分辨率图像。这保证了渲染的完整性（不会因为遮挡或视角外推导致空白区域），同时将计算量控制在实时范围内。(3)~遮挡感知渲染精修：由于隐空间渲染在低分辨率进行，直接上采样会丢失细节。Tele-Aloha 利用高质量输入图像的加权混合机制精修解码图像——根据源视图中人体表面的可见性计算权重，将高分辨率 RGB 输入融合到目标视角，最终输出 2K 分辨率渲染图像。

该系统还结合低延迟虹膜追踪和自动立体显示屏，使用户无需佩戴头显即可体验强烈的三维立体感。

\textbf{设计动机}：采用隐空间而非直接 RGB 空间渲染的动机是：在极稀疏视图下，直接 RGB 预测容易产生空白和伪影，而隐特征空间能编码更丰富的语义信息，经解码后保证视觉完整性。使用加权混合而非纯神经渲染的动机是：输入图像本身已经是 4K 高质量的，充分利用这些真实像素比纯神经预测能获得更清晰的结果。

\textbf{公式与参数}：整个系统总成本约 15,000 美元，实现 2048 $\times$ 2048 分辨率、30 FPS、低于 150~ms 端到端延迟的双向实时通信（SIGGRAPH 2024）。

\begin{takeaway}
身体与运动建模面临数量级挑战：全身化身需处理上百关节和非刚性衣物变形，前馈高斯预测和隐空间渲染策略在遮挡和稀疏视角条件下仍能实现高保真重建，Tele-Aloha 等工程化系统已将成本降至消费级。
\end{takeaway}

%=====================================================================

\subsection{近期方法脉络：可动画生物与化身重建、动物与生物体重建、Compositional human-object Gaussian scenes等}

这一组近期工作延续 3DGS、Gaussian splats、Dynamic Gaussian Splatting，主要把改动集中在扩展、增量改进、应用、压缩等环节。下面只保留相对既有路线的变化，便于和前文方法直接对照。

\paragraph{CORGI}
CORGI（2026.07）~\cite{2607_00321} 延续 3DGS 在“可动画生物与化身重建”上的基本设定，主要改变是：CORGI：Canonical-Driven Orbital Generation生成一致的环绕观测，结合Consistency-aware Deformable 3DGS在D-SMAL先验上建模跨视角形变。高斯基元仍承担可渲染动态场景表示。
其设计动机是：提出面向单图狗重建的无3D监督框架，将姿态规范化、多视图生成和可变形3DGS结合起来，显式缓解生成视角不一致和动物关节形变带来的几何/外观错位。
具体而言，方法首先用Canonical和Orbit LoRA将任意输入姿态规范化并合成360度视频观测；随后用CA-3DGS模块锚定D-SMAL动物先验，在每个视角上建模局部非刚性偏移并保持跨视角一致；最后结合扩散式修复网络迭代去除生成伪影，使单图输入能够形成可渲染、可动画的3DGS动物表示。
适用边界是：依赖专门的犬类先验和LoRA生成模型；对非犬类动物、严重遮挡和生成模型失败案例的泛化仍有限；分类基于摘要和源码章节片段。

\paragraph{BEAST3D}
BEAST3D（2026.06）~\cite{2606_02937} 延续 Gaussian splats 在“动物与生物体重建”上的基本设定，主要改变是：BEAST3D：ViT预测3D Gaussian splats并通过可微渲染重建held-out views，同时分割动物前景。高斯基元仍承担可渲染动态场景表示。
其设计动机是：提出面向实验动物多视角视频的自监督预训练框架，用3D高斯重建和前景分割共同学习结构特征，减少人工姿态标注依赖。
具体而言，模型输入校准多视角图像，预测3D Gaussian splats以重建留出视角，并同时学习动物/背景分割；已知相机几何帮助少视角条件下稳定学习。
适用边界是：依赖校准多视角实验设置；主要面向实验动物，开放场景泛化有限。

\paragraph{CoGS}
CoGS（2026.06）~\cite{2606_28820} 延续 Dynamic Gaussian Splatting 在“Compositional human-object Gaussian scenes”上的基本设定，主要改变是：组合式 Gaussian Splatting，分解人体、物体和背景三条分支。高斯基元仍承担可渲染动态场景表示。
其设计动机是：将复杂动态场景分解为 articulated human、manipulated object 和 static background，分别建模并协调优化。
具体而言，人体分支由完整 canonical prior 初始化，对象和背景使用不同运动模型与 Gaussian 表示，通过组合渲染处理共享像素和遮挡。
适用边界是：单目视频约束不足，少见区域和复杂遮挡仍难恢复；依赖人体先验和对象分割质量。

\paragraph{SignNet-1M}
SignNet-1M（2026.06）~\cite{2606_24361} 延续 3DGS 在“Sign-language and motion dataset augmentation”上的基本设定，主要改变是：使用 3DGS 做新视角渲染，并结合扩散模型进行背景和身份编辑。高斯基元仍承担可渲染动态场景表示。
其设计动机是：提出覆盖 ASL、CSL 和 DGS 的大规模增强手语数据集，并通过新视角、背景、身份和成像扰动提升鲁棒性。
具体而言，3DGS 用于从有限视频中合成旋转、缩放等新视角，扩散模型用于背景替换和 signer 替换，后处理模拟真实采集扰动。
适用边界是：3DGS 是数据增强手段而非论文的主要方法创新；未从摘要看到对 3DGS 训练或表示的改进。

\paragraph{RigPAPR}
RigPAPR（2026.06）~\cite{2606_06685} 延续 Gaussian splat LBS 在“点云与高斯动画表示”上的基本设定，主要改变是：RigPAPR：基于Proximity Attention Point Rendering的rig动画方法，避免Gaussian splat线性蒙皮的关节边界伪影。高斯基元仍承担可渲染动态场景表示。
其设计动机是：指出对单个splat做刚性LBS会破坏canonical tiling并产生关节缝隙/尖刺，转向proximity attention point rendering来改善动画。
具体而言，方法从静态点云重建和单目驱动视频中恢复rig和可重定位资产，用邻近注意力而非逐splat刚性蒙皮处理关节区域。
适用边界是：不是标准3DGS方法；与高斯表示的关系主要是对比和替代。

\paragraph{RAGA}
RAGA（2026.06）~\cite{2606_29329} 延续 3DGS avatar 在“Avatar-scene lighting and shadows”上的基本设定，主要改变是：基于精确 ray-Gaussian line integrals 的 Gaussian 空间阴影计算。高斯基元仍承担可渲染动态场景表示。
其设计动机是：提出 Ray-Traced Gaussian Shadow Casting，用沿阴影光线的高斯不透明度积分估计遮挡权重。
具体而言，对每个遮挡 Gaussian 沿 shadow ray 积分其 opacity profile，并用理论最大积分归一化，从而为 avatar-scene、multi-avatar 和 object interaction 产生物理更合理的阴影。
适用边界是：阴影准确性依赖 Gaussian opacity 与几何的对应关系；复杂半透明材质和全局光照仍未完全覆盖。

\paragraph{PointSplat}
PointSplat（2026.06）~\cite{2606_32036} 延续 3DGS avatar 在“Compact human Gaussian representations”上的基本设定，主要改变是：PointSplat：在 3D 空间进行 human-centric prediction 以减少多视图冗余。高斯基元仍承担可渲染动态场景表示。
其设计动机是：将预测从 view-centric 转为 3D 空间中的 human-centric representation，以同时保持高保真和低带宽。
具体而言，网络直接学习并生成紧凑 Gaussian/point-based 人体表示，减少跨视角冗余，适配有限计算和传输带宽的 live streaming。
适用边界是：依赖人体先验和输入视角质量；宽松衣物、遮挡和快速动作可能增加难度。

\paragraph{Flex4DHuman}
Flex4DHuman（2026.06）~\cite{2606_13655} 延续 4D Gaussian splats 在“动态人体多视角生成”上的基本设定，主要改变是：Flex4DHuman：多视角视频扩散模型，仅用相对相机位姿条件生成可送入4D Gaussian splats重建的视图。高斯基元仍承担可渲染动态场景表示。
其设计动机是：用相对相机位姿编码替代骨架、深度、法线或目标视角几何先验，使动态人物多视角补全更灵活。
具体而言，模型基于Wan视频扩散骨干，通过五轴位置编码扩展时空RoPE以注入视角索引和连续SE(3)位姿，输出同步多视角视频。
适用边界是：本身不是3DGS重建方法；下游4D Gaussian质量依赖生成视频一致性。

\paragraph{TryOnCrafter}
TryOnCrafter（2026.06）~\cite{2606_26092} 延续 4DGS 在“Camera-controllable virtual try-on”上的基本设定，主要改变是：构建可渲染 4D 试衣代理以支持任意视角动态渲染。不改变 3DGS 主线的表示、渲染或优化管线。
其设计动机是：突破传统 VVT 只能跟随源相机轨迹的问题，让服装试穿视频可在任意相机轨迹下合成。
具体而言，摘要未直接确认代理是否为 4D Gaussian 表示，但任务与动态可渲染人体代理高度接近 4DGS 应用；方法强调结构同步、纹理幻化和背景一致性。
适用边界是：3DGS 证据来自关键词召回和可渲染 4D proxy 描述，若全文未使用 Gaussian 表示则相关性需下调。

\paragraph{Hand-4DGS}
Hand-4DGS（2026.06）~\cite{2606_19156} 延续 4D Gaussian Splatting 在“手部与第一视角重建”上的基本设定，主要改变是：Hand-4DGS：mesh-guided表示和temporal convolutions结合的前馈4D手部Gaussian Splatting框架。高斯基元仍承担可渲染动态场景表示。
其设计动机是：将mesh结构先验和时间卷积引入前馈4DGS手部重建，实现约60 FPS推理和跨场景泛化。
具体而言，方法使用手部mesh指导高斯布局和结构约束，并通过时序卷积建模快速手部运动，直接从第一视角视频输出动态高斯手部。
适用边界是：依赖手部mesh先验；严重遮挡和手物交互仍困难。
\section{语音驱动、局部建模与扩展应用}
\label{sec:16_3}

语音驱动的说话人头像是数字化身最活跃的应用方向之一。此外，手部、局部人体部位的建模以及服装试穿与编辑也是 3DGS 在人体领域的重要应用。本节从说话人头像、手部与局部建模、服装试穿与编辑三个角度展开讨论，它们与前两节形成互补：如果说前两节关注的是"如何重建和驱动一个完整的化身"，那么本节更多关注"如何让化身说话、交互和换装"。

\subsection{基础方法：从音频到高斯变形}

本节讨论~8~条技术路线：（1）实时与高效方法；（2）唇形同步与情绪控制；（3）少样本与跨身份泛化；（4）可动画手部化身；（5）手-物交互；（6）服装-人体分离与重建；（7）虚拟试穿；（8）服装重建。

\paragraph{GaussianTalker}

GaussianTalker（2024.04）~\cite{2404_16012} 是将 3DGS 应用于说话人头像的先驱工作之一。

\textbf{Baseline 与问题}：基于 NeRF 的说话人头像方法（如 AD-NeRF、GeneFace）虽然能生成高质量面部动画，但体渲染的射线步进机制导致推理速度极慢（每帧数秒），且在动态面部区域容易产生扭曲伪影——根本原因在于直接在动态区域生成逐点外观时，外观空间的跳变和不连续性使得连续平滑的神经场难以精确拟合。

\textbf{改了什么}：GaussianTalker 提出基于变形而非外观生成的范式：维护一组持久的规范 3D 高斯表示静态平均面部形状，通过音频驱动的变形场预测逐帧属性偏移量。

\textbf{实现原理}：核心架构包含两个关键模块。(1)~三平面编码的规范高斯：为了建立高斯之间的空间关系（相邻面部区域应具有相似的变形模式），该方法使用多分辨率三平面表示编码规范 3D 位置的潜特征 $f(x_c)$，从中解码高斯属性（缩放、旋转、球谐系数、不透明度），而非直接存储独立的高斯参数。这使得空间上邻近的高斯共享相似的特征嵌入。

(2)~空间-音频交叉注意力：音频特征（由 DeepSpeech 编码器提取）和高斯空间特征通过交叉注意力模块融合。注意力机制允许每个高斯根据其空间位置关注音频特征的不同部分，例如嘴部高斯更关注语音的辅音特征，而额头高斯几乎不受音频影响。融合后的特征通过 MLP 预测每个高斯的位置、旋转和球谐偏移量。

\textbf{设计动机}：变形范式相比外观生成具有天然的平滑性和连续性优势：面部运动是空间上连续的变形，而非逐像素独立生成的颜色值。三平面编码利用了面部区域的局部相似性，交叉注意力实现了音频信号到空间位置的精确映射。

\textbf{公式与参数}：该方法以 120 FPS 的渲染速度实现了高面部保真度和唇形同步精度，使用分段式训练策略和光度损失、感知损失、唇形同步损失联合优化。

\paragraph{TalkingGaussian}

TalkingGaussian（2024.04）~\cite{2404_15264} 同样采用了基于变形的范式，但提出了与 GaussianTalker 不同的解耦策略。

\textbf{Baseline 与问题}：直接在动态面部区域生成逐点外观（appearance generation）时，外观空间的跳变和不连续性使得连续平滑的神经场难以精确拟合，导致面部扭曲伪影。

\textbf{改了什么}：TalkingGaussian 纯粹使用变形（deformation）来表示面部运动，通过平滑连续的变形驱动持久的高斯基元，而非直接修改点的外观颜色。

\textbf{实现原理}：关键发现是面部运动和口腔运动在空间和时间上具有不同的特性：面部表情变化是渐进且空间连续的，而口腔（嘴唇开合、牙齿暴露）是快速且空间不连续的。为此，该方法将模型分解为面部和口腔两个独立分支，每个分支有独立的变形场。面部分支专注于全局表情变化和唇形同步，口腔分支专注于嘴内部的几何变化。

\textbf{设计动机}：面部和口腔的解耦避免了单一变形场需要同时学习两种截然不同的运动模式，显著提高了各自的拟合精度。

\paragraph{GSTalker}

GSTalker（2024.04）~\cite{2404_19040} 学习音频驱动的高斯变形场，与 GaussianTalker 的三平面编码和交叉注意力策略不同，GSTalker 集成多分辨率哈希网格三平面（将 3D 空间特征编码为三个正交 2D 平面之和的紧凑表示）来编码规范高斯的空间邻域关系，并引入时间平滑模块。该模块通过最小化相邻帧高斯位置的加速度，学习精细面部细节的准确时序连贯变形。

\paragraph{GaussianSpeech}

GaussianSpeech（2024.11）~\cite{2411_18675} 提出了紧凑高效的 3DGS 化身表示，与 GaussianTalker 使用球谐系数编码颜色不同，GaussianSpeech 使用轻量级两层颜色 MLP 生成表情依赖的颜色，使面部皱纹和微表情随音频内容动态变化。此外引入体积感知剪枝策略（选择不透明度与体积乘积最大的 25K 高斯）和皱纹/感知损失来合成面部细节，并发布了新的大规模多视角音频-视觉序列数据集。

\subsubsection{实时与高效方法}

本路线介绍面向实时推理和高效部署的说话人头像方法。

\paragraph{EGSTalker}

EGSTalker（2025.10）~\cite{2510_08587} 从三个维度系统性地优化了 3DGS 说话人头像的效率与质量权衡。读者已从上一节了解到 GaussianTalker 的变形范式——维护规范 3D 高斯并通过音频驱动变形场预测逐帧偏移量。EGSTalker 继承了这一变形范式，但精准定位了其中三个效率瓶颈并逐一重新设计：用 Kolmogorov-Arnold 网络替代 MLP 以更少的参数实现更强的非线性映射能力，用轻量级代理注意力替代全交叉注意力以降低计算开销，并引入两阶段训练策略使训练数据需求从 10 分钟缩短到 3--5 分钟。

\textbf{Baseline 与问题}。GaussianTalker 等 3DGS 说话人方法使用标准 MLP 将空间特征映射到高斯参数、用完整交叉注意力融合音频与空间特征。这两处设计存在效率与质量的双重瓶颈。在静态高斯构建阶段，标准 MLP 的高维非线性映射存在信息损失——面部特征到高斯参数之间的复杂关系难以被常规 MLP 的分段线性激活函数（如 ReLU）精确捕捉。在动态变形阶段，标准交叉注意力的计算复杂度为 $O(N^2 d)$（$N$ 为空间点数，$d$ 为特征维度），当面部分辨率提高时计算量急剧膨胀，阻碍实时推理。此外，大多数方法需要 10 分钟以上的训练视频才能构建高质量的静态高斯。

\textbf{改了什么}。EGSTalker 将静态构建和动态变形两个环节中的 MLP 映射全部替换为 Kolmogorov-Arnold 网络（KAN），将音频-空间融合机制从全交叉注意力替换为高效空间-音频注意力（Efficient Spatial-Audio Attention, ESAA），并采用两阶段训练策略——第一阶段仅构建静态高斯，第二阶段再学习音频驱动的变形。

\textbf{实现原理}。整个框架分为静态初始化和动态变形两个阶段。

\textit{静态高斯初始化}。面部运动具有区域特性：空间上相邻的点遵循相似的运动轨迹，不同区域则呈现截然不同的运动模式。标准 3DGS 虽能有效建模静态头部，但难以捕获区域内和区域间的空间关系。为此，EGSTalker 设计了空间结构编码器，包含两个关键组件。

第一是\textbf{多分辨率哈希三平面编码器}。它将头部区域通过三个正交 2D 网格分层编码：每个 3D 空间点 $x = (x, y, z)$ 投影到这三个网格上，在多分辨率层级上插值后通过 Hadamard 积（逐元素乘法）融合，生成紧凑特征向量 $f_v(x)$。直觉上，三平面表示将 3D 空间分解为三个 2D 平面——稀疏区域（如头部边缘）的计算冗余被自动削减，密集区域（如嘴唇、眼睛）的细节得到保留。

第二是 \textbf{KAN 映射网络}。传统 MLP 使用固定激活函数（如 ReLU），本质上是用分段线性函数近似复杂的面部特征到高斯参数的映射，在高维空间中容易产生信息瓶颈。KAN（Kolmogorov-Arnold 网络）利用多项式近似，在网络的边上而非节点上学习可优化的 B 样条激活函数，能更精确地捕捉面部特征与高斯参数之间的复杂非线性关系。静态高斯参数的生成公式为：

\begin{equation}
    G_{\text{static}} = \{\mu_s, s_s, r_s, \text{SH}_s, \alpha_s\} = \text{KAN}(f_v(x))
\end{equation}

其中 $f_v(x)$ 为三平面编码的空间特征，输出包含位置 $\mu$、缩放 $s$、旋转 $r$、球谐系数 $\text{SH}$ 和不透明度 $\alpha$ 五组静态高斯参数。

\textit{高效音频驱动变形}。静态高斯构建完成后，第二阶段通过音频信号驱动动态变形。关键创新是\textbf{高效空间-音频注意力（ESAA）}模块。

现有方法的音频-空间融合存在两类问题：基于元素级调制（element-wise modulation）的方法假设静态 3D 坐标始终对应一致的面部区域，但面部是动态变化的，该假设经常不成立；GaussianTalker 使用标准交叉注意力虽能解决这个问题，但 $O(N^2 d)$ 的计算复杂度无法支持实时推理。ESAA 引入\textbf{代理令牌（agent token）}机制来解决这一矛盾：首先，一个轻量级 MLP 将空间特征 $f_v$ 投影为少量代理令牌 $A = \text{MLP}(f_v)$；然后，这些代理令牌分两步完成信息聚合与广播：

\begin{equation}
    \text{ACA}(f_v, f_a) = \text{Attn}\big(f_v, A, \text{SDP}(A, f_a, f_a)\big)
\end{equation}

其中 $\text{SDP}(Q, K, V) = \text{Softmax}(QK^T / \sqrt{d_k})V$ 为缩放点积注意力。直觉上，代理令牌 $A$ 先作为查询从音频特征 $f_a$ 中聚合关键信息（代理聚合），再将聚合后的信息作为值广播回所有空间点（代理广播）。代理令牌数量被限制为空间序列长度的 0.5\%，将复杂度从 $O(N^2 d)$ 降低至 $O(N \cdot n \cdot d)$（$n \ll N$），其中 0.5\% 的比例在消融实验中被证明是同步精度与推理速度的最佳平衡点。

此外，音频驱动的运动具有周期性特征（如重复的音节节奏），因此 EGSTalker 借鉴 FaceFormer 的思路，采用周期位置编码（Periodic Positional Encoding, PPE）来编码时间步信息：

\begin{equation}
    \text{PPE}(t, 2i) = \sin\!\Big(\frac{t}{p \cdot 10000^{2i/d}}\Big), \quad \text{PPE}(t, 2i\!+\!1) = \cos\!\Big(\frac{t}{p \cdot 10000^{2i/d}}\Big)
\end{equation}

其中 $t$ 为时间步，$d$ 为特征维度，$p$ 为周期参数。相比标准位置编码的固定周期，PPE 通过可调周期 $p$ 显式建模音频驱动运动的周期性模式，增强了时序连贯性。

除音频外，眨眼和头部姿态等非音频因素也影响面部运动。EGSTalker 将眼部特征 $f_e$、姿态嵌入 $f_p$ 和全局零向量 $f_z$（用于捕获与两者均无关的运动残差）投影到共享潜空间后与音频特征 $f_a$ 拼接，再通过 ESAA 与空间特征融合。

融合后的音频感知空间特征 $f_d$ 送入 KAN 预测逐帧变形偏移量：

\begin{equation}
    G_{\text{deform}} = \Delta\{\mu_n, s_n, r_n, \text{SH}_n, \alpha_n\} = \text{KAN}(f_d(x_n))
\end{equation}

最终第 $i$ 帧通过泼溅静态与动态高斯参数之和渲染：$I_i = \text{Splat}(G_{\text{static}} + G_{\text{deform}})$。

\textbf{设计动机}。EGSTalker 的三项核心设计各有明确的效率考量：（1）用 KAN 替代 MLP 并非为了增大模型容量，而是因为 KAN 的可学习激活函数能以更少的参数量达到同等甚至更好的非线性映射精度——消融实验显示，仅将 MLP 换为 KAN 就显著提升了重建质量（PSNR 提升），同时不增加推理延迟；PPE 的加入进一步增强了时序连贯性（在消融中与 ESAA 联合评估时对唇形同步有独立贡献）；（2）ESAA 的代理注意力机制的核心洞察是：并非所有空间点都需要直接与所有音频特征交互，少量代理令牌即可近似全局音频-空间依赖关系，从而在保留融合质量的同时大幅降低计算量；（3）两阶段训练策略将静态结构学习与动态变形学习解耦，使每个阶段可以独立优化——第一阶段仅需少量帧即可建立高质量的静态高斯分布，消融实验表明去除第一阶段会导致嘴唇、眼睛等关键区域的点密度显著降低，几何精度与音视频对齐质量随之退化。

\textbf{公式与参数}。EGSTalker 仅需约 6500 帧（约 4.3 分钟的 25 FPS 视频）即可完成训练，远少于大多数需要 10 分钟以上训练数据的方法。两个阶段的训练损失各有侧重。静态初始化阶段采用复合损失：

\begin{equation}
    L_{\text{static}} = L_{L_1} + \lambda_{\text{D-SSIM}} L_{\text{D-SSIM}} + \lambda_{\text{lpips}} L_{\text{lpips}}
\end{equation}

其中 $\lambda_{\text{D-SSIM}}$ 和 $\lambda_{\text{lpips}}$ 平衡结构相似性与感知质量。变形阶段在静态损失基础上额外引入唇部重建损失 $L_{\text{lip}}$（基于面部关键点裁剪的唇部区域计算）以提升唇形同步精度：

\begin{equation}
    L_{\text{def}} = L_{L_1} + \lambda_{\text{D-SSIM}} L_{\text{D-SSIM}} + \lambda_{\text{lpips}} L_{\text{lpips}} + \lambda_{\text{lip}} L_{\text{lip}}
\end{equation}

在单张 NVIDIA Tesla T4 GPU（16 GB 显存）上评测，EGSTalker 在 self-driven 设置下的 PSNR、LPIPS、FID 等图像质量指标和 LMD、LSE-D、LSE-C 等唇形同步指标均优于或持平 GaussianTalker 基线，同时达到了所有方法中最高的渲染帧率（FPS）。ESAA 模块中代理令牌占比的消融实验揭示了精度-速度的清晰折线：移除 ESAA（0\%）时 FPS 最高但唇形同步显著退化；引入少量代理令牌后 FPS 大幅提升；0.5\% 是同步精度与推理速度的最佳平衡点；进一步提升到 1.0\% 仅略微改善 LMD 但降低 FPS。

\paragraph{PGSTalker}

PGSTalker（2025.09）~\cite{2509_16922} 提出像素感知密度控制策略和轻量级多模态门控融合模块，与 EGSTalker 的代理注意力策略互补。

\textbf{Baseline 与问题}：标准 3DGS 的自适应密度控制仅基于屏幕空间梯度均值判断是否分裂或克隆高斯，这一全局统一的准则未考虑不同面部区域的运动幅度差异。嘴唇和眉毛在说话时运动幅度远大于额头和头皮，但标准致密化对两者施加相同的梯度阈值，导致动态区域（嘴唇）高斯不足而静态区域（额头）高斯冗余。

\textbf{改了什么}：PGSTalker 从密度控制和音频-空间特征融合两个维度改进说话人头像：(1) 引入像素感知密度控制策略，根据面部区域的运动需求自适应分配点密度；(2) 设计轻量级多模态门控融合模块（Multimodal Gated Fusion Module）替代简单的特征拼接来融合音频和空间特征，提升高斯变形预测的精度。

\textbf{实现原理}：(1)~像素感知密度控制：在归一化设备坐标（NDC，Normalized Device Coordinate）空间中计算多视角平均梯度，并引入像素感知准则——对不同图像区域赋予差异化的致密化权重。在嘴唇、眉毛等动态面部区域，降低梯度阈值使致密化更积极，增强细节表达；在额头、头皮等静态区域，提高阈值抑制冗余高斯的生成。这一策略使有限的高斯预算集中在视觉上最重要的动态区域。(2)~多模态门控融合模块：音频特征和空间特征通过门控机制融合。门控单元学习每个特征维度对最终变形预测的贡献权重，动态抑制不相关的信息通道，相比简单的特征拼接或加法融合更精确地建模音频-视觉的对应关系。

\textbf{设计动机}：说话人头像的渲染质量瓶颈往往不在高斯总数，而在高斯的空间分布。关键动态区域（嘴部、眉毛）的高斯密度不足是唇形模糊的主因，而静态区域的高斯冗余则浪费计算资源。像素感知密度控制通过重新分配高斯密度解决这一不匹配。

\textbf{公式与参数}：在公开数据集上的实验表明 PGSTalker 在渲染质量、唇形同步精度和推理速度上均优于现有 NeRF 和 3DGS 方法，同时展现出对未见身份的强泛化能力。

\paragraph{EmbedTalk}

EmbedTalk（2026.03）~\cite{2603_07604} 用学习嵌入（learnable Gaussian embeddings）替代传统的三平面编码来驱动时间变形，与 EGSTalker 的三平面+KAN 框架不同，EmbedTalk 为每个高斯学习一个独立的嵌入向量，通过音频条件化的 MLP 将嵌入映射为高斯属性偏移量。这种设计避免了三平面的 3D 到 2D 投影信息损失，在移动 GPU（RTX 2060 6 GB）上实现超过 60 FPS 的渲染，同时模型更加紧凑。

\paragraph{Audio-Plane}

Audio-Plane（2025.03）~\cite{2503_22605} 将 4D 体积分解为音频无关的空间平面和音频依赖的调制平面，形成紧凑可解释的音频平面表示（Audio Factorization Plane），并通过音频引导显著性泼溅机制——利用音频信号计算逐高斯的运动显著性权重，在渲染时自适应强调嘴部等高动态区域——提升唇形同步精度。

\subsubsection{唇形同步与情绪控制}

本路线介绍提升说话人头像唇形同步精度和情绪控制能力的方法。

唇形同步的准确性直接决定说话人头像的真实感。然而，说话人头像中长期存在一个被忽视的瓶颈，即\textbf{音素-视素对齐歧义}（phoneme-viseme alignment ambiguity）：不同的音素（如 /t/ 和 /d/、/k/ 和 /g/）对应极其相似的唇形（视素），仅靠渲染模型内部的弱对齐（特征相加或注意力机制）难以区分它们，导致唇形模糊或错位。

\paragraph{PASE}

PASE（2025.04）~\cite{2504_05803} 提出音素感知语音编码器（Phoneme-Aware Speech Encoder），首次将音素级别的语义先验显式引入音频-视觉对齐，将唇形同步精度提升了 13.7\%（NeRF 框架，LSE-C 指标）和 14.2\%（3DGS 框架，LSE-D 指标）。

\textbf{Baseline 与问题}。读者已经了解到 GaussianTalker 和 SyncTalk 等方法使用 DeepSpeech、HuBERT、Wav2Vec 2.0 或 Whisper 等通用声学特征作为音频条件驱动渲染模型。这些特征是为语音识别或说话人识别设计的，聚焦于声学域中音素的区分能力。当它们被直接用于说话人头像时，音素-视素之间的对齐完全依赖渲染模型内部的弱对齐机制（简单的特征拼接或注意力），而非显式的音素级别约束。这种弱对齐无法解决根本问题：不同音素可能映射到极其相似的唇形（例如 /d/ 和 /t/ 的唇部姿态几乎一致，区别仅在于声带是否振动），而通用声学特征并不编码这种视觉区分性。结果是唇形不准确（音素被混淆时唇形错位）、唇形模糊（模型无法学到语音信号与唇部动态之间的精确映射关系），甚至嘴部完全闭合时仍露出牙齿。

\textbf{改了什么}。PASE 将传统的「通用声学特征提取 + 渲染模型内隐式对齐」范式，替换为「STFT 全频段频谱 $\to$ GRU（Gate Recurrent Unit，门控循环单元）时序编码 $\to$ 音素嵌入增强的交叉注意力 $\to$ 对比学习」的显式音素级别对齐框架。具体而言，旧范式中 Mel 频谱图（压制 $>$4\,kHz 高频）+ 通用预训练编码器被替换为 STFT 频谱图（保留 0--8\,kHz 全频段）+ 8 层 GRU 编码器；旧范式中渲染模型内部的隐式对齐被替换为独立的音素级别对齐模块，该模块将音素嵌入融入音频特征作为查询向量，通过交叉注意力融合视觉特征，并用对比损失显式拉开正负样本对之间的距离。

\textbf{实现原理}。PASE 的核心思路是：利用语音和视觉模态共享的语义层——音素——作为跨模态对齐的锚点，既利用时间关系约束又利用稳定的语义结构。整个框架包含四个协作模块。

(1)~\textbf{STFT 音频编码器}。直觉上，唇形由口腔内发音器官的运动决定，而许多发音器官的运动线索藏在高频段——例如摩擦音 /s/ 和 /sh/ 的能量分布集中在 4--8\,kHz，但 Mel 频谱图为模拟人耳感知特性而压制了高频分量，恰好在视觉线索最丰富的频段丢失了信息。PASE 改用 STFT 频谱图（参数：n-fft $=$ 512，win-length $=$ 512，hop-length $=$ 128），保留完整的线性频率分辨率。频谱图输入 8 层 GRU（门控循环单元）网络，GRU 的更新门 $z_t$ 和重置门 $r_t$ 自适应调节时间建模的粒度——当音素切换时门控开度增大以捕获快速变化，稳态时门控收紧以保持稳定——比固定感受野的 CNN 更适合捕获音素级别的时序动态：

\begin{align}
    z_t &= \sigma(W_z [h_{t-1}, x_t]), \quad r_t = \sigma(W_r [h_{t-1}, x_t]) \\
    \tilde{h}_t &= \tanh(W_h [r_t \odot h_{t-1}, x_t]) \\
    h_t &= (1 - z_t)\, h_{t-1} + z_t\, \tilde{h}_t
\end{align}

其中 $\sigma$ 为 Sigmoid 函数，$h_{t-1}$ 为前一时间步的隐状态，$W_z$、$W_r$、$W_h$ 为可学习权重矩阵。最后一个 GRU 层的输出即为包含动态时序信息的语音嵌入 $A_p$。

(2)~\textbf{多通道 CNN 视觉编码器}。受 Wav2Lip 和 SyncNet 启发，设计多通道 CNN 从逐帧裁剪的唇部图像中提取时空视觉特征 $V_p$，输出与语音特征同维度的嵌入向量。唇部图像通过面部关键点（第 48--68 号关键点）检测后裁剪获得。

(3)~\textbf{预测与重构模块}。为增强模型对特征噪声和遮挡的鲁棒性，对音频特征 $A \in \mathbb{R}^{T \times D}$ 和视觉特征 $V \in \mathbb{R}^{T \times D}$ 独立生成二值掩码矩阵，随机选取部分时间步用可学习向量替换原始特征，训练时要求模型预测并重构被掩码的部分。重构损失为：

\begin{equation}
    \mathcal{L}_\text{rec} = \|V_\text{masked} - V_\text{orig}\|_2^2 + \|A_\text{masked} - A_\text{orig}\|_2^2
\end{equation}

该辅助损失有效缓解过拟合，并在某一模态缺失或被遮挡时保持模型的推理能力。

(4)~\textbf{音素级别对齐模块}。这是 PASE 的核心创新。对于给定音素 $p$，以音频特征 $A_p$ 为锚点（anchor），将音素标签嵌入为特征向量 $F_p$ 后融入音频特征生成查询向量：

\begin{equation}
    Q_p = A_p + F_p
\end{equation}

这一设计的直觉是：查询向量同时编码了「听起来是什么」（音频内容）和「属于哪个音素」（语义标签），使模型在融合时显式感知音素身份。随后以视觉特征 $V_p$ 为键和值，通过交叉注意力得到融合特征：

\begin{equation}
    F_\text{fusion} = \mathrm{CrossAttn}(Q_p,\, V_p,\, V_p)
\end{equation}

最后，对于正样本对 $(A_p, F^+_\text{fusion})$（同一音素的音频-视觉配对）和负样本对 $(A_p, F^-_\text{fusion})$（不同音素的配对），使用 InfoNCE 式对比损失显式拉开正负对距离：

\begin{equation}
    \mathcal{L}_\text{con} = -\log \frac{\exp\big(s(A_p, F^+_\text{fusion}) / \tau\big)}{\exp\big(s(A_p, F^+_\text{fusion}) / \tau\big) + \sum_{F^-_\text{fusion}} \exp\big(s(A_p, F^-_\text{fusion}) / \tau\big)}
\end{equation}

其中 $s(\cdot)$ 为相似度函数，$\tau$ 为温度系数。总体目标函数为 $\mathcal{L} = \mathcal{L}_\text{con} + \lambda\, \mathcal{L}_\text{rec}$，其中 $\lambda$ 为平衡重构损失的权重因子。

\textbf{设计动机}。PASE 的设计围绕三个核心观察。（1）音素-视素对齐歧义是说话人头像中唇形不准确的根本原因——不同音素（/t/ vs.\ /d/）在声学上截然不同但在视觉上几乎无法区分，通用声学特征不编码这种视觉区分性。（2）Mel 频谱图对高频的压制恰好丢弃了对唇形区分最有价值的频段——摩擦音等辅音的高频能量分布是唇形差异的关键线索。（3）音素是语音和视觉模态共享的语义层——利用这一稳定的语义结构作为对齐锚点，比仅依赖时间关系约束更可靠。STFT+GRU 的组合保留了全频段信息和音素级时序动态，音素嵌入+对比学习显式区分了视觉上易混淆的音素对。

\textbf{公式与参数}。在 LRS2 数据集上训练，使用 Montreal Forced Aligner 获取精确的音素-时间戳对齐标注。GRU 层数为 8，学习率 $5 \times 10^{-5}$，窗口大小 5，批大小 16。在 NeRF 框架（SyncTalk）上训练 60\,000 步（微调至 88\,000 步），在 3DGS 框架（TalkingGaussian）上训练 20\,000 步，均在单张 RTX 4090 上完成。相比最优基线（HuBERT），PASE 在 NeRF 模型上 LSE-C 提升 13.7\%、LSE-D 提升 14.2\%，在 3DGS 模型上优势进一步放大，且两项指标均接近真实视频水平。STFT+GRU 组合相比 Mel+GRU 在 NeRF 上 LMD 降低 8.0\%，相比 STFT+CNN 在 3DGS 上 LMD 降低 2.2\%。模型参数量仅 37.7M（Wav2Vec 2.0 为 317M，减少 8.8 倍），GPU 显存消耗降低 57.3\%（2051.5\,MB vs.\ 4805.3\,MB），推理时间仅比最轻量的 DeepSpeech 长 38.4\%（4.72\,s vs.\ 3.41\,s）。

\paragraph{VisualSpeaker}

VisualSpeaker（2025.07）~\cite{2507_06060} 提出感知唇读损失（perceptual lip-reading loss），与 PASE 从音频编码器端改进唇形同步不同，VisualSpeaker 从渲染监督端引入唇读知识。通过在大型唇读数据集上预训练视觉自动语音识别模型（Visual Speech Recognition），在训练期间用该模型评估 3DGS 化身渲染图像的"可读性"——即渲染的唇形能否被正确识别为对应的语音——作为辅助监督信号，将唇顶点误差降低了 56.1\%。

\paragraph{SyncTalk++}

SyncTalk++（2025.06）~\cite{2506_14742} 系统性地解决了说话人头像中的同步问题。与其他方法只关注唇形同步不同，SyncTalk++ 从三个维度统一建模同步：动态肖像渲染器（基于 3DGS 的身份保持渲染）确保身份一致性，面部同步控制器（使用在 LRS2 数据集上训练的音频-视觉同步编码器而非通用 ASR 特征）对齐唇形与语音并使用 3D 面部混合形状模型重建准确表情，头部同步稳定器通过平滑头部姿态估计消除抖动。此外还引入分布外（OOD）音频表情生成器和躯干修复器处理训练数据中未见过的语音内容。

\paragraph{TaoAvatar}

TaoAvatar（2025.03）~\cite{2503_17032} 实现了在 Apple Vision Pro 上以 90 FPS 运行的全身体说话人化身。与上述纯 3DGS 方法不同，TaoAvatar 采用教师-学生框架：教师模型使用 StyleUNet 生成高保真非刚性变形，然后通过知识蒸馏将 StyleUNet 的变形"烘焙"为轻量级 MLP 网络（非刚性变形烘焙+双层混合形状补偿），结合扩展的 SMPLX++ 参数化模型（包含衣物和头发的完整几何）和绑定在网格三角面片上的 3D 高斯实现高效移动端渲染。

在情绪控制方面，除唇形同步外，说话人头像还需具备根据情绪状态调整面部表情的能力。

\paragraph{EmoTalkingGaussian}

EmoTalkingGaussian（2025.02）~\cite{2502_00654} 通过唇对齐情绪面部生成器训练模型，使化身能根据连续情绪值（效价 valence 和唤醒度 arousal）操纵面部情绪。该方法继承 TalkingGaussian 的变形管线，关键改进在于用情绪条件化的图像生成器创建唇形同步且带有目标情绪的训练数据——通过法线图损失提升渲染质量，并通过同步损失约束渲染图像与输入音频的唇形一致性，避免情绪编辑时唇形失步。

\paragraph{EmoDiffTalk}

EmoDiffTalk（2025.12）~\cite{2512_05991} 引入动作单元（AU）提示的高斯扩散过程实现细粒度面部动画器，与 EmoTalkingGaussian 在效价-唤醒度空间操作不同，EmoDiffTalk 先用三平面预测高斯颜色属性替代球谐系数以获得更好的颜色归因，然后将音频编码为 AU 码驱动高斯位置等属性，配合文本到 AU 的情绪控制器建立从文本情绪描述到 AU 激活向量的直接映射，实现可扩展的动态情绪编辑。

\subsubsection{少样本与跨身份泛化}

本路线介绍实现说话人头像跨身份泛化和少样本适配的方法。

大多数说话人头像方法需要每人重新训练，限制了实用化。如何让一个模型同时为多个说话者服务，甚至泛化到训练中未见过的身份，是该方向的重要课题。

\paragraph{GenSync}

GenSync（2025.05）~\cite{2505_01928} 修改了 GaussianTalker 的训练与身份建模环节——用身份感知解耦模块（Identity-Aware Disentanglement Module）替代单身份隐式编码——3DGS 变形范式和空间-音频交叉注意力机制保持不变，使得单一模型能同时为多个说话者合成唇形同步视频，训练速度比每人单独训练的 GaussianTalker 快 6.8 倍。

\textbf{Baseline 与问题}。读者已经了解到 GaussianTalker 采用变形范式驱动 3DGS 面部动画，但其每个模型只能训练一个身份（10 个身份需训练约 62 小时）。先前多身份方法如 MI-NeRF 虽然尝试了多身份合成，但依赖逐帧 3DMM 表情提取（计算开销大且需逐身份优化），无法实现实时推理。核心瓶颈在于：不同说话者的唇形运动模式各不相同（有人说话嘴张得大、有人嘴张得小），如果不将"说什么"（音频驱动的运动）和"谁在说"（身份特征）显式分离，网络无法在多身份间共享运动知识。

\textbf{改了什么}。GenSync 引入身份感知解耦模块（受 MI-NeRF 的因子化身份-音频表示启发，但将条件从逐帧 3DMM 表情替换为音频嵌入，消除逐帧优化开销），将音频嵌入和身份向量映射到共享潜空间后做 Hadamard 积（逐元素相乘）交互，取代 GaussianTalker 中隐式绑定单一身份的音频-运动映射。3DGS 规范表示、三平面编码和空间-音频交叉注意力模块均保持不变。此外，GenSync 采用两阶段训练策略：第一阶段为每个身份独立优化 8000 次迭代的规范高斯；第二阶段将所有身份联合训练 50000 次迭代，共享解耦模块和注意力网络。

\textbf{实现原理}。身份感知解耦模块的设计直觉是：说话时的唇形运动由"说什么"和"谁在说"共同决定，两者之间不是简单的线性叠加，而是存在非线性耦合——同一个音素，不同人的嘴形幅度和节奏各不相同。因此模块将音频嵌入 $a$ 和可学习身份向量 $i$ 分别通过投影矩阵 $U_1$、$U_2$ 映射到共享潜空间，再做 Hadamard 积捕获非线性交互，同时保留各模态的独立贡献：

\begin{equation}
    M(a, i) = C\big[(U_1 a) \odot (U_2 i)\big] + W_2 a + W_3 i
\end{equation}

其中 $\odot$ 表示 Hadamard 积（逐元素乘法），$C$、$W_2$、$W_3$ 为可学习权重矩阵。直觉上，$C[(U_1 a) \odot (U_2 i)]$ 项让音频和身份在潜空间中"对话"——每个音频特征维度都被身份特征调制，产生个性化的运动指令；$W_2 a + W_3 i$ 则保留音频和身份各自不依赖对方的独立贡献，避免信息丢失。

解耦模块的输出随后送入融合空间-音频注意力网络。该网络继承 GaussianTalker 的空间-音频交叉注意力结构，但输入从纯音频特征变为解耦后的身份条件化音频特征。具体地，以规范高斯特征 $f(\mu_c)$ 为查询（query），以解耦模块输出 $M(a,i)$、眼部特征 $e$ 和视角 $v$ 的拼接为键值（key/value），通过交叉注意力计算融合的空间-音频嵌入 $h_n$，再由变形 MLP $F_d$ 预测每个高斯各属性（位置、旋转、缩放、颜色等）的变形偏移。其中眼部特征 $e$ 由 OpenFace 提取，面部几何先验由 Basel Face Model 提供。

\textbf{设计动机}。将身份和音频解耦而非混合编码有三个好处：(1) 模型可以在身份间共享音频驱动的通用唇形运动知识，只需学习每个身份的紧凑向量（而非完整模型），大幅减少训练开销；(2) 解耦使身份嵌入可以跨身份交换——实验表明，将身份 A 的嵌入（说话嘴张大）替换到身份 B 的渲染中，身份 B 的唇形幅度会增大，验证了身份嵌入确实捕获了说话风格；此外，跨音频分布测试显示，用身份 A 的音频驱动其他身份时合成唇形依然合理，即使驱动音频来自不同性别，模型也能生成同步的唇部运动，说明解耦有助于泛化；(3) 直接以音频嵌入为条件（而非逐帧 3DMM 表情）避免了昂贵的表情提取和逐帧优化，实现实时推理。

\textbf{公式与参数}。在 10 个身份的实验中，GenSync 联合训练总耗时约 9 小时，而 GaussianTalker 为每个身份单独训练需约 62 小时（6.8 倍加速）。每个身份的规范高斯在第一阶段独立优化 8000 次迭代，第二阶段联合训练 50000 次迭代共享网络。在 NVIDIA A30 GPU 上评测，GenSync 的 LPIPS 为 0.078（GaussianTalker 0.073）、FID 为 21.59（20.51），视觉质量与基线相当；唇形同步分数（SyncNet）为 11.98，优于 GaussianTalker 的 12.26（分数越低越好）。视频裁剪至 512$\times$512 分辨率，训练/测试帧比例 9:1。

\paragraph{MGGTalk}

MGGTalk（2025.04）~\cite{2504_00665} 仅需单目数据集即可泛化到未见身份，与 GenSync 的多身份联合训练策略不同，MGGTalk 利用深度信息增强几何特征并通过面部对称性补充外观特征——当一侧面部被遮挡时，镜像另一侧的特征来推断完整外观，实现从野外单目数据到未见身份的零样本泛化。

\paragraph{SDTalk}

SDTalk（2026.05）~\cite{2605_09956} 是一次拍摄的 3DGS 说话人框架，通过结构化面部先验（可见分支精确拟合可观察表面，完成分支利用 DINOv2 全局特征预测被遮挡的内部结构如牙齿和口腔——这一"结构化面部先验"将人脸重建分为可见区域和遮挡区域两条路径，前者依赖光度损失精确重建，后者借助基础模型的语义特征从上下文推断不可见结构）和双分支运动场（可见区域的精确重建 + 未观察区域的先验补全）实现跨身份泛化。

\paragraph{InsTaG}

InsTaG（2025.02）~\cite{2502_20387} 提出身份无关预训练和运动对齐适应策略，与 MGGTalk 的单目泛化不同，InsTaG 采用预训练+微调范式：先在大规模多身份数据上训练通用的音频-运动映射网络（身份无关预训练），再用几秒目标身份视频通过运动对齐适应快速个性化，将个性化时间压缩到秒级。

\paragraph{EmoTaG}

EmoTaG（2026.03）~\cite{2603_21332} 在 FLAME 参数空间而非直接在 3D 高斯上预测运动，与 EmoDiffTalk 的 AU（Action Unit，面部动作单元）驱动策略不同，EmoTaG 引入门控残差运动网络（GRMN），通过身份条件编码器融合音频和情绪信号的 AdaIN 调制，专家运动解码器的三个互补分支（残差、门控和语义情绪引导）协同预测 FLAME 表情参数和下颌姿态。

\paragraph{Splat-Portrait}

Splat-Portrait（2026.01）~\cite{2601_18633} 自动学习将单张肖像图像解耦为静态 3DGS 重建和预测的 2D 背景，无需运动先验即可生成自然的唇形运动——采用基于 Splatter-Image 的静态生成器从单目视频帧预测像素对齐的高斯参数，音频条件动态解码器在微调阶段注入语音驱动的运动，同时用 SDS 从极端视角精修外观。

手部是人体中自由度最高、最具表现力的部位之一。3DGS 的灵活表示能力使其适合建模手部的复杂关节运动和精细外观。从本节开始，讨论重点从头部和身体转向手部等局部部位的建模。

\subsubsection{可动画手部化身}

本路线介绍基于~3DGS~的高保真可动画手部化身建模方法。

\paragraph{HandSCS}

HandSCS（2025.03）~\cite{2503_14736} 提出结构坐标空间（Structural Coordinate Space, SCS）。

\textbf{Baseline 与问题}：手部具有极高的关节自由度（29 个自由度），且不同手指之间具有解剖学关联（如无名指弯曲时小指倾向于跟随）。标准 3DGS 的逐高斯独立优化忽略了这种结构化运动模式，导致姿态泛化能力差。

\textbf{改了什么}：HandSCS 为每个高斯提供来自姿态内和姿态间的显式结构引导，弥合稀疏骨骼关节和密集高斯基元之间的差距。

\textbf{实现原理}：SCS 包含两个互补组件。在姿态内层面，结构坐标基由静态骨骼和动态骨骼混合构成。静态骨骼从 MANO 运动学模型导出，并增加跨手指连接以捕获相邻手指间的解剖学关系；动态骨骼通过在静态骨骼段上进行注意力加权的插值构建，自适应建模姿态依赖的结构变化。每个高斯在 SCS 中具有角度-径向描述符，编码其相对于最近骨骼的位置。在姿态间层面，姿态间一致性损失要求相似关节配置下的高斯属性保持一致：如果两个手姿态中食指的弯曲角度相同，则食指区域的高斯应具有相似的属性。

\textbf{设计动机}：稀疏骨骼提供了全局结构约束，但无法直接指导密集高斯的行为。SCS 通过静态+动态骨骼的混合坐标系统，将骨骼的结构信息"传播"到每个高斯，使其在优化时具有明确的解剖学先验。

\paragraph{JGHand}

JGHand（2025.01）~\cite{2501_19088} 提出一种仅由 3D 关键点驱动的手部化身方法，无需依赖参数化手部模型。其核心是基于关键点的可微分骨架变换：通过层级化的旋转-缩放-对齐流程（$B = PK'F'RFK$），从规范骨架出发依次执行骨骼单位化、局部坐标系旋转、目标骨骼长度恢复和手掌平面对齐，实现对任意骨骼长度和姿态的精确变形。在规范空间中，高斯用 $(u,v,d)$ 归一化坐标参数化并绑定到最近骨骼，通过 Fast-SNARF 权重场插值获得连续的蒙皮权重，再经线性混合蒙皮（LBS）变换到姿态空间。此外，JGHand 借鉴 Screen Space Ambient Occlusion 思想，基于渲染深度图设计可微分阴影计算层，通过卷积核采样模拟手指自遮挡产生的阴影，无需额外训练数据即可提升真实感。

\paragraph{WildGHand}

WildGHand（2026.02）~\cite{2602_20556} 针对野外视频中手部化身重建面临的遮挡、运动模糊和光照变化等扰动问题，提出动态扰动解耦与扰动感知优化两大策略。动态扰动解耦（DPD）模块将扰动显式建模为 3D 高斯属性上的时变偏差 $\Delta g_l = \sigma_l \cdot \phi([z_l, g_l])$，其中时间嵌入 $z_l$ 经帧索引的位置编码获得，全局缩放因子 $\sigma_l \in (-1, 1)$ 控制偏差影响强度——扰动越强的帧获得更大的 $\sigma_l$，训练时以 30\% 概率将 $\sigma_l$ 置零以鼓励主网络捕获手部外观的主成分。扰动感知优化（PAO）策略利用 SAM 分割输入图像为多个区域，根据每个区域的重建误差 $E_u$ 和手部前景比例 $\eta_u$ 自适应生成各向异性加权掩码，使受扰动影响的区域获得更低的损失权重，从而在 PSNR 上实现高达 15.8\% 的相对提升。

\subsubsection{手-物交互}

本路线介绍基于~3DGS~的手-物交互重建与建模方法。

\paragraph{InteractAvatar}

InteractAvatar（2025.04）~\cite{2504_07949} 首次建模手-脸交互的光真实外观。其采用混合 mesh-Gaussian 表示，将高斯绑定到 FLAME（面部）和 MANO（手部）网格面上，随网格变形而刚性变换，兼顾 mesh 的变形可控性和 3DGS 的细节表现力。针对手部的高自由度运动，动态精修模块通过两组 MLP 分别建模几何动态（位置、缩放、旋转偏移）和外观动态（颜色、透明度偏移），捕获关节弯曲时的细皱纹和自投射阴影等姿态依赖变化。手-脸交互模块则利用 DECAF 的基于位置动力学（PBD）的碰撞求解器获取顶点偏移，并通过 PointNet++ 提取手-脸全局几何特征，结合相对位姿和面部表情参数，预测交互引起的非刚性变形和外观变化；距离加权的偏移调制确保远离接触点的区域不受交互模块影响。

\paragraph{GHOST}

GHOST（2026.03）~\cite{2603_18912} 修改了 3DGS 在手-物联合重建中的表示与优化环节——用 2D 高斯盘替代 3D 高斯椭球表示物体与手部，引入几何先验一致性损失和抓取感知对齐替代纯光度优化——多视角渲染管线保持不变。

\textbf{Baseline 与问题}：从单目 RGB 视频重建手-物交互面临三重挑战。首先，标准 3DGS 使用 3D 高斯椭球表示物体，在稀疏视角和频繁遮挡下，被手指遮挡的物体区域无法获得有效监督，导致重建出现空洞或手指形状的"洞"。其次，SfM 恢复的物体点云与 HaMeR 估计的手部平移存在尺度不一致和空间漂移，导致手与物体在 3D 空间中并不真正接触。最后，手部是高度变形的非刚体，标准 3DGS 的全局高斯集合无法跟随 MANO 网格的姿态变化。

\textbf{改了什么}：GHOST 从表示、几何和空间对齐三个层面进行改进。表示层面，物体和手部均采用 2D 高斯盘（2DGS）而非 3D 高斯椭球，提供更紧凑的表面建模。几何层面，引入基于大规模 3D 物体数据库（Objaverse）的几何先验检索与对齐机制，以及几何一致性损失 $L_{geo}$ 替代仅依赖光度监督的范式。空间对齐层面，提出抓取感知（grasp-aware）联合优化，同时精修手部平移 $T_t^h$ 和物体全局尺度 $S_o$。

\textbf{实现原理}：整个管线分为三个阶段。

\emph{第一阶段：预处理与几何先验检索。}给定输入视频 $V = \{I_t\}_{t=1}^T$，先用 SAM2 分割物体掩码 $M_t^o$，再通过 SfM 估计相机内参 $K$、相机位姿 $(R_t^c, T_t^c)$ 和稀疏物体点云 $P_{sfm}$。为弥补遮挡导致的不完整，GHOST 利用视觉语言模型（如 InternVL）从采样帧提取物体文本描述 $d$，通过 OpenShape 的 CLIP 模型计算文本嵌入 $\phi_{txt}(d)$，在 Objaverse 的预计算物体嵌入中执行最近邻搜索，检索 $k$ 个候选网格。然后对每个候选网格优化仿射变换（旋转 $R_p$、平移 $T_p$、缩放 $S_p$），使投影轮廓与物体掩码的 IoU 最大：

\begin{equation}
L_{IoU} = 1 - \text{IoU}\bigl(\pi_t([R_p|T_p][S_p \cdot p]), M_t^o\bigr),
\end{equation}

其中 $\pi_t$ 是相机投影函数。最终选取 $L_{IoU}$ 最小的对齐网格作为物体几何先验 $O \in \mathbb{R}^{N \times 3}$。

\emph{第二阶段：抓取感知的手-物对齐。}GHOST 首先检测"抓取帧"：计算物体运动 $\Delta T^o$ 和每只手运动 $\Delta T^h$ 在 $xy$ 平面上的余弦相似度 $c_h$，当 $c_h > \tau_{sim}=0.5$ 时标记为抓取帧。在抓取帧上，GHOST 联合优化手部平移 $T_t^h$ 和物体尺度 $S_o$，目标函数包含三个互补项：
\begin{itemize}
    \item 接触损失 $L_{contact}$——Chamfer 距离度量手部网格与缩放后几何先验之间的距离，确保手与物体真正接触：
    $L_{contact} = \text{CD}(T_t^h + V_t^h,\; [R_t^c | T_t^c][S_o \cdot O])$；
    \item 2D 投影损失 $L_{proj}$——将平移后的手部关节重投影回图像，与 HaMeR 预测的 2D 关节对齐，保持 2D 一致性：$L_{proj} = \|\pi_t(T_t^h + J_t^h) - J_{2D,t}^h\|_1$；
    \item 时序平滑损失 $L_{temp}$——约束相邻帧的手部平移变化量：$L_{temp} = \frac{1}{T-1}\sum_{t=1}^{T-1}\|T_{t+1}^h - T_t^h\|_2^2$。
\end{itemize}
三者加权组合为 $L = \lambda_1 L_{contact} + \lambda_2 L_{proj} + \lambda_3 L_{temp}$，经验设定 $\lambda_1=10^3,\; \lambda_2=10^{-1},\; \lambda_3=10$。

\emph{第三阶段：手-物联合 Gaussian Splatting 优化。}物体首先用 2D 高斯盘表示为 $\mathcal{G}_o = \{c_o, r_o, s_o, \alpha_o, SH_o\}$，并引入两个关键损失。其一是手感知背景损失 $L_{bkg,h}$：标准背景损失 $L_{bkg}$ 在手遮挡物体的帧中会错误地将被手指覆盖的物体区域当作背景删除，$L_{bkg,h}$ 改用手+物体联合掩码区分前景与背景，避免手指遮挡造成空洞。其二是几何一致性损失 $L_{geo}$：

\begin{align}
L_{out} &= \text{ReLU}(\|c_o - O\|_2 - \tau_{out})^2, \quad L_{fill} = \text{ReLU}(\|O - c_o\|_2 - \tau_{fill})^2, \\
L_{geo} &= \lambda_{geo}(L_{out} + L_{fill}),
\end{align}

直觉上，$L_{out}$ 惩罚离先验表面过远的高斯（防止漂移），$L_{fill}$ 惩罚先验表面上无高斯覆盖的大间隙（鼓励填补遮挡空洞），其中 $\tau_{out}$、$\tau_{fill}$ 分别为偏离和空洞的距离阈值，$\lambda_{geo}=5$。手部则将 2D 高斯盘绑定到 MANO 网格面上，随网格变形而刚性变换，采用两阶段策略：先优化物体高斯，再冻结物体、仅优化手部高斯的外观参数（缩放 $s_h$、不透明度 $\alpha_h$、球谐 $SH_h$），网格驱动的位置和旋转保持固定。

\textbf{设计动机}：GHOST 的核心洞察是"先验引导 + 分阶段优化"的解耦策略。几何先验 $O$ 提供了物体在遮挡区域的"形状猜想"，使 $L_{geo}$ 能够引导高斯填满真实表面而非依赖稀疏的光度监督。抓取感知对齐利用了物理常识——物体运动必然由手部驱动——将接触约束集中在物体实际被操控的帧上，避免了在非交互帧引入错误接触信号。两阶段优化将非刚体的手部与刚体物体分离，让物体先在无手干扰的条件下收敛，再联合精修，避免了手-物高斯之间的梯度干扰。

\textbf{公式与参数}。GHOST 在 ARCTIC、HO3D 和 in-the-wild 三个数据集上进行评估。在 ARCTIC 数据集上，GHOST 在物体 Chamfer Distance 上比 HOLD（NeRF 方法，每序列优化约数小时）降低约 30\%，在 PSNR 和 SSIM 渲染指标上同样取得最优。与同为 Gaussian-based 的 BIGS 相比，GHOST 在物体几何精度和手部重建质量上均有显著提升，尤其在严重遮挡场景下优势更为明显——BIGS 缺乏几何先验，被手指遮挡的物体区域容易出现空洞，而 GHOST 的 $L_{fill}$ 损失有效引导高斯填满这些区域。GHOST 的每序列优化时间约为 BIGS 的 $1/13$，在保持更高精度的同时实现了数量级的加速。运行效率方面，几何先验检索阶段在 Objaverse 中执行 $k$ 近邻搜索约需 1--2 分钟，抓取感知对齐和 Gaussian Splatting 优化共需约 10--15 分钟（单张 RTX 3090），整个流程无需多视角监督，仅依赖单目 RGB 输入。

\paragraph{BIGS}

BIGS（2025.04）~\cite{2504_09097} 解决了从单目 RGB 视频进行双手类别无关交互重建的问题。其核心设计是双手共享高斯：在规范空间中仅优化单手（右手）的 3D 高斯，通过对 $x$ 轴翻转将其映射到左手，使两只手的 3D 信息在同一规范空间中有效积累，缓解双手严重遮挡下的信息不足。手部和物体各自使用 TriplaneNet 提取特征，分别通过外观、几何和变形 MLP 预测高斯参数，再经 LBS 或刚体变换映射到图像空间。对于被遮挡的物体不可见区域，BIGS 引入基于预训练扩散模型的分数蒸馏采样（SDS）损失，结合文本反转学习物体描述和 PiDi 边界条件 ControlNet 引导，在不依赖物体形状先验的情况下重建完整物体几何。

\paragraph{Interaction-Aware 4D Gaussian Splatting}

Interaction-Aware 4D Gaussian Splatting（2025.11）~\cite{2511_14540} 提出交互感知的手-物高斯表示，将动态 HOI 场景分解为手部、物体和背景三组独立高斯分别建模。其关键改进是为每个高斯引入可优化的权重 $w$ 和半径 $o$：$w$ 通过 RBF 核函数建模邻近高斯间的空间权重，实现类 LBS 的分段线性变换以获得更清晰的轮廓；$o$ 控制局部影响范围，小半径对应更锐利的几何轮廓。物体场 $F_O$ 首次将手部位置作为额外输入，构建交互感知的动态变形场，使物体运动能够感知并响应手部操作。配合渐进式五阶段优化策略（初始化、预热、HOI 精修、背景优化、协同重建），仅依赖轻量级手部追踪器提供粗 3D 引导即可实现高质量重建。

\paragraph{HOGSA}

HOGSA（2025.01）~\cite{2501_02845} 提出基于 3DGS 的双手手-物交互数据增强框架，通过基于网格的 3DGS 分别建模左手、右手和物体，并扩展单手抓取姿态优化模块为双手生成多样化的手-物交互姿态。

\textbf{Baseline 与问题}。双手手-物交互理解（姿态估计、接触估计）依赖大规模标注数据集，但由于双手与物体之间存在严重遮挡和高度自由度的运动，收集和标注高质量的大规模数据集极为困难。已有的数据增强方法存在三类局限：（1）传统渲染管线（如 Corona et al.）需要复杂的 3D 扫描和后处理来获取高质量形状与纹理贴图，还需要手部模型的蒙皮权重，专业性要求高且渲染结果缺乏真实感；（2）基于 NeRF 的方法（如 HO-NeRF）虽能生成多样化数据，但需要离线建模，渲染过程耗时，无法用于大规模数据增强；（3）已有的 3DGS 数据增强方法（如 MANUS）仅针对单手-物体交互，无法扩展到双手场景——双手场景中的遮挡更为严重，模型更难学习有效的姿态先验。此外，由于多分辨率图像输入导致的渲染模糊和标注不精确，合成数据与真实数据之间的差距会降低模型性能。

\textbf{改了什么}。HOGSA 从建模、姿态扩展和渲染质量三个层面进行改进。建模层面，采用基于网格的 3DGS（mesh-based 3DGS）将高斯核定义在 MANO-HD 手部网格和物体网格的三角面上，而非使用全局自由高斯，使高斯随网格变形自动移动，支持任意姿态驱动渲染。姿态扩展层面，将单手抓取姿态优化模块扩展为双手版本，生成多样化的双手交互姿态以扩大数据集的姿态分布。渲染质量层面，设计超分辨率模块提升 3DGS 粗渲染图像的质量，弥合合成数据与真实数据之间的差距。

\textbf{实现原理}。整个流程分为三个模块。

\emph{第一模块：Hand-Object Gaussian Splatting（HOGS）。}HOGSA 为每个交互序列分别构建左手、右手和物体的 3DGS 模型。高斯核的初始化方式借鉴 GaMeS，将高斯中心定义在网格三角面的重心坐标上：$\mathbf{x}_c = \beta_1 \mathbf{v}_1 + \beta_2 \mathbf{v}_2 + \beta_3 \mathbf{v}_3$，其中 $\{\beta_i\}$ 为可训练参数且 $\beta_1 + \beta_2 + \beta_3 = 1$，$\{\mathbf{v}_i\}$ 为三角面顶点。旋转矩阵 $\mathbf{R}$ 由三角面的法线和切向量构造，缩放矩阵 $\mathbf{S}$ 由三角面的几何尺寸决定。手部运动通过 MANO-HD 的蒙皮权重和骨骼变换矩阵驱动：$\mathbf{x}_t^h = (\sum_{i=1}^{n} w_i \mathbf{B}_i) \cdot \mathbf{x}_c^h$，物体则通过 $SE(3)$ 刚体变换驱动：$\mathbf{x}_t^o = \mathbf{T} \times \mathbf{x}_c^o$。通过改变姿态参数和相机参数，HOGS 可以渲染任意交互姿态的任意视角图像。

\emph{第二模块：姿态优化模块（Pose Optimization Module, POM）。}为扩大数据集的姿态多样性，HOGSA 将已有的单手抓取姿态优化方法扩展到双手场景。POM 以原始数据集中的手-物姿态为种子，在保持抓取物理合理性的前提下，优化生成多样化的双手交互姿态。这些新姿态被输入 HOGS 模型，渲染出大量新的交互图像。

\emph{第三模块：超分辨率模块（Super-Resolution Module, SRM）。}由于输入图像分辨率不一致，3DGS 渲染出的图像常常模糊。SRM 对粗渲染图像进行超分辨率处理，提升图像的视觉质量，使其更接近真实照片。渲染后的新图像进一步与 COCO2017 背景图像融合并裁剪至 $224 \times 224$ 分辨率，构成最终的增强数据。

\textbf{设计动机}。HOGSA 的设计围绕双手手-物交互数据的三个核心瓶颈展开。数据量瓶颈通过 3DGS 的高效渲染解决——相比 NeRF 的离线建模和慢速渲染，3DGS 可以快速生成大量新视角图像。姿态多样性瓶颈通过 POM 解决——原始数据集中的姿态分布有限，POM 在物理约束下扩展姿态空间，使增强数据覆盖更多交互模式。真实感瓶颈通过 SRM 解决——3DGS 的渲染质量受多分辨率输入限制，超分辨率模块弥合了合成图像与真实图像之间的视觉差距。三者组合形成端到端的增强管线：输入原始交互序列，输出大规模、多姿态、高真实感的双手交互数据。

\textbf{公式与参数}。HOGSA 在 H2O 和 ARCTIC 两个基准数据集上进行评估。增强数据与原始数据混合后用于微调双手交互理解的基线模型。实验表明，使用 HOGSA 增强数据后，基线模型在手-物姿态估计精度和接触估计准确性上均有提升，在 H2O 和 ARCTIC 上达到最优水平。消融分析系统性地验证了三个模块各自的贡献：POM 带来的姿态多样性对模型泛化能力影响最大，SRM 对缩小合成-真实差距贡献显著，而背景融合策略进一步提升了模型的鲁棒性。该方法已被 AAAI 2025 接收。

虚拟试穿和服装编辑是电子商务、游戏和社交媒体中的核心应用。3DGS 的显式表示使其天然支持服装的分离、替换和重新组合。与前文讨论的"从视频中重建并分解已有服装"不同，本节面向更实用的场景：给定目标服装的图片，将其"穿"到已重建的人体上。

\subsubsection{服装-人体分离与重建}

本路线介绍将服装与人体几何分离并进行独立~3DGS~重建的方法。

\paragraph{GGAvatar}

GGAvatar（2024.11）~\cite{2411_09952} 修改了 3DGS 在穿衣人体重建中的表示初始化与训练策略——为每件服装构建独立的参数化模板替代单层点偏移策略，用分阶段训练（隔离训练 + 联合训练）替代端到端训练——3DGS 的高斯泼溅渲染管线保持不变。

\textbf{Baseline 与问题}。前述 GART、HUGS、GaussianAvatar 等 3DGS 化身模型均采用"点偏移"（points offset）策略：在 SMPL 网格表面叠加一层偏移量来近似外层服装的位置。这种单层表示面临两个根本局限。第一，穿衣人体至少包含两层几何——皮肤层和服装层，仅用一层偏移无法同时表示两者，更无法将它们分离。第二，服装具有与身体截然不同的拓扑结构和运动模式（宽松衣袖的摆动、裙摆的飘逸），单层偏移将身体和服装"绑死"在一起，使得服装编辑和虚拟试穿等应用无法实现。即使是已有的分离式模型（如 SCARF、DELTA），也需要大量 3D 扫描数据进行多视角训练，且不同服装之间的分离能力有限。本节从服装-人体分离、虚拟试穿和服装重建三个角度展开讨论，它们与前文讨论的"整体化身重建"形成互补。

\textbf{改了什么}。GGAvatar 对 3DGS 化身的表示和训练做出三处关键修改。(1)~在表示层面，为每件服装 $c \in \{1, 2, \ldots, n_c\}$ 构建独立的参数化模板 $\mathbf{T}^{(c)}(\beta, \theta)$，而非在身体模板上叠加单层偏移。身体和各件服装各自拥有一组独立的高斯集合 $\mathcal{V}^{(b)}$ 和 $\{\mathcal{V}^{(c)}\}$。(2)~在变形层面，为服装高斯引入可学习的蒙皮权重偏移 $\Delta w$，使服装的运动不完全依赖身体的蒙皮权重，能够表达宽松服装的独立运动。(3)~在训练层面，采用两阶段策略：先"隔离训练"各层高斯（仅优化本层、冻结其他层），再"联合训练"全部高斯（不再增删高斯、仅优化属性），防止身体和服装高斯在训练过程中相互穿透。

\textbf{实现原理}。整个流程分为服装模板估计、高斯表示与变形、分阶段渲染三个阶段。

第一阶段是服装模板估计（Garment Templates Estimation），解决"如何获得每件服装的初始 3D 形状"这一起点问题。直觉上，如果能为每件衣服分别建立一个 3D 网格模板，那么身体和各件服装就有了各自的"骨架"，后续分离就水到渠成。具体地，首先用 FrankMocap 估计人体姿态参数 $\beta$（体型）和 $\theta$（姿态），选择正面视角与 SCHP（Self Correction for Human Parsing）的语义分割结果对齐，将分割后的服装区域送入 ISP（Implicit Surface Prediction）模型。ISP 通过两个 MLP 分别预测服装正面和背面的 SDF 值与 3D 形状，再按缝合顺序编号将前后片合并为完整的服装网格 $\mathbf{T}^{(c)}(\beta, \theta)$。所有服装模板与身体 SMPL 模型存储在统一的规范空间（canonical space）中，共享相同的 $\beta$、$\theta$ 参数和关节对应 $J(\beta)$，但拥有各自独立的蒙皮权重 $\mathcal{W}^{(c)}$，确保变形时身体和服装保持空间对齐。服装模板的顶点提取后作为对应服装高斯的初始位置。

第二阶段是高斯表示与可学习变形。身体和各件服装分别拥有一组 3D 高斯，每个高斯包含位置 $\mu^{(i)}$、旋转 $R^{(i)}$、缩放 $s^{(i)}$、不透明度 $\eta^{(i)}$ 和球谐颜色 $f^{(i)}$ 五种属性（身体高斯还额外附加可学习的顶点偏移 $O \in \mathbb{R}^{|\mathcal{V}^{(b)}| \times 3}$，以捕捉 SMPL 模板无法表达的局部几何细节）。关键创新在于可学习蒙皮权重：标准 3DGS 化身使用固定的 SMPL 蒙皮权重驱动高斯变形，而 GGAvatar 为每个高斯增加一个可学习的偏移量 $\Delta w_j$，使得服装区域的高斯能够学习到与身体不同的运动模式。具体公式为：

\begin{equation}
\tilde{\mathcal{W}}_k(\mu^{(i)}) = \mathcal{W}_k(\mu^{(i)}) + \Delta w_j, \quad i \in \mathcal{V}, \ k \in \{1, 2, \ldots, n_k\}
\end{equation}

其中 $\mathcal{W}_k(\mu^{(i)})$ 是 SMPL 提供的第 $k$ 个关节的初始蒙皮权重，$\Delta w_j$ 是可学习的偏移量，$n_k = 23$ 是关节数量。利用调整后的蒙皮权重和骨骼变换 $\{B_1, B_2, \ldots, B_{n_k}\}$，规范空间中的高斯通过线性混合蒙皮（LBS）变形到观测空间：

\begin{equation}
\mathcal{T}^{(i)} = \sum_{k=1}^{n_k} \tilde{\mathcal{W}}_k(\mu^{(i)}) B_k, \quad \mu' = \tilde{\mathcal{T}}^{(i)}_{(1:3,1:3)} \mu^{(i)} + \tilde{\mathcal{T}}^{(i)}_{(4,1:3)}
\end{equation}

即每个高斯的观测位置是所有骨骼变换的蒙皮权重加权平均，其中权重矩阵 $\mathcal{T}^{(i)}$ 的左上子矩阵对应旋转、右侧向量对应平移。值得注意的是，旋转同样通过该权重矩阵变形：$R' = \tilde{\mathcal{T}}_{(1:3,1:3)} R$，而缩放、不透明度和球谐颜色在变形过程中保持不变。这意味着宽松袖口处的高斯可以通过学习较大的 $\Delta w$ 来实现与身体骨骼不同的摆动轨迹，同时旋转也随蒙皮权重偏移而自适应调整。

第三阶段是分阶段渲染与损失函数。在隔离训练阶段，身体和各件服装的高斯分别独立优化，使用标准的 3DGS 致密化-修剪策略，每层的重建损失仅在对应 SCHP 分割区域内计算，防止不同层的高斯相互入侵。隔离阶段的损失包括 $\mathcal{L}_1$ 重建损失、掩码损失 $\mathcal{L}_{mask}$（确保各层高斯仅在对应分割区域内渲染）、S3IM 损失和高斯正则化 $\mathcal{L}_{G\_reg}$。其中 S3IM（Stochastic Structural Similarity）替代传统 SSIM：将渲染图像随机裁剪为多个 patch 后分别计算 SSIM，以捕获全局结构信息，而非仅在局部窗口内计算相似性。在联合训练阶段，所有高斯集合同时渲染但不增删，损失函数在隔离阶段损失基础上增加等距正则化（isometric regularization）损失和碰撞损失 $\mathcal{L}_{col}$（防止身体和服装高斯在空间上重叠）：

\begin{equation}
\mathcal{L}_{iso} = \sum_{i=1}^{|\mathcal{V}|} \sum_{j \in \mathcal{N}(i)} \left(\lambda_\mu |d(\mu^{(i)}, \mu^{(j)}) - d(\mu_o^{(i)}, \mu_o^{(j)})| + \lambda_\Sigma |d(\boldsymbol{\Sigma}^{(i)}, \boldsymbol{\Sigma}^{(j)}) - d(\boldsymbol{\Sigma}_o^{(i)}, \boldsymbol{\Sigma}_o^{(j)})|\right)
\end{equation}

其中 $d$ 为 $L_2$ 距离，下标 $o$ 表示规范姿态下的属性值，$\mathcal{N}(i)$ 为高斯 $i$ 的邻居集合。直觉上，等距正则化要求"变形前后，相邻高斯之间的距离和协方差差异应尽量保持不变"，这防止了高斯在变形过程中产生撕裂或堆叠等伪影。整个管线基于共享骨架驱动，身体和各件服装的高斯通过同一个骨架变形到目标姿态，但由于各自拥有独立的蒙皮权重，可以在保持空间对齐的同时实现独立的运动模式。

\textbf{设计动机}。GGAvatar 的三个关键设计各有其动机。(1)~独立服装模板的动机是：单层点偏移本质上是在身体表面"画"一层衣服，无法表达服装与身体之间的空间间隙（如宽松外套与身体之间的空隙），更无法将服装完整剥离。独立模板为每件服装提供了自己的几何基础。(2)~可学习蒙皮权重的动机是：SMPL 的蒙皮权重是为身体表面设计的，直接应用于服装会迫使服装高斯完全跟随身体骨骼运动，无法表达宽松服装的自由摆动。(3)~分阶段训练的动机是：如果同时优化身体和服装的高斯，由于 3DGS 的 $\alpha$-混合 渲染机制，身体高斯可能"长到"服装区域内部（因为渲染时前景高斯遮挡后景，身体高斯可以通过占据服装位置来降低重建损失），隔离训练通过冻结其他层来防止这种"入侵"行为。

\textbf{公式与参数}。GGAvatar 使用 SMPL 模型（$n_k = 23$ 个关节、$\beta \in \mathbb{R}^{10}$ 体型参数、$\theta \in \mathbb{R}^{(n_k+1) \times 3}$ 姿态参数）作为身体基础模板，通过 FrankMocap 估计姿态参数，SCHP 进行人体语义分割，ISP 模型生成服装网格模板。训练仅需单目视频输入。在 People Snapshot 数据集上，GGAvatar（约 20 分钟训练，单张 RTX 3080 Ti）在 male-3-casual 上达到 31.01 dB PSNR / 0.9812 SSIM，整体性能与 GaussianAvatar（约 30 分钟训练）等非分离模型相当，在多个序列上取得最高 PSNR。相比之下，同为分离式模型的 SCARF 在相同 20 分钟训练条件下仅达到约 18--20 dB PSNR，而完整训练需要约 40 小时（V100），质量仍不及 GGAvatar。NeRF 类方法（如 Neural Body）需要 14 小时至 1 天的训练时间，GGAvatar 的训练速度是其数百倍。渲染速度约 80 FPS（$540 \times 540$ 分辨率）。GGAvatar 实现了完全的服装分离能力——支持颜色编辑、单品服装转移（如单独替换上衣或裤子）和动画驱动等应用。

\paragraph{LayerGS}

LayerGS（2026.01）~\cite{2601_05853} 将任意姿态的人体分解为可动画的多层化身，每层编码为一组 2D 高斯实现准确的几何和光真实渲染。该方法采用三阶段训练策略：先训练单层规范化身重建外层服装的粗略形状，再冻结外层优化内层身体，最后冻结内层精修外层以解决层间交叠伪影。通过预训练 2D 扩散模型（Stable Diffusion V1.5 + ControlNet）的分数蒸馏采样修复被遮挡区域。训练完成后从每层 2D 高斯提取 3D 网格并将高斯绑定到对应三角面片，利用线性混合蒙皮驱动网格变形并附加 Laplacian 变形解决层间穿透，实现可动画的 3D 虚拟试穿。

\paragraph{Cloth-HUGS}

Cloth-HUGS（2026.04）~\cite{2604_15875} 在共享规范空间中使用独立的高斯层分别表示身体和服装，通过 SMPL 驱动的关节和学习的线性混合蒙皮权重变形规范体积。身体和服装各自拥有独立的 triplane 特征编码，经三个 MLP 解码器分别预测外观、几何校正和变形属性（包括可学习的 LBS 权重和姿态依赖偏移）。引入物理启发的约束——利用 SNUG 生成的布料网格提供双向 Chamfer 距离模拟一致性损失、ARAP 正则化保持局部布料结构、掩码监督确保轮廓一致性——以及深度感知的两通道渲染策略（先渲染身体和场景，再根据深度遮罩合成服装层），实现鲁棒的身体-服装-场景合成。

\subsubsection{虚拟试穿}

本路线介绍基于~3DGS~的三维虚拟服装试穿方法。

\paragraph{GaussianVTON}

GaussianVTON（2024.05）~\cite{2405_07472} 首次提出仅使用图像作为编辑提示的 3D 虚拟试穿管线，集成了高斯泼溅编辑和 2D 虚拟试穿。三阶段精修策略逐步解决面部模糊、服装不准确和视角质量下降的问题。

\textbf{Baseline 与问题}。前述 GGAvatar、LayerGS 等方法已经实现了服装与人体在 3DGS 中的分离表示，但它们聚焦于"从视频中重建并分解已有的服装"。虚拟试穿面临的是一个不同的问题：给定一件目标服装的图片，如何将其"穿"到一个已重建的 3D 人体场景上。已有的 3D 场景编辑方法（如 InstructNeRF2NeRF、GaussianEditor）依赖文本提示驱动编辑（例如"把短袖变成黑色"），但文本描述无法精确传达服装的具体款式、纹理和图案，导致编辑结果不准确。GaussianVTON 的核心洞察是：用图像而非文本作为编辑提示，可以精确控制试穿结果。本节介绍从二维编辑蒸馏到三维、从文本驱动到图像驱动的虚拟试穿方法。

\textbf{改了什么}。本方法修改了 3DGS 的场景编辑环节，用图像提示（image prompting）驱动的逐视角编辑替代文本提示驱动编辑，3DGS 的场景表示和高斯泼溅渲染管线保持不变。具体而言，GaussianVTON 将预训练的 2D 虚拟试穿扩散模型 LaDI-VTON 作为逐视角编辑器，对已重建 3DGS 场景的多视角训练图像逐一进行服装替换，然后通过重新优化 3D 高斯参数将编辑效果融入 3D 场景。值得注意的是，LaDI-VTON 本身需要多种辅助输入（姿态关键点、人体解析图、inpainting 掩码、dense pose），GaussianVTON 在管线中自动提取这些辅助数据，用户只需提供目标人物图像和目标服装图像即可完成全部操作。这种"2D 编辑 + 3D 重建"的范式避免了设计专用的 3D 试穿网络，任何基于扩散的 2D VTON 模型都可以即插即用。

\textbf{实现原理}。整个流程分为编辑和精修两个阶段。

第一阶段是核心编辑策略 Edit Recall Reconstruction（ERR）。在标准 3DGS 编辑流程中，常用的是 IterativeDU 策略：逐视角依次编辑、替换训练图像并立即渲染更新。问题在于，早期视角的编辑结果会影响后续视角的高斯参数，导致不同视角的编辑不一致——例如某个视角"凭空"增加了原本不存在的几何结构。ERR 的直觉很简单：先把所有视角的图片都编辑完，再统一更新数据集和渲染。具体地，设原始数据集为 $\{I_0^v\}$（$v$ 为视角索引），这些原始图像被独立缓存，并在所有后续编辑阶段中作为扩散模型的条件输入持续使用。ERR 用 LaDI-VTON 对所有视角同时进行编辑，得到 $\{I_1^v\}$，然后才统一更新数据集并重新渲染。与 IterativeDU 相比，ERR 的关键区别在于"全部编辑完再更新"——避免了编辑顺序不同导致的累积误差和不一致几何变化。同时，ERR 相比 SDS、DDS、PDS 等基于分数蒸馏的编辑策略，不会导致编辑视角相对于输入视图发生根本性变化，更适合需要精确区域控制的虚拟试穿任务。

第二阶段是三阶段精修策略（Three-Stage Refinement），逐步解决扩散模型引入的典型缺陷：

(1)~面部一致性精修（Face Consistency）。2D VTON 模型在替换服装时常导致面部模糊或失真。GaussianVTON 用 MediaPipe FaceMesh 提取编辑前面部关键点 $F^v = K(I_0^v)$，在编辑后将原始面部信息回传到 $I_1^v$，得到面部修正后的 $I_2^v$。

(2)~层级稀疏编辑（Hierarchical Sparse Editing）。多视角编辑的另一个痛点是部分视角的服装编辑出现错误（如颜色偏差、图案丢失）。GaussianVTON 用 Lang-Segment-Anything 分割出每个视角的服装区域 $M_W^v = S(I_2^v)$，然后对每三个相邻视角 $i, j, k$ 计算 dHash 相似度，找出与另外两个视角最不一致的视角 $\hat{v}$：

\begin{equation}
\hat{v} = \min_{i \neq j \neq k} \left(\max\left(\text{dHash}(I_{M_W}^i, I_{M_W}^j),\ \text{dHash}(I_{M_W}^i, I_{M_W}^k),\ \text{dHash}(I_{M_W}^j, I_{M_W}^k)\right)\right)
\end{equation}

直觉上，如果三个视角中有一个的服装外观与其余两个显著不同，说明该视角的编辑可能出错，需要用 LaDI-VTON 重新编辑。这种稀疏策略只重编有问题的视角，而非全部重做，兼顾了效率和准确性。

(3)~渲染优化（Rendering Optimization）。经过前两步精修的图像可能在质量上仍有退化。GaussianVTON 采用 DDNM（Denoising Diffusion Null-Space Model）去噪恢复图像质量，核心模型为 $y = Ax + n$，其中 $x$ 为原始图像，$A$ 为退化矩阵，$n$ 为噪声。最终得到精修视角图像 $I_4^v$，用于更新数据集并重新优化 3D 高斯参数。

\textbf{设计动机}。GaussianVTON 的设计围绕一个核心矛盾展开：2D 扩散模型编辑能力强但缺乏 3D 一致性，3DGS 渲染灵活但不具备编辑能力。通过"2D 编辑 + ERR 统一更新 + 三阶段精修"的级联设计，将 2D 模型的编辑能力逐步蒸馏到 3D 高斯中。图像提示而非文本提示的选择使得编辑精度从"语义级别"提升到"像素级别"。这对于虚拟试穿至关重要：一件碎花连衣裙的每个花纹细节都无法用文本精确描述。此外，ERR 策略避免了 IterativeDU 因编辑顺序不同导致累积误差的问题，而层级稀疏编辑则提供了高效的错误检测和修正机制。

\textbf{公式与参数}。GaussianVTON 在所有实验中使用自适应尺寸编辑（adaptive sizing editing），而非此前工作常用的固定 $512 \times 512$ 裁剪，确保自定义数据的编辑效果和可控区域不受影响。2D VTON 模型 LaDI-VTON 在 DressCode 数据集上预训练后直接使用，无需在多视角数据上进一步微调。高斯渲染器采用 GaussianEditor 的优化实现。在量化评估中，GaussianVTON 在 CLIP 文本-图像方向相似性、CLIP 图像-图像方向相似性、FID、SSIM、PSNR 和 LPIPS 六项指标上均优于基于文本提示的基线方法（InstructN2N 和 GaussianEditor）。

\paragraph{GS-VTON}

GS-VTON（2024.10）~\cite{2410_05259} 提出图像提示的 3D 虚拟试穿方法，核心分为个性化扩散模型适配和人物感知 3DGS 编辑两部分。首先通过 LoRA 微调 Stable Diffusion Inpainting 模型，利用参考驱动的多图联合编辑策略（将首张图的注意力特征注入其余图的去噪过程）生成跨视角一致的训练图像集。然后在 3DGS 编辑阶段引入人物感知自注意力机制，将 LoRA 微调模型的编辑方向与已编辑图像集的注意力特征按比例融合，保持不同视角下服装外观的一致性。该方法不限制图像分辨率（突破此前 $512 \times 512$ 的约束），单张 V100 GPU 上 LoRA 微调约 30 分钟、3DGS 编辑约 25 分钟。

\paragraph{LAGA}

LAGA（2024.05）~\cite{2405_12663} 实现了从文本描述创建可分解的高保真化身，每层对应特定的服装或人体本身。生成阶段采用由粗到精策略：先用稀疏高斯（5000 点）配合密度引导优化服装粗略形状，再通过带扰动的复制上采样和 SDS 损失细化高频细节；同时提出 dual-SDS 损失，同时优化单层局部图像和多层合成全局图像，确保各层服装之间的视觉连贯性。在服装转移阶段，引入三个正则化损失引导高斯运动：人体拟合损失使服装贴合目标体型轮廓，相似性损失保持转移前后服装的整体形状，可见性损失确保服装不被内层遮挡。该方法支持不同体型之间的服装自由转移，突破了此前仅限相似体型间转移的限制。

\subsubsection{服装重建}

本路线介绍基于~3DGS~的服装几何重建与网格提取方法。

\textbf{GarmentGS}（2025.05）~\cite{2505_02126}。本方法修改了 3DGS 的高斯基元优化与几何重建环节，用密集点云先验替代无引导的自适应密度控制，渲染损失与球谐函数保持不变。其核心目标是实现\emph{非水密}（non-watertight）的单层服装网格重建，使提取的网格直接可用于布料模拟和虚拟试穿。

\textbf{Baseline 做法与问题}。标准 3DGS 通过自适应密度控制（clone/split）逐步调整高斯基元的分布，但高斯基元的位置和形状完全由渲染损失驱动，缺乏对真实几何表面的约束。这导致两个突出问题：（1）3D 高斯椭球体是"胖"的体积元素，与服装这类薄片状表面对齐不佳，提取的网格常常出现内-外双层结构和碎片面；（2）重建结果多为水密网格，无法直接生成具有领口、袖口等功能性开口的可穿戴服装。已有的正则化方法（如 2DGS）虽能提升表面重建质量，但无法实现非水密重建；SDF 联合方法虽能部分实现非水密重建，但 SDF 训练速度慢，丧失了 3DGS 快速渲染的优势。

\textbf{改了什么}。GarmentGS 在 3DGS 训练过程中引入三步点云引导操作，将高斯基元精确约束到服装表面上，并在网格后处理阶段利用点云去除碎片。

\textbf{实现原理}。整体流程分为密集点云重建、高斯基元优化和网格去噪三个阶段。

（1）\emph{密集点云重建}。首先从多视角服装图像快速重建密集点云作为几何先验。传统 MVS 方法处理 100 张多视角图像需要约 2 小时，GarmentGS 对管线做了三项简化——仅保留位置和法线等几何信息（去掉颜色）、降低深度图分辨率、减少多视角一致性优化迭代次数——将重建时间压缩至约 10 分钟，同时保留足够的几何精度。

（2）\emph{高斯基元优化}。这是本方法的核心，包含三个按顺序执行的操作：先将高斯拉到表面、再将椭球压扁为薄片、最后旋转薄片使其法线对齐表面法线——三步协同将自由散布的 3D 高斯转化为紧贴服装表面的扁平圆盘。

\emph{移动高斯（Moving Gaussians）}。直觉上，标准 3DGS 的高斯基元散落在空间中，并不一定贴在物体表面上。GarmentGS 利用密集点云作为"锚点"，把每个高斯基元拉到最近的点云位置上。具体地，先用密集点云构建 K-D 树，对每个高斯基元查询欧氏距离最近的点云点，然后将高斯基元的均值直接移到该点：

\begin{equation}
P_{\text{Gaussian}} = Q_{\text{Point}}
\end{equation}

其中 $Q_{\text{Point}}$ 是点云中距 $P_{\text{Gaussian}}$ 最近的点。这确保高斯基元始终贴附在服装表面上。

\emph{展平高斯（Flattening Gaussians）}。直觉上，3D 高斯椭球是胖的体积，而服装是薄片。GarmentGS 将每个椭球沿最短轴方向压缩，使其变为"扁平的圆盘"，更好地贴合服装表面。3DGS 中协方差矩阵 $\Sigma_i = R_i S_i S_i^T R_i^T$ 的缩放矩阵 $S_i = \text{diag}(s_1, s_2, s_3)$ 控制椭球形状。将 $S_i$ 中最小尺度分量 $s_3$（即最短轴方向的尺度）设为趋近于零，即可将椭球压扁为椭圆盘：

\begin{equation}
S_{\text{thin}} = \text{diag}(s_1, s_2, s_3) \rightarrow \text{diag}(s_1, s_2, 0)
\end{equation}

对应的展平损失为 $\mathcal{L}_{\text{thin}} = (s_1, s_2, s_3) \cdot \mathbf{e}_1$，鼓励最小尺度分量不断缩小。

\emph{旋转高斯（Rotating Gaussians）}。展平后的椭圆盘还需要旋转使其法线对齐服装表面的法线方向，否则倾斜的圆盘仍无法精确贴合表面。GarmentGS 将每个椭圆盘的法线 $\mathbf{m}_i$ 与其对应点云点的法线 $\mathbf{n}_i$ 对齐，旋转损失为：

\begin{equation}
\mathcal{L}_{\text{normal}} = 1 - |\mathbf{m}_i \cdot \mathbf{n}_i|
\end{equation}

当两个法线完全平行时，内积为 $\pm 1$，损失为零。

最终的训练损失为三部分之和：

\begin{equation}
\mathcal{L} = \mathcal{L}_{\text{RGB}} + \lambda_{\text{thin}} \mathcal{L}_{\text{thin}} + \lambda_{\text{normal}} \mathcal{L}_{\text{normal}}
\end{equation}

其中 $\mathcal{L}_{\text{RGB}}$ 是标准 3DGS 的 RGB 重建损失（$\mathcal{L}_1$ + D-SSIM），$\lambda_{\text{thin}} = 100$，$\lambda_{\text{normal}} = 0.1$。展平损失的权重远大于法线损失，反映了"先压扁、再对齐"的优先级策略。

（3）\emph{网格去噪}。虽然高斯基元优化实现了非水密重建，但 TSDF 融合提取的网格仍存在内-外双层结构，内层主要由碎片面构成。GarmentGS 利用密集点云（严格分布在服装表面）训练 LOF（Local Outlier Factor，局部异常因子）异常检测模型，遍历网格顶点，移除被判定为异常的碎片面，最终得到单层的、可直接用于布料模拟的服装网格。

\textbf{设计动机}。GarmentGS 的设计围绕一个关键洞察展开：服装重建与一般物体重建有本质区别——服装必须是\emph{非水密}的（有领口、袖口等开口），且网格必须是\emph{单层}的（而非封闭体）。现有方法要么追求水密表面（SDF 类方法），要么虽能产生开放表面但网格质量差（正则化方法）。GarmentGS 选择密集点云而非 SDF 作为几何先验，兼顾了三点：点云本身天然是非水密的（只覆盖表面存在的区域），重建速度远快于 SDF 训练，且可以作为网格后处理的参考基准。三步优化（移动$\to$展平$\to$旋转）的顺序也有讲究：先将高斯拉到表面上（移动），再压扁为薄片（展平），最后调整朝向（旋转），每一步都为下一步创造条件。本节从服装-人体分离、虚拟试穿和服装重建三个角度，展现了 3DGS 显式表示在服装领域从分离到编辑再到重建的完整技术链。

\textbf{公式与参数}。GarmentGS 在 DeepFashion3D-v2 数据集上进行评估，每个场景训练 30,000 次迭代，所有实验在单张 RTX 3090 上完成。与 3DGS、GOF、GSDF、2DGS、PGSR 五个基线方法相比，GarmentGS 在 Chamfer Distance（几何精度）和 PSNR/SSIM/LPIPS（渲染质量）上均取得最优，并且是唯一同时实现非水密重建、单层网格提取和纹理映射的方法。消融实验表明，移动操作对 Chamfer Distance 的改善最为显著，旋转操作对表面法线一致性贡献最大，三个操作缺一不可。

\paragraph{PGC}

PGC（2025.03）~\cite{2503_20779} 仅需要单个静态帧的多视角捕获即可重建可模拟的服装，采用混合网格嵌入 3D 高斯泼溅表示。高斯的位置定义为相对三角面片的偏移量，使其随网格变形自动移动，捕获飞散纤维等近场高频细节。同时提出基于 Disney BRDF 的布料专用 PBR 着色模型（含织物特有的光泽散射分量），从输入图像中分解出空间变化的反照率和粗糙度等材质参数。渲染时采用高斯-PBR 混合策略：低频成分由 PBR 着色提供姿态依赖的光影效果，高频成分由 3DGS 补充织物纹理细节，使单姿态重建即可泛化到新姿态和新光照。

\paragraph{ReWeaver}

ReWeaver（2026.01）~\cite{2601_16672} 从稀疏多视角 RGB 图像重建拓扑准确的 3D 服装和缝纫图案，仅需四个视角即可工作。该方法采用双路径 Transformer 架构，通过可学习的面片查询和曲线查询预测 3D 接缝曲线、3D 表面面片及其拓扑连接关系，再经面内注意力模块将 3D 几何隐式展平为对应的 2D 缝纫面板。所有几何预测通过超网络实例化为连续参数化函数（MLP），支持任意密度的可微采样。生成的接缝和面板与多视角图像精确对齐，输出完整的 2D-3D 服装表示，可直接用于高保真物理模拟和机器人操作。

\begin{takeaway}
语音驱动、手部建模和服装重建将 3DGS 化身从被动重建推向主动交互：音频-面部映射实现实时说话人头像，混合网格嵌入高斯表示实现可模拟服装重建，这些扩展应用正持续丰富 3DGS 在人体领域的适用范围。
\end{takeaway}

%=====================================================================

\section{近期研究（2026 年 8--11 月）}
\paragraph{G-Skin：以生成姿态图像为高斯骨架绑定提供监督}
2026.08~\cite{2608_01726}。G-Skin 接收已有静态高斯和给定骨架，将二维生成模型产生的多姿态图像用作蒙皮监督，在测试对象上优化高斯随骨骼运动的权重。几何正则约束相邻位置的变形，减弱生成图像不准确时的绑定漂移。

作者用 50 个类人和 35 个非人形的生成资产评价。关联部位已经粘连或道具与身体拓扑融合时，局部连续性可能促成错误绑定。此流程仍需逐对象优化，不能称为一步预测骨架或蒙皮权重。

\paragraph{PD-GS：把离散音素和连续语音融合到高斯头像动画}
2026.08~\cite{2608_05218}。PD-GS 先建立规范姿态的高斯头部，再用连续音频和对齐的音素共同生成形变。融合模块在嘴唇闭合等发音关键区段更重视离散语言线索，连续音频仍提供上下文。

作者在 HDTF 的主要评价协议中使用两个常见对象，分别比较唇部关键点、同步、渲染质量和模型侧延迟。这属于个体特定头像实验，少样本人脸推广及舌等其他发音器官仍是后续问题。

\paragraph{4DAnyone：用骨架条件与跨视图合成支持单目人体四维重建}
2026.08~\cite{2608_20335}。4DAnyone 从单目视频恢复人体骨架，用骨架作为稀疏几何条件生成多视角视频。参考上下文压缩减少视角增加时的外观条件开销，目标视图的分组在去噪中动态关联，以减小各组之间的结构漂移。生成结果随后用于高斯人体重建。

测试覆盖 DNA-Rendering 和 DyMVHumans，并检查多视角合成与下游四维高斯结果。论文的失败例显示，远离骨架的宽松衣物可能跨视图不一致，错误骨架也可能被一致地传播到全部新视角。

\paragraph{ManifoldSplat：用自然语言改变绑定到高斯头像的局部几何}
2026.10~\cite{2610_03599}。ManifoldSplat 从单目视频建立可动画头像，根据文本生成局部 FLAME 网格形状增量，并传递给绑定于网格的高斯。几何编辑和渲染保持不同角色，表面一致性约束用于抑制形状改变后原语与网格脱节。

作者在 NeRSemble 子集保留正面相机，并用提示一致性、身份和渲染指标评价多种局部编辑。大形变会放大 FLAME 与原高斯的残余不匹配，实际编辑空间也受预定义形变原型限制。

\paragraph{HumanSkill-4D：用持续语义区域和关节轨迹回答头像查询}
2026.10~\cite{2610_24452}。HumanSkill-4D 把人体部位与衣物标签关联到已有码本中的持久高斯区域，让固定多模态模型能按语义取出原语。涉及动作时，模型把自然语言条件组合为关节轨迹上的运动和时间算子，返回匹配区域及时间段。流程不对单个头像另训练。

作者在扩展已有动作和人体数据的三个指代任务上评价。冻结解析器的类别和动作模板有明确边界，例如未登记的小饰品或动作无法仅凭自由文本可靠查询。

\paragraph{EGGHead：在固定原语预算内重建单图头部高斯}
2026.10~\cite{2610_26898}。EGGHead 学习按头部语义区域分配有限高斯，并借助 Voronoi 划分推断稀疏原语的局部形状。轻量对称模块补全遮挡区域，推断无需为每个新对象再执行逐个优化。

论文在不重合的 FFHQ 和 CelebA-HQ 图像上评价。模型未将纹理拆成反照率与照明，不能用于重打光；不均匀光线会影响对称补全，完全正脸时后脑仍可能不完整。

\paragraph{World2HOSI：让长流程人体与物体交互共享可更新状态}
2026.10~\cite{2610_27725}。World2HOSI 将静态高斯外观、可交互几何和持久物体身份组织成统一场景状态。导航、人与物体交互和人体与场景互动由不同专家合成，每段动作完成后的物体变化再写入状态，为下一段确定位置与目标。

论文在十个真实重建或合成场景的二十条日常指令上度量任务完成和几何质量，并展示有跨段依赖的动作。可选物理修正以专家给出的参考动作为前提，难以纠正本身不可行的大偏差计划。

\paragraph{复杂度先验引导通用人体高斯的分配与局部监督}
2026.10~\cite{2610_35152}。该方法从衣着人体网格的局部几何与纹理变化估计顶点复杂度，据此把更多高斯分配给褶皱或细节区，并在训练时更经常对这些区域抽取局部图像监督。整体图像监督仍保留，以防只优化高频区域。

作者用 THuman2.0、2K2K 和 X-Humans 的已标定渲染视图比较。这是静态多视图人体重建方法，逐帧应用到视频并不会自带时间一致性，也不能与单张无标定头像重建共享同一输入口径。


\section*{本章小结}
\addcontentsline{toc}{section}{本章小结}

\begin{futurebox}{本章总结与展望}{ch16_summary}
本章系统梳理了 3DGS 在人体化身与动画领域的研究进展，覆盖 258 篇论文和三个核心方向。

在面部与头部重建方面，3DGS 已经建立了从基础架构（基于参数化模型的混合表示）到高级功能（单图像重建、生成模型、重光照、头发建模）的完整技术栈。HyperGaussians 的高维扩展、ScaffoldAvatar 的局部表情建模和 GRMM 的残余几何建模代表了当前的技术前沿。

在身体与运动方面，研究重点从最初的可行性验证（3DGS-Avatar、HUGS）发展到精细化的变形建模、物理可信的衣物动画以及极具挑战性的野外单目视频处理。铰接物体重建、人体-物体交互和人体新视角合成为身体级 3DGS 提供了丰富的研究问题。

在语音驱动、局部建模与扩展应用方面，3DGS 的实时渲染能力使其成为说话人头像的理想选择。手部建模、服装试穿等方向则展现了 3DGS 显式表示在部件级解耦和可编辑性方面的独特优势。

展望未来，该领域的关键挑战包括：(1) 从单张图像实现全身+面部+手部的统一高质量化身创建；(2) 在移动设备上的实时高保真渲染；(3) 物理可信的动态建模（衣物、头发、软组织的交互）；(4) 多人场景中的实时化身渲染和交互。
\end{futurebox}
